% !TEX program = pdflatex
% =====================================================================
%  SAPHO — Referência Técnica Completa
%  A plataforma AURORA + YANC para concepção de soft-processors
%
%  Documento gerado a partir do ESTUDO DO CÓDIGO-FONTE dos repositórios
%  da organização nipscernlab (não da documentação existente).
%
%  Compilação:  pdflatex main && pdflatex main   (duas passadas p/ TOC)
%  Requer:      TeX Live 2021+ ou MiKTeX. Pacotes todos padrão.
% =====================================================================
\documentclass[11pt,a4paper,oneside]{report}

% ---------- Codificação e idioma ----------
\usepackage[utf8]{inputenc}
\usepackage[T1]{fontenc}
\usepackage[brazil]{babel}
\usepackage{lmodern}
\usepackage{textcomp}
\usepackage{microtype}

% ---------- Layout ----------
\usepackage[a4paper,margin=2.5cm,headheight=14pt]{geometry}
\usepackage{graphicx}
\usepackage{xcolor}
\usepackage{booktabs}
\usepackage{array}
\usepackage{tabularx}
\usepackage{longtable}
\usepackage{multirow}
\usepackage{enumitem}
\usepackage{csquotes}
\usepackage{listings}
\usepackage{fancyhdr}
\usepackage{titlesec}
\usepackage{parskip}
\usepackage[hidelinks]{hyperref}
\usepackage{url}

\hypersetup{
  pdftitle={SAPHO — Referência Técnica Completa (AURORA + YANC)},
  pdfauthor={NIPSCERN / UFJF},
  pdfsubject={Plataforma SAPHO de soft-processors},
  colorlinks=true,
  linkcolor=auroraDeep,
  citecolor=auroraDeep,
  urlcolor=auroraTeal
}

% ---------- Paleta AURORA (aurora borealis) ----------
\definecolor{auroraDeep}{HTML}{0B3D5B}   % azul profundo
\definecolor{auroraTeal}{HTML}{0E8F8F}   % verde-água
\definecolor{auroraMint}{HTML}{2FB89B}   % menta
\definecolor{auroraInk}{HTML}{1A2330}    % quase-preto
\definecolor{auroraGray}{HTML}{5A6675}
\definecolor{auroraLight}{HTML}{EEF4F6}
\definecolor{codeBg}{HTML}{F6F8FA}
\definecolor{codeKw}{HTML}{0B5FA5}
\definecolor{codeCom}{HTML}{6A8759}
\definecolor{codeStr}{HTML}{A31515}

% ---------- Colunas tabularx ----------
\newcolumntype{Y}{>{\raggedright\arraybackslash}X}
\newcolumntype{Z}{>{\raggedright\arraybackslash\small}X}

% ---------- Estilo de código (listings) ----------
\lstset{
  basicstyle=\ttfamily\small,
  keywordstyle=\color{codeKw}\bfseries,
  commentstyle=\color{codeCom}\itshape,
  stringstyle=\color{codeStr},
  backgroundcolor=\color{codeBg},
  frame=single,
  rulecolor=\color{auroraLight},
  framesep=4pt,
  breaklines=true,
  breakatwhitespace=false,
  postbreak=\mbox{\textcolor{auroraGray}{$\hookrightarrow$}\space},
  columns=fullflexible,
  keepspaces=true,
  showstringspaces=false,
  tabsize=2,
  upquote=true,
  extendedchars=true,
  literate=%
    {á}{{\'a}}1 {à}{{\`a}}1 {â}{{\^a}}1 {ã}{{\~a}}1
    {é}{{\'e}}1 {ê}{{\^e}}1 {í}{{\'i}}1
    {ó}{{\'o}}1 {ô}{{\^o}}1 {õ}{{\~o}}1
    {ú}{{\'u}}1 {ü}{{\"u}}1 {ç}{{\c c}}1
    {Á}{{\'A}}1 {Ã}{{\~A}}1 {É}{{\'E}}1 {Ó}{{\'O}}1 {Ç}{{\c C}}1
    {±}{{$\pm$}}1 {→}{{$\rightarrow$}}1 {≈}{{$\approx$}}1
}

% Linguagem auxiliar p/ a notação C± (CMM)
\lstdefinelanguage{CMM}{
  morekeywords={processor, in, out, void, comp, float, int, fixed, return,
    if, else, while, for, const, include, datatype, numBits, MABits,
    MIBits, gain},
  morecomment=[l]{//},
  morecomment=[s]{/*}{*/},
  morestring=[b]"
}
\lstdefinestyle{cmm}{language=CMM}

% ---------- Cabeçalhos ----------
\pagestyle{fancy}
\fancyhf{}
\fancyhead[L]{\small\itshape SAPHO — Referência Técnica}
\fancyhead[R]{\small\thepage}
\renewcommand{\headrulewidth}{0.4pt}
\renewcommand{\headrule}{\hbox to\headwidth{\color{auroraTeal}\leaders\hrule height \headrulewidth\hfill}}

% ---------- Títulos de capítulo ----------
\titleformat{\chapter}[display]
  {\normalfont\huge\bfseries\color{auroraDeep}}
  {\filright\Large\color{auroraTeal}Capítulo \thechapter}{8pt}{\Huge}
\titlespacing*{\chapter}{0pt}{10pt}{24pt}
\titleformat{\section}{\Large\bfseries\color{auroraDeep}}{\thesection}{0.6em}{}
\titleformat{\subsection}{\large\bfseries\color{auroraInk}}{\thesubsection}{0.6em}{}

% ---------- Macros de conveniência ----------
\newcommand{\cmm}{C\textpm{}}                      % nome da linguagem C±
\newcommand{\repo}[1]{\texttt{nipscernlab/#1}}
\newcommand{\file}[1]{\path{#1}}
\newcommand{\tool}[1]{\textsc{#1}}
\newcommand{\code}[1]{\lstinline[basicstyle=\ttfamily\normalsize]|#1|}

% Caixa de nota neutra (sem dependências externas)
\newenvironment{nota}
  {\par\medskip\noindent\begin{tabularx}{\linewidth}{!{\color{auroraTeal}\vrule width 3pt}>{\hspace{6pt}}X}}
  {\end{tabularx}\par\medskip}

\setlist{itemsep=2pt,topsep=4pt}

% =====================================================================
\begin{document}

% ----------------------------- CAPA -----------------------------
\begin{titlepage}
\centering
\vspace*{2.5cm}
{\color{auroraTeal}\rule{\linewidth}{2pt}}\\[10pt]
{\Huge\bfseries\color{auroraDeep} SAPHO}\\[6pt]
{\LARGE\color{auroraInk} Referência Técnica Completa}\\[14pt]
{\large\color{auroraGray} A plataforma \textbf{AURORA} (IDE) + \textbf{YANC} (compiladores)\\
para concepção, compilação, simulação e visualização de \emph{soft-processors}}\\[10pt]
{\color{auroraTeal}\rule{\linewidth}{2pt}}\\[2cm]

{\large\itshape Documento construído a partir do estudo direto do código-fonte\\
dos repositórios da organização \texttt{nipscernlab}.}\\[2cm]

{\large\bfseries NIPSCERN}\\[2pt]
{\normalsize Núcleo de pesquisa --- Universidade Federal de Juiz de Fora (UFJF)}\\[1.5cm]

{\normalsize Versão da suíte documentada: \textbf{SAPHO/AURORA 6.3.2}}\\[2pt]
{\normalsize Junho de 2026}

\vfill
{\footnotesize\color{auroraGray} Escopo: núcleo SAPHO (AURORA, YANC, \textit{toolchain bundle},
canal de distribuição \texttt{sapho}, \textit{forks} \texttt{gtkwave-nipscern} e
\texttt{netlistsvg-aurora}, e o subsistema \textit{Aurora Intelligence}).\par}
\end{titlepage}

% --------------------------- SUMÁRIO ---------------------------
\cleardoublepage
\tableofcontents
\cleardoublepage

% --------------------------- CAPÍTULOS ---------------------------
% !TEX root = ../main.tex
\chapter{Contexto e visão geral}
\label{cap:01-contexto-visao-geral}

Este capítulo situa a plataforma \tool{SAPHO} no seu contexto institucional,
descreve o que ela é, apresenta o seu modelo conceitual (a IDE \tool{AURORA},
os compiladores \tool{YANC} e o \emph{toolchain bundle}, mais os subsistemas
\tool{PRISM} e \tool{Aurora Intelligence}), identifica o público-alvo e oferece
um panorama de alto nível do fluxo de trabalho. Ele termina com um mapa dos 18
capítulos deste documento. Os detalhes de cada assunto ficam reservados aos
capítulos próprios; aqui o objetivo é dar a visão de conjunto que permite ler o
restante com orientação.

\begin{nota}
\textbf{Fonte da verdade.} As afirmações deste documento foram verificadas
diretamente no código-fonte dos repositórios da organização
\texttt{nipscernlab} (arquivos \path{package.json}, \path{CITATION.cff},
\path{README.md}, scripts de \emph{bootstrap}, módulos do processo principal do
Electron e os compiladores do \tool{YANC}). Quando a documentação histórica
diverge do código, o código prevalece.
\end{nota}

\section{O laboratório NIPS-CERN e a UFJF}
\label{sec:01-nipscern}

\tool{SAPHO} é desenvolvida pelo \textbf{NIPS-CERN}, um núcleo de pesquisa
sediado na \textbf{Universidade Federal de Juiz de Fora (UFJF)}, em Minas
Gerais, Brasil. O laboratório se localiza no prédio do PPEE --- Programa de
Pós-Graduação em Engenharia Elétrica da UFJF --- e mantém colaboração direta
com o CERN, contribuindo para o experimento ATLAS do LHC (instrumentação de
detectores, teste de PMTs e eletrônica do TileCal).

O grupo passou por uma evolução de identidade: foi originalmente conhecido como
\emph{LAPTEL} (Laboratório de Processamento de Sinais e Telecomunicações),
depois \emph{NIPS} (Núcleo de Instrumentação e Processamento de Sinais) e,
após a formalização da colaboração com o CERN, \textbf{NIPS-CERN}. A coordenação
é do Prof. Dr. Luciano Manhães de Andrade Filho, e a equipe reúne pesquisadores,
mestres, doutores e estudantes de iniciação científica --- um traço marcante do
laboratório é o envolvimento direto de graduandos em experimentos reais.

A autoria registrada da IDE \tool{AURORA} (arquivo \file{CITATION.cff} do
repositório \repo{aurora}) aponta para a afiliação \enquote{NIPS-CERN ---
Universidade Federal de Juiz de Fora}, e o endereço institucional é
\url{https://www.nipscern.com}. As instituições parceiras que aparecem no sítio
do laboratório incluem FAPEMIG, CAPES, CNPq, a própria UFJF e o CERN.

\subsection{Soft-processors em FPGA: o terreno técnico}
\label{sec:01-soft-processor}

\tool{SAPHO} pertence ao domínio do projeto de \emph{soft-processors}. Um
soft-processor é um processador descrito inteiramente em uma linguagem de
descrição de hardware (HDL, aqui \emph{Verilog}) e implementado sobre a malha
lógica reconfigurável de um FPGA, em vez de ser fabricado como um circuito
integrado fixo. Essa abordagem permite que a arquitetura do processador ---
largura de palavra, conjunto de instruções, periféricos, organização de memória
--- seja parametrizada e modificada livremente, recompilada e reimplementada em
minutos. É justamente essa maleabilidade que torna o soft-processor uma
ferramenta natural para pesquisa e ensino de arquitetura de computadores: o
aluno ou pesquisador projeta o hardware, escreve software para ele e observa o
comportamento, tudo dentro de um único ciclo iterativo.

\tool{SAPHO} adota um modelo de processador próprio, projetado para ser
modificado por quem o estuda. O fluxo completo --- da descrição em alto nível
até a imagem de memória e o \emph{testbench} --- roda localmente, sem dependência
de nuvem. A arquitetura do núcleo \tool{SAPHO} é tratada em detalhe no
\autoref{cap:02-arquitetura-soft-processor-sapho}.

\section{O que é a plataforma SAPHO}
\label{sec:01-o-que-e-sapho}

\textbf{SAPHO} é a suíte integrada para concepção, compilação, simulação e
visualização de soft-processors mantida pelo NIPS-CERN. O nome é um acrônimo cuja
expansão aparece de forma consistente no código e nos metadados do projeto:

\begin{nota}
\textbf{SAPHO = \emph{Scalable-Architecture Processor for Hardware Optimization}.}
Esta expansão é declarada no campo \code{description} do \file{package.json} do
repositório \repo{aurora} (\enquote{Scalable-Architecture Processor for Hardware
Optimization}), no \file{README.md} do repositório de distribuição \repo{sapho}
e na página de projeto do sítio do laboratório. Não foram encontradas no
código-fonte expansões alternativas concorrentes.
\end{nota}

Vale distinguir três usos do nome, que convivem no código:

\begin{description}[leftmargin=1.4em]
  \item[A plataforma/suíte.] \tool{SAPHO} como o ecossistema completo que une a
        IDE, os compiladores e o toolchain (o sentido predominante neste
        documento).
  \item[O canal de distribuição.] O repositório \repo{sapho} é o \emph{canal de
        release} estável de onde os usuários finais baixam o instalador. O
        repositório \repo{aurora}, por sua vez, é onde a IDE é
        \emph{desenvolvida}. Essa separação em dois repositórios é intencional
        (ver \autoref{sec:01-nomes} e o \autoref{cap:16-build-distribuicao}).
  \item[O produto empacotado.] O \file{package.json} do repositório \repo{aurora}
        nomeia o pacote como \code{sapho} e o \code{productName} como
        \code{SAPHO}; o instalador gerado chama-se
        \path{sapho-aurora-Setup-vX.Y.Z.exe}.
\end{description}

A versão da suíte documentada neste material é a \textbf{6.3.2} (campo
\code{version} do \file{package.json} e do \file{CITATION.cff}), e a distribuição
é licenciada sob \textbf{MIT}.

\section{O modelo conceitual: AURORA + YANC + toolchain bundle}
\label{sec:01-modelo-conceitual}

Conceitualmente, \tool{SAPHO} é a composição de três peças principais, mais dois
subsistemas que operam dentro da IDE. A separação de responsabilidades é nítida:
o \tool{YANC} compila, o \emph{toolchain bundle} simula e sintetiza, e a
\tool{AURORA} orquestra tudo por trás de uma única interface.

\begin{tabularx}{\linewidth}{l Y}
\toprule
\textbf{Peça} & \textbf{Papel} \\
\midrule
\tool{AURORA} & A IDE desktop (Electron + Monaco) onde o usuário escreve o
código, dispara as compilações, simula, vê ondas e inspeciona o RTL. É o
repositório vivo do desenvolvimento. \\
\tool{YANC} & Os compiladores que traduzem o código de alto nível até Verilog
sintetizável, imagens de memória e \emph{testbench}. \\
\emph{toolchain bundle} & O \emph{snapshot} portátil msys/mingw64 que reúne os
simuladores e sintetizadores de terceiros (Icarus Verilog, Verilator, Yosys,
GCC/G++, Python+cocotb) que a \tool{AURORA} baixa e empacota. \\
\bottomrule
\end{tabularx}

\subsection{AURORA --- a IDE}
\label{sec:01-aurora}

\tool{AURORA} é a IDE desktop da plataforma. Conforme o \file{package.json} e o
\file{README.md} do repositório \repo{aurora}, ela é construída sobre
\textbf{Electron} (versão 39 nas dependências de desenvolvimento) e o editor
\textbf{Monaco} (o mesmo motor de edição do VS Code). Ela unifica, sob uma
interface única: a escrita de \cmm{}, a montagem (\emph{assembling}), a síntese
de Verilog, a simulação, a visualização de formas de onda e a inspeção de RTL.

A IDE associa-se ao tipo de arquivo de projeto \textbf{\code{.spf}} (\emph{SAPHO
Project File}), declarado em \code{fileAssociations} no \file{package.json}, e
organiza o trabalho de forma orientada a processadores: cada projeto contém um
ou mais processadores, com codificação por cor e uma árvore de arquivos
unificada. A arquitetura interna da \tool{AURORA} é detalhada no
\autoref{cap:05-aurora-arquitetura}, e seus subsistemas funcionais (editor,
árvore, compilação, simulação) nos capítulos seguintes.

\subsection{YANC --- os compiladores}
\label{sec:01-yanc}

\textbf{YANC} (\enquote{Yet Another Compiler}) é a espinha dorsal de compilação
da plataforma. Segundo o \file{README.md} do repositório \repo{yanc}, ele recebe
um programa de alto nível --- escrito em \cmm{} (uma linguagem do tipo C, com
operadores de ponto fixo, ponto flutuante, números complexos e notação de Dirac
como cidadãos de primeira classe) ou em \emph{C++} --- e o compila por completo
até um núcleo \tool{SAPHO} sintetizável, as imagens de memória de programa e de
dados, e um \emph{testbench} pronto para executar.

Os binários do \tool{YANC} incluem, entre outros, \code{cmmcomp} (\cmm{} para
\emph{assembly}), \code{asmcomp} (\emph{assembly} para binário/hex),
\code{appcomp} e \code{comp2gtkw}; eles são distribuídos com bibliotecas HDL e
macros. O \tool{YANC} é desenvolvido em C, com Flex e Bison, e pode ser usado de
forma autônoma a partir de scripts de \emph{shell}, embora o uso principal seja
dirigido pela \tool{AURORA}. A linguagem \cmm{} é o tema do
\autoref{cap:03-linguagem-cmm}, e o interior dos compiladores, do
\autoref{cap:04-yanc-compiladores}.

\subsection{O toolchain bundle}
\label{sec:01-toolchain}

O \emph{toolchain bundle} é um \emph{snapshot} portátil de ferramentas de
terceiros que a \tool{AURORA} não traz no seu repositório fonte, mas baixa
durante o \emph{bootstrap}. Conforme o script
\file{components/Scripts/download-toolchain.js}, o bundle (arquivo
\code{aurora-msys-v1.zip}, distribuído pelo repositório \repo{aurora-toolchain})
reúne, num único diretório msys/mingw64, o Icarus Verilog (\code{iverilog},
\code{vvp}), o Verilator, o Yosys, GCC/G++, \code{make}, Perl e um Python com
cocotb (incluindo a VPI estática \path{libcocotbvpi_verilator.a}). O
\tool{YANC}, o \emph{fork} \repo{gtkwave-nipscern} e demais componentes são
baixados separadamente por scripts irmãos. O toolchain é descrito no
\autoref{cap:14-toolchain-bundle}, e o catálogo dos componentes de terceiros, no
\autoref{cap:15-catalogo-terceiros}.

\subsection{PRISM --- visualizador de RTL}
\label{sec:01-prism}

\textbf{PRISM} é o subsistema de visualização de RTL da \tool{AURORA}. Ele
renderiza a estrutura do processador e os caminhos de dados como diagramas de
netlist, a partir da síntese do \textbf{Yosys} combinada com o
\textbf{netlistsvg}. O \file{package.json} aponta a dependência
\code{@silimate/netlistsvg} para o \emph{fork} \repo{netlistsvg-aurora}, e o
netlistsvg roda em processo (a partir de \path{node_modules}). A interface do
PRISM vive em \path{html/prism/} (arquivos \path{prism.html}, \path{prism.js} e
\path{prism.css}). O subsistema é tratado em detalhe no
\autoref{cap:11-aurora-prism-rtl}.

\subsection{Aurora Intelligence --- subsistema de IA}
\label{sec:01-aurora-intelligence}

\textbf{Aurora Intelligence} é o subsistema de inteligência artificial integrado
à IDE. Ele está \textbf{completamente implementado}: o módulo
\file{main/ai/provider.js} mostra a integração via Vercel AI SDK, com suporte aos
provedores \code{openai}, \code{anthropic}, \code{google}, \code{deepseek},
\code{groq} e \code{ollama}, mais um servidor MCP (\emph{Model Context Protocol})
próprio (\file{main/ai/aurora_mcp_server.js}), um \emph{tool bridge}, gestão de
chaves (\code{keystore}), conversas, e a ponte com CLIs de agente
(\file{main/ai/claude_code.js}, \file{main/ai/codex_cli.js}).

\begin{nota}
\textbf{Status do provedor próprio.} Embora o subsistema esteja completo, o
\emph{provedor} e os \emph{modelos próprios} do laboratório ainda \textbf{não
foram treinados} --- esse treinamento é um item de \emph{roadmap}. Hoje, o
Aurora Intelligence opera por meio dos provedores de terceiros listados acima,
com as chaves de API fornecidas pelo usuário. O subsistema é detalhado no
\autoref{cap:12-aurora-intelligence-ia}.
\end{nota}

\section{Público-alvo}
\label{sec:01-publico}

\tool{SAPHO} destina-se a \textbf{pesquisa e ensino em arquitetura de
computadores e projeto de soft-processors}. A vertente de ensino é concreta: a
plataforma é empregada na disciplina de \textbf{DLP --- Dispositivos Lógicos
Programáveis} da UFJF, na qual os estudantes projetam, compilam e simulam
arquiteturas de processador customizadas usando o ecossistema. Isso aproxima a
plataforma de dois perfis complementares:

\begin{itemize}
  \item \textbf{Estudantes e docentes} de cursos de engenharia/computação que
        precisam de um caminho concreto entre os conceitos teóricos de
        arquitetura de computadores e a prática de projeto de hardware ---
        escrever um programa, compilá-lo para um processador customizado,
        simulá-lo e ver as formas de onda, sem montar um \emph{toolchain} de
        EDA manualmente.
  \item \textbf{Pesquisadores} que desejam experimentar variações de
        arquitetura (largura, ISA, periféricos) de forma rápida e iterativa,
        com todo o fluxo executando localmente.
\end{itemize}

\section{Panorama de alto nível do fluxo}
\label{sec:01-fluxo}

Em alto nível, o trabalho na plataforma percorre o caminho \emph{projeto $\to$
\cmm{} $\to$ compilação $\to$ simulação $\to$ ondas/RTL}. O usuário cria um
projeto (arquivo \code{.spf}) com um ou mais processadores, escreve o código em
\cmm{}, dispara a compilação a partir da barra de ferramentas, e então simula e
inspeciona os resultados. O diagrama a seguir resume o encadeamento conceitual;
os detalhes de cada etapa ficam para os capítulos correspondentes, e o pipeline
completo, ponta a ponta, é o tema do \autoref{cap:17-pipeline-ponta-a-ponta}.

\begin{lstlisting}[caption={Fluxo conceitual de alto nível da plataforma SAPHO.},captionpos=b]
  Projeto (.spf)            Cap. 7  — arvore de projeto / processadores
       |
       v
  Codigo C+- / C++          Cap. 3  — linguagem C+-    | Cap. 6 — editor
       |
       v
  Compilacao  (YANC)        Cap. 4  — cmmcomp -> asmcomp
   cmmcomp -> asmcomp       Cap. 8  — orquestracao da compilacao na AURORA
       |
       +--> Verilog sintetizavel + imagens de memoria (.mif) + testbench
       |
       v
  Simulacao                 Cap. 9  — Icarus Verilog / Verilator + cocotb
   (iverilog / verilator)
       |
       +--> formas de onda (.vcd / .fst)
       |
       v
  Visualizacao
   - Ondas  (GTKWave/Surfer)  Cap. 10
   - RTL    (PRISM)           Cap. 11  — Yosys + netlistsvg
\end{lstlisting}

A correspondência entre os passos do fluxo e os botões da barra de ferramentas
de compilação da \tool{AURORA}, conforme o \file{README.md} do repositório
\repo{aurora}, é:

\begin{tabularx}{\linewidth}{l l Y}
\toprule
\textbf{Passo} & \textbf{Botão} & \textbf{Ferramenta de apoio} \\
\midrule
\cmm{}  & \code{cmmcomp}  & \code{cmmcomp} + \code{asmcomp} (por processador) \\
Veri    & \code{vericomp} & Icarus Verilog (sintetiza o \emph{top-level}) \\
Wave    & \code{wavecomp} & GTKWave (abre o \code{.vcd}/\code{.fst} gerado) \\
PRISM   & \code{prismcomp}& Yosys + netlistsvg (visualizador de RTL) \\
\bottomrule
\end{tabularx}

A IDE distingue ainda os modos \emph{Processor} (um único processador ativo,
ciclo completo \cmm{} $\to$ ASM $\to$ IVL $\to$ Wave) e \emph{Project} (todos os
processadores do \code{.spf}), com a chave \emph{Compile-\&-Simulate} escolhendo
entre simulação dirigida por \emph{testbench} e síntese somente-Verilog. Esses
modos são detalhados no \autoref{cap:08-aurora-compilacao}.

\section{Convenção de nomes}
\label{sec:01-nomes}

Os nomes próprios da plataforma têm significados precisos, resumidos a seguir:

\begin{tabularx}{\linewidth}{l Y}
\toprule
\textbf{Nome} & \textbf{Significado} \\
\midrule
\tool{SAPHO} & \emph{Scalable-Architecture Processor for Hardware Optimization} --- a
plataforma/suíte completa; também o nome do canal de distribuição
(repositório \repo{sapho}) e do produto empacotado. \\
\tool{AURORA} & A IDE desktop (Electron + Monaco); repositório de
desenvolvimento \repo{aurora}. \\
\tool{YANC} & \emph{Yet Another Compiler} --- a suíte de compiladores
(\code{cmmcomp}, \code{asmcomp}, \code{appcomp}, \code{comp2gtkw});
repositório \repo{yanc}. \\
\tool{PRISM} & O visualizador de RTL embutido na IDE (Yosys + netlistsvg). \\
\tool{Aurora Intelligence} & O subsistema de IA da IDE. \\
\cmm{} & A linguagem de alto nível principal aceita pelo \tool{YANC}. \\
\bottomrule
\end{tabularx}

\begin{nota}
\textbf{Dois repositórios, de propósito.} O repositório \repo{aurora} é onde a
GUI é \emph{desenvolvida}; o repositório \repo{sapho} é o \emph{canal de
release} de onde os usuários finais baixam o instalador estável (a suíte =
\tool{YANC} + \tool{AURORA}). A divisão é intencional, não um defeito.
\end{nota}

\textbf{POLARIS} foi um sucessor planejado da IDE que está \textbf{descontinuado};
a IDE viva e mantida da plataforma é a \tool{AURORA}, e este documento não
descreve o POLARIS.

\section{Plataformas suportadas}
\label{sec:01-plataformas}

O suporte de plataforma difere entre as peças, conforme verificado no código:

\begin{itemize}
  \item \textbf{AURORA (IDE):} hoje, somente \textbf{Windows}. A configuração do
        \code{electron-builder} no \file{package.json} declara apenas o alvo
        \code{nsis} para \code{win} na arquitetura \code{x64}; não há alvos para
        macOS ou Linux. O \file{README.md} reforça que a IDE empacota um
        \emph{toolchain} específico de Windows e não é suportada em macOS ou
        Linux no momento.
  \item \textbf{YANC (compiladores):} \textbf{Windows e Linux}, com scripts de
        \emph{setup} próprios para cada sistema (\path{Scripts/setup.bat} no
        Windows, \path{Scripts/setup.sh} no Linux), conforme o \file{README.md}
        do repositório \repo{yanc}.
\end{itemize}

\section{Mapa do documento}
\label{sec:01-mapa}

Esta referência técnica está organizada em 18 capítulos, mais apêndices. O mapa
abaixo orienta a leitura:

\begin{longtable}{c >{\raggedright\arraybackslash}p{0.30\linewidth} >{\raggedright\arraybackslash}X}
\toprule
\textbf{\#} & \textbf{Título} & \textbf{Escopo} \\
\midrule
\endhead
1  & Contexto e visão geral & Este capítulo: laboratório, plataforma, modelo conceitual, fluxo, nomes e mapa. \\
2  & Arquitetura do soft-processor SAPHO & O modelo de processador \tool{SAPHO}: organização, parametrização, datapath. \\
3  & A linguagem \cmm{} & Sintaxe, tipos, diretivas de hardware e semântica da linguagem de alto nível. \\
4  & YANC --- os compiladores & Interior dos compiladores (\code{cmmcomp}, \code{asmcomp} e demais), Flex/Bison, bibliotecas HDL. \\
5  & AURORA --- arquitetura & Estrutura interna da IDE: processo principal (Electron), \emph{renderer}, IPC, módulos. \\
6  & AURORA --- editor & O editor Monaco, modos de edição, \emph{split}, introspecção de código. \\
7  & AURORA --- projeto e árvore & Projetos \code{.spf}, processadores, árvore de arquivos e hierarquia. \\
8  & AURORA --- compilação & Orquestração da compilação: modos, especificações de comando, \emph{builders}. \\
9  & AURORA --- simulação e ondas & Simulação com Icarus Verilog/Verilator e cocotb; geração de traços. \\
10 & AURORA --- visualização (GTKWave / Surfer) & Visualizadores de forma de onda integrados. \\
11 & AURORA --- PRISM (RTL) & O visualizador de RTL via Yosys + netlistsvg. \\
12 & AURORA --- Aurora Intelligence (IA) & O subsistema de IA: provedores, MCP, ferramentas, status do provedor próprio. \\
13 & AURORA --- produtividade & Recursos de produtividade da IDE (backups, projetos recentes, integrações). \\
14 & O toolchain bundle & O \emph{snapshot} msys/mingw64: conteúdo, \emph{bootstrap} e empacotamento. \\
15 & Catálogo de componentes de terceiros & Inventário das ferramentas FOSS embutidas e suas licenças. \\
16 & Build e distribuição & \emph{Build} com electron-builder, instalador NSIS, auto-atualização, split de repositórios. \\
17 & Pipeline ponta a ponta & O fluxo completo, do código à onda/RTL, integrando todas as peças. \\
18 & Engenharia e evolução & Práticas de engenharia, versionamento e evolução do projeto. \\
\bottomrule
\end{longtable}

\noindent Os apêndices reúnem um glossário de termos e uma lista de referências
externas.

% !TEX root = ../main.tex
\chapter{A arquitetura do soft-processor SAPHO}
\label{cap:02-arquitetura-soft-processor-sapho}

Este capítulo descreve o processador (\emph{soft-processor}) que o \emph{toolchain}
\tool{YANC} emite ao final da compilação. A fonte da verdade é o conjunto de módulos
Verilog em \path{yanc/HDL/} --- \file{processor.v}, \file{core.v}, \file{ula.v},
\file{instr\_dec.v}, \file{addr\_dec.v} e \file{myFIFO.v} --- e o \emph{back-end} do
montador (\emph{ASMComp}) em \path{yanc/Compilers/ASMComp/Sources/}, que parametriza e
instancia esses módulos, gera as imagens de memória (\path{.mif}) e produz o
\emph{testbench}. A geração do código \cmm/C++ até a montagem em Assembly é tratada em
\autoref{cap:03-pipeline-de-compilacao-yanc}; aqui o foco é o \emph{hardware} resultante.

\section{Visão geral do estilo de arquitetura}
\label{sec:arq-visao-geral}

O processador SAPHO é um núcleo de \textbf{datapath baseado em acumulador} (\emph{single
accumulator}) com \textbf{arquitetura Harvard}: memória de programa e memória de dados são
fisicamente separadas, cada uma com sua própria largura e seu próprio barramento de
endereço. O coração do datapath é um único registrador acumulador (\lstinline|racc|),
declarado em \file{core.v} como

\begin{lstlisting}[language=Verilog]
reg signed [NUBITS-1:0] racc;
// ...
always @ (posedge clk or posedge rst) if (rst) racc <= 0; else racc <= ula_out;
\end{lstlisting}

A cada ciclo o acumulador captura a saída da Unidade Lógico-Aritmética (a \tool{ULA}). A
maioria das instruções segue o padrão \emph{acumulador-com-memória}: um operando vem da
memória de dados (entrada \lstinline|in1| da ULA) e o outro é o próprio acumulador
(entrada \lstinline|in2|), e o resultado realimenta o acumulador. Há ainda uma
\textbf{pilha de dados} (\emph{data stack}) que fornece o segundo operando para as variantes
de instrução com prefixo \lstinline|S_| (\emph{stack}), permitindo expressões aninhadas sem
variáveis temporárias em memória.

Três características definem o estilo do núcleo:

\begin{description}[leftmargin=2.2em]
  \item[Totalmente parametrizável] Toda largura --- palavra, mantissa, expoente, operando,
    endereços, profundidade de pilhas, número de portas --- é um \lstinline|parameter|
    Verilog. O mesmo RTL gera desde um processador de 8 bits inteiros até um de dezenas de
    bits com ponto flutuante customizado.
  \item[Pague apenas pelo que usar] Cada instrução do conjunto é um \lstinline|parameter|
    booleano (0/1). O montador liga (\lstinline|.OPCODE(1)|) somente as instruções que o
    programa efetivamente usa; tudo o mais permanece em 0 e é eliminado pelo
    \lstinline|generate ... if| que envolve cada bloco. Assim, o circuito sintetizado contém
    apenas as unidades aritméticas e os caminhos de controle realmente exercitados.
  \item[Datapath de ciclo curto] Não há \emph{pipeline} de execução clássico com
    \emph{forwarding}; o fluxo é busca de instrução, decodificação registrada, seleção de
    operandos e operação combinacional na ULA, com o resultado registrado no acumulador.
    Registradores de \emph{prefetch}/decodificação introduzem latência fixa entre busca e
    escrita.
\end{description}

\begin{nota}
Os nomes dos parâmetros são abreviações em maiúsculas fixadas no RTL: \lstinline|NUBITS|
(largura da palavra), \lstinline|NBMANT| (bits de mantissa), \lstinline|NBEXPO| (bits de
expoente), \lstinline|NBOPER| (largura do campo operando), \lstinline|NBOPCO| (largura do
opcode, fixa em 7), \lstinline|MDATAS|/\lstinline|MINSTS| (tamanho das memórias de dados e
de instruções) e \lstinline|MDATAW|/\lstinline|MINSTW| (bits de endereço correspondentes,
derivados via \lstinline|$clog2|).
\end{nota}

\subsection{Hierarquia de módulos}
\label{sec:arq-hierarquia}

O topo sintetizável é o módulo \lstinline|processor| (\file{processor.v}), que conecta o
\lstinline|core| às duas memórias. O \lstinline|core| (\file{core.v}) contém o datapath e o
controle. A tabela a seguir resume a hierarquia confirmada no RTL.

\begin{longtable}{l l}
\toprule
\textbf{Módulo (arquivo)} & \textbf{Função} \\
\midrule
\endhead
\lstinline|processor| (\file{processor.v}) & Topo: instancia \lstinline|core|, \lstinline|mem_instr| e \lstinline|mem_data|. \\
\lstinline|mem_instr| (\file{processor.v}) & Memória de programa (Harvard), carregada por \lstinline|$readmemb|. \\
\lstinline|mem_data| (\file{processor.v})  & Memória de dados (Harvard), 1 leitura + 1 escrita por ciclo. \\
\lstinline|core| (\file{core.v})           & Datapath + controle: acumulador, pilha de dados, ULA. \\
\lstinline|instr_fetch| (\file{core.v})    & Busca de instrução: PC, \emph{prefetch}, pilha de subrotina. \\
\lstinline|pc| (\file{core.v})             & Contador de programa. \\
\lstinline|prefetch| (\file{core.v})       & Separa opcode/operando, decide saltos (JMP/JIZ/CAL/RET) e interrupção. \\
\lstinline|stack| (\file{core.v})          & Pilha genérica (usada para dados e para retorno de subrotina). \\
\lstinline|instr_dec| (\file{instr\_dec.v})& Decodificador: opcode $\rightarrow$ sinais de controle e \lstinline|ula_op|. \\
\lstinline|ula_in1_ctrl| (\file{core.v})   & Seleção do operando 1 (memória ou topo da pilha). \\
\lstinline|ula_in2_ctrl| (\file{core.v})   & Seleção do operando 2 (acumulador ou entrada de I/O). \\
\lstinline|mem_ctrl|/\lstinline|rel_addr| (\file{core.v}) & Endereçamento indireto e relativo (arrays, FFT, ponteiros). \\
\lstinline|io_ctrl| (\file{core.v})        & Controle das portas de entrada e saída. \\
\lstinline|ula| (\file{ula.v})             & Unidade Lógico-Aritmética (todas as operações). \\
\lstinline|addr_dec| (\file{addr\_dec.v})  & Decodificador de endereço de porta (1-de-N), no nível do \emph{wrapper}. \\
\lstinline|myFIFO| (\file{myFIFO.v})       & FIFO genérica (substitui IP de FIFO do fabricante). \\
\bottomrule
\end{longtable}

\subsection{Fluxo de dados do núcleo}
\label{sec:arq-fluxo}

O diagrama ASCII a seguir resume o caminho confirmado em \file{core.v}, da busca de
instrução até a realimentação do acumulador.

\begin{lstlisting}
       +----------------+      instr (NBOPCO+NBOPER bits)
       |  mem_instr     |----------------------------+
       | (programa)     |<--- instr_addr            |
       +----------------+                           v
              ^                            +--------------------+
              |                            |  instr_fetch       |
       +------+------+  pc_addr            |  (pc, prefetch,    |
       |  pc (PC)    |<-------- pc_load    |   stack subrotina) |
       +-------------+                     +----+----------+----+
                                               |opcode    |operand
                                               v          |
                                       +---------------+   |
                                       |  instr_dec    |   |  (endereco de dados)
                                       | -> ula_op,    |   |
                                       |    push/pop,  |   v
                                       |    wr,ldi,sti |  +--------------+
                                       +-------+-------+  |  mem_data    |
                                               |          | (dados)      |
            +--------------+   in1   +----------+----+     +------+-------+
            | data stack   |-------->| ula_in1_ctrl |<-----------+ mem_data_out
            | (sp)         |         +-------+-------+
            +------^-------+                 | in1
                   |                         v
                   |  sp_in           +------------+   out
                   +------------------|    ULA     |--------+
                   |                  +-----+------+        |
            +------+-------+  in2           ^               v
            | ula_in2_ctrl |---------------/         +-------------+
            | (acc / io_in)|                         |  racc (ACC) |
            +------^-------+<------------------------+-------------+
                   | io_in                            (realimenta in2)
\end{lstlisting}

\section{Memórias: arquitetura Harvard}
\label{sec:arq-memorias}

As duas memórias são instanciadas no topo (\file{processor.v}) e têm larguras e
profundidades independentes.

\subsection{Memória de instruções}
\label{sec:arq-mem-instr}

A memória de programa é o módulo \lstinline|mem_instr|. Cada palavra tem
\lstinline|NBOPCO + NBOPER| bits (opcode + operando) e é apenas de leitura em operação. O
conteúdo é carregado em simulação por \lstinline|$readmemb| a partir do arquivo
\path{.mif} de instruções; sob Yosys (síntese) o \lstinline|$readmemb| é ignorado por uma
guarda \lstinline|`ifdef YOSYS|.

\begin{lstlisting}[language=Verilog]
module mem_instr
#( parameter NADDRE = 8, parameter NBDATA = 12, parameter FNAME = "instr.mif" )
( input clk, input [$clog2(NADDRE)-1:0] addr, output reg [NBDATA-1:0] data = 0 );
  reg [NBDATA-1:0] mem [0:NADDRE-1];
  `ifdef YOSYS
  `else
    initial $readmemb(FNAME, mem);
  `endif
  wire wr = 0; // avoid unnecessary warnings
  always @ (posedge clk) begin
    data <= mem[addr];
    if (wr) mem[addr] <= 0;
  end
endmodule
\end{lstlisting}

\subsection{Memória de dados}
\label{sec:arq-mem-dados}

A memória de dados (\lstinline|mem_data|) tem palavra de \lstinline|NUBITS| bits, é
\textbf{signed} e oferece uma porta de leitura e uma de escrita simultâneas no mesmo ciclo
(endereços \lstinline|addr_rd| e \lstinline|addr_wr| independentes). A escrita é condicionada
ao sinal \lstinline|wr| produzido pelo decodificador.

\begin{lstlisting}[language=Verilog]
module mem_data
#( parameter NADDRE = 8, parameter NBDATA = 32, parameter FNAME = "data.mif" )
( input clk, input wr,
  input  [$clog2(NADDRE)-1:0] addr_rd, addr_wr,
  input  signed [NBDATA-1:0] data_in,
  output reg signed [NBDATA-1:0] data_out );
  reg [NBDATA-1:0] mem [0:NADDRE-1];
  `ifdef YOSYS
  `else
    initial $readmemb(FNAME, mem);
  `endif
  always @ (posedge clk) begin
    if (wr) mem[addr_wr] <= data_in;
    data_out <= mem[addr_rd];
  end
endmodule
\end{lstlisting}

No topo, as duas memórias são ligadas ao núcleo:

\begin{lstlisting}[language=Verilog]
mem_instr #(.NADDRE(MINSTS), .NBDATA(NBOPCO+NBOPER), .FNAME(IFILE))
          minstr(clk, instr_addr, instr);
mem_data  #(.NADDRE(MDATAS), .NBDATA(NUBITS), .FNAME(DFILE))
          mdata(clk, mem_wr, mem_addr_rd, mem_addr_wr, mem_data_out, mem_data_in);
\end{lstlisting}

\section{Busca de instruções, PC e saltos}
\label{sec:arq-fetch}

\subsection{Contador de programa}
\label{sec:arq-pc}

O contador de programa (módulo \lstinline|pc|) é um registrador de \lstinline|MINSTW| bits
que, a cada ciclo, ou incrementa o endereço atual ou carrega um novo valor (\lstinline|load|).
O incremento é sempre por um.

\begin{lstlisting}[language=Verilog]
wire [NBITS-1:0] um  = {{NBITS-1{1'b0}}, {1'b1}};
wire [NBITS-1:0] val = (load) ? data : addr;
always @ (posedge clk or posedge rst) begin
  if (rst) addr <= 0;
  else     addr <= val + um;
end
\end{lstlisting}

\subsection{Prefetch, opcode e operando}
\label{sec:arq-prefetch}

O módulo \lstinline|prefetch| fatia a palavra de instrução em opcode e operando, e centraliza
a lógica de salto. A instrução tem o opcode nos bits mais significativos e o operando nos
menos significativos:

\begin{lstlisting}[language=Verilog]
assign opcode  = mem_instr[NBOPCO+NBOPER-1:NBOPER];
assign operand = mem_instr[NBOPER       -1:     0];
\end{lstlisting}

Os saltos são decididos diretamente sobre o opcode (sem passar pela ULA). O \lstinline|JMP|
(opcode 15) é incondicional; o \lstinline|JIZ| (opcode 16, \emph{jump if zero}) salta quando o
acumulador é zero --- o sinal \lstinline|is_um| indica acumulador não nulo:

\begin{lstlisting}[language=Verilog]
// JMP                                       JIZ
assign wJMP = (opcode == 5'd15) | ((opcode == 5'd16) & ~is_um);
\end{lstlisting}

Quando a chamada de subrotina (\lstinline|CAL|) está habilitada, o \lstinline|prefetch|
reconhece também \lstinline|CAL| (opcode 17) e \lstinline|RET| (opcode 18), gerando os
sinais de \emph{push}/\emph{pop} para a pilha de subrotina:

\begin{lstlisting}[language=Verilog]
wire wCAL = (opcode == 5'd17);
wire wRET = (opcode == 5'd18);
assign pc_load  = wJMP | wCAL | wRET;
assign isp_push = wCAL;
assign isp_pop  = wRET;
\end{lstlisting}

\subsection{Interrupção e marcador \texttt{\#TOAQUI}}
\label{sec:arq-itr-toaqui}

Dois mecanismos opcionais são parametrizados por endereço. Quando \lstinline|ITRADD > 0|, a
entrada \lstinline|itr| força o PC a saltar para o endereço de interrupção; quando
\lstinline|TOAQUIADDR > 0|, a saída \lstinline|cheguei| pulsa no ciclo em que o endereço
buscado iguala o marcador (\lstinline|instr_addr == TOAQUIADDR|), permitindo instrumentar o
ponto exato de execução.

\subsection{Pilha de subrotina e pilha de dados}
\label{sec:arq-stacks}

Ambas reusam o mesmo módulo genérico \lstinline|stack|, parametrizado por largura
(\lstinline|NBITS|) e profundidade (\lstinline|DEPTH|). A pilha de subrotina
(\lstinline|isp|, profundidade \lstinline|SDEPTH|) guarda endereços de retorno e só é
sintetizada quando \lstinline|CAL| está presente. A pilha de dados (\lstinline|sp|,
profundidade \lstinline|DDEPTH|, largura \lstinline|NUBITS|) fornece o segundo operando às
instruções de variante \lstinline|S_|/\lstinline|SF_|.

\begin{lstlisting}[language=Verilog]
reg  [NADDR-1:0] pointer = 0;
wire [NADDR-1:0] pmaisum = pointer + um;
wire [NADDR-1:0] pmenoum = pointer - um;
always @ (posedge clk or posedge rst) begin
  if      (rst ) pointer <= zero;
  else if (push) pointer <= pmaisum;
  else if (pop ) pointer <= pmenoum;
end
always @ (posedge clk) if (push) mem[pointer] <= in;
assign                     out = mem[pmenoum];
\end{lstlisting}

\section{Decodificação de instruções}
\label{sec:arq-decod}

O decodificador (\file{instr\_dec.v}) recebe o opcode de 7 bits e produz: os sinais de
\emph{push}/\emph{pop} das pilhas, o código de operação da ULA (\lstinline|ula_op|, 6 bits),
o \emph{enable} de escrita em memória (\lstinline|mem_wr|), os controles de I/O
(\lstinline|req_in|, \lstinline|out_en|) e os controles de endereçamento indireto
(\lstinline|ldi|, \lstinline|sti|, \lstinline|fft|, \lstinline|lda|, \lstinline|sta|).

\subsection{Reconhecimento condicional de cada opcode}
\label{sec:arq-decod-cond}

Cada instrução só gera lógica de reconhecimento quando seu parâmetro está ligado. O padrão,
repetido para todos os opcodes, é um \lstinline|generate|:

\begin{lstlisting}[language=Verilog]
wire wLOD; generate if ((LOD) != 0)
  begin : dec_wlod assign wLOD = opcode == 7'd00; end
  else begin : dec_wlod assign wLOD = 1'b0; end endgenerate
\end{lstlisting}

Os sinais agregados (por exemplo, \lstinline|mem_wr|, \lstinline|push|, \lstinline|pop|) são
ORs dos reconhecedores das instruções que os disparam. Por exemplo, a escrita em memória é
ativada por \lstinline|SET|, \lstinline|SET_P|, \lstinline|STI|, \lstinline|ISI| e
\lstinline|STA|:

\begin{lstlisting}[language=Verilog]
assign mem_wr = wSET | wSET_P | wSTI | wISI | wSTA;
\end{lstlisting}

\subsection{Geração do código da ULA (\lstinline|ula_op|)}
\label{sec:arq-decod-ula}

O campo \lstinline|ula_op| (6 bits) é montado bit a bit. Cada bit (\lstinline|b0|..\lstinline|b5|)
é o OR dos reconhecedores das instruções cujo código de ULA tem aquele bit em 1. O valor
final é registrado (atraso de um ciclo) antes de chegar à ULA:

\begin{lstlisting}[language=Verilog]
always @ (posedge clk) ula_op <= {b5,b4,b3,b2,b1,b0};
\end{lstlisting}

\section{Conjunto de instruções (ISA)}
\label{sec:arq-isa}

O opcode tem \textbf{7 bits} (\lstinline|NBOPCO = 7|, fixo em \file{processor.v} e replicado
como \lstinline|NBITS_OPC| no montador, em \file{eval.c}). A associação \emph{mnemônico
$\rightarrow$ opcode} está no analisador léxico do montador
(\file{ASMComp.l}); a associação \emph{opcode $\rightarrow$ \lstinline|ula_op|} está no
decodificador (\file{instr\_dec.v}). As convenções de sufixo são:

\begin{description}[leftmargin=2.6em]
  \item[\texttt{S\_}] segundo operando vindo da \textbf{pilha de dados} (em vez da memória).
  \item[\texttt{F\_}] versão em \textbf{ponto flutuante}.
  \item[\texttt{SF\_}] ponto flutuante com operando vindo da pilha.
  \item[\texttt{P\_}] faz \textbf{PUSH} antes da operação.
  \item[\texttt{\_M}] opera sobre o valor de \textbf{memória} (em vez do acumulador).
\end{description}

\subsection{Controle de fluxo, memória e I/O}
\label{sec:arq-isa-ctrl}

\begin{longtable}{l l l}
\toprule
\textbf{Op.} & \textbf{Mnemônico} & \textbf{Descrição} \\
\midrule
\endhead
0  & \lstinline|LOD|    & Carrega o acumulador com dado da memória. \\
1  & \lstinline|P_LOD|  & PUSH seguido de LOD. \\
2  & \lstinline|LDI|    & Load com endereçamento indireto (índice no acumulador). \\
3  & \lstinline|ILI|    & Load indireto com reversão de bits (FFT). \\
4  & \lstinline|SET|    & Armazena o acumulador na memória. \\
5  & \lstinline|SET_P|  & SET seguido de POP. \\
6  & \lstinline|STI|    & Store com endereçamento indireto (índice na pilha). \\
7  & \lstinline|ISI|    & Store indireto com reversão de bits (FFT). \\
8  & \lstinline|PSH|    & Empurra o acumulador na pilha de dados. \\
9  & \lstinline|POP|    & Retira da pilha de dados para o acumulador. \\
10 & \lstinline|INN|    & Entrada de dado (porta de I/O). \\
11 & \lstinline|F_INN|  & Entrada em ponto flutuante (aplica I2F). \\
12 & \lstinline|P_INN|  & PUSH seguido de INN. \\
13 & \lstinline|PF_INN| & PUSH seguido de F\_INN. \\
14 & \lstinline|OUT|    & Saída de dado (porta de I/O). \\
15 & \lstinline|JMP|    & Salto incondicional. \\
16 & \lstinline|JIZ|    & Salta se o acumulador for zero. \\
17 & \lstinline|CAL|    & Chamada de subrotina (empilha retorno). \\
18 & \lstinline|RET|    & Retorno de subrotina (desempilha). \\
96 & \lstinline|NOP|    & Nenhuma operação. \\
102& \lstinline|LDA|    & \lstinline|acc = mem[acc]| (indireto base-less). \\
103& \lstinline|STA|    & \lstinline|mem[topo_pilha] = acc|, com POP. \\
\bottomrule
\end{longtable}

\begin{nota}
O mnemônico \lstinline|LEA| (\emph{load effective address}) não é um opcode próprio: no
montador ele se expande para um \lstinline|LOD| de uma constante igual ao endereço de uma
variável (\file{eval.c}, função \lstinline|instr_lea|). Isso permite passar o endereço-base
de um array para uma subrotina, que então o desreferencia em tempo de execução com
\lstinline|ADD| + \lstinline|LDA|/\lstinline|STA|.
\end{nota}

\subsection{Operações aritméticas, lógicas e de comparação}
\label{sec:arq-isa-arit}

A tabela a seguir lista os opcodes de cálculo (variantes \lstinline|S_|, \lstinline|F_|,
\lstinline|SF_|, \lstinline|P_| e \lstinline|_M| compartilham a unidade aritmética da base).
A coluna \lstinline|ula_op| é o código que o decodificador entrega à ULA (de
\file{instr\_dec.v}).

\begin{longtable}{l l l l}
\toprule
\textbf{Op.} & \textbf{Mnemônico} & \textbf{ula\_op} & \textbf{Operação} \\
\midrule
\endhead
19--22  & \lstinline|ADD|/\lstinline|S_ADD|/\lstinline|F_ADD|/\lstinline|SF_ADD| & 2 / 3 & Soma (inteira / float). \\
23--26  & \lstinline|MLT| .. \lstinline|SF_MLT| & 4 / 5 & Multiplicação (inteira / float). \\
27--30  & \lstinline|DIV| .. \lstinline|SF_DIV| & 6 / 7 & Divisão (inteira / float). \\
31--32  & \lstinline|MOD|/\lstinline|S_MOD|     & 8     & Resto da divisão (inteiro). \\
33--36  & \lstinline|SGN| .. \lstinline|SF_SGN| & 9 / 10& Copia o sinal de um operando ao outro. \\
37--42  & \lstinline|NEG| .. \lstinline|PF_NEG_M|& 11--14& Negação (inteira / float). \\
43--48  & \lstinline|ABS| .. \lstinline|PF_ABS_M|& 15--18& Valor absoluto (inteiro / float). \\
49--54  & \lstinline|PST| .. \lstinline|PF_PST_M|& 19--22& Zera se negativo (\emph{positive set}). \\
55--57  & \lstinline|NRM|/\lstinline|NRM_M|/\lstinline|P_NRM_M|& 23 / 24& Divisão por constante \lstinline|NUGAIN|. \\
58--60  & \lstinline|I2F|/\lstinline|I2F_M|/\lstinline|P_I2F_M|& 25 / 26& Inteiro $\rightarrow$ float. \\
61--63  & \lstinline|F2I|/\lstinline|F2I_M|/\lstinline|P_F2I_M|& 27 / 28& Float $\rightarrow$ inteiro. \\
64--69  & \lstinline|AND| .. \lstinline|S_XOR|  & 29--31& Bit a bit: AND, OR, XOR. \\
70--72  & \lstinline|INV|/\lstinline|INV_M|/\lstinline|P_INV_M|& 32 / 33& Inversão bit a bit (NOT). \\
73--76  & \lstinline|LAN| .. \lstinline|S_LOR|  & 34 / 35& Lógico booleano AND/OR. \\
77--79  & \lstinline|LIN|/\lstinline|LIN_M|/\lstinline|P_LIN_M|& 36 / 37& Inversão lógica (NOT booleano). \\
80--83  & \lstinline|LES| .. \lstinline|SF_LES| & 38 / 39& Menor que (inteiro / float). \\
84--87  & \lstinline|GRE| .. \lstinline|SF_GRE| & 40 / 41& Maior que (inteiro / float). \\
88--89  & \lstinline|EQU|/\lstinline|S_EQU|     & 42    & Igualdade (inteiro e float). \\
90--91  & \lstinline|SHL|/\lstinline|S_SHL|     & 43    & Deslocamento à esquerda. \\
92--93  & \lstinline|SHR|/\lstinline|S_SHR|     & 44    & Deslocamento lógico à direita. \\
94--95  & \lstinline|SRS|/\lstinline|S_SRS|     & 45    & Deslocamento aritmético à direita. \\
97      & \lstinline|F_ROT|                     & 46    & Aproximação de raiz quadrada (potência de 2). \\
98--101 & \lstinline|F_SU1| .. \lstinline|SF_SU2|& 47 / 48& Subtração em float (no operando 1 ou 2). \\
104--105& \lstinline|F_SCL|/\lstinline|SF_SCL|  & 49    & Escala float por $2^{k}$. \\
106--107& \lstinline|XPO|/\lstinline|XPO_M|     & 50 / 51& Expoente base-2 do float como inteiro. \\
\bottomrule
\end{longtable}

\section{A Unidade Lógico-Aritmética (ULA)}
\label{sec:arq-ula}

A ULA (\file{ula.v}) é organizada como um banco de submódulos especializados (um por
operação) cujas saídas convergem para um multiplexador final selecionado por
\lstinline|ula_op|. Cada submódulo só é instanciado quando seu parâmetro está ligado; caso
contrário, sua saída recebe \lstinline|{NUBITS{1'bx}}| (don't-care), e a síntese descarta o
caminho. A interface do módulo é simples:

\begin{lstlisting}[language=Verilog]
module ula
#( parameter NUBITS = 32, parameter NBMANT = 23, parameter NBEXPO = 8,
   parameter signed [NUBITS-1:0] NUGAIN = 64, /* ... flags por operacao ... */ )
( input [5:0] op,
  input  signed [NUBITS-1:0] in1, in2,
  output signed [NUBITS-1:0] out );
\end{lstlisting}

\subsection{Aritmética inteira}
\label{sec:arq-ula-int}

As operações inteiras são diretas e tratam a palavra como \lstinline|signed [NUBITS-1:0]|.
Exemplos confirmados no RTL:

\begin{lstlisting}[language=Verilog]
// ADD            // MLT             // DIV             // MOD
assign out = in1 + in2;   assign out = in1 * in2;
assign out = in1 / in2;   assign out = in1 % in2;
\end{lstlisting}

A negação é \lstinline|-in|; o valor absoluto é \lstinline|(in[NUBITS-1]) ? -in : in|; o
\emph{positive set} (\lstinline|PST|) zera valores negativos. A normalização (\lstinline|NRM|)
divide por uma constante de hardware \lstinline|NUGAIN|:

\begin{lstlisting}[language=Verilog]
module ula_nrm #( parameter NUBITS = 32, parameter signed NUGAIN = 1 )
( input signed [NUBITS-1:0] in, output signed [NUBITS-1:0] out );
  assign out = in/NUGAIN;
endmodule
\end{lstlisting}

\subsection{Operações lógicas, condicionais, comparações e deslocamentos}
\label{sec:arq-ula-logic}

As operações bit a bit são \lstinline|&|, \lstinline^|^, \lstinline|^| e \lstinline|~|. As
condicionais booleanas tratam ``zero'' como falso: \lstinline|LAN| produz 1 se ambos forem
diferentes de zero, \lstinline|LOR| produz 1 se algum for diferente de zero, \lstinline|LIN|
inverte a condição. As comparações (\lstinline|LES|, \lstinline|GRE|, \lstinline|EQU|)
devolvem a palavra estendida com zeros à esquerda e 1 no bit menos significativo quando
verdadeiro:

\begin{lstlisting}[language=Verilog]
// LES (menor que, com sinal)
assign out = {{(NUBITS-1){1'b0}}, (in1 < in2)};
\end{lstlisting}

Os deslocamentos são \lstinline|<<| (\lstinline|SHL|), \lstinline|>>| lógico (\lstinline|SHR|)
e \lstinline|>>>| aritmético com sinal (\lstinline|SRS|).

\subsection{Ponto flutuante customizado}
\label{sec:arq-ula-float}

O SAPHO usa um formato de ponto flutuante \textbf{próprio}, parametrizável (não o IEEE-754).
A palavra de \lstinline|NUBITS| bits é dividida em:

\begin{itemize}[leftmargin=1.6em]
  \item \textbf{1 bit de sinal} (o mais significativo);
  \item \textbf{\lstinline|NBEXPO| bits de expoente} com sinal (em complemento de dois);
  \item \textbf{\lstinline|NBMANT| bits de mantissa}, com o bit ``1'' líder
    \emph{armazenado explicitamente} (não há bit implícito como no IEEE).
\end{itemize}

A consistência \lstinline|NUBITS == NBMANT + NBEXPO + 1| é verificada pelo montador em
\file{eval.c} (\lstinline|eval_finish|); se violada, a compilação aborta. A conversão de e
para o formato é feita em software pelo montador (rotinas \lstinline|f2mf|/\lstinline|mf2f|
em \file{t2t.c}): \lstinline|f2mf| desempacota um \lstinline|float| IEEE de 32 bits, reescala
o expoente para \lstinline|NBEXPO| e a mantissa para \lstinline|NBMANT| (com arredondamento e
renormalização), e empacota no formato SAPHO.

O hardware opera sobre esse formato. A \textbf{denormalização} (\lstinline|ula_denorm|)
alinha os expoentes de dois operandos antes da soma/comparação, deslocando a mantissa de
menor expoente:

\begin{lstlisting}[language=Verilog]
// equaliza expoentes: o menor expoente desloca para casar com o maior
wire signed [EXP:0] eme    =  e1_in-e2_in;
wire                ege    =  eme[EXP];
wire        [EXP:0] shift2 = (ege) ? {EXP+1{1'b0}} : eme;
wire        [EXP:0] shift1 = (ege) ? -eme          : {EXP+1{1'b0}};
always @ (*) e_out = (ege) ? e2_in : e1_in; // expoente final = o maior
\end{lstlisting}

A \textbf{normalização} pós-operação (\lstinline|ula_norm|) desloca a mantissa à esquerda até
que o bit mais significativo seja 1, ajustando o expoente. A multiplicação em float
(\lstinline|ula_fmlt|) soma os expoentes e multiplica as mantissas; a divisão
(\lstinline|ula_fdiv|) subtrai os expoentes e divide as mantissas estendidas. A subtração é
realizada pelo próprio somador, com inversão de sinal do operando indicado
(\lstinline|F_SU1|/\lstinline|F_SU2|):

\begin{lstlisting}[language=Verilog]
wire su1 = (F_SU1 != 0) & (op == 6'd47); // inverte o sinal de in1
wire su2 = (F_SU2 != 0) & (op == 6'd48); // inverte o sinal de in2
\end{lstlisting}

Há ainda operações de ``cirurgia de expoente'': \lstinline|F_SCL| escala um float por $2^{k}$
somando \lstinline|k| ao expoente (com saturação à faixa do expoente), e \lstinline|XPO|
extrai o expoente base-2 como inteiro. A conversão inteiro$\leftrightarrow$float
(\lstinline|I2F|/\lstinline|F2I|) e a aproximação de raiz quadrada (\lstinline|F_ROT|, potência
de dois mais próxima) completam o repertório de ponto flutuante.

\subsection{Números complexos}
\label{sec:arq-ula-complex}

Não há uma unidade de hardware dedicada a números complexos no datapath. Um valor complexo é
representado como um \textbf{par de variáveis} float (parte real e parte imaginária); o
compilador e a biblioteca os manipulam por operações sobre cada componente. No nível de
simulação, o gerador de HDL (\file{hdl.c}) reconhece os pares (sufixos de parte real/imaginária)
e os junta em um sinal \lstinline|comp_<nome>| de largura dupla, apenas para visualização de
forma de onda --- não é parte do circuito sintetizado.

\subsection{O multiplexador final}
\label{sec:arq-ula-mux}

O submódulo \lstinline|ula_mux| seleciona, por \lstinline|op|, qual resultado vai à saída. As
operações de ponto flutuante que exigem normalização (soma, multiplicação, divisão,
\lstinline|I2F|, \lstinline|F_ROT|, subtrações) compartilham a saída normalizada
\lstinline|smx|, produzida por \lstinline|norm_mux|:

\begin{lstlisting}[language=Verilog]
always @ (*) case (op)
  6'd0 : out = in2 ;   // NOP
  6'd1 : out = in1 ;   // LOD
  6'd2 : out = add ;   // ADD
  6'd3 : out = smx ;   // F_ADD (normalizada)
  6'd4 : out = mlt ;   // MLT
  // ...
  6'd42: out = equ ;   // EQU
  6'd43: out = shl ;   // SHL
  default: out = {NUBITS{1'bx}};
endcase
\end{lstlisting}

\section{Entrada e saída}
\label{sec:arq-io}

\subsection{Portas, \emph{request} e \emph{enable}}
\label{sec:arq-io-portas}

A interface de I/O do núcleo expõe um barramento de entrada (\lstinline|io_in|), um de saída
(\lstinline|io_out|, que espelha \lstinline|mem_data_out|), endereços de porta
(\lstinline|addr_in|, \lstinline|addr_out|) e sinais de controle (\lstinline|req_in|,
\lstinline|out_en|). O módulo \lstinline|io_ctrl| só ativa esses sinais quando há
instruções de I/O presentes. A entrada é injetada no operando 2 da ULA via
\lstinline|ula_in2_ctrl|: quando \lstinline|req_in| está ativo, o dado da porta substitui o
acumulador.

\begin{lstlisting}[language=Verilog]
reg req_inr; always @ (posedge clk) req_inr <= req_in;
reg [NUBITS-1:0] ior; always @ (posedge clk) ior <= io_in;
assign out = (req_inr) ? ior : acc;
\end{lstlisting}

\subsection{Decodificação de múltiplas portas}
\label{sec:arq-io-decod}

No nível do \emph{wrapper} gerado (\file{hdl.c}), quando há mais de uma porta de entrada ou
saída, um decodificador 1-de-N (\lstinline|addr_dec|) converte o índice de porta no sinal de
\emph{enable} correspondente. Com uma única porta, o sinal é ligado diretamente.

\begin{lstlisting}[language=Verilog]
module addr_dec #( parameter NPORT = 4 )
( input valid_in, input [$clog2(NPORT)-1:0] index,
  output reg [NPORT-1:0] valid_out );
  integer i;
  always @* for (i = 0; i < NPORT; i = i+1)
    valid_out[i] = valid_in & (index == i[$clog2(NPORT)-1:0]);
endmodule
\end{lstlisting}

\subsection{FIFO genérica}
\label{sec:arq-io-fifo}

Para fluxos de dados, o repositório fornece uma FIFO genérica (\lstinline|myFIFO|,
\file{myFIFO.v}) parametrizada por largura da palavra (\lstinline|WORD|), profundidade
(\lstinline|LENGTH|) e limiar de ``quase vazia'' (\lstinline|ALMOST|). Ela expõe os sinais
usuais --- \lstinline|empty|, \lstinline|full|, \lstinline|almost_empty|, \lstinline|usedw| ---
e foi escrita para substituir o IP de FIFO de fabricantes de FPGA.

\section{Endereçamento indireto: arrays, ponteiros e FFT}
\label{sec:arq-indireto}

O endereçamento indireto é tratado por \lstinline|mem_ctrl| e \lstinline|rel_addr|. Há dois
modos:

\begin{description}[leftmargin=2.2em]
  \item[Indireto relativo] (\lstinline|LDI|/\lstinline|STI| e variantes \lstinline|ILI|/\lstinline|ISI|):
    o endereço efetivo é a base (operando da instrução) somada a um deslocamento. Para FFT,
    o deslocamento passa por \textbf{reversão de bits} (\lstinline|FFTSIZ| bits), implementada
    em \lstinline|rel_addr| quando \lstinline|USEFFT| está ligado.
  \item[Indireto \emph{base-less}] (\lstinline|LDA|/\lstinline|STA|): o endereço vem
    diretamente do acumulador (leitura) ou do topo da pilha de dados (escrita), sem a base do
    operando. Isso viabiliza ponteiros em tempo de execução e a passagem de arrays como
    parâmetros de função.
\end{description}

\begin{lstlisting}[language=Verilog]
generate
  if (LDA) begin : rd_addr assign mem_addr_rd = lda ? ula[MDATAW-1:0] : rel_rd; end
  else     begin : rd_addr assign mem_addr_rd = rel_rd;                          end
endgenerate
generate
  if (STA) begin : wr_addr assign mem_addr_wr = sta ? stk_ofst : rel_wr; end
  else     begin : wr_addr assign mem_addr_wr = rel_wr;                  end
endgenerate
\end{lstlisting}

A reversão de bits da FFT é feita invertendo a ordem dos \lstinline|FFTSIZ| bits menos
significativos do deslocamento:

\begin{lstlisting}[language=Verilog]
integer i;
always @ (*) for (i = 0; i < FFTSIZ; i = i+1) aux[i] = offset[FFTSIZ-1-i];
wire [MDATAW-1:0] add = (fft) ? {offset[MDATAW-1:FFTSIZ], aux} : offset;
assign out = (use_oft) ? add + addr : addr;
\end{lstlisting}

\section{Parametrização de hardware a partir do \cmm}
\label{sec:arq-param}

Os parâmetros de hardware declarados no programa \cmm (largura de palavra, mantissa, etc.)
chegam ao montador via \emph{logs} intermediários e são lidos em \file{eval.c} pela rotina
\lstinline|eval_init|, que os busca em \path{app_log.txt}:

\begin{lstlisting}[language=C]
eval_get("app_log.txt","nubits", aux); nubits = atoi(aux); // largura da palavra
eval_get("app_log.txt","nbmant", aux); nbmant = atoi(aux); // bits de mantissa
eval_get("app_log.txt","nbexpo", aux); nbexpo = atoi(aux); // bits de expoente
\end{lstlisting}

Outros parâmetros vêm de diretivas no próprio Assembly (\file{ASMComp.l}), processadas por
\lstinline|eval_opernd| em \file{eval.c}: \lstinline|#NDSTAC| (profundidade da pilha de
dados), \lstinline|#SDEPTH| (pilha de subrotina), \lstinline|#NUIOIN|/\lstinline|#NUIOOU|
(número de portas de entrada/saída), \lstinline|#NUGAIN| (constante de normalização) e
\lstinline|#FFTSIZ| (bits da FFT).

A largura do campo operando (\lstinline|nbopr|) é derivada do maior entre a memória de
instruções e a de dados:

\begin{lstlisting}[language=C]
nbopr = (n_ins > n_dat) ? ceil(log2(n_ins)) : ceil(log2(n_dat));
\end{lstlisting}

Por fim, o gerador de HDL (\file{hdl.c}, \lstinline|hdl_vv_file|) emite a instância do
\lstinline|processor| com todos esses valores e liga, uma a uma, somente as instruções
usadas:

\begin{lstlisting}[language=C]
fprintf(f_veri, "processor#(.NUBITS(%d),\n", nubits);
fprintf(f_veri,            ".NBMANT(%d),\n", nbmant);
fprintf(f_veri,            ".NBEXPO(%d),\n", nbexpo);
// ... MDATAS, MINSTS, SDEPTH, DDEPTH, NBIOIN, NBIOOU, FFTSIZ, NUGAIN ...
for (int i = 0; i < opc_cnt(); i++) fprintf(f_veri, ".%s(1),\n", opc_get(i));
\end{lstlisting}

\section{Imagens de memória (\path{.mif}) geradas}
\label{sec:arq-mif}

O montador escreve duas imagens de memória em texto binário, uma por linha, à medida que
varre o Assembly (\file{eval.c}): \path{<proc>_inst.mif} (programa) e \path{<proc>_data.mif}
(dados). Ambas são abertas em \lstinline|eval_init|:

\begin{lstlisting}[language=C]
snprintf(aux, sizeof(aux), "%s/Hardware/%s_data.mif", proc_dir, prname); f_data  = fopen(aux, "w");
snprintf(aux, sizeof(aux), "%s/Hardware/%s_inst.mif", proc_dir, prname); f_instr = fopen(aux, "w");
\end{lstlisting}

\subsection{Imagem de instruções}
\label{sec:arq-mif-instr}

Cada instrução é uma linha com \lstinline|NBITS_OPC| (7) bits de opcode concatenados a
\lstinline|nbopr| bits de operando, ambos convertidos por \lstinline|itob| (inteiro para
string binária). Para uma operação aritmética com variável:

\begin{lstlisting}[language=C]
fprintf(f_instr, "%s%s\n", itob(opc_idx, NBITS_OPC), itob(var_find(va), nbopr));
\end{lstlisting}

Esse formato (opcode nos bits altos, operando nos baixos) é exatamente o que o
\lstinline|prefetch| espera ao fatiar \lstinline|mem_instr[NBOPCO+NBOPER-1:NBOPER]| e
\lstinline|mem_instr[NBOPER-1:0]| (ver \autoref{sec:arq-prefetch}). Saltos gravam o endereço
do rótulo (\lstinline|lab_find|); entradas e saídas gravam o índice de porta.

\subsection{Imagem de dados}
\label{sec:arq-mif-dados}

Cada variável ocupa uma linha de \lstinline|nubits| bits. Variáveis inteiras gravam o valor
diretamente; variáveis float são inicializadas com o equivalente binário de \lstinline|0.0|
no formato SAPHO via \lstinline|f2mf|:

\begin{lstlisting}[language=C]
if (!is_const && type > 1)
    fprintf(f_data, "%s\n", itob(f2mf("0.0", NULL), nubits)); // float -> 0.0
else
    fprintf(f_data, "%s\n", itob(var_val(va), nubits));        // inteiro/constante
\end{lstlisting}

Arrays são preenchidos por \lstinline|arr_add| (de \file{array.c}), opcionalmente a partir de
um arquivo de inicialização (diretiva \lstinline|#arrays|).

\section{Geração do \emph{testbench} (\path{_tb.v})}
\label{sec:arq-tb}

Ao final da varredura, \lstinline|eval_finish| invoca \lstinline|hdl_vv_file| (o \path{.v} de
topo do processador) e \lstinline|hdl_tb_file| (o \emph{testbench}), ambos em \file{hdl.c}.
O \emph{testbench} gerado (\path{<proc>_tb.v}) é autônomo e contém:

\begin{description}[leftmargin=2.2em]
  \item[Relógio e \emph{reset}] um período derivado da frequência configurada
    (\lstinline|T = 1000.0/sim_clk()| ns), com pulso de \emph{reset} inicial e geração de
    clock por \lstinline|always #(T/2) clk = ~clk|.
  \item[Instância do processador] conectada conforme as portas existentes (entrada, saída,
    interrupção, \lstinline|cheguei|).
  \item[Estímulo de entrada] cada porta de entrada lê de um arquivo
    \path{Simulation/input_<i>.txt} via \lstinline|$fscanf| na borda de descida do clock.
  \item[Captura de saída] cada porta de saída grava em \path{Simulation/output_<i>.txt} via
    \lstinline|$fdisplay| + \lstinline|$fflush| quando seu \emph{enable} está ativo.
  \item[Dump de forma de onda] \lstinline|$dumpfile|/\lstinline|$dumpvars| produzem o
    \path{.vcd}, incluindo as variáveis e arrays do usuário (espelhados no \path{.v} de topo)
    e os \emph{flags} de pilha e os erros de arredondamento da ULA.
  \item[Término] um bloco de progresso imprime o avanço, e o \emph{testbench} encerra com
    \lstinline|$finish| ao detectar o marcador de fim de programa (\lstinline|proc.valr10|).
\end{description}

\begin{lstlisting}[language=C]
double T = 1000.0/sim_clk(); // periodo em ns (frequencia em MHz)
fprintf(f_veri, "    clk = 0;\n    rst = 1;\n    #%f;\n    rst = 0;\n", T);
fprintf(f_veri, "always #%f clk = ~clk;\n\n", T/2.0);
// ...
fprintf(f_veri, "always @ (posedge clk) if (proc.valr10 == %d) begin\n", sim_get_fim());
fprintf(f_veri, "    $display(\"Info: end of program!\");\n    $finish;\nend\n");
\end{lstlisting}

\subsection{Visibilidade de simulação versus síntese}
\label{sec:arq-tb-vis}

Sinais de instrumentação (espelhos de variáveis, \emph{taps} do PC e da memória, erros de
arredondamento em \lstinline|real|) ficam atrás da macro \lstinline|YANC_SIM_VIS|, definida
automaticamente para Icarus Verilog (que predefine \lstinline|__ICARUS__|) e para Verilator
quando se passa \lstinline|+define+YANC_TRACE|. Nunca são incluídos na síntese. No \path{.v}
de topo, comentários \lstinline|/* verilator tracing_off */| e \lstinline|/* verilator
tracing_on */| delimitam quais sinais entram na forma de onda sob Verilator (escolhido por
velocidade), mantendo os \emph{taps} caros (como o erro de arredondamento real da ULA)
disponíveis apenas sob Icarus.

\begin{nota}
A interação entre as imagens \path{.mif}, o \emph{testbench} e as ferramentas de simulação
(Icarus, Verilator, GTKWave) e síntese (Yosys), bem como o empacotamento dessas ferramentas
no \emph{bundle} portátil, são detalhados em \autoref{cap:03-pipeline-de-compilacao-yanc} e
nos capítulos sobre o \emph{toolchain}.
\end{nota}

% !TEX root = ../main.tex
\chapter{A linguagem C± (CMM) e o caminho C++}
\label{cap:03-linguagem-cmm}

Este capítulo descreve a linguagem \cmm{} (CMM) tal como ela é definida pela
gramática real do compilador \tool{cmmcomp}, parte do projeto \repo{yanc}. A
descrição parte das duas fontes canônicas da sintaxe: o scanner Flex
\file{Compilers/CMMComp/Sources/CMMComp.l} e o parser Bison
\file{Compilers/CMMComp/Sources/CMMComp.y}. As ações semânticas residem nos
módulos C de \file{Compilers/CMMComp/Sources/} (entre eles
\file{diretivas.c}, \file{data\_declar.c}, \file{stdlib.c}, \file{oper.c},
\file{saltos.c}, \file{funcoes.c}, \file{array\_index.c} e \file{messages.c}).
O catálogo de mensagens vem de \file{Compilers/CMMComp/Headers/messages.h}, e os
exemplos vêm de \file{Compilers/CMMComp/Tests/}. A última seção documenta o
caminho C++ (\tool{cpppp} e \tool{cppcomp}), cuja gramática vive em
\file{Compilers/CPPComp/Sources/}.

\begin{nota}
A linguagem \cmm{} é um front-end de software para o soft-processor SAPHO. Ela
produz \emph{assembly} simbólico (\path{.asm}), que é posteriormente convertido
em imagem de memória e em Verilog pelos estágios seguintes do toolchain (ver
\autoref{cap:02-yanc-pipeline}). O foco deste capítulo é a \emph{linguagem-fonte}:
sintaxe, tipos, operadores, biblioteca padrão e diretivas.
\end{nota}

\section{Estrutura de um programa \cmm{}}
\label{sec:cmm-estrutura}

Pela regra de partida do parser, um programa \cmm{} é uma sequência de
elementos de topo de nível, sem ordem rígida entre eles:

\begin{lstlisting}[language=C]
fim           : prog {prog_emit();}
prog          : prog_elements | prog prog_elements
prog_elements : direct | declar | funcao
\end{lstlisting}

Ou seja, cada elemento de topo de nível é uma das três coisas:

\begin{description}[leftmargin=2.2em]
  \item[\code{direct}] uma \emph{diretiva} de hardware (linhas iniciadas por
  \texttt{\#}), que configura parâmetros do processador-alvo.
  \item[\code{declar}] uma \emph{declaração} de variável ou de \emph{array}
  global.
  \item[\code{funcao}] a definição de uma \emph{função}.
\end{description}

Não existe a noção sintática de um bloco \enquote{\code{processor \{ ... \}}};
o que individualiza um processador é o conjunto de diretivas no início do arquivo
(em especial \lstinline|#PRNAME|, que dá o nome do processador). O ponto de
entrada da execução é a função \lstinline|main()|: a ausência dela é um erro
detectado ao final da compilação, com a mensagem \lstinline|MSG_ERR_NO_MAIN|.

Um arquivo \cmm{} mínimo e completo, retirado de
\file{Tests/cmm\_define/Software/cmm\_define.cmm}, ilustra a forma canônica:
diretivas no topo, depois a função \lstinline|main()|, normalmente com um laço
infinito \lstinline|while (1)| que modela o comportamento contínuo do hardware:

\begin{lstlisting}[style=cmm]
#PRNAME cmm_define
#NUBITS 16
#NDSTAC 4
#SDEPTH 4
#NUIOIN 1
#NUIOOU 1
#NBMANT 10
#NBEXPO 5
#NUGAIN 128

#define SCALE  3
#define OFFSET 7

void main()
{
    int x;
    while (1)
    {
        x = in(0);
        x = x + SCALE;
        out(0, x);
        x = x + OFFSET;
        out(0, x);
    }
}
\end{lstlisting}

\section{Diretivas de hardware}
\label{sec:cmm-diretivas}

As diretivas são reconhecidas pelo scanner como \emph{tokens} próprios (regras
\lstinline|"#PRNAME"|, \lstinline|"#NUBITS"|, etc. em \file{CMMComp.l}) e
aparecem na gramática nas regras \code{direct} (configuração) e \code{dire\_inter}
(comportamentais). Cada diretiva ocupa uma linha e recebe um argumento: um nome
(em \lstinline|#PRNAME|) ou um inteiro (nas demais).

\subsection{Diretivas de configuração}

A regra \code{direct} mapeia cada diretiva para uma chamada de
\lstinline|dire_exec()| (em \file{diretivas.c}), que atualiza o estado de
compilação e anexa um \code{STMT\_DIRECTIVE} ao programa. A tabela a seguir lista
as diretivas exatamente como aparecem no parser, com os comentários do próprio
código.

\begin{longtable}{@{}p{0.16\linewidth}p{0.12\linewidth}p{0.62\linewidth}@{}}
\toprule
\textbf{Diretiva} & \textbf{Argumento} & \textbf{Significado (do parser)} \\
\midrule
\endhead
\lstinline|#PRNAME| & nome & nome do processador \\
\lstinline|#NUBITS| & inteiro & largura da palavra da ULA \\
\lstinline|#NBMANT| & inteiro & largura da mantissa (bits) \\
\lstinline|#NBEXPO| & inteiro & largura do expoente (bits) \\
\lstinline|#NDSTAC| & inteiro & profundidade da pilha de dados \\
\lstinline|#SDEPTH| & inteiro & profundidade da pilha de sub-rotinas \\
\lstinline|#NUIOIN| & inteiro & número de portas de entrada \\
\lstinline|#NUIOOU| & inteiro & número de portas de saída \\
\lstinline|#NUGAIN| & inteiro & constante de divisão usada por \lstinline|norm(.)| \\
\lstinline|#FFTSIZ| & inteiro & tamanho da FFT (na forma $2^{\texttt{FFTSIZ}}$) \\
\bottomrule
\end{longtable}

A função \lstinline|dire_exec()| guarda em variáveis globais do compilador apenas
o que o front-end \cmm{} precisa consultar em tempo de compilação: o nome
(\lstinline|prname|), a largura da mantissa (\lstinline|nbmant|), a do expoente
(\lstinline|nbexpo|), e a contagem de portas de entrada e de saída
(\lstinline|nuioin|, \lstinline|nuioou|). As demais diretivas são repassadas
adiante no pipeline via o \code{STMT\_DIRECTIVE} emitido. Os valores-padrão,
definidos em \file{diretivas.c}, são \lstinline|nbmant = 16|,
\lstinline|nbexpo = 6|, \lstinline|nuioin = 1| e \lstinline|nuioou = 1|.

\begin{nota}
As contagens \lstinline|#NUIOIN| e \lstinline|#NUIOOU| são checadas em tempo de
compilação contra os índices de porta usados em \lstinline|in(.)|/\lstinline|out(.,.)|.
Um índice fora da faixa produz \lstinline|MSG_ERR_NO_IN_PORT| ou
\lstinline|MSG_ERR_NO_OUT_PORT|.
\end{nota}

\subsection{Diretivas comportamentais}

Duas diretivas adicionais marcam pontos no fluxo de execução, em vez de
configurar parâmetros. Elas são reconhecidas pelo scanner com nomes que diferem
do \emph{token} interno (ver \file{CMMComp.l}):

\begin{itemize}[leftmargin=1.5em]
  \item \lstinline|#PRACA| (token \code{ITRADD}) --- marca o ponto de retorno
  para o \emph{reset} acionado pelo pino de interrupção (\code{itr}). É tratada
  pela regra \code{dire\_inter}, que emite \lstinline|stmt_dire_inter()|.
  \item \lstinline|#TOAQUI| (token \code{TOAQUI}) --- marca um endereço que o
  pino \code{cheguei} reflete (verdadeiro quando \texttt{PC == addr}). Emite
  \lstinline|stmt_dire_toaqui()|.
\end{itemize}

Diferentemente das diretivas de configuração, estas são \emph{statements}: podem
aparecer dentro do corpo de funções (a regra \code{dire\_inter} é uma das
alternativas de \code{stmt\_full}).

\section{Tipos de dados}
\label{sec:cmm-tipos}

O scanner reconhece exatamente quatro palavras-chave de tipo, cada uma associada
a um código numérico interno (ver \file{CMMComp.l}):

\begin{lstlisting}[language=C]
"comp"  {type_tmp = 3; yylval.ival = 3; return TYPE;}
"float" {type_tmp = 2; yylval.ival = 2; return TYPE;}
"int"   {type_tmp = 1; yylval.ival = 1; return TYPE;}
"void"  {              yylval.ival = 0; return TYPE;}
\end{lstlisting}

Os códigos de tipo são usados em toda a tabela de símbolos e nas decisões de
codegen. A tabela abaixo consolida os códigos, confirmados em
\file{data\_declar.c} (campo \lstinline|v_table[id].type|):

\begin{longtable}{@{}p{0.10\linewidth}p{0.20\linewidth}p{0.60\linewidth}@{}}
\toprule
\textbf{Código} & \textbf{Tipo} & \textbf{Descrição} \\
\midrule
\endhead
0 & \lstinline|void| & ausência de tipo (também: variável ainda não declarada) \\
1 & \lstinline|int|  & inteiro de \lstinline|#NUBITS| bits, em complemento de dois \\
2 & \lstinline|float| & ponto flutuante (mantissa \lstinline|#NBMANT|, expoente \lstinline|#NBEXPO|) \\
3 & \lstinline|comp| & número complexo (par de valores em ponto flutuante) \\
4 & --- & parte imaginária de um \lstinline|comp| (gerada automaticamente) \\
5 & --- & literal complexo (\code{CNUM}), usado nas reduções de expressão \\
\bottomrule
\end{longtable}

\subsection{Inteiro (\texttt{int})}

O tipo \lstinline|int| é a palavra inteira da ULA, com largura dada por
\lstinline|#NUBITS| e aritmética em complemento de dois. É o tipo dos índices de
\emph{array}, dos resultados lógicos e dos operandos de operações bit-a-bit e de
deslocamento. Literais inteiros são reconhecidos pela regra \code{INNUM}
(\lstinline|[0-9]+|) e produzem o \emph{token} \code{INUM}.

\subsection{Ponto flutuante (\texttt{float})}

O tipo \lstinline|float| representa números fracionários no formato de ponto
flutuante do SAPHO, parametrizado por \lstinline|#NBMANT| (mantissa) e
\lstinline|#NBEXPO| (expoente). Não existe na linguagem um tipo separado
\enquote{ponto fixo}: o único tipo fracionário é \lstinline|float|, e a precisão
efetiva é configurada pelas diretivas. Literais com ponto decimal são
reconhecidos pela regra \code{FLNUM} (\lstinline|(0|[1-9]+[0-9]*)\.?[0-9]*|) e
produzem \code{FNUM}.

As funções transcendentais (raiz, trigonometria, exponencial, logaritmo) e os
arredondamentos sempre operam sobre, ou produzem, valores \lstinline|float|.

\subsection{Número complexo (\texttt{comp})}

O tipo \lstinline|comp| é a marca distintiva da \cmm{}. Cada variável
\lstinline|comp| ocupa, na tabela de símbolos, dois registros: o da parte real
(código de tipo 3) e o da parte imaginária (código 4), criado automaticamente em
\lstinline|declar_var()| por \lstinline|get_img_id()| sempre que
\lstinline|type_tmp > 2|. Literais complexos são reconhecidos pela regra
\code{CONUM} do scanner:

\begin{lstlisting}[language=C]
CONUM  [-]*{WS}*{FLNUM}{WS}*[+|-]{WS}*{FLNUM}{WS}*[i]
\end{lstlisting}

ou seja, a forma \lstinline|a+bi| (com \lstinline|i| ao final), produzindo o
\emph{token} \code{CNUM}. O símbolo \lstinline|i| é \emph{reservado} para indicar
a parte imaginária; usá-lo como identificador gera \lstinline|MSG_ERR_RESERVED_I|.

Exemplo de uso de literais complexos (de \file{Tests/cmm\_comp\_arith/Software/cmm\_comp\_arith.cmm}):

\begin{lstlisting}[style=cmm]
void main() {
    comp x; comp y; comp r;
    while (1) {
        r = (1.0+2.0i) * (3.0+4.0i);   fout(0, real(r)); fout(0, imag(r));
        r = (4.0+2.0i) / (1.0+1.0i);   fout(0, real(r)); fout(0, imag(r));
        r = (1.0+2.0i) + (3.0+4.0i);   fout(0, real(r)); fout(0, imag(r));
        x = 1.0+2.0i; y = 3.0+4.0i; r = x * y;  fout(0, real(r)); fout(0, imag(r));
    }
}
\end{lstlisting}

A aritmética \lstinline|comp| dá suporte a \lstinline|+|, \lstinline|-|,
\lstinline|*| e \lstinline|/| entre operandos complexos, com o tratamento das
duas componentes feito pela camada de operadores (\file{oper.c}).

\section{Declaração de variáveis e \emph{arrays}}
\label{sec:cmm-declaracao}

A regra \code{declar} cobre todas as formas de declaração. As formas e suas ações
semânticas, lidas diretamente do parser:

\begin{longtable}{@{}p{0.52\linewidth}p{0.40\linewidth}@{}}
\toprule
\textbf{Forma sintática} & \textbf{Ação} \\
\midrule
\endhead
\lstinline|TYPE id_list ;| & lista de variáveis sem inicialização \\
\lstinline|TYPE ID = exp ;| & declaração com inicialização \\
\lstinline|TYPE ID [INUM] STRING ;| & \emph{array} 1D inicializado por arquivo \\
\lstinline|TYPE ID [INUM][INUM] STRING ;| & \emph{array} 2D inicializado por arquivo \\
\lstinline+TYPE ID [INUM] # | ID | ID >;+ & \emph{array} 1D inicializado por produto matriz-vetor \\
\lstinline+TYPE ID [INUM] # exp | ID >;+ & \emph{array} 1D inicializado por produto constante-vetor \\
\bottomrule
\end{longtable}

A \code{id\_list} permite declarar várias variáveis e/ou \emph{arrays} de uma vez:

\begin{lstlisting}[language=C]
id_list : IID | id_list ',' IID
IID    : ID                           {declar_var   ($1         );}
       | ID '[' INUM ']'              {declar_arr_1d($1,$3   ,-1);}
       | ID '[' INUM ']' '[' INUM ']' {declar_arr_2d($1,$3,$6,-1);}
\end{lstlisting}

\subsection{Variáveis escalares}

Uma variável escalar é declarada por \lstinline|declar_var()|, que rejeita
redeclaração com \lstinline|MSG_ERR_VAR_EXISTS| (\enquote{a variável já existe})
e, no caso de \lstinline|comp|, registra automaticamente a parte imaginária. Há
ainda a forma com inicialização imediata, \lstinline|TYPE ID = exp ;|, que
declara e em seguida atribui (\lstinline|stmt_assign|).

\subsection{\emph{Arrays} 1D e 2D}

\emph{Arrays} podem ser uni- ou bidimensionais. O tamanho é um literal inteiro
(\code{INUM}). Há três formas de inicialização:

\begin{description}[leftmargin=1.6em]
  \item[Sem inicialização] na \code{id\_list}, com índice \texttt{-1} indicando
  ausência de arquivo.
  \item[Por arquivo] um literal de string após as dimensões dá o nome de um
  arquivo de texto cujos valores preenchem a memória do \emph{array} em tempo de
  compilação. O parser emite \lstinline|MSG_INFO_ARRAY_FILE_INIT| ao reconhecer
  essa forma. Exemplo de uso: \lstinline|int x[128] "values.txt";|.
  \item[Por notação de Dirac] formas \lstinline+# |M|v>+ e
  \lstinline+# c|v>+ inicializam o \emph{array} como produto matriz-vetor ou
  constante-vetor (ver \autoref{sec:cmm-dirac}).
\end{description}

\subsection{\emph{Array} com índice invertido (\emph{bit-reversal})}

A \cmm{} oferece uma sintaxe especial de indexação para a FFT: a forma
\lstinline|x[i)| (colchete de abertura, parêntese de fechamento) acessa o
elemento cujo índice tem os \emph{bits} de \texttt{i} invertidos. Isso aparece na
gramática tanto em expressões quanto em atribuições:

\begin{lstlisting}[language=C]
| ID '[' exp ')'  '='     exp ';'  {stmt_append(stmt_array_assign($1, $3, NULL, $6, 1));} // reversed
| ID '[' exp ')'                   {$$ = expr_array_index(v_table[$1].type, $1, 1, $3, NULL);}
\end{lstlisting}

O quinto argumento de \lstinline|stmt_array_assign| (\texttt{1}) distingue o
acesso invertido do normal. O uso típico aparece em \file{Tests/proc\_fft}:

\begin{lstlisting}[style=cmm]
temp    = wpv[sind]*data[j);
data[j) =   data[k) - temp;
data[k) =   data[k) + temp;
\end{lstlisting}

\section{Operadores}
\label{sec:cmm-operadores}

A precedência e a associatividade dos operadores são declaradas no parser e
seguem a convenção de C (\enquote{mais abaixo na lista $\Rightarrow$ maior
prioridade}):

\begin{lstlisting}[language=C]
%left LOR
%left LAN
%left '|'
%left '^'
%left '&'
%left EQU DIF
%left '>' '<' GREQU LESEQ
%left SHIFTL SHIFTR SSHIFTR
%left '+' '-'
%left '*' '/' '%'
%left '!' '~' PPLUS
\end{lstlisting}

\subsection{Aritméticos, bit-a-bit e de deslocamento}

A tabela a seguir lista cada operador binário com a operação interna que a regra
de expressão produz (\lstinline|expr_binop(OP_...)| no parser):

\begin{longtable}{@{}p{0.16\linewidth}p{0.18\linewidth}p{0.56\linewidth}@{}}
\toprule
\textbf{Símbolo} & \textbf{Operação} & \textbf{Descrição} \\
\midrule
\endhead
\lstinline|+| & \code{OP\_ADD} & soma \\
\lstinline|-| & \code{OP\_SUB} & subtração \\
\lstinline|*| & \code{OP\_MUL} & multiplicação \\
\lstinline|/| & \code{OP\_DIV} & divisão \\
\lstinline|%| & \code{OP\_MOD} & resto (somente inteiros) \\
\lstinline|&| & \code{OP\_AND} & E bit-a-bit \\
\lstinline+|+ & \code{OP\_OR} & OU bit-a-bit \\
\lstinline|^| & \code{OP\_XOR} & OU-exclusivo bit-a-bit \\
\lstinline|<<| & \code{OP\_SHL} & deslocamento à esquerda \\
\lstinline|>>| & \code{OP\_SHR} & deslocamento à direita \\
\lstinline|>>>| & \code{OP\_SSHR} & deslocamento à direita com sinal (complemento de dois) \\
\bottomrule
\end{longtable}

O operador \lstinline|>>>| (\enquote{\emph{arithmetic right shift}}) preserva o
bit de sinal --- é uma extensão da \cmm{} em relação a C. As operações bit-a-bit
e de deslocamento são restritas a operandos inteiros: aplicá-las a um
\lstinline|comp| produz \lstinline|MSG_ERR_BITWISE_COMPLEX| ou
\lstinline|MSG_ERR_SHIFT_COMPLEX|; o resto (\lstinline|%|) sobre não-inteiro
produz \lstinline|MSG_ERR_MOD_NON_INT|.

\subsection{Relacionais e lógicos}

Os operadores relacionais e lógicos produzem resultados inteiros (0/1):

\begin{tabularx}{\linewidth}{Y Y Y}
\toprule
\textbf{Símbolo} & \textbf{Operação} & \textbf{Descrição} \\
\midrule
\lstinline|<|   & \code{OP\_LT} & menor que \\
\lstinline|>|   & \code{OP\_GT} & maior que \\
\lstinline|<=|  & \code{OP\_LE} & menor ou igual \\
\lstinline|>=|  & \code{OP\_GE} & maior ou igual \\
\lstinline|==|  & \code{OP\_EQ} & igual \\
\lstinline|!=|  & \code{OP\_NE} & diferente \\
\lstinline|&&|  & \code{OP\_LAN} & E lógico \\
\lstinline+||+ & \code{OP\_LOR} & OU lógico \\
\bottomrule
\end{tabularx}

Comparar um valor complexo dispara avisos de que o compilador usará o módulo
(\lstinline|MSG_WARN_CMP_COMPLEX| e variantes); uma operação lógica só entre
inteiros é o caso esperado, e um operando \lstinline|float| ou \lstinline|comp|
gera aviso (\lstinline|MSG_WARN_LOGIC_FLOAT|, \lstinline|MSG_WARN_LOGIC_COMP|).

\subsection{Unários e incremento}

Os operadores unários são reconhecidos por regras dedicadas de expressão:

\begin{itemize}[leftmargin=1.5em]
  \item \lstinline|-exp| --- negação aritmética (\code{OP\_NEG});
  \item \lstinline|!exp| --- negação lógica (\code{OP\_LIN});
  \item \lstinline|~exp| --- inversão bit-a-bit (\code{OP\_INV}, só inteiros ---
  caso contrário \lstinline|MSG_ERR_INV_NON_INT|);
  \item \lstinline|+exp| --- operador nulo (retorna o próprio operando).
\end{itemize}

O incremento \lstinline|++| (token \code{PPLUS}) existe em duas posições: como
\emph{statement} de atribuição (\lstinline|x++;|, \lstinline|x[i]++;|,
\lstinline|x[i][j]++;|) e como expressão (\lstinline|expr_pplus|). Incrementar um
número complexo é rejeitado com \lstinline|MSG_ERR_INCR_COMPLEX|.

\section{Atribuições}
\label{sec:cmm-atribuicoes}

A regra \code{assignment} cobre a atribuição escalar, o incremento e os acessos a
\emph{array} (incluindo a forma de índice invertido e a 2D):

\begin{lstlisting}[language=C]
assignment : ID '=' exp ';'                          {stmt_assign}
           | ID PPLUS ';'                            {stmt_pplus}
           | ID '[' exp ']' PPLUS ';'                {stmt_pplus (1D)}
           | ID '[' exp ']' '[' exp ']' PPLUS ';'    {stmt_pplus (2D)}
           | ID '[' exp ']' '=' exp ';'              {stmt_array_assign (normal)}
           | ID '[' exp ')' '=' exp ';'              {stmt_array_assign (reversed)}
           | ID '[' exp ']' '[' exp ']' '=' exp ';'  {stmt_array_assign (2D)}
\end{lstlisting}

Atribuições entre tipos diferentes não são bloqueadas, mas emitem avisos de
conversão, por exemplo \lstinline|MSG_WARN_INT_RECV_FLOAT| (\enquote{variável é
int, mas recebe float}) e suas variantes para todas as combinações
int/float/comp.

A regra \code{assignment} ainda contém todo o bloco de atribuições em notação de
Dirac, descrito em \autoref{sec:cmm-dirac}.

\section{Controle de fluxo}
\label{sec:cmm-fluxo}

A \cmm{} oferece os comandos de controle de fluxo de C. Cada um corresponde a
uma alternativa de \code{stmt\_full} no parser.

\subsection{\texttt{if} / \texttt{else}}

A regra resolve a ambiguidade clássica do \enquote{\emph{dangling else}} com as
declarações \lstinline|%nonassoc THEN| e \lstinline|%nonassoc ELSE|:

\begin{lstlisting}[language=C]
if_else_stmt : if_exp stmt_full else_intro stmt_full {stmt_append(if_fim());}  // com else
             | if_exp stmt_full %prec THEN           {stmt_append(if_stmt());} // sem else
if_exp       : IF '(' exp ')'                        {if_exp($3);}
else_intro   : ELSE                                  {else_stmt();}
\end{lstlisting}

\subsection{\texttt{while}, \texttt{do/while} e \texttt{for}}

O \lstinline|while| e o \lstinline|do/while| são tratados em \file{saltos.c}. O
\lstinline|for| é \emph{açúcar sintático}: o parser o reescreve em tempo de
\emph{parsing} para um \lstinline|while| equivalente, conforme o comentário da
gramática:

\begin{lstlisting}[language=C]
// for (init; cond; step) body  desugars at parse-time into:
//     init;
//     while (cond) { body; step; }
\end{lstlisting}

A cláusula de inicialização do \lstinline|for| aceita a forma declarativa
(\lstinline|TYPE ID = exp|) ou a de atribuição (\lstinline|ID = exp|); a de passo
aceita \lstinline|ID = exp| ou \lstinline|ID++|.

\subsection{\texttt{break} e \texttt{continue}}

\lstinline|break;| e \lstinline|continue;| são \emph{statements} próprios
(\lstinline|exec_break()| / \lstinline|exec_continue()|). Pelo comentário do
parser, o \lstinline|break| liga-se ao laço \emph{ou} ao \lstinline|switch| mais
interno por meio de uma pilha de alvos em \file{saltos.c}; um \lstinline|break|
solto gera \lstinline|MSG_ERR_BREAK_LOST| e um \lstinline|continue| fora de laço
gera \lstinline|MSG_ERR_CONTINUE_LOST|.

\subsection{\texttt{switch} / \texttt{case} / \texttt{default}}

O corpo do \lstinline|switch| é uma sequência plana de rótulos
(\lstinline|case|, \lstinline|default|) e \emph{statements}, no modelo de C ---
um rótulo é só um marcador, não um contêiner. Isso permite \emph{fall-through}
real, blocos \lstinline|{ }| dentro de um \lstinline|case| e
\lstinline|switch| aninhados. Os rótulos de \lstinline|case| aceitam tanto
constantes inteiras quanto \emph{float}:

\begin{lstlisting}[language=C]
case_label    : CASE INUM ':' {case_test($2,1);}
              | CASE FNUM ':' {case_test($2,2);}
default_label : DEFAULT   ':' {defaut_test();}
\end{lstlisting}

O exemplo \file{Tests/ctrl\_flow/Software/ctrl\_flow.cmm} exercita
sistematicamente todas as combinações de \lstinline|if|/\lstinline|else|/%
\lstinline|while|/\lstinline|break|; um trecho:

\begin{lstlisting}[style=cmm]
void main()
{
    while (1)
    {
        if (1)        { out(0, 11); }   // if sem else, tomado
        if (0)        { out(0, 12); }   // if sem else, nao tomado
        if (1)        { out(0, 21); }
        else          { out(0, 22); }

        int bc = 5;
        while (bc)
        {
            if (bc == 3) { break; }     // break liga-se a este while
            out(0, 100 + bc);
            bc = bc - 1;
        }
        out(0, 0);
    }
}
\end{lstlisting}

\section{Funções}
\label{sec:cmm-funcoes}

A definição de função segue a regra \code{funcao}:

\begin{lstlisting}[language=C]
funcao     : func_intro par_list ')' '{' stmt_list '}' {func_ret($1);}
func_intro : TYPE ID '('                  {declar_fun($1,$2); $$ = $2;}
par_items  : TYPE ID                       {declar_par($1,$2);}
           | par_items ',' TYPE ID         {declar_par($3,$4);}
\end{lstlisting}

Características confirmadas no parser e nos módulos de apoio:

\begin{itemize}[leftmargin=1.5em]
  \item O tipo de retorno é um dos quatro \code{TYPE}; \lstinline|void| indica
  função que não retorna valor.
  \item A lista de parâmetros pode ser vazia. Os parâmetros são escalares: pelo
  comentário do parser, \emph{arrays} não são permitidos como parâmetros.
  \item O retorno é feito por \lstinline|return exp;| (função com valor) ou
  \lstinline|return;| (void). Retornar valor de função \lstinline|void| produz
  \lstinline|MSG_ERR_VOID_RETURN_VALUE|; usar como valor uma função que não
  retorna nada produz \lstinline|MSG_ERR_VOID_FUNC_USE|.
  \item Conversões de tipo nos parâmetros e no retorno geram avisos dedicados
  (família \lstinline|MSG_WARN_CONV_*_PARAM| e \lstinline|MSG_WARN_RET_*|).
  \item A contagem de parâmetros na chamada é verificada
  (\lstinline|MSG_ERR_PARAM_COUNT|); uma função inexistente gera
  \lstinline|MSG_ERR_FUNC_WHERE|.
\end{itemize}

Chamadas aparecem em dois contextos: como \emph{statement}
(\lstinline|ID(exp_list);|, regra \code{void\_call}) e como expressão
(\lstinline|ID(exp_list)|, regra \code{func\_call}).

\section{Entrada e saída}
\label{sec:cmm-io}

A entrada e a saída usam portas indexadas por inteiro, com variantes inteiras e
de ponto flutuante. As regras de I/O da gramática:

\begin{longtable}{@{}p{0.34\linewidth}p{0.58\linewidth}@{}}
\toprule
\textbf{Forma} & \textbf{Descrição} \\
\midrule
\endhead
\lstinline|in(INUM)| & lê dado externo da porta (expressão; \code{OP\_STD\_IN}) \\
\lstinline|fin(INUM)| & lê dado externo convertendo para \emph{float} (\code{OP\_STD\_FIN}) \\
\lstinline|out(INUM, exp);| & escreve para fora do processador (\code{stmt\_out}) \\
\lstinline|fout(INUM, exp);| & escreve convertendo para \emph{float} (\code{stmt\_out}) \\
\bottomrule
\end{longtable}

As portas em \lstinline|in|/\lstinline|fin| e \lstinline|out|/\lstinline|fout|
são literais inteiros (\code{INUM}), validados contra \lstinline|#NUIOIN| e
\lstinline|#NUIOOU|. Há ainda a saída em notação de Dirac (\code{std\_vout}),
\lstinline+out(p, c|v>);+, que exporta um vetor ponderado por uma constante
(ver \autoref{sec:cmm-dirac}). Avisos \lstinline|MSG_WARN_USE_FOUT| e
\lstinline|MSG_WARN_USE_OUT| orientam a escolha entre as variantes inteira e
\emph{float}.

\section{Biblioteca padrão (\emph{stdlib})}
\label{sec:cmm-stdlib}

A biblioteca padrão da \cmm{} é composta por funções intrínsecas reconhecidas
como \emph{tokens} pelo scanner e mapeadas, no parser, para nós de expressão
\lstinline|expr_stdlib(OP_STD_...)| (ou para \emph{statements}, no caso das que
não retornam valor). A tabela cataloga todas elas com a descrição do próprio
parser.

\subsection{Funções especiais}

\begin{longtable}{@{}p{0.18\linewidth}p{0.16\linewidth}p{0.56\linewidth}@{}}
\toprule
\textbf{Função} & \textbf{Operação} & \textbf{Descrição} \\
\midrule
\endhead
\lstinline|norm(x)| & \code{OP\_STD\_NRM} & divide o argumento por \lstinline|#NUGAIN| (evita o divisor da ULA); só inteiro \\
\lstinline|pset(x)| & \code{OP\_STD\_PST} & retorna zero se o argumento for negativo \\
\lstinline|abs(x)| & \code{OP\_STD\_ABS} & valor absoluto; para \lstinline|comp|, o módulo \\
\lstinline|sign(x,y)| & \code{OP\_STD\_SIGN} & retorna \lstinline|y| com o sinal de \lstinline|x| \\
\lstinline|copy(x,y)| & \lstinline|stmt_copy| & copia \lstinline|x| em \lstinline|y| (sem checagem de tipo) \\
\bottomrule
\end{longtable}

A \lstinline|norm(.)| sobre não-inteiro produz \lstinline|MSG_ERR_NORM_NON_INT|;
\lstinline|sign(.,.)| e \lstinline|pset(.)| com complexos produzem
\lstinline|MSG_ERR_SIGN_COMPLEX| e \lstinline|MSG_ERR_PSET_COMPLEX|.

\subsection{Funções não-lineares e de arredondamento}

\begin{longtable}{@{}p{0.18\linewidth}p{0.18\linewidth}p{0.54\linewidth}@{}}
\toprule
\textbf{Função} & \textbf{Operação} & \textbf{Descrição} \\
\midrule
\endhead
\lstinline|sqrt(x)| & \code{OP\_STD\_SQRT} & raiz quadrada (\emph{float}; raiz principal se complexo) \\
\lstinline|atan(x)| & \code{OP\_STD\_ATAN} & arco-tangente \\
\lstinline|sin(x)| & \code{OP\_STD\_SIN} & seno \\
\lstinline|cos(x)| & \code{OP\_STD\_COS} & cosseno \\
\lstinline|tan(x)| & \code{OP\_STD\_TAN} & tangente (minimax dedicado, não sin/cos) \\
\lstinline|cosh(x)| & \code{OP\_STD\_COSH} & cosseno hiperbólico \\
\lstinline|sinh(x)| & \code{OP\_STD\_SINH} & seno hiperbólico \\
\lstinline|tanh(x)| & \code{OP\_STD\_TANH} & tangente hiperbólica \\
\lstinline|exp(x)| & \code{OP\_STD\_EXP} & exponencial $e^x$ \\
\lstinline|log(x)| & \code{OP\_STD\_LOG} & logaritmo natural (ln) \\
\lstinline|pow(x,y)| & \code{OP\_STD\_POW} & $x^y$ \\
\lstinline|floor(x)| & \code{OP\_STD\_FLOOR} & maior \emph{float} inteiro $\le x$ \\
\lstinline|ceil(x)| & \code{OP\_STD\_CEIL} & menor \emph{float} inteiro $\ge x$ \\
\lstinline|round(x)| & \code{OP\_STD\_ROUND} & \emph{float} inteiro mais próximo (\emph{ties away from zero}) \\
\bottomrule
\end{longtable}

Pelo comentário do parser, \lstinline|pow(x,y)| escolhe a estratégia conforme o
expoente: constante inteira $\to$ \emph{square-and-multiply}; variável inteira
$\to$ laço em tempo de execução; caso geral $\to$ $\exp(y \cdot \ln x)$.
Diversas dessas funções ainda não têm versão complexa, e o uso com
\lstinline|comp| é rejeitado (\lstinline|MSG_ERR_SQRT_COMPLEX|,
\lstinline|MSG_ERR_EXP_COMPLEX|, \lstinline|MSG_ERR_LOG_COMPLEX|,
\lstinline|MSG_ERR_POW_COMPLEX|).

\subsection{Funções de números complexos}

\begin{longtable}{@{}p{0.20\linewidth}p{0.18\linewidth}p{0.52\linewidth}@{}}
\toprule
\textbf{Função} & \textbf{Operação} & \textbf{Descrição} \\
\midrule
\endhead
\lstinline|real(z)| & \code{OP\_STD\_REAL} & parte real de um complexo \\
\lstinline|imag(z)| & \code{OP\_STD\_IMAG} & parte imaginária de um complexo \\
\lstinline|fase(z)| & \code{OP\_STD\_FASE} & fase (argumento) de um complexo \\
\lstinline|mod2(z)| & \code{OP\_STD\_MOD2} & módulo ao quadrado de um complexo \\
\lstinline|complex(x,y)| & \code{OP\_STD\_COMP} & cria um complexo a partir de dois reais \\
\lstinline|conj(z)| & \code{OP\_STD\_CONJ} & conjugado complexo ($a - bi$); um real vira $x + 0i$ \\
\bottomrule
\end{longtable}

As funções \lstinline|real|, \lstinline|imag|, \lstinline|mod2| e
\lstinline|fase| exigem argumento complexo (\lstinline|MSG_ERR_REAL_ARG_COMP| e
análogas); \lstinline|complex(.,.)| não aceita argumentos já complexos
(\lstinline|MSG_ERR_COMPLEX_OF_COMPLEX|); \lstinline|conj| de um real emite
\lstinline|MSG_WARN_CONJ_REAL|.

\section{Álgebra linear em notação de Dirac (bra-ket)}
\label{sec:cmm-dirac}

A \cmm{} incorpora um conjunto de \emph{statements} em notação de Dirac
(bra-ket) para álgebra linear sobre vetores e matrizes. Os delimitadores são
tratados no scanner:

\begin{lstlisting}[language=C]
"⟩"   return BRA;
"⟨"   return KET;
"|I|" return EYE;     // matriz identidade
"|0⟩" return VZERO;   // vetor nulo
\end{lstlisting}

\begin{nota}
A sintaxe usa os caracteres Unicode de bra-ket: \lstinline|⟨| (token \code{KET})
e \lstinline|⟩| (token \code{BRA}). O caractere \lstinline|#| separa o vetor ou
matriz de destino da expressão. O operador \lstinline+|+ delimita matrizes e
vetores conforme a forma.
\end{nota}

\subsection{Produto interno como expressão}

A única forma de Dirac que aparece como \emph{expressão} (e não como
\emph{statement}) é o produto interno $\langle a | b \rangle$:

\begin{lstlisting}[language=C]
exp : KET ID '|' ID BRA   {$$ = expr_inner(expr_var(...,$2), expr_var(...,$4));}
\end{lstlisting}

ou seja, \lstinline+<a|b>+ devolve o produto interno entre os vetores
\lstinline|a| e \lstinline|b|. Vetores de tamanhos diferentes geram
\lstinline|MSG_ERR_VECTOR_SIZE_DIFF|; operandos que não são vetores geram
\lstinline|MSG_ERR_INNER_NEEDS_VECTORS|.

\subsection{Atribuições em notação de Dirac}

As demais formas são \emph{statements} de atribuição (parte da regra
\code{assignment}). A tabela lista cada forma com a função de codegen invocada,
exatamente como no parser:

\begin{longtable}{@{}p{0.42\linewidth}p{0.50\linewidth}@{}}
\toprule
\textbf{Forma} & \textbf{Significado} \\
\midrule
\endhead
\lstinline+A # |B|a>;+ & $A \leftarrow B \cdot a$ (matriz-vetor); \lstinline|stmt_dirac_Mv| \\
\lstinline+a # c|b>;+ & $a \leftarrow c \cdot b$ (constante-vetor); \lstinline|stmt_dirac_cv| \\
\lstinline+a # |b> + c|d>;+ & $a \leftarrow b + c \cdot d$ (soma de vetores); \lstinline|stmt_dirac_apcb| \\
\lstinline+A # |a><b|;+ & $A \leftarrow a \cdot b^{T}$ (produto externo); \lstinline|stmt_dirac_vvt| \\
\lstinline+A # |B| - |a><b|;+ & $A \leftarrow B - a \cdot b^{T}$; \lstinline|stmt_dirac_Mmvvt| \\
\lstinline+A # c|B|;+ & $A \leftarrow c \cdot B$ (constante-matriz); \lstinline|stmt_dirac_cM| \\
\lstinline+A # c|I|;+ & $A \leftarrow c \cdot I$ (identidade ponderada); \lstinline|stmt_dirac_cI| \\
\lstinline+a # |0>;+ & zera todos os elementos de \lstinline|a|; \lstinline|stmt_dirac_v0| \\
\lstinline+a # c|in(p)>;+ & preenche \lstinline|a| com a entrada da porta \lstinline|p| ponderada por \lstinline|c|; \lstinline|stmt_dirac_cvin| \\
\lstinline+a # c -> |a>;+ & registrador de deslocamento (\emph{shift register}); \lstinline|stmt_dirac_shift| \\
\bottomrule
\end{longtable}

A saída de um vetor ponderado, \lstinline+out(p, c|v>);+ (regra \code{std\_vout},
\lstinline|stmt_vout|), também usa a notação. Erros de dimensão e de tipo nessas
operações têm mensagens dedicadas: \lstinline|MSG_ERR_NOT_A_VECTOR|,
\lstinline|MSG_ERR_NOT_A_MATRIX|, \lstinline|MSG_ERR_MATRIX_ROW_MISMATCH|,
\lstinline|MSG_ERR_DIM_MISMATCH|, entre outras.

\subsection{Exemplo: RLS (\emph{Recursive Least Squares})}

O exemplo \file{Tests/proc\_rls/Software/proc\_rls.cmm} reúne a maioria das
formas, inclusive a declaração de \emph{array} inicializada por Dirac
(\lstinline+float Px[4] # |P|x>;+):

\begin{lstlisting}[style=cmm]
void rls_update(float d)
{
    float e = d - ⟨w|x⟩;             // produto interno como expressao

    float Px[4] # |P|x⟩;             // declaracao: Px = P * x
    float g     = 1.0/(0.99 + ⟨x|Px⟩);
    float K [4] # g|Px⟩;             // declaracao: K = g * Px

    w # |w⟩ + e|K⟩;                  // soma de vetores
    P # |P| - |K⟩⟨Px|;               // matriz - produto externo
    P # 1.010101|P|;                 // constante * matriz
}

void main()
{
    P # 1000.0|I|;                   // P = 1000 * I
    w # |0⟩;                         // zera w
    int j = 0;
    while (j < 20)
    {
        x # fin(0)*0.001 -> |x⟩;     // shift register a partir de entrada float
        rls_update(in(1)*0.001);
        out(0,1000.0|w⟩);            // saida de vetor ponderado
        j++;
    }
}
\end{lstlisting}

\section{Pré-processamento: \texttt{\#define} e \texttt{\#include}}
\label{sec:cmm-preproc}

\subsection{\texttt{\#define} (macro de objeto)}

O \lstinline|#define| é tratado \emph{inteiramente no scanner}, conforme o
cabeçalho de \file{CMMComp.l}. É suportada apenas a forma \enquote{de objeto}
(\emph{object-like}): \lstinline|#define NAME body| registra o corpo, e cada uso
posterior de \lstinline|NAME| é substituído reanalisando o corpo a partir de um
\emph{buffer} empilhado. Um sinalizador \lstinline|active| por macro impede que
uma definição auto-referente se expanda indefinidamente. Macros \emph{de função},
\lstinline|#ifdef| e \lstinline|#include| estão fora do escopo do scanner da
\cmm{}.

\begin{lstlisting}[style=cmm]
#define SCALE  3
#define OFFSET 7
// ... SCALE e OFFSET sao expandidos para 3 e 7 antes de o parser ve-los
\end{lstlisting}

O teste negativo \file{NegTests/define\_recursive.cmm} verifica que uma macro que
se referencia recursivamente termina como variável não declarada (a expansão não
entra em laço infinito).

\subsection{Macros de \emph{assembly} (includes da stdlib)}

As funções transcendentais são implementadas como \emph{macros de assembly},
disponíveis em \file{Compilers/CMMComp/Includes/}. Esses arquivos contêm rotinas
prontas que o codegen injeta quando a função correspondente é usada (emitindo as
mensagens \lstinline|MSG_INFO_*_MACRO|). Os arquivos disponíveis são:

\begin{tabularx}{\linewidth}{Y Y}
\toprule
\textbf{Arquivo} & \textbf{Função} \\
\midrule
\file{float\_sqrt.asm} & raiz quadrada \\
\file{float\_atan.asm} & arco-tangente \\
\file{float\_sin.asm} & seno \\
\file{float\_tan.asm} & tangente \\
\file{float\_exp.asm} & exponencial \\
\file{float\_log.asm} & logaritmo natural \\
\bottomrule
\end{tabularx}

Como ilustração, \file{float\_sqrt.asm} implementa a raiz por iterações de
Newton-Raphson sobre uma estimativa inicial (potência de 2 mais próxima):

\begin{lstlisting}[language=bash]
@float_sqrt     SET sqrt_num       // obtem o parametro
              F_ROT                // primeira estimativa
                PSH
              F_DIV sqrt_num       // iteracao 1
             SF_ADD
              F_MLT 0.5
                ...                // mais iteracoes
                RET
\end{lstlisting}

Alguns exemplos (como \file{Tests/Seno} e \file{Tests/ArcTan}) também usam
tabelas de \emph{look-up} (\path{*_LUT.txt}) carregadas via inicialização de
\emph{array} por arquivo.

\section{Catálogo de mensagens do compilador}
\label{sec:cmm-mensagens}

O compilador é bilíngue: \file{messages.h} define cada mensagem com a macro
\lstinline|M(pt, en)|, e a opção de linha de comando \lstinline|-en|/\lstinline|-pt|
seleciona o idioma (padrão: português). As mensagens dividem-se em erros
(abortam a compilação), avisos (\enquote{atenção}, não abortam) e informações.
A tabela a seguir resume as principais categorias e mensagens representativas;
todas têm o símbolo correspondente em \file{messages.h}.

\begin{longtable}{@{}p{0.30\linewidth}p{0.30\linewidth}p{0.32\linewidth}@{}}
\toprule
\textbf{Categoria} & \textbf{Símbolo} & \textbf{Condição} \\
\midrule
\endhead
Geral & \lstinline|MSG_ERR_SYNTAX| & erro de sintaxe (parser) \\
Geral & \lstinline|MSG_ERR_NO_MAIN| & ausência de \lstinline|main()| \\
Geral & \lstinline|MSG_ERR_CANT_OPEN_FILE| & arquivo não pôde ser aberto \\
Variável & \lstinline|MSG_ERR_VAR_EXISTS| & redeclaração de variável \\
Variável & \lstinline|MSG_ERR_VAR_NOT_FOUND| & variável inexistente \\
Variável & \lstinline|MSG_ERR_DECLARE_VAR_PLEASE| & uso de variável não declarada \\
Número & \lstinline|MSG_ERR_RESERVED_I| & \lstinline|i| usado como identificador \\
Número & \lstinline|MSG_ERR_INT_MAX_OVERFLOW| & inteiro fora da faixa \\
\emph{Array} & \lstinline|MSG_ERR_ARRAY_NEEDS_IDX| & \emph{array} usado sem índice \\
\emph{Array} & \lstinline|MSG_ERR_NOT_ARRAY| & indexação de não-\emph{array} \\
\emph{Array} & \lstinline|MSG_ERR_ARRAY_2D| / \lstinline|_1D| & número de dimensões incompatível \\
Complexo & \lstinline|MSG_ERR_INCR_COMPLEX| & incremento de complexo \\
Complexo & \lstinline|MSG_ERR_SQRT_COMPLEX| & \lstinline|sqrt| de complexo (não implementado) \\
Complexo & \lstinline|MSG_ERR_REAL_ARG_COMP| & argumento de \lstinline|real| não é complexo \\
I/O & \lstinline|MSG_ERR_NO_IN_PORT| & porta de entrada inexistente \\
I/O & \lstinline|MSG_ERR_NO_OUT_PORT| & porta de saída inexistente \\
Operador & \lstinline|MSG_ERR_MOD_NON_INT| & \lstinline|%| com não-inteiro \\
Operador & \lstinline|MSG_ERR_BITWISE_COMPLEX| & bit-a-bit com complexo \\
Operador & \lstinline|MSG_ERR_SHIFT_COMPLEX| & deslocamento de complexo \\
Operador & \lstinline|MSG_ERR_LOGIC_NON_INT| & operação lógica com não-inteiro \\
Álgebra & \lstinline|MSG_ERR_INNER_NEEDS_VECTORS| & produto interno exige vetores \\
Álgebra & \lstinline|MSG_ERR_VECTOR_SIZE_DIFF| & vetores de tamanhos diferentes \\
Álgebra & \lstinline|MSG_ERR_DIM_MISMATCH| & dimensões incompatíveis \\
Função & \lstinline|MSG_ERR_PARAM_COUNT| & número de parâmetros errado \\
Função & \lstinline|MSG_ERR_FUNC_WHERE| & função inexistente \\
Função & \lstinline|MSG_ERR_VOID_RETURN_VALUE| & retorno com valor em \lstinline|void| \\
Salto & \lstinline|MSG_ERR_BREAK_LOST| & \lstinline|break| fora de laço/switch \\
Salto & \lstinline|MSG_ERR_CONTINUE_LOST| & \lstinline|continue| fora de laço \\
Interrupção & \lstinline|MSG_ERR_DUP_INTERRUPT| & duas interrupções definidas \\
\bottomrule
\end{longtable}

Além dos erros, há famílias de avisos de conversão de tipo
(\lstinline|MSG_WARN_INT_RECV_FLOAT|, \lstinline|MSG_WARN_CONV_*_PARAM|,
\lstinline|MSG_WARN_RET_*|), de arredondamento de índice de \emph{array}
(\lstinline|MSG_WARN_IDX_FLOAT| e variantes), de comparação envolvendo complexos
(\lstinline|MSG_WARN_CMP_*|) e de variáveis/funções não usadas
(\lstinline|MSG_WARN_UNUSED_*|). Ao final de uma compilação bem-sucedida o
compilador imprime \lstinline|MSG_OK_CMM_DONE|.

\subsection{Testes negativos}

O diretório \file{NegTests/} contém fixtures malformadas que o compilador deve
rejeitar de forma limpa (saída não-zero, sem \emph{crash}). O manifesto
\file{NegTests/manifest.txt} associa cada fixture à substring esperada na
mensagem:

\begin{tabularx}{\linewidth}{Y Y}
\toprule
\textbf{Fixture} & \textbf{Mensagem esperada} \\
\midrule
\path{syntax.cmm} & \enquote{Syntax error} \\
\path{no_main.cmm} & \enquote{main() function} \\
\path{garbage.cmm} & \enquote{Syntax error} \\
\path{undeclared.cmm} & \enquote{declared the variable} \\
\path{define_recursive.cmm} & \enquote{declare the variable} \\
\bottomrule
\end{tabularx}

\section{Interface de linha de comando do \texttt{cmmcomp}}
\label{sec:cmm-cli}

O ponto de entrada está em \lstinline|main()| no próprio \file{CMMComp.y} e usa o
parser de argumentos de \file{args.c}. Pelas mensagens de ajuda em
\file{messages.h} (\lstinline|MSG_CLI_HELP|), a invocação é:

\begin{lstlisting}[language=bash]
cmmcomp [opcoes] -i <arq.cmm> -n <nome> -p <dir> -m <dir> -t <dir>
\end{lstlisting}

\begin{longtable}{@{}p{0.26\linewidth}p{0.66\linewidth}@{}}
\toprule
\textbf{Opção} & \textbf{Descrição} \\
\midrule
\endhead
\lstinline|-i|, \lstinline|--input| & arquivo-fonte \cmm{} (obrigatória) \\
\lstinline|-n|, \lstinline|--name| & nome do processador / base do \path{.asm} (obrigatória) \\
\lstinline|-p|, \lstinline|--proc-dir| & diretório do processador (obrigatória) \\
\lstinline|-m|, \lstinline|--macros-dir| & diretório de Macros (obrigatória) \\
\lstinline|-t|, \lstinline|--temp-dir| & diretório Temp (obrigatória) \\
\lstinline|-A|, \lstinline|--array| & expande cada elemento de \emph{array} como sinal individual no GTKWave \\
\lstinline|-en|, \lstinline|-pt| & idioma das mensagens (padrão: \lstinline|-pt|) \\
\lstinline|-h|, \lstinline|--help| & mostra ajuda e sai \\
\lstinline|-V|, \lstinline|--version| & mostra a versão e sai \\
\bottomrule
\end{longtable}

\section{O caminho C++}
\label{sec:cmm-cpp}

Além da \cmm{}, o YANC oferece um caminho alternativo a partir de C/C++ para o
mesmo soft-processor, composto por dois binários em
\file{Compilers/CPPComp/Sources/}:

\begin{description}[leftmargin=2em]
  \item[\tool{cpppp}] o pré-processador C autônomo (\file{cpppp.c}). Trata
  \lstinline|#include "..."| e \lstinline|#include <...>| (com diretórios
  \lstinline|-I|), \lstinline|#define| de objeto e de função (incluindo
  variádicas \lstinline|...|/\lstinline|__VA_ARGS__|, \emph{stringize} \lstinline|#|
  e \emph{paste} \lstinline|##|), \lstinline|#undef|,
  \lstinline|#ifdef|/\lstinline|#ifndef|/\lstinline|#else|/\lstinline|#endif|,
  \lstinline|#pragma once| e repassa \lstinline|#pragma yanc| verbatim.
  \item[\tool{cppcomp}] o compilador propriamente dito (\file{main.c},
  \file{CPPComp.l}, \file{CPPComp.y}, \file{codegen.c}). Consome o código já
  pré-processado e emite \path{.asm} no mesmo layout de pipeline do
  \tool{cmmcomp} (\path{<proc_dir>/Software/<prname>.asm}).
\end{description}

\subsection{Subconjunto da linguagem aceito}

O scanner \file{CPPComp.l} tokeniza um subconjunto de C99 com extensões de C++.
As palavras-chave reconhecidas incluem, além das de fluxo de controle
(\lstinline|if|, \lstinline|else|, \lstinline|while|, \lstinline|for|,
\lstinline|do|, \lstinline|switch|, \lstinline|case|, \lstinline|default|,
\lstinline|break|, \lstinline|continue|, \lstinline|return|, \lstinline|goto|),
os tipos \lstinline|void|, \lstinline|int|, \lstinline|float|, \lstinline|char|,
\lstinline|short|, \lstinline|long|, \lstinline|double|, \lstinline|bool| /
\lstinline|_Bool|, \lstinline|unsigned| e \lstinline|signed|. Entre os recursos
de C++ reconhecidos estão \lstinline|struct|, \lstinline|class|, \lstinline|this|,
\lstinline|virtual|, \lstinline|override|, \lstinline|final|, os \emph{casts} de
C++ (\lstinline|static_cast|, \lstinline|reinterpret_cast|, \lstinline|const_cast|,
\lstinline|dynamic_cast|), \lstinline|operator|, \lstinline|template|,
\lstinline|typename|, \lstinline|auto|, \lstinline|public| /
\lstinline|private| / \lstinline|protected|, \lstinline|namespace|,
\lstinline|using|, \lstinline|new| / \lstinline|delete| e
\lstinline|sizeof|. Literais \lstinline|true|, \lstinline|false| e
\lstinline|nullptr| são tratados como inteiros.

Para desambiguar nomes de tipo de nomes de expressão, o scanner usa o
\enquote{\emph{typedef-name hack}}: a cada identificador, consulta a tabela de
símbolos e emite \code{TYPEDEF\_NAME} ou \code{TEMPLATE\_NAME} no lugar de
\code{IDENT} quando o nome resolve para um \emph{typedef} ou para um
\emph{template} de função.

Diversos especificadores são reconhecidos e ignorados no alvo
(\lstinline|inline|, \lstinline|register|, \lstinline|volatile|,
\lstinline|restrict|, \lstinline|explicit|, \lstinline|mutable|,
\lstinline|noexcept|, \lstinline|_Noreturn|); \lstinline|constexpr| é
tratado como \lstinline|const| (token \code{KW\_CONST}).

\subsection{O alvo e o \texttt{\#pragma yanc}}

C/C++ não carrega diretivas de processador na sintaxe; em vez disso, os
parâmetros do alvo têm \emph{padrões fixos} em \file{Headers/config.h}, que
\emph{são} o alvo YANC: palavra de 32 bits, mantissa 23 / expoente 8 (formato
estilo IEEE-754 de precisão simples), pilhas de dados e de retorno com 128
posições. Um programa pode sobrescrever esses padrões por arquivo com
\lstinline|#pragma yanc <chave> <valor>| no escopo de arquivo, processado
diretamente no scanner. As chaves reconhecidas espelham as diretivas \cmm{}:
\lstinline|prname|, \lstinline|nubits|, \lstinline|nbmant|, \lstinline|nbexpo|,
\lstinline|nugain|, \lstinline|ndstac|, \lstinline|sdepth|, \lstinline|nuioin|,
\lstinline|nuioou|, \lstinline|fftsiz| e \lstinline|itradd|. A tabela mostra os
padrões de \file{config.h}:

\begin{tabularx}{\linewidth}{Y Y}
\toprule
\textbf{Parâmetro} & \textbf{Padrão} \\
\midrule
\lstinline|CFG_NUBITS| & 32 \\
\lstinline|CFG_NBMANT| & 23 \\
\lstinline|CFG_NBEXPO| & 8 \\
\lstinline|CFG_NUGAIN| & 128 \\
\lstinline|CFG_NDSTAC| & 128 \\
\lstinline|CFG_SDEPTH| & 128 \\
\lstinline|CFG_NUIOIN| & 1 \\
\lstinline|CFG_NUIOOU| & 1 \\
\lstinline|CFG_FFTSIZ| & 3 \\
\bottomrule
\end{tabularx}

\subsection{Intrínsecas de I/O e \emph{assembly} inline}

A I/O do alvo é exposta como funções intrínsecas reconhecidas pelo codegen
(\file{codegen.c}): \lstinline|in(porta)| emite a instrução \code{INN} (a porta
deve ser literal inteiro) e \lstinline|out(porta, valor)| emite \code{OUT}. Como
a saída do hardware é uma palavra inteira, enviar um \lstinline|float| por
\lstinline|out| gera um aviso orientando a converter para inteiro. Há ainda
intrínsecas de heap (\lstinline|malloc|, \lstinline|free|) quando o programa usa
\lstinline|new|/\lstinline|delete| ou alocação dinâmica.

A linguagem aceita \emph{assembly} embutido pela forma
\lstinline|asm("...");| (regra \code{S\_ASM} no parser; sinônimos
\lstinline|asm| e \lstinline|__asm__| no scanner), emitido verbatim e nunca
fundido pelo \emph{peephole}.

\subsection{Cabeçalhos de biblioteca (\texttt{Includes})}

O diretório \file{Compilers/CPPComp/Includes/} fornece versões enxutas de
cabeçalhos da biblioteca padrão adaptadas ao alvo:

\begin{tabularx}{\linewidth}{Y Y}
\toprule
\textbf{Cabeçalho} & \textbf{Conteúdo} \\
\midrule
\file{cmath} & subconjunto de funções matemáticas em software \\
\file{cstdint} & inteiros de largura fixa \\
\file{cstddef} & definições básicas (\lstinline|size_t|, etc.) \\
\file{cstring} & funções de manipulação de memória/strings \\
\file{vector} & contêiner dinâmico \\
\file{array} & contêiner de tamanho fixo \\
\file{bit} & utilitários de bits \\
\file{limits} & limites numéricos \\
\bottomrule
\end{tabularx}

Como exemplo, \file{cmath} implementa \lstinline|sqrt| por Newton-Raphson em
software (24 iterações) e expõe \lstinline|isfinite|/\lstinline|isnan|/%
\lstinline|isinf| como \emph{no-ops}, pois o formato de \emph{float} do alvo não
tem padrões de bits para NaN ou infinito.

\subsection{Recursos de C++ exercitados pelos testes}

A suíte \file{Compilers/CPPComp/Tests/} contém dezenas de programas que vão de um
\emph{smoke test} em C (\file{test1/Software/test1.cpp}) a uma mini-STL escrita na
própria linguagem (\file{test21/Software/test21.cpp}, com um \lstinline|Vector|
templatizado com RAII, crescimento sob demanda e \lstinline|operator[]|, além de
um \lstinline|unique_ptr| dono de recurso). Os testes mais avançados
(\file{test44} a \file{test50}) implementam pipelines de DSP
(\emph{blind deconvolution}, FISTA, FIR inverso) em múltiplos arquivos
\path{.cpp}/\path{.hpp}, demonstrando que o caminho C++ suporta projetos
multi-arquivo de porte realista para o mesmo soft-processor SAPHO.

\begin{nota}
O \tool{cppcomp} compartilha o número de versão (\file{yanc\_version.h}) e o
layout de saída com o \tool{cmmcomp}: ambos produzem \path{.asm} em
\path{<proc_dir>/Software/<prname>.asm}, consumido pelos estágios seguintes do
pipeline (ver \autoref{cap:02-yanc-pipeline}).
\end{nota}

% !TEX root = ../main.tex
\chapter{YANC: os compiladores e o pipeline de tradução}
\label{cap:04-yanc-compiladores}

O \tool{YANC} (\enquote{Yet Another Compiler}) é o conjunto de programas de linha de
comando que transforma o programa de alto nível do usuário no hardware
sintetizável de um \emph{soft-processor} \tool{SAPHO}, junto com as imagens de
memória e o \emph{testbench} necessários para simulá-lo. Diferentemente de um
compilador convencional, cuja saída é código de máquina para uma CPU já
existente, o \tool{YANC} emite Verilog: o produto final é a descrição RTL de um
processador customizado mais a memória de instruções (\file{.mif}) que esse
processador executa. A IDE \tool{AURORA} (\autoref{cap:05-aurora-arquitetura}) é a
camada que orquestra esses binários; este capítulo documenta os binários em si e
a cadeia de tradução que eles compõem.

Todo o material descrito aqui vive no repositório \repo{yanc}, em
\path{Compilers/}, \path{Scripts/} e no \file{Makefile} na raiz. A versão da
suíte é definida em um único lugar, \file{Compilers/yanc_version.h}
(\lstinline|#define YANC_VERSION "5.2"|), incluído por todos os binários para que
a saída de \lstinline|--version| seja consistente.

\section{Os sete binários e o modelo de construção}
\label{sec:yanc-binarios}

O \file{Makefile} é a fonte única de verdade da construção: ele define, em um só
lugar, como cada binário é compilado, e é consumido por \file{Scripts/setup.sh},
\file{Scripts/setup.bat}, \file{Scripts/aurora.bat} e pelos fluxos de CI/release.
São sete binários produzidos em \path{bin/}:

\begin{longtable}{l l l}
\toprule
\textbf{Binário} & \textbf{Entrada} & \textbf{Saída} \\
\midrule
\endhead
\tool{cmmcomp}   & \cmm{} (\file{.cmm})            & \file{.asm} + tabelas de log \\
\tool{cppcomp}   & C++ pré-processado (\file{.c})  & \file{.asm} + tabelas de log \\
\tool{cpppp}     & C++ (\file{.cpp}, com diretivas)& C++ pré-processado \\
\tool{appcomp}   & \file{.asm}                     & \file{app\_log.txt} (1ª passada) \\
\tool{asmcomp}   & \file{.asm} (+ logs)            & \file{.v} + \file{.mif} + \file{\_tb.v} \\
\tool{comp2gtkw} & padrão de bits (stdin)          & número complexo legível (stdout) \\
\tool{gen\_gtkw} & cabeçalho de \file{.vcd}        & \file{.gtkw} formatado \\
\bottomrule
\end{longtable}

\subsection{Tecnologias de construção: Flex, Bison e GCC}
\label{sec:yanc-flexbison}

Os binários se dividem por tecnologia, conforme as receitas do \file{Makefile}:

\begin{description}[leftmargin=1.2em]
  \item[Flex + Bison] \tool{cmmcomp} e \tool{cppcomp}. A receita roda
    \lstinline|bison -y -d <gram>.y| (gera \file{y.tab.c} e \file{y.tab.h}) e
    \lstinline|flex <lexer>.l| (gera \file{lex.yy.c}); em seguida o GCC compila
    esses arquivos gerados junto com os \file{.c} escritos à mão.
  \item[Apenas Flex] \tool{appcomp} (de \file{app.l}) e \tool{asmcomp} (de
    \file{ASMComp.l}). Não há gramática \file{.y}: a análise é dirigida por
    estados dentro do próprio analisador léxico.
  \item[Apenas GCC] \tool{cpppp} (uma única unidade de tradução, \file{cpppp.c}),
    \tool{comp2gtkw} e \tool{gen\_gtkw} (\file{Scripts/comp2gtkw.c} e
    \file{Scripts/gen\_gtkw.c}).
\end{description}

A lista de fontes por binário é fixada no \file{Makefile}. Por exemplo,
\tool{cmmcomp} compila \lstinline|ast.c data_assign.c oper.c stdlib.c| (entre
outros) mais \file{CMMComp.l} e \file{CMMComp.y}; \tool{asmcomp} compila
\lstinline|eval.c labels.c opcodes.c hdl.c simulacao.c array.c| mais
\file{ASMComp.l}.

\subsection{Compilação cruzada para Windows sem dependência de DLL}
\label{sec:yanc-crosscompile}

O \file{Makefile} detecta a plataforma por \lstinline|$(OS)|: quando é
\lstinline|Windows_NT| (inclusive sob MSYS2), o sufixo de saída vira \file{.exe};
em POSIX não há sufixo. O compilador é parametrizável por \lstinline|CC|, e o
cabeçalho da receita documenta explicitamente o uso de um \emph{toolchain} de
compilação cruzada:

\begin{lstlisting}[language=bash]
make CC=x86_64-w64-mingw32-gcc        # Windows standalone (sem DLLs do MSYS2)
\end{lstlisting}

Construir com \tool{x86\_64-w64-mingw32-gcc} produz \file{.exe} que não dependem
das DLLs de runtime do MSYS2, ou seja, binários autocontidos que a \tool{AURORA}
pode empacotar e distribuir. No fluxo nativo do Windows, \file{setup.bat}
localiza preferencialmente \lstinline|x86_64-w64-mingw32-gcc.exe| no \lstinline|PATH|
ou em \path{<msys2>/mingw64/bin} (\file{Scripts/setup.bat}). Os alvos
\lstinline|make <nome>| (curtos, sem sufixo) e \lstinline|make install DESTDIR=<dir>|
completam o conjunto de operações de build.

\section{O front end C+-: \tool{cmmcomp}}
\label{sec:yanc-cmmcomp}

O \tool{cmmcomp} é o compilador da linguagem \cmm{}, descrita em detalhe no
\autoref{cap:03-linguagem-cmm}; aqui o foco é seu papel como binário no pipeline.
Ele lê um arquivo \file{.cmm} e produz um \file{.asm} mais um conjunto de tabelas
auxiliares no diretório temporário.

\subsection{Interface de linha de comando}
\label{sec:yanc-cmmcomp-cli}

A análise de argumentos vive em \file{Compilers/CMMComp/Sources/args.c}. Todas as
opções que carregam caminhos são obrigatórias; a falta de qualquer uma encerra o
programa com mensagem de uso. As opções (de \file{args.h}, \lstinline|struct cli_args|):

\begin{longtable}{l l}
\toprule
\textbf{Opção} & \textbf{Significado} \\
\midrule
\endhead
\lstinline|-i / --input|       & arquivo-fonte \cmm{} (dentro de \path{<proc-dir>/Software}) \\
\lstinline|-n / --name|        & nome do processador (base do \file{.asm} de saída) \\
\lstinline|-p / --proc-dir|    & diretório do processador \\
\lstinline|-m / --macros-dir|  & diretório de macros \\
\lstinline|-t / --temp-dir|    & diretório temporário \\
\lstinline|-A / --array|       & emite cada elemento de array como sinal próprio no GTKWave \\
\lstinline|-h / --help|        & imprime ajuda e sai \\
\lstinline|-V / --version|     & imprime versão e sai \\
\bottomrule
\end{longtable}

Há ainda um \emph{flag} de idioma processado antes de tudo
(\lstinline|parse_lang_flag(&argc, argv)| em \lstinline|main|), que seleciona o
idioma das mensagens (\lstinline|-en|/\lstinline|-pt|) e o remove de
\lstinline|argv|.

\subsection{Léxico, diretivas e \texttt{\#define} no analisador}
\label{sec:yanc-cmmcomp-lexer}

O analisador léxico (\file{CMMComp.l}) reconhece as diretivas de configuração do
processador como \emph{tokens} próprios: \lstinline|#PRNAME|, \lstinline|#NUBITS|,
\lstinline|#NBMANT|, \lstinline|#NBEXPO|, \lstinline|#NDSTAC|, \lstinline|#SDEPTH|,
\lstinline|#NUIOIN|, \lstinline|#NUIOOU|, \lstinline|#NUGAIN|, \lstinline|#FFTSIZ|,
mais as diretivas comportamentais \lstinline|#PRACA| (ponto de retorno da
interrupção) e \lstinline|#TOAQUI| (marcador de PC). Reconhece também as palavras
da biblioteca padrão (\lstinline|in|, \lstinline|out|, \lstinline|norm|,
\lstinline|sqrt|, \lstinline|sin|, \lstinline|complex|, etc.), os tipos
(\lstinline|int|, \lstinline|float|, \lstinline|comp|, \lstinline|void|), os
operadores de dois símbolos (\lstinline|<<|, \lstinline|>>>|, \lstinline|==|,
\lstinline|&&|) e a notação de Dirac em UTF-8 (\lstinline|BRA|/\lstinline|KET|).

Um \lstinline|#define| do tipo objeto é tratado inteiramente no léxico: a regra
captura \lstinline|#define NAME corpo|, registra o corpo em uma tabela
(\lstinline|g_def|, capacidade \lstinline|CMM_MAX_DEFINES = 256|) e, ao reencontrar
\lstinline|NAME|, re-analisa o corpo a partir de um buffer empilhado
(\lstinline|yy_scan_string| + \lstinline|yypush_buffer_state|), desempilhado no
\lstinline|<<EOF>>|. Um \emph{flag} \lstinline|active| por macro impede a expansão
infinita de uma definição auto-referente. Macros do tipo função, \lstinline|#ifdef|
e \lstinline|#include| estão fora de escopo do léxico do \cmm{}.

\subsection{Gramática dirigida por AST e a entrada do programa}
\label{sec:yanc-cmmcomp-ast}

A gramática (\file{CMMComp.y}) constrói uma AST (\emph{abstract syntax tree})
durante o \emph{parse}: cada regra produtora monta a subárvore correspondente
(\lstinline|expr_node *| na pilha do Bison) sem emitir código inline. Quando o
\emph{parse} termina (regra \lstinline|fim : prog|), dispara-se
\lstinline|prog_emit()|, que percorre a AST inteira e gera o assembly em uma única
varredura. O \lstinline|main| (em \file{CMMComp.y}) encadeia
\lstinline|parse_lang_flag| $\rightarrow$ \lstinline|cli_parse| $\rightarrow$
\lstinline|parse_init| $\rightarrow$ \lstinline|yyparse| $\rightarrow$
\lstinline|parse_end|.

As diretivas têm efeito duplo (em \file{diretivas.c}, \lstinline|dire_exec|):
atualizam estado global de compilação lido pelas demais regras
(p.\,ex.\ \lstinline|nbmant|, \lstinline|nbexpo|, \lstinline|nuioin|,
\lstinline|nuioou|, \lstinline|prname|) e, ao mesmo tempo, anexam um nó
\lstinline|STMT_DIRECTIVE| à lista de \emph{statements}, para que sua emissão
aconteça no percurso final da AST.

\subsection{A inicialização e os arquivos gerados}
\label{sec:yanc-cmmcomp-init}

\lstinline|parse_init| (em \file{global.c}) abre o \file{.cmm} de entrada em
\path{<proc-dir>/Software/<input>} e cria o \file{.asm} de saída em
\path{<proc-dir>/Software/<name>.asm}. Cria ainda dois arquivos auxiliares no
diretório temporário: \file{cmm_log.txt} (informações para o assembler e para o
GTKWave) e \file{pc\_<name>\_mem.txt} (a \emph{side-table} de linhas, descrita na
\autoref{sec:yanc-sidetable}). Ao final, \lstinline|parse_end| anexa as macros
necessárias ao \file{.asm} (\lstinline|mac_copy|), verifica consistência das
variáveis (\lstinline|check_var|), grava \lstinline|num_ins| no
\file{cmm_log.txt} e gera \file{trad_cmm.txt}, o mapa de número de linha para
texto-fonte.

\section{O front end C++: \tool{cpppp} e \tool{cppcomp}}
\label{sec:yanc-cpp}

O caminho C++ usa dois binários em sequência, espelhando a disposição de
diretórios do caminho \cmm{}.

\subsection{\tool{cpppp}: o pré-processador C}
\label{sec:yanc-cpppp}

O \tool{cpppp} (uma única unidade, \file{Compilers/CPPComp/Sources/cpppp.c}) roda
antes do \tool{cppcomp} e trata as diretivas do pré-processador, conforme o
cabeçalho do arquivo: \lstinline|#include "nome"| (relativo ao arquivo atual e aos
caminhos \lstinline|-I|), \lstinline|#include <nome>| (apenas caminhos
\lstinline|-I|), \lstinline|#define NAME corpo| (macros tipo objeto e tipo função,
incluindo variádicas com \lstinline|...|/\lstinline|__VA_ARGS__|, \emph{stringize}
com \lstinline|#| e colagem com \lstinline|##|), \lstinline|#undef|,
\lstinline|#ifdef|/\lstinline|#ifndef|/\lstinline|#else|/\lstinline|#endif|,
\lstinline|#pragma once| e \lstinline|#pragma yanc ...| (repassado verbatim). Para
o \lstinline|#pragma once|, os caminhos são canonizados (\lstinline|path_canon|:
absolutos, com barras normais e em minúsculas no Windows) para que o mesmo arquivo
incluído por grafias diferentes compare igual. A invocação é
\lstinline|cpppp -i input.c [-o out.c] [-I dir]*|.

\subsection{\tool{cppcomp}: o compilador C}
\label{sec:yanc-cppcomp}

O \tool{cppcomp} (Flex + Bison: \file{CPPComp.l}, \file{CPPComp.y}, mais
\file{codegen.c}, \file{symtab.c}, \file{types.c}, \file{ast.c}, \file{main.c})
recebe o \file{.c} já pré-processado e emite o \file{.asm}. Seu \lstinline|main|
(em \file{main.c}) lê \lstinline|-i|, \lstinline|-p|, \lstinline|-n| e
\lstinline|-t|; constrói a AST (\lstinline|ast_unit|, \lstinline|st_init|,
\lstinline|yyparse|); deriva o nome do processador do basename quando
\lstinline|-n| não é dado; e grava a saída em
\path{<proc-dir>/Software/<prname>.asm} (criando \path{Software/} se faltar).

O ponto crucial para a integração é que \tool{cppcomp} produz \emph{as mesmas}
tabelas auxiliares que \tool{cmmcomp}, de modo que o restante do pipeline
(\tool{appcomp}, \tool{asmcomp}, \tool{gen\_gtkw}) é idêntico para os dois front
ends. Em \file{codegen.c}, \lstinline|write_cmm_log| grava \file{cmm_log.txt}
(uma linha por variável, no formato \lstinline|func var tipo [tamanho]|, mais
\lstinline|num_ins|) e \lstinline|write_trad_cmm| grava \file{trad_cmm.txt} com os
mesmos três cabeçalhos de \emph{scaffolding} negativos do caminho \cmm{}
(\lstinline|-1 INTERNAL|, \lstinline|-2 void main();|, \lstinline|-3 END|)
seguidos de cada linha-fonte numerada a partir de 1. Como o
\lstinline|src_path| é o arquivo efetivamente compilado (o \file{pp.cpp}
pré-processado), os números de linha batem com os carimbos de linha da AST.

\section{O formato do assembly intermediário (\file{.asm})}
\label{sec:yanc-asm}

O \file{.asm} é a representação intermediária que une os dois front ends ao back
end. É texto, uma instrução por linha, no formato \emph{mnemônico} seguido (quando
houver) de um operando, com comentários estilo \lstinline|//| e rótulos de salto
prefixados por \lstinline|@|. As diretivas de configuração aparecem como linhas
\lstinline|#...|. O trecho abaixo, extraído de
\file{Compilers/CMMComp/Includes/float_sqrt.asm}, é um exemplo real de corpo de
rotina:

\begin{lstlisting}
@float_sqrt     SET sqrt_num       // pega o parametro
              F_ROT                // estimativa inicial (potencia de 2 mais proxima)
                PSH                // atualiza x
              F_DIV sqrt_num       // iteracao 1
             SF_ADD
              F_MLT 0.5
                RET
\end{lstlisting}

\subsection{O conjunto de mnemônicos}
\label{sec:yanc-asm-mnemonics}

Os mnemônicos válidos são definidos identicamente nos analisadores de
\tool{appcomp} (\file{app.l}) e \tool{asmcomp} (\file{ASMComp.l}). Cada mnemônico
é mapeado a um índice de opcode e a um \enquote{próximo estado} que determina como
o operando seguinte será interpretado (sem operando; endereço em memória de dados;
endereço em memória de instruções; porta de entrada; porta de saída; ou
endereçamento indireto fixo de pseudo-instruções). As famílias incluem:

\begin{description}[leftmargin=1.2em]
  \item[Memória/acumulador] \lstinline|LOD|, \lstinline|P_LOD|, \lstinline|LDI|,
    \lstinline|ILI| (carga, com variantes de \emph{push} e endereçamento indireto/
    invertido para FFT); \lstinline|SET|, \lstinline|SET_P|, \lstinline|STI|,
    \lstinline|ISI| (armazenamento).
  \item[Pilha] \lstinline|PSH|, \lstinline|POP|.
  \item[E/S] \lstinline|INN|, \lstinline|F_INN|, \lstinline|P_INN|,
    \lstinline|PF_INN|, \lstinline|OUT|.
  \item[Saltos/sub-rotinas] \lstinline|JMP|, \lstinline|JIZ|, \lstinline|CAL|,
    \lstinline|RET|.
  \item[Aritmética] \lstinline|ADD|, \lstinline|MLT|, \lstinline|DIV|,
    \lstinline|MOD| e seus variantes de ponto flutuante (\lstinline|F_...|) e de
    pilha (\lstinline|S_...|, \lstinline|SF_...|).
  \item[Funções especiais] \lstinline|SGN|, \lstinline|NEG|, \lstinline|ABS|,
    \lstinline|PST|, \lstinline|NRM|, \lstinline|I2F|, \lstinline|F2I|,
    \lstinline|F_ROT| (raiz por potência de 2 mais próxima), \lstinline|F_SCL|,
    \lstinline|XPO|.
  \item[Lógicos/bit a bit] \lstinline|AND|, \lstinline|ORR|, \lstinline|XOR|,
    \lstinline|INV|, \lstinline|LAN|, \lstinline|LOR|, \lstinline|LIN|;
    comparações \lstinline|LES|, \lstinline|GRE|, \lstinline|EQU|;
    deslocamentos \lstinline|SHL|, \lstinline|SHR|, \lstinline|SRS|.
  \item[Ponteiros/pseudo-instruções] \lstinline|LDA|, \lstinline|STA|,
    \lstinline|LEA| (\emph{load-effective-address}) e a família \lstinline|..._V|
    (acesso com deslocamento constante, p.\,ex.\ \lstinline|LOD_V|,
    \lstinline|ADD_V|).
\end{description}

O \tool{asmcomp} associa a cada mnemônico um índice numérico de 7 bits
(\lstinline|NBITS_OPC = 7|, que deve casar com o Verilog em \file{proc.v}); o
operando ocupa \lstinline|nbopr| bits, calculado em \lstinline|eval_init| como
\lstinline|ceil(log2(max(n_ins, n_dat)))|.

\section{A primeira passada: \tool{appcomp}}
\label{sec:yanc-appcomp}

O \tool{appcomp} (somente Flex, \file{app.l} + \file{eval.c}) faz a primeira
passada sobre o \file{.asm}. Seu papel é \emph{coletar} os parâmetros do
processador e \emph{resolver} os endereços de variáveis e rótulos, gravando tudo
em \file{app_log.txt} (no diretório \lstinline|-t|). A invocação é
\lstinline|appcomp -i <arquivo.asm> -t <temp-dir>|.

A passada é necessária porque o \file{.asm} é único, mas referencia rótulos antes
de defini-los (saltos para frente). Contando instruções e registrando o índice
onde cada rótulo aparece, a primeira passada permite que a segunda passada já
conheça o endereço final de cada salto. Concretamente, em \file{eval.c}:

\begin{itemize}[leftmargin=1.2em]
  \item As diretivas \lstinline|#PRNAME|, \lstinline|#NUBITS|, \lstinline|#NBMANT|,
    \lstinline|#NBEXPO| têm seu argumento capturado e gravado
    (\lstinline|eval_direct| + \lstinline|eval_opernd|, estados 1 a 4).
  \item \lstinline|#ITRAD| e \lstinline|#TOAQUI| registram o índice da instrução
    corrente como \lstinline|itr_addr| e \lstinline|toaqui_addr|
    (\lstinline|eval_itrad|, \lstinline|eval_toaqui|).
  \item Cada opcode incrementa o contador de instruções \lstinline|n_ins| no
    momento certo (no operando, para instruções com operando; imediatamente, para
    instruções sem operando, em \lstinline|eval_opcode|).
  \item Cada operando que é variável é adicionado à tabela
    (\lstinline|var_add|), construindo a memória de dados; arrays com e sem
    inicialização são reconhecidos por estados dedicados.
  \item Cada rótulo \lstinline|@nome| é gravado com seu índice de instrução
    (\lstinline|eval_label| escreve \lstinline|@nome n_ins|).
  \item A pseudo-instrução \lstinline|LEA <nome>| conta como uma variável (o alvo)
    mais uma constante sintética \lstinline|__addr_<nome>| (o endereço
    materializado).
\end{itemize}

Ao final (\lstinline|eval_finish|), o log recebe \lstinline|n_ins| (número de
instruções) e \lstinline|n_dat| (número de dados, \lstinline|var_cnt()|). Há uma
verificação de sanidade: se a memória de dados teria duas células ou menos, o
processador é considerado inútil e a compilação é abortada.

\section{A segunda passada: \tool{asmcomp}}
\label{sec:yanc-asmcomp}

O \tool{asmcomp} (Flex, \file{ASMComp.l} + \file{eval.c} + \file{hdl.c} +
\file{simulacao.c} + \file{array.c} + \file{labels.c} + \file{opcodes.c}) é o back
end. Ele faz a segunda passada sobre o \emph{mesmo} \file{.asm}, agora com os
endereços já resolvidos pela primeira passada, e emite os três produtos finais:
o Verilog do processador, as imagens de memória \file{.mif} e o \emph{testbench}.

\subsection{Interface e inicialização}
\label{sec:yanc-asmcomp-cli}

As opções (de \file{ASMComp/Headers/args.h}) acrescentam, ao conjunto do
\tool{cmmcomp}, o diretório de HDL e os parâmetros de simulação:

\begin{longtable}{l l}
\toprule
\textbf{Opção} & \textbf{Significado} \\
\midrule
\endhead
\lstinline|-i / --input|       & arquivo \file{.asm} (caminho completo) \\
\lstinline|-p / --proc-dir|    & diretório do processador (saída em \path{Hardware/}) \\
\lstinline|-d / --hdl-dir|     & diretório de HDL \\
\lstinline|-m / --macros-dir|  & diretório de macros \\
\lstinline|-t / --temp-dir|    & diretório temporário (lê os logs aqui) \\
\lstinline|-f / --freq|        & frequência de operação em MHz \\
\lstinline|-c / --clocks|      & número de clocks a simular \\
\bottomrule
\end{longtable}

\lstinline|eval_init| lê \file{app_log.txt} para recuperar \lstinline|prname|,
\lstinline|n_ins|, \lstinline|n_dat|, \lstinline|nubits|, \lstinline|nbmant|,
\lstinline|nbexpo| e, se existirem, \lstinline|itr_addr| e \lstinline|toaqui_addr|.
Registra os rótulos (\lstinline|lab_reg| lê \file{app_log.txt} e, ao encontrar
\lstinline|@fim|, fixa o endereço de fim em \file{simulacao.c}). Calcula
\lstinline|nbopr| (\lstinline|(n_ins > n_dat) ? ceil(log2(n_ins)) : ceil(log2(n_dat))|)
e abre, com \lstinline|fopen|, as duas imagens de memória em
\path{<proc-dir>/Hardware/} (\file{<prname>\_data.mif} e
\file{<prname>\_inst.mif}). Os arquivos de simulação por porta
(\path{<proc-dir>/Simulation/input_<i>.txt} e \path{output_<i>.txt}) são abertos
mais tarde pelo \emph{testbench} gerado, via \lstinline|$fopen| no
\file{<prname>\_tb.v} (\file{hdl.c}).

\subsection{Geração das memórias \file{.mif}}
\label{sec:yanc-asmcomp-mif}

São duas memórias, gravadas como \file{.mif} (uma palavra binária por linha) em
\path{<proc-dir>/Hardware/}:

\begin{description}[leftmargin=1.2em]
  \item[\file{<prname>\_data.mif}] memória de dados. Cada variável vista pela
    primeira vez é registrada (\lstinline|instr_ula|) e uma célula é emitida:
    floats são inicializados com zero pela conversão \enquote{my float}
    (\lstinline|f2mf("0.0", ...)|), inteiros e constantes recebem seu valor
    (\lstinline|var_val|). Arrays e a pseudo-instrução \lstinline|LEA| também
    materializam suas células aqui, em ordem de primeiro aparecimento.
  \item[\file{<prname>\_inst.mif}] memória de instruções. Cada instrução vira a
    concatenação do índice de opcode em \lstinline|NBITS_OPC| bits com o operando
    em \lstinline|nbopr| bits (\lstinline|itob(opc_idx, NBITS_OPC)| seguido de
    \lstinline|itob(operando, nbopr)|). O operando é o endereço de dados
    (\lstinline|var_find|), o endereço de instrução do salto
    (\lstinline|lab_find|), o índice da porta de E/S, ou o deslocamento.
\end{description}

A conversão de inteiro para string binária é \lstinline|itob| (em
\file{ASMComp/Sources/t2t.c}); a conversão de float IEEE-754 para o formato de
ponto flutuante interno (\enquote{my float}, parametrizado por \lstinline|nbmant|
e \lstinline|nbexpo|) é \lstinline|f2mf| no mesmo arquivo. Ao final
(\lstinline|eval_finish|), há a checagem de consistência
\lstinline|nubits == nbmant + nbexpo + 1|.

\subsection{Geração do Verilog e do \emph{testbench}}
\label{sec:yanc-asmcomp-hdl}

Encerrado o léxico, \lstinline|eval_finish| chama \lstinline|hdl_vv_file| e
\lstinline|hdl_tb_file| (em \file{hdl.c}):

\begin{description}[leftmargin=1.2em]
  \item[\file{<prname>.v}] o Verilog de topo do processador: instancia o núcleo
    \lstinline|p_<prname>| com clock, reset, portas de E/S, o pino de interrupção
    e, quando há \lstinline|#TOAQUI|, o pino \lstinline|cheguei|. Inclui também o
    bloco de visibilidade de simulação (sinais \lstinline|valr2| e
    \lstinline|linetabs|) descrito na \autoref{sec:yanc-sidetable}.
  \item[\file{<prname>\_tb.v}] o \emph{testbench} gerado, escrito no diretório
    temporário. Ele gera clock e reset, instancia o processador, lista os sinais
    a despejar (\lstinline|$dumpvars| de clock, reset, portas de E/S,
    \lstinline|valr2|, \lstinline|linetabs| e variáveis/arrays do usuário),
    suporta um modo \lstinline|+HEADER_ONLY| (que despeja apenas o cabeçalho do
    VCD, para acelerar a coleta de sinais por \tool{gen\_gtkw}) e termina a
    simulação (\lstinline|$finish|) ao atingir o endereço \lstinline|@fim|.
\end{description}

\subsection{A tabela de simulação e o tradutor de opcodes}
\label{sec:yanc-asmcomp-sim}

Em paralelo às memórias, o \tool{asmcomp} grava \file{trad_opcode.txt} (em
\file{simulacao.c}, via \lstinline|sim_add|): uma linha por instrução, no formato
\lstinline|<idx> <opcode> <operando>|. Esse arquivo é o que o GTKWave usa, via
filtro de arquivo, para mostrar o assembly legível na trilha \enquote{Assembly}.
Variáveis e arrays declarados pelo usuário no \file{.cmm} (descobertos consultando
\file{cmm_log.txt} em \lstinline|sim_is_var|/\lstinline|sim_is_arr|) são
registrados com nomes mangleados (p.\,ex.\ \lstinline|me1_f_global_v_x_e_|) para
aparecerem corretamente no waveform.

\section{A \emph{side-table} de linhas e o trace em lockstep}
\label{sec:yanc-sidetable}

O recurso central que liga a execução em hardware ao código-fonte é a
\emph{side-table} de linhas: para cada valor do contador de programa (PC), ela
mapeia a instrução de volta à linha-fonte de alto nível que a originou. Esse mapa
é a base da visualização \emph{em lockstep} no GTKWave, onde se acompanha
simultaneamente a instrução de assembly e a linha de \cmm{}/C++ que o processador
está executando.

\subsection{O arquivo \file{pc\_<name>\_mem.txt}}
\label{sec:yanc-sidetable-pc}

O \tool{cmmcomp} grava esse arquivo durante a emissão (em \file{global.c}).
A cada instrução normal adicionada (\lstinline|add_instr|), ele escreve uma linha
com a linha-fonte corrente codificada em 20 bits binários
(\lstinline|itob(emit_line, 20)|). Para construções de \emph{scaffolding}
(\lstinline|add_sinst|), usa códigos negativos sentinela: \lstinline|-1|
(\lstinline|INTERNAL|), \lstinline|-2| (\lstinline|void main();|) e \lstinline|-3|
(\lstinline|END|). Assim, o arquivo tem uma entrada por instrução, na mesma ordem
da memória de instruções, e cada entrada é o número da linha-fonte (ou um
sentinela). A variável \lstinline|emit_line| é mantida pelo percurso da AST e lida
por \lstinline|add_instr|, de modo que cada instrução carrega a linha do nó que a
gerou. O caminho C++ produz exatamente o mesmo arquivo, com os mesmos sentinelas,
via \lstinline|cg_line| / \file{codegen.c}.

\subsection{O arquivo \file{trad_cmm.txt}}
\label{sec:yanc-sidetable-trad}

O complemento é \file{trad_cmm.txt}, que mapeia número de linha para o
\emph{texto} daquela linha. Os três primeiros registros são os códigos negativos
de \emph{scaffolding} (\lstinline|-1 INTERNAL|, \lstinline|-2 void main();|,
\lstinline|-3 END|); em seguida, cada linha-fonte numerada a partir de 1. Antes de
gravar, o \tool{cmmcomp} substitui os caracteres de Dirac em UTF-8 por
\lstinline|<| e \lstinline|>| (\lstinline|substituir_braket| em \file{global.c})
para que sejam exibíveis no GTKWave.

\subsection{Como o hardware consome a tabela}
\label{sec:yanc-sidetable-hw}

O \tool{asmcomp} embute no \file{<prname>.v} a leitura dessa tabela
(\file{hdl.c}). Uma memória \lstinline|reg [19:0] min| é carregada no início com
\lstinline|$readmemb("pc_<prname>_mem.txt", min)|. A cada \lstinline|posedge clk|,
o índice do PC (\lstinline|pc_sim_val|) indexa essa memória produzindo
\lstinline|linetab|, registrado em \lstinline|linetabs| (sinal \emph{traceado}). O
valor do PC propaga por uma cadeia de dez registradores (\lstinline|valr1| a
\lstinline|valr10|), dos quais \lstinline|valr2| é o estágio exibido (a trilha
\enquote{Assembly}). No GTKWave, o filtro de arquivo \file{trad_opcode.txt}
traduz \lstinline|valr2| em mnemônico, e o filtro \file{trad_cmm.txt} traduz
\lstinline|linetabs| na linha-fonte: as duas trilhas avançam em \emph{lockstep}
com a execução.

\section{Os auxiliares de GTKWave: \tool{comp2gtkw} e \tool{gen\_gtkw}}
\label{sec:yanc-gtkw}

\subsection{\tool{comp2gtkw}: tradutor de números complexos}
\label{sec:yanc-comp2gtkw}

O \tool{comp2gtkw} (\file{Scripts/comp2gtkw.c}) é um filtro de processo: lê de
\lstinline|stdin| o padrão de bits de um sinal complexo e escreve em
\lstinline|stdout| a forma legível \lstinline|<real> <imag>i|. O formato de
entrada é autodescritivo: os primeiros 16 bits codificam \lstinline|nbm| (8 bits)
e \lstinline|nbe| (8 bits); os \lstinline|nbits = nbm + nbe + 1| bits seguintes
são a parte real, e os próximos \lstinline|nbits| são a parte imaginária. Cada
metade é decodificada do formato \enquote{my float} (\lstinline|b2mf|) para
\lstinline|float| e impressa com três casas. O GTKWave invoca esse executável como
filtro de processo nos sinais do tipo \lstinline|comp| (ver
\autoref{sec:yanc-gen-gtkw}), exibindo o número complexo em vez do binário cru.

\subsection{\tool{gen\_gtkw}: gerador do \file{.gtkw}}
\label{sec:yanc-gen-gtkw}

O \tool{gen\_gtkw} (\file{Scripts/gen\_gtkw.c}) lê o \emph{cabeçalho} de um
\file{.vcd} e escreve um arquivo de salvamento \file{.gtkw} já formatado, com
todos os sinais agrupados, coloridos, com aliases legíveis e com os filtros de
tradução conectados. Substitui o antigo formatador em Tcl em tempo de execução: o
GTKWave v4 usado no fluxo nipscern ignora \lstinline|--script|, então a formatação
precisa viver no próprio arquivo de salvamento, aberto com
\lstinline|gtkwave <vcd> -a <out.gtkw>|. A invocação é:

\begin{lstlisting}[language=bash]
gen_gtkw <in.vcd> <out.gtkw> <tmp_base> <comp2gtkw_exe>
\end{lstlisting}

O programa faz \emph{parsing} do cabeçalho do VCD (escopos \lstinline|$scope|,
sinais \lstinline|$var|, parando em \lstinline|$enddefinitions| para não ler um
VCD de gigabytes inteiro), detecta cada instância de processador (todo escopo que
possui ao mesmo tempo um sinal \lstinline|valr2| e um \lstinline|linetabs|), e
emite uma seção formatada por instância. Funciona tanto para um quanto para
vários processadores, sendo o caso de um único processador apenas o $N = 1$.
Para cada processador, conecta a trilha
\enquote{Assembly} (\lstinline|valr2|) ao filtro de arquivo
\path{<tmp_base>/<tipo>/trad_opcode.txt}, a trilha \enquote{C+-}
(\lstinline|linetabs|) ao filtro \path{<tmp_base>/<tipo>/trad_cmm.txt}, e os
sinais \lstinline|comp_me3_| ao filtro de processo \tool{comp2gtkw}. A lógica de
formatação é portada de \file{js/wave/gtkw_proc_writer.js} da \tool{AURORA}; os
\emph{flags} de formato (\lstinline|@<hex>|) foram calibrados contra um arquivo de
salvamento real do GTKWave.

\section{Os scripts de orquestração}
\label{sec:yanc-scripts}

Os scripts em \path{Scripts/} demonstram e automatizam o pipeline ponta a ponta.
Os exemplos POSIX (\file{.sh}) têm contrapartidas Windows (\file{.bat}).

\subsection{\file{env}: resolução do ambiente}
\label{sec:yanc-env}

\file{env.sh} (e \file{env.bat}) é carregado por \lstinline|source| pelos demais
scripts. Resolve a raiz do repositório a partir da própria localização do arquivo,
define \lstinline|YANC_BIN = <repo>/bin| e localiza as ferramentas de simulação
(\tool{iverilog}, \tool{vvp}, \tool{verilator}, \tool{gtkwave}, \tool{fst2vcd}).
Primeiro carrega o cache \file{Scripts/tools.local.sh} (gravado por \file{setup};
ignorado pelo git, específico da máquina) e, para o que faltar, recorre ao
\lstinline|PATH|. Nenhum caminho é codificado fixamente.

\subsection{\file{setup}: preparação única}
\label{sec:yanc-setup}

\file{setup.sh}/\file{setup.bat} preparam o ambiente uma vez após o clone:
localizam o \emph{toolchain} de build (gcc/bison/flex; no Windows,
preferencialmente \tool{x86\_64-w64-mingw32-gcc}), produzem os binários em
\path{bin/} (mantendo os existentes, ou compilando da fonte, ou baixando os
pré-compilados do release mais recente), localizam os simuladores e o GTKWave, e
gravam todos os caminhos resolvidos em \file{tools.local.sh}. Aceitam
\lstinline|--rebuild| (recompilar) e \lstinline|--download| (baixar
pré-compilados).

\subsection{\file{single\_proc}, \file{single\_proc\_cpp} e \file{multi\_proc}}
\label{sec:yanc-runners}

\begin{description}[leftmargin=1.2em]
  \item[\file{single\_proc.sh}] roda o pipeline \cmm{} completo de um processador:
    \tool{cmmcomp} $\rightarrow$ \tool{appcomp} $\rightarrow$ \tool{asmcomp}
    $\rightarrow$ simulação $\rightarrow$ GTKWave. O exemplo embutido é
    \lstinline|proc_fft| (uma FFT de ordem 8 usando endereçamento por bits
    invertidos do \cmm{}).
  \item[\file{single\_proc\_cpp.sh}] o equivalente para o front end C++:
    \tool{cpppp} $\rightarrow$ \tool{cppcomp} $\rightarrow$ \tool{appcomp}
    $\rightarrow$ \tool{asmcomp} $\rightarrow$ simulação $\rightarrow$ GTKWave.
    Aceita também \lstinline|--no-view| (roda o pipeline sem abrir o GTKWave).
  \item[\file{multi\_proc.sh}] demonstra um circuito de topo com dois
    processadores \tool{SAPHO} internos (\lstinline|ZeroCross| e
    \lstinline|ProcDTW|). Cada processador passa pelo pipeline de compilação até
    seu Verilog; o \emph{testbench} é o \file{top_level_tb.v} escrito pelo
    usuário (o \emph{testbench} auto-gerado só conduz um processador isolado).
\end{description}

Todos os \emph{runners} escrevem apenas em \path{Teste/} (ignorado pelo git) e
leem HDL, macros e binários direto do repositório.

\subsection{O flag \texttt{-{}-sim iverilog|verilator}}
\label{sec:yanc-simflag}

Os \emph{runners} aceitam \lstinline|--sim iverilog| (padrão) ou
\lstinline|--sim verilator|, selecionando o simulador:

\begin{description}[leftmargin=1.2em]
  \item[Icarus (\tool{iverilog} + \tool{vvp})] compila o \emph{testbench}
    auto-gerado (\file{<proc>_tb.v}) junto com o \file{<proc>.v} e os arquivos de
    HDL (\file{addr_dec.v}, \file{instr_dec.v}, \file{processor.v}, \file{core.v},
    \file{ula.v}). Roda primeiro com \lstinline|+HEADER_ONLY| (VCD de cabeçalho,
    para \tool{gen\_gtkw}) e depois a simulação real com saída FST.
  \item[Verilator] reusa o \file{<proc>_tb.v} como topo, com
    \lstinline|--binary --timing --trace +define+YANC_TRACE| (a flag
    \lstinline|YANC_TRACE| habilita o \emph{harness} de visibilidade) e vários
    \lstinline|-Wno-...| para tolerar o estilo do Verilog gerado.
\end{description}

Em ambos os casos, o passo final é idêntico: \tool{gen\_gtkw} gera o layout
\file{.gtkw} a partir do cabeçalho do VCD e o GTKWave abre o waveform com esse
layout.

\subsection{\file{regress.sh}: suíte de regressão}
\label{sec:yanc-regress}

\file{regress.sh} é a suíte de regressão única do repositório. Constrói os cinco
binários de compilação uma vez e roda três fases sobre eles: a fase \cmm{} (para
cada teste em \path{Compilers/CMMComp/Tests/}, roda \tool{cmmcomp} e compara o
\file{.asm} com um \file{golden.asm}; para o subconjunto com simulação, roda
também \tool{appcomp} + \tool{asmcomp} + iverilog + vvp e compara as saídas); a
fase C++ (para cada teste em \path{Compilers/CPPComp/Tests/}, roda o pipeline C++
completo e compara contra um \file{golden.txt}); e a fase negativa \cmm{} (para
cada fixture em \path{Compilers/CMMComp/NegTests/}, assegura que \tool{cmmcomp}
rejeita o programa malformado com saída não-zero limpa e o diagnóstico esperado).
Aceita \lstinline|--update| (regenerar goldens), \lstinline|--no-sim|,
\lstinline|--cmm-only| e \lstinline|--cpp-only|, entre outros.

\section{Diagrama do pipeline}
\label{sec:yanc-diagrama}

O fluxo completo, dos dois front ends ao GTKWave, é o seguinte:

\begin{lstlisting}
  C+- (.cmm)                       C++ (.cpp)
     |                                |
     |                          [ cpppp ]   (pre-processador: #include/#define/#ifdef)
     |                                |
     |                          C++ pre-proc. (pp.cpp)
     |                                |
 [ cmmcomp ]                     [ cppcomp ]
     |   \                            |   \
     |    +--> cmm_log.txt            |    +--> cmm_log.txt
     |    +--> pc_<n>_mem.txt         |    +--> pc_<n>_mem.txt   (side-table de linhas)
     |    +--> trad_cmm.txt           |    +--> trad_cmm.txt     (linha -> texto)
     v                                v
  <name>.asm  <---- (mesmo formato intermediario) ---->  <prname>.asm
     |
     v
 [ appcomp ]  (1a passada)
     |
     +--> app_log.txt   (prname/nubits/nbmant/nbexpo, rotulos@->addr,
     |                    variaveis, itr_addr, toaqui_addr, n_ins, n_dat)
     v
 [ asmcomp ]  (2a passada)
     |
     +--> Hardware/<prname>_inst.mif   (memoria de instrucoes: opcode|operando)
     +--> Hardware/<prname>_data.mif   (memoria de dados)
     +--> Hardware/<prname>.v          (Verilog sintetizavel + $readmemb da side-table)
     +--> Temp/<prname>_tb.v           (testbench)
     +--> Temp/trad_opcode.txt         (idx -> mnemonico, para a trilha Assembly)
     v
 [ iverilog+vvp  |  verilator ]   (--sim iverilog|verilator)
     |
     +--> <proc>_tb.vcd / .fst   (waveform)
     +--> <proc>_tb_hdr.vcd      (apenas cabecalho, via +HEADER_ONLY)
     v
 [ gen_gtkw ]   (le o cabecalho do VCD -> escreve .gtkw formatado)
     |           filtros: trad_opcode.txt (Assembly), trad_cmm.txt (C+-),
     |           comp2gtkw (sinais complexos)
     v
 [ GTKWave ]   trilhas Assembly + C+- avancando em lockstep com a execucao
\end{lstlisting}

\begin{nota}
O \file{.asm} é o ponto de convergência: os dois front ends (\cmm{} e C++)
produzem exatamente o mesmo formato intermediário e as mesmas tabelas auxiliares,
de modo que todo o back end (\tool{appcomp}, \tool{asmcomp}, \tool{gen\_gtkw}) é
compartilhado. A linguagem \cmm{} em si é detalhada no
\autoref{cap:03-linguagem-cmm}; o RTL gerado e o conjunto de instruções do
processador, no \autoref{cap:02-arquitetura-soft-processor-sapho}.
\end{nota}

% !TEX root = ../main.tex
\chapter{AURORA IDE: arquitetura geral}
\label{cap:05-aurora-arquitetura}

A \tool{AURORA} é a IDE desktop da plataforma \tool{SAPHO}: uma aplicação
\href{https://www.electronjs.org/}{Electron} na qual o usuário projeta
soft-processors, edita código em \cmm, Verilog e SystemVerilog, compila pela
toolchain \tool{YANC} (\autoref{cap:07-yanc-pipeline}), simula, visualiza o RTL
no \tool{PRISM} e conversa com o subsistema Aurora Intelligence
(\autoref{cap:12-aurora-intelligence}). Este capítulo documenta a
\emph{arquitetura do processo Electron}: a separação entre processo principal
(\emph{main}) e processo de renderização (\emph{renderer}), as janelas criadas,
a topologia de IPC (\emph{inter-process communication}) e os módulos de
infraestrutura do \emph{main} (estado, ciclo de vida, caminhos, rastreio de
processos filhos, log, carregamento do bundle do renderer).

O escopo aqui é o processo \emph{main} e a fronteira IPC. A interface em si
(componentes \href{https://lit.dev/}{Lit}, editor \tool{Monaco}, terminal,
painéis) é tratada em \autoref{cap:06-aurora-renderer}; a toolchain, o pipeline
de compilação e o \tool{PRISM} têm capítulos próprios. Quando um assunto pertence
a outro capítulo, ele é apenas mencionado de passagem.

\begin{nota}
Toda afirmação deste capítulo foi verificada lendo o código-fonte em
\path{C:/Users/chrys/Documents/GitHub/aurora}. Os arquivos centrais são
\file{main.js} (ponto de entrada), \file{main/windows.js} (fábricas de janelas),
\file{main/lifecycle.js} (ciclo de vida), os módulos \file{main/ipc/*.js}
(handlers IPC) e os preloads em \path{js/app/preload*.js}.
\end{nota}

\section{Modelo de dois processos do Electron}
\label{sec:05-dois-processos}

O Electron herda do Chromium um modelo multiprocessos. A \tool{AURORA} usa duas
classes de processo:

\begin{description}[leftmargin=*]
  \item[Processo \emph{main} (Node.js)] É o processo privilegiado. Tem acesso a
    Node, ao sistema de arquivos, ao \lstinline|child_process| (spawn de
    compiladores/simuladores), à rede e às APIs nativas do Electron
    (\lstinline|app|, \lstinline|BrowserWindow|, \lstinline|dialog|,
    \lstinline|session|, \lstinline|safeStorage|). Existe exatamente um por
    instância da aplicação. Todo o código sob \file{main/} roda aqui.
  \item[Processo \emph{renderer} (Chromium)] É a página web que desenha a UI.
    Roda \emph{sandboxed}, sem acesso direto a Node nem ao filesystem. Cada
    \lstinline|BrowserWindow| tem o seu próprio renderer. Para alcançar qualquer
    capacidade do \emph{main}, o renderer chama uma API curada exposta por um
    \emph{preload} via \lstinline|contextBridge|.
\end{description}

\subsection{O padrão central: ``renderer orquestra, main executa''}
\label{sec:05-padrao-central}

A linha divisória de toda a \tool{AURORA} é: \emph{o renderer decide o que fazer;
o main faz o que pode causar efeito colateral}. O renderer monta a intenção
(abrir projeto, ler arquivo, compilar uma etapa, fazer commit, gerar um SVG) e a
envia por um canal IPC; o \emph{main} valida a entrada, executa a operação
privilegiada (filesystem, \lstinline|spawn|, rede, keychain) e devolve o
resultado. Nenhuma capacidade de Node, \lstinline|fs| ou \lstinline|spawn| vive
no renderer --- todas estão no \emph{main}, atrás de um canal IPC.

Esse padrão aparece de forma explícita no executor de compilação. O renderer
constrói um \lstinline|CommandSpec| estruturado (\lstinline|{binary, args}|) e o
manda por \lstinline|exec-spec|; o \emph{main} valida o binário contra uma
allowlist da toolchain, confere \emph{protected flags} e só então dispara
\lstinline|spawn(binary, args, { shell: false })|. Como descrito no cabeçalho de
\file{main/compile/executor.js}, com \lstinline|shell:false| cada argumento vai
para o filho exatamente como está --- não há parsing de shell, portanto não há
bug de \emph{quoting} nem vetor de injeção. É exatamente essa propriedade que
permite que o sistema de \emph{command override} da Aurora Intelligence exista
com segurança (\autoref{cap:12-aurora-intelligence}).

\section{O ponto de entrada: \file{main.js}}
\label{sec:05-mainjs}

O arquivo \file{main.js} é deliberadamente magro: toda a lógica real vive em
\file{main/}, e este arquivo apenas \emph{conecta os módulos} numa ordem de boot
determinística. Suas responsabilidades, em ordem de execução, são:

\begin{enumerate}[leftmargin=*]
  \item \textbf{Configurar o logger antes de tudo.} A primeira instrução após os
    \lstinline|require| é \lstinline|configureLogger()| (de \file{main/logger.js}),
    para que toda chamada de log subsequente já use a configuração definitiva.
  \item \textbf{Rede de segurança do processo main.} Inicia o
    \lstinline|crashReporter| em modo local (\lstinline|uploadToServer: false|,
    coletando minidumps) e registra \lstinline|process.on('uncaughtException')|
    e \lstinline|process.on('unhandledRejection')| para que nada que escape de um
    callback derrube o processo silenciosamente. Tudo \emph{best-effort}: a rede
    de segurança nunca pode ser o que bloqueia o boot.
  \item \textbf{Ajuste de GPU / compositor.} As \emph{command-line switches} têm
    de ser anexadas \emph{antes} de o app ficar pronto, pois o processo de GPU as
    lê no lançamento. Veja \autoref{sec:05-gpu}.
  \item \textbf{AppUserModelID (Windows).} Chama
    \lstinline|app.setAppUserModelId('com.nipscern.sapho')| o mais cedo possível
    --- antes de existir qualquer \lstinline|BrowserWindow| --- para que o Windows
    associe o processo a uma identidade estável de \emph{jumplist}/taskbar.
  \item \textbf{Registro de IPC e ciclo de vida.} Adquire o \emph{single-instance
    lock} via \lstinline|lifecycle.register()| e, se obtido, registra todos os
    módulos de IPC (\autoref{sec:05-ipc}).
  \item \textbf{\lstinline|app.whenReady()|.} Quando o Electron está pronto:
    cria o diretório \path{/tmp} do MSYS, roda a higiene de temporários
    (\file{main/temp_gc.js}), reconstrói a jumplist, instala a
    \emph{Content-Security-Policy} (\autoref{sec:05-csp}) e finalmente chama
    \lstinline|windows.createSplashScreen()| --- a splash agenda a criação da
    janela principal.
\end{enumerate}

\subsection{Switches de GPU}
\label{sec:05-gpu}

Antes de \lstinline|app.whenReady|, \file{main.js} anexa duas switches:

\begin{lstlisting}[language=Java]
app.commandLine.appendSwitch('enable-gpu-rasterization');
app.commandLine.appendSwitch('enable-zero-copy');
\end{lstlisting}

A combinação rasteriza os \emph{tiles} na GPU e os carrega sem cópia pela CPU ---
um ganho para painéis de muito \emph{scroll}/\emph{paint} (terminal, editor,
configuração de waveform). O código documenta explicitamente que \emph{não} se
usa \lstinline|disable-frame-rate-limit| (removeria o \emph{vsync} do Chromium e
poderia saturar o processo de GPU durante os \emph{loops} de
\lstinline|requestAnimationFrame| da splash). A aceleração de hardware é deixada
\textbf{ligada} --- \lstinline|app.disableHardwareAcceleration()| nunca é chamado
--- e o \lstinline|backgroundThrottling| fica no padrão.

\subsection{Content-Security-Policy}
\label{sec:05-csp}

Dentro de \lstinline|app.whenReady|, \file{main.js} instala uma CSP via
\lstinline|session.defaultSession.webRequest.onHeadersReceived|, entregue como
\emph{header} de resposta para que cubra tanto o carregamento empacotado
(\lstinline|file://|) quanto o servidor de desenvolvimento. O \emph{handler}
\emph{substitui} (não acrescenta) qualquer CSP que um servidor de dev tenha
definido, deixando a da \tool{AURORA} como autoritativa. As diretivas incluem,
entre outras, \lstinline|default-src 'self'|, \lstinline|object-src 'none'|,
\lstinline|frame-ancestors 'none'| e um \lstinline|connect-src| que, em produção,
só permite a origem própria mais a detecção local do \tool{Ollama}
(\lstinline|localhost:11434|) --- os provedores de IA em nuvem e o \tool{MCP}
rodam no \emph{main}, não no renderer.

\section{Janelas: \file{main/windows.js}}
\label{sec:05-janelas}

As fábricas de janela vivem em \file{main/windows.js}, com exceção da janela do
\tool{PRISM}, que mora em \file{main/ipc/prism.js} por carregar muita lógica de
compilação. Ao todo a \tool{AURORA} cria \textbf{cinco} tipos de
\lstinline|BrowserWindow|: principal, splash, atualização, \tool{PRISM} e o
\emph{Design Lab}. A tabela abaixo resume cada uma e suas opções de segurança ---
em todas elas \lstinline|contextIsolation| está ligado, \lstinline|sandbox| está
ligado e \lstinline|nodeIntegration| está desligado.

\begin{longtable}{@{}p{0.16\linewidth} p{0.30\linewidth} p{0.46\linewidth}@{}}
  \toprule
  \textbf{Janela} & \textbf{Preload dedicado} & \textbf{Características} \\
  \midrule
  \endhead
  Principal & \file{js/app/preload.js} &
    \lstinline|frame:false|, \lstinline|thickFrame:true|, barra de título
    customizada; 1280\texttimes820 (mín.\ 720\texttimes480); carrega
    \path{index.html}. Bloqueio de navegação: \lstinline|setWindowOpenHandler|
    nega \lstinline|window.open| e \lstinline|will-navigate| é cancelado.
    Múltiplas janelas principais podem coexistir no mesmo processo. \\
  Splash & \file{js/app/preload_splash.js} &
    720\texttimes480, \lstinline|frame:false|, \lstinline|transparent:true|,
    \lstinline|alwaysOnTop|, \lstinline|skipTaskbar|; carrega
    \path{html/splash.html}. Exibe progresso \emph{real} de boot. \\
  Atualização & \file{js/app/preload_update.js} &
    540\texttimes660, \lstinline|frame:false|, \lstinline|transparent:true|,
    \lstinline|alwaysOnTop|; janela única que percorre os estados
    \emph{available} / \emph{downloading} / \emph{downloaded}
    (\autoref{cap:13-updater-build}). Não fecha durante um download em curso. \\
  \tool{PRISM} & \file{js/app/preload_prism.js} &
    1400\texttimes900 (mín.\ 1000\texttimes700), \lstinline|frame:false|,
    \lstinline|thickFrame:true|; carrega \path{html/prism/prism.html}.
    Criada em \file{main/ipc/prism.js} (\autoref{cap:08-prism-rtl}). \\
  Design Lab & (sem preload) &
    1100\texttimes820, janela emoldurada (controles nativos do SO), singleton.
    Galeria interna de componentes, aberta pela paleta de comandos; é uma página
    estática que não fala com nada. \\
  \bottomrule
  \caption{As cinco janelas da \tool{AURORA} e suas opções principais.}
\end{longtable}

\subsection{Janela principal: \lstinline|createMainWindow|}
\label{sec:05-main-window}

A janela principal é criada sem moldura nativa (\lstinline|frame:false|) porque a
\tool{AURORA} desenha a própria barra de título; \lstinline|thickFrame:true|
preserva o \emph{Aero snap}, o redimensionamento de borda e as animações em todas
as versões do Windows. As \lstinline|webPreferences| materializam o modelo de
segurança: \lstinline|contextIsolation:true|, \lstinline|sandbox:true|,
\lstinline|nodeIntegration:false|, \lstinline|nodeIntegrationInSubFrames:false|,
\lstinline|webviewTag:false|, \lstinline|enableWebSQL:false|,
\lstinline|allowRunningInsecureContent:false| e
\lstinline|experimentalFeatures:false|. O único caminho do renderer para o
\emph{main} é o \lstinline|electronAPI| exposto pelo preload.

A função também:

\begin{itemize}[leftmargin=*]
  \item registra o protocolo \lstinline|sapho:| e a extensão \path{.spf} no
    Windows e reconstrói a \emph{jumplist} a cada criação de janela;
  \item notifica o renderer sobre maximizar/restaurar/fullscreen (evento
    \lstinline|window-state|) para que o ícone da barra de título acompanhe;
  \item ao detectar um \path{.spf} passado na linha de comando, envia
    \lstinline|open-spf-file| ao renderer após \lstinline|did-finish-load|;
  \item no \lstinline|close| da \emph{última} janela principal, dispara
    \lstinline|stopAllToolchain()| e fecha a janela auxiliar do \tool{PRISM},
    garantindo que nenhum compile/simulação/síntese sobreviva ao fechamento da
    IDE.
\end{itemize}

\subsection{Splash com progresso real}
\label{sec:05-splash}

A splash não é um \emph{timer} fixo: a barra avança conforme a janela principal
cruza marcos genuínos de boot. O fluxo, codificado em
\lstinline|createSplashScreen|, é:

\begin{lstlisting}[language=bash]
splash carrega (did-finish-load)         -> progresso 10%  "boot"
  cria janela principal OCULTA (deferShow)
                                         -> progresso 24%  "window"
                                         -> progresso 36%  "loading"
  webContents da principal:
    dom-ready                            -> progresso 56%  "dom"
    did-finish-load                      -> progresso 80%  "resources"
    ready-to-show                        -> progresso 90%  "editor"
  renderer sinaliza app:renderer-ready   -> progresso 96%  -> handoff 100%
  splash responde splash:filled          -> espera 2 s     -> reveal
\end{lstlisting}

O \lstinline|reveal| maximiza e mostra a janela principal e fecha a splash de
imediato (sem \emph{fade}, para não haver \emph{flash} de desktop vazio entre
elas). Um teto absoluto de \textbf{15\,s} força a transferência caso algum marco
nunca dispare. O coordenador usa \lstinline|ipcMain.once| para os sinais
\lstinline|splash:filled| e \lstinline|app:renderer-ready|, de modo que uma
segunda janela criada depois não re-dispare o \emph{handoff}.

\subsection{Controles de janela roteados pelo emissor}
\label{sec:05-window-controls}

Como vários janelas principais podem coexistir no mesmo processo (ver
\autoref{sec:05-lifecycle}), \lstinline|registerWindowControls| roteia cada IPC
de janela para \emph{a janela que o enviou}, via
\lstinline|BrowserWindow.fromWebContents(event.sender)|, e não para o singleton
\lstinline|state.mainWindow|. Os canais são: \lstinline|window:minimize|,
\lstinline|window:maximize-toggle|, \lstinline|window:close| e
\lstinline|window:get-state| (mais zoom: \lstinline|zoom-in|,
\lstinline|zoom-out|, \lstinline|zoom-reset|, e \lstinline|open-design-lab|).

\section{Ciclo de vida da aplicação: \file{main/lifecycle.js}}
\label{sec:05-lifecycle}

O módulo \file{main/lifecycle.js} concentra o ciclo de vida do app via a função
\lstinline|register()|, registrada por \file{main.js} no boot.

\subsection{Single-instance lock e múltiplas janelas}

No início, \lstinline|register()| detecta um \path{.spf} na linha de comando
(\lstinline|process.argv.find(arg => arg.endsWith('.spf'))|) e tenta adquirir o
\emph{single-instance lock} com \lstinline|app.requestSingleInstanceLock()|. A
variável de ambiente \lstinline|SAPHO_SKIP_SINGLE_INSTANCE=1| ignora o lock (uso
exclusivo dos testes). Se o lock não é obtido, a função loga o motivo em
\emph{ambos} o arquivo de log e o \lstinline|stderr| e chama
\lstinline|app.quit()|, retornando \lstinline|false| --- o que faz \file{main.js}
abortar o registro dos handlers num processo que está morrendo.

Quando uma segunda instância é lançada, o evento \lstinline|second-instance|
\textbf{cria uma nova janela} no mesmo processo (em vez de só focar a existente),
permitindo vários projetos lado a lado sem duplicar processos Electron nem correr
risco de corrida sobre a pasta compartilhada \path{components/Temp}. Se essa
segunda invocação trouxe um \path{.spf}, ele é enviado à nova janela.

\subsection{Encerramento: limpeza em duas fases}
\label{sec:05-before-quit}

O \emph{handler} \lstinline|before-quit| é a limpeza \emph{autoritativa} e roda em
duas fases serializadas, para evitar \lstinline|EBUSY| no \lstinline|rmdir| de
\path{Temp}:

\begin{description}[leftmargin=*]
  \item[Fase 1 --- soltar handles.] Fecha todos os \emph{watchers} de arquivo e
    de diretório (chokidar, que no Windows usa \lstinline|ReadDirectoryChangesW|)
    e chama \lstinline|stopAllToolchain()| (\autoref{sec:05-process-registry}),
    que mata todos os processos filhos da toolchain. Também encerra o servidor
    \tool{MCP} da Aurora (a ponte HTTP local). Tudo isso corre em paralelo, com
    um teto de \textbf{5\,s} via \lstinline|Promise.race| para não travar o quit.
  \item[Fase 2 --- apagar \path{Temp}.] Só depois de a fase 1 liberar os handles,
    remove recursivamente \path{components/Temp} (\lstinline|maxRetries:3|) e a
    recria vazia.
\end{description}

Há ainda \lstinline|window-all-closed| (sai do app fora do macOS) e
\lstinline|activate| (recria a janela principal se não houver nenhuma, padrão
macOS).

\section{Caminhos e estado do main}
\label{sec:05-paths-state}

\subsection{\file{main/paths.js}}

Centraliza os caminhos estáticos. \lstinline|appRoot| é
\lstinline|app.getAppPath()| (a raiz do projeto em dev, ou
\path{resources/app.asar} empacotado). O \lstinline|componentsPath| --- onde mora
o toolchain bundle (\autoref{cap:09-toolchain-bundle}) --- depende do modo: em dev
é \path{<appRoot>/components}; empacotado, fica ao lado do executável
(\lstinline|path.dirname(app.getPath('exe'))/components|), pois os componentes são
distribuídos \emph{fora} do asar. O sinalizador \lstinline|isDev| deriva de
\lstinline|process.env.NODE_ENV === 'development'|.

\subsection{\file{main/state.js}}
\label{sec:05-state}

Estado mutável compartilhado do \emph{main}. Todos os módulos importam a
\emph{mesma} referência de objeto, de modo que escritas feitas por um handler
ficam visíveis a todos. Substitui as antigas variáveis \lstinline|let| de escopo
de arquivo. Os campos incluem:

\begin{itemize}[leftmargin=*]
  \item \textbf{Janelas:} \lstinline|mainWindow|, \lstinline|splashWindow|,
    \lstinline|updateWindow|, \lstinline|prismWindow|.
  \item \textbf{Updater:} \lstinline|isQuitting|,
    \lstinline|downloadInProgress|, \lstinline|updateAvailable|,
    \lstinline|updateInfo| etc.
  \item \textbf{Projeto:} \lstinline|currentOpenProjectPath| (a fonte única da
    verdade sobre qual projeto está aberto, lida por \file{git.js} e
    \file{search.js}) e \lstinline|fileToOpen|.
  \item \textbf{Processos de simulação:} \lstinline|currentVvpProcess|,
    \lstinline|vvpProcessPid|, \lstinline|currentGtkwaveProcesses| e
    \lstinline|childProcesses| --- o \emph{Set} central de todo filho da
    toolchain, \emph{tree-killed} ao fechar a janela/sair.
  \item \textbf{Watchers:} \lstinline|activeWatchers|,
    \lstinline|activeDirectoryWatchers| e seus caches de \emph{stats}.
\end{itemize}

\section{Rastreio e \emph{tree-kill} de processos filhos}
\label{sec:05-process-registry}

O módulo \file{main/process_registry.js} é o registro central de todo processo
filho da toolchain que a \tool{AURORA} dispara (compiladores \cmm/ASM,
\tool{iverilog}, \tool{vvp}, \tool{Verilator} com seus \tool{g++}/\tool{make}/
\tool{ccache}, \tool{yosys} para o \tool{PRISM}, \tool{gtkwave}, \tool{cocotb}/
Python) mais a rotina única que força a parada de todos.

\subsection{Registro e ponto de spawn único}

\begin{description}[leftmargin=*]
  \item[\lstinline|trackChild(child)|] adiciona o filho a
    \lstinline|state.childProcesses| e o remove automaticamente em
    \lstinline|exit|/\lstinline|close|/\lstinline|error| --- o \emph{Set} só
    contém processos vivos. Marca também \lstinline|toolchainEverRan = true|.
  \item[\lstinline|spawnTracked(cmd, args, opts)|] é o ponto de \emph{spawn}
    obrigatório da toolchain: equivale a \lstinline|trackChild(spawn(...))|, com a
    mesma assinatura e o mesmo retorno de \lstinline|child_process.spawn|. Spawnar
    por aqui torna automático o registro para \emph{tree-kill}.
\end{description}

Os CLIs de agente de IA (\file{claude_code}, \file{codex_cli})
\emph{deliberadamente} não passam por aqui --- eles gerenciam as próprias árvores
de processo e são derrubados por seus \lstinline|killAll()| em
\lstinline|stopAllToolchain()|.

\subsection{\lstinline|stopAllToolchain()|}

É \emph{best-effort} e idempotente (memoizada numa única \lstinline|stopPromise|),
chamável tanto no \lstinline|close| da janela quanto no \lstinline|before-quit|.
Cobre, em ordem:

\begin{enumerate}[leftmargin=*]
  \item todo filho rastreado, \emph{tree-killed} por PID
    (\lstinline|taskkill /F /T /PID|);
  \item o processo de simulação \emph{legacy} de slot único, caso não estivesse
    rastreado;
  \item \textbf{apenas se algum filho realmente rodou nesta sessão}
    (\lstinline|toolchainEverRan|): varreduras de \emph{backstop} por nome
    (\lstinline|vvp.exe|, \lstinline|gtkwave.exe|) e por prefixo de caminho
    (os \lstinline|V<top>.exe| do Verilator sob \path{components/Temp});
  \item os CLIs de IA (Claude Code / Codex) via \lstinline|killAll()|;
  \item as gerações de chat de IA em curso, via \lstinline|abortAll()| (aborta os
    \emph{streams} HTTP).
\end{enumerate}

A guarda \lstinline|toolchainEverRan| é o que faz uma sessão só de edição fechar
instantaneamente: as varreduras externas (que custam um \lstinline|taskkill| frio
e um \lstinline|Get-CimInstance Win32_Process| via PowerShell) só rodam se houve
provadamente o que matar. As primitivas de \emph{kill} (\lstinline|killProcessSilently|,
\lstinline|killProcessesByName|, \lstinline|killProcessesByPathPrefix|) vivem em
\file{main/utils.js}.

\section{Higiene de temporários e log}
\label{sec:05-tempgc-log}

\subsection{\file{main/temp_gc.js}}

Higiene universal de temporários, \emph{best-effort} e não-bloqueante, rodada uma
vez no startup (enfileirada para fora do caminho crítico via
\lstinline|setImmediate|). \lstinline|runStartupGC()| (1) poda os arquivos
\path{aurora-mcp-<pid>.json} obsoletos que o Claude Code escreve no diretório
temporário do SO e que nunca eram limpos, e (2) limpa as imagens de anexo da IA
(\lstinline|cleanupTempImages(0)|), que são \emph{one-shot} por turno.

\subsection{\file{main/logger.js}}

Configuração central do
\href{https://github.com/megahertz/electron-log}{\tool{electron-log}}, que é um
singleton: como toda chamada \lstinline|require('electron-log')| no código retorna
a mesma instância, basta configurar transportes/níveis/formato uma vez aqui e todo
o \emph{main} herda. O transporte de arquivo fica no nível \lstinline|info|, com
rotação a $\sim$5\,MB; o transporte de console é \lstinline|debug| em dev e
\lstinline|warn| em produção. Os caminhos de log por SO (padrão do
\tool{electron-log}) são:

\begin{tabularx}{\linewidth}{Y Z}
  \toprule
  \textbf{Sistema} & \textbf{Caminho de \path{main.log}} \\
  \midrule
  Windows & \path{%APPDATA%/Aurora-IDE/logs/main.log} \\
  macOS   & \path{~/Library/Logs/Aurora-IDE/main.log} \\
  Linux   & \path{~/.config/Aurora-IDE/logs/main.log} \\
  \bottomrule
\end{tabularx}

\section{Carregando o renderer empacotado pelo Vite}
\label{sec:05-render-loader}

\subsection{\file{main/render_loader.js}}

Toda janela de primeira parte (principal, splash, atualização, \tool{PRISM})
carrega sua página por \lstinline|loadPage(win, relPath)|, centralizando a
seleção dev/prod num só lugar:

\begin{itemize}[leftmargin=*]
  \item \textbf{Dev (Vite):} quando \lstinline|AURORA_RENDERER_URL| está definida
    \emph{e} o app não está empacotado, a página é carregada do servidor de dev do
    Vite (com HMR), anexando o caminho relativo à origem do servidor.
  \item \textbf{Produção:} carrega a página empacotada de \path{dist/} via
    \lstinline|file://|.
\end{itemize}

Não há \emph{fallback} para a fonte crua: após a migração para \tool{Lit}, a
\path{index.html} crua carrega \lstinline|import|s de \lstinline|lit| que só
resolvem pelo \emph{bundler}, então um \path{dist/} ausente é um erro de build
--- e é exibido explicitamente como uma página de erro, não como uma UI
silenciosamente quebrada.

\subsection{\file{scripts/launch-electron.js}}

O \lstinline|npm start| não chama \lstinline|electron .| diretamente, e sim
\file{scripts/launch-electron.js}. O motivo: alguns shells pais (o terminal
integrado do VS Code, o Claude Code) exportam
\lstinline|ELECTRON_RUN_AS_NODE=1|; herdada por \lstinline|electron .|, ela faria
o Electron bootar como Node puro (\lstinline|app| viraria \lstinline|undefined| e
o app quebraria em \lstinline|app.getAppPath()| dentro de \file{main/paths.js}). O
script \emph{deleta} essa variável do ambiente e então faz
\lstinline|spawn(electron, ['.', ...argv])|.

\section{Topologia de IPC}
\label{sec:05-ipc}

A fronteira IPC é o coração da arquitetura: é por ela que o renderer alcança o
\emph{main}. No total existem \textbf{157} registros de \lstinline|ipcMain| no
processo principal (contagem por \lstinline|grep| de
\lstinline|ipcMain.handle/on/once| em \file{main/} e \file{main.js}). A maioria
usa \lstinline|ipcMain.handle| (\emph{request/response} via
\lstinline|invoke|/\lstinline|Promise|); alguns usam \lstinline|ipcMain.on|
(sinais unidirecionais, como controles de janela e a conclusão de \emph{tool
call}).

\subsection{Os módulos \file{main/ipc/*}}

\begin{longtable}{@{}p{0.18\linewidth} r p{0.58\linewidth}@{}}
  \toprule
  \textbf{Módulo} & \textbf{Handlers} & \textbf{Papel} \\
  \midrule
  \endhead
  \file{ipc/files.js}    & 27 &
    Sistema de arquivos, diálogos e \emph{watchers} (chokidar):
    \lstinline|read-file|, \lstinline|write-file|, \lstinline|copy-file|,
    \lstinline|mkdir|, \lstinline|file:delete|, \lstinline|watch-file|,
    \lstinline|watch-directory|, \lstinline|dialog:showOpen|,
    \lstinline|show-save-dialog|, \lstinline|create-backup| etc. \\
  \file{ipc/git.js}      & 28 &
    Controle de versão via \href{https://github.com/steveukx/git-js}{simple-git}
    sobre o diretório do projeto aberto: \lstinline|git:status|,
    \lstinline|git:diff|, \lstinline|git:commit|, \lstinline|git:push|,
    \lstinline|git:pull|, \lstinline|git:clone|, \lstinline|git:branches|,
    \lstinline|git:stash| etc. (\autoref{cap:11-git-github}). \\
  \file{ipc/ai.js}       & 22 &
    Aurora Intelligence: gestão de chaves (\lstinline|ai:set-key|,
    \lstinline|ai:get-key-status| --- a chave \emph{nunca} sai do main),
    \lstinline|ai:chat-start|/\lstinline|ai:chat-abort|, \emph{bridges} de
    Claude Code/Codex, manifesto de \emph{tools}, histórico de conversas
    (\autoref{cap:12-aurora-intelligence}). \\
  \file{ipc/project.js}  & 11 &
    Lifecycle de projeto (\lstinline|project:open|, \lstinline|project:close|,
    \lstinline|project:createStructure|) e CRUD de processadores
    (\lstinline|create-processor-project|, \lstinline|delete-processor|,
    \lstinline|rename-processor|); reescreve o \path{.spf} nesses eventos. \\
  \file{ipc/compile.js}  & 7 &
    Execução de compilação/simulação \emph{legacy} e abertura de visualizadores
    de waveform: \lstinline|launch-gtkwave-only|, \lstinline|launch-surfer|,
    \lstinline|cancel-vvp-process|, \lstinline|kill-current-spec-process|,
    \lstinline|check-process-running|, \lstinline|decode-complex|. \\
  \file{ipc/prism.js}    & 6 &
    \tool{PRISM}: dono da janela do visualizador de RTL e de todos os passos de
    síntese (Yosys $\to$ netlistsvg): \lstinline|prism-compile-with-paths|,
    \lstinline|prism-recompile|, \lstinline|generate-svg-from-module|,
    \lstinline|prism:build-digitaljs| (\autoref{cap:08-prism-rtl}). \\
  \file{ipc/github_auth.js} & 7 &
    Conexão de conta GitHub (PAT validado + \emph{Device Flow} OAuth), com o
    token cifrado por \lstinline|safeStorage| (DPAPI no Windows):
    \lstinline|github:connect|, \lstinline|github:status|,
    \lstinline|github:list-repos|, \lstinline|github:oauth-login| etc. \\
  \file{ipc/system.js}   & 6 &
    Utilidades de sistema: \lstinline|get-components-path|,
    \lstinline|toolchain:python-status|, \lstinline|join-path|,
    \lstinline|path-dirname|, \lstinline|app:reload| e
    \lstinline|renderer:error| (a \emph{error boundary} do renderer encaminha
    erros não tratados para o log persistente). \\
  \file{ipc/search.js}   & 2 &
    \emph{Find in Files} no projeto: \lstinline|search:in-project|, uma varredura
    recursiva síncrona limitada por profundidade/tamanho/matches, ancorada em
    \lstinline|state.currentOpenProjectPath|. \\
  \bottomrule
  \caption{Os nove módulos sob \file{main/ipc/} (somam 116 handlers).}
\end{longtable}

\subsection{Outros registradores de IPC}

Além de \file{main/ipc/*}, registram canais: \file{main/compile/executor.js}
(4 --- \lstinline|exec-spec|, \lstinline|exec-spec-streamed|,
\lstinline|exec-spec-allowed-binaries|, \lstinline|exec-spec-protected-flags|),
\file{main/lsp/verible_lsp.js} (9, \lstinline|lsp:*|),
\file{main/lsp/slang_lsp.js} (6, \lstinline|slang:*|),
\file{main/format/clang_format.js} (2, \lstinline|format:*|),
\file{main/treesitter/grammars.js} (2, \lstinline|treesitter:*| --- serve bytes
WASM ao realce semântico), \file{main/updater.js} (8, \lstinline|update*|) e
\file{main/windows.js} (10 --- os oito controles de janela/zoom mais os dois
sinais da splash \lstinline|splash:filled| e \lstinline|app:renderer-ready|,
registrados via \lstinline|ipcMain.once| dentro de \lstinline|createSplashScreen|).
Os LSP, \emph{clang-format} e \emph{tree-sitter} são todos
\emph{best-effort}: se o binário/grammar não estiver instalado, o lado \emph{main}
vira \emph{no-op}.

\section{Os preloads como ponte segura}
\label{sec:05-preloads}

Cada janela tem um preload que executa num contexto isolado e usa
\lstinline|contextBridge.exposeInMainWorld| para expor \emph{apenas} as funções
IPC que aquele renderer realmente consome --- uma allowlist de canais. Nenhum
preload expõe um \lstinline|send|/\lstinline|invoke| genérico, então um renderer
não pode alcançar um canal arbitrário.

\subsection{\file{js/app/preload.js} (janela principal)}

O preload principal expõe \emph{oito} \emph{namespaces} no
\lstinline|window| via \lstinline|contextBridge|:

\begin{tabularx}{\linewidth}{Y Z}
  \toprule
  \textbf{Objeto global} & \textbf{Conteúdo} \\
  \midrule
  \lstinline|electronAPI| & Arquivos, \emph{watchers}, projeto, compilação/
    simulação (\lstinline|execSpec|, \lstinline|execSpecStreamed|),
    \tool{PRISM}, diálogos, UI/janela, terminal, updater, busca,
    \lstinline|getPathForFile| e \lstinline|logRendererError|. \\
  \lstinline|terminalAPI| & Terminal integrado (canais do terminal expostos
    também como \emph{namespace} próprio, além do \emph{spread} em
    \lstinline|electronAPI|). \\
  \lstinline|aiAPI|       & Aurora Intelligence (provedores, chaves, chat,
    \emph{tools}, conversas). \\
  \lstinline|gitAPI|      & Git + integração GitHub. \\
  \lstinline|lspAPI|      & LSP Verilog (Verible). \\
  \lstinline|slangAPI|    & LSP semântico SystemVerilog (Slang). \\
  \lstinline|clangFormatAPI| & Formatação C/C++/\cmm. \\
  \lstinline|treeSitterAPI| & Bytes WASM de \emph{grammars} tree-sitter. \\
  \bottomrule
\end{tabularx}

Cada método é um \emph{wrapper} fino sobre \lstinline|ipcRenderer.invoke| ou
\lstinline|ipcRenderer.send|, mantendo o mapeamento 1:1 com os
\lstinline|ipcMain.handle/on| do \emph{main}. As chaves de API em texto puro
trafegam \emph{só} no sentido renderer $\to$ main: não existe canal
\lstinline|getKey|, pois ler bytes decifrados de volta para o renderer anularia o
propósito de manter o \emph{keystore} no \emph{main}. \lstinline|getPathForFile|
usa \lstinline|webUtils.getPathForFile| para resolver o caminho de um \lstinline|File|
em \emph{drag\&drop} (necessário no Electron 32+, que removeu \lstinline|file.path|).

\subsection{Preloads auxiliares}

\begin{description}[leftmargin=*]
  \item[\file{preload_splash.js}] expõe \lstinline|splashAPI| com apenas
    \lstinline|onProgress|, \lstinline|notifyFilled| e \lstinline|getAppVersion|
    --- a splash não precisa de mais nada.
  \item[\file{preload_update.js}] expõe \lstinline|updateAPI| espelhando 1:1 os
    canais de \file{main/updater.js}: \lstinline|onState|, \lstinline|onProgress|,
    \lstinline|onError|, \lstinline|onShake| (main $\to$ renderer) e
    \lstinline|download|, \lstinline|install|, \lstinline|dismiss|
    (renderer $\to$ main).
  \item[\file{preload_prism.js}] expõe um \lstinline|electronAPI| reduzido com
    cada canal que a janela do \tool{PRISM} usa (\lstinline|readFile|,
    \lstinline|prismCompileWithPaths|, \lstinline|buildDigitalJS|,
    \lstinline|openSourceFile| etc.), e comenta explicitamente que \emph{não} há
    \lstinline|send|/\lstinline|removeAllListeners| genéricos, para preservar a
    allowlist de canais.
\end{description}

\section{Diagrama da topologia}
\label{sec:05-diagrama}

O diagrama abaixo resume processos, janelas e a direção do IPC.

\begin{lstlisting}
+===========================================================================+
|                       PROCESSO MAIN (Node.js)                             |
|                                                                           |
|  main.js  -> lifecycle.register() (single-instance lock, before-quit)     |
|          -> <modulo>.register() para cada IPC abaixo                       |
|          -> windows.createSplashScreen() em app.whenReady                  |
|                                                                           |
|  Estado/infra:  state.js   paths.js   logger.js   temp_gc.js              |
|                 process_registry.js (trackChild / spawnTracked /          |
|                                      stopAllToolchain -> taskkill /T)      |
|                                                                           |
|  IPC (157 ipcMain):                                                        |
|    files  project  compile  git  github_auth  prism  ai  search  system   |
|    + executor(exec-spec)  lsp(verible/slang)  format  treesitter  updater |
|                                                                           |
|    spawn(shell:false) -> YANC / iverilog / vvp / verilator / yosys /      |
|                          gtkwave / cocotb   (todos rastreados)            |
+=================================|=========================================+
            contextBridge (preload)  |  ipcRenderer.invoke / .send  /  eventos
   +--------------+--------------+----+------------+-------------------+
   |              |              |                 |                   |
+--v---------+ +--v--------+ +---v---------+ +-----v--------+ +--------v------+
| RENDERER   | | RENDERER  | | RENDERER    | | RENDERER     | | RENDERER      |
| principal  | | splash    | | update      | | PRISM        | | Design Lab    |
| index.html | | splash.   | | update-     | | prism.html   | | design-lab.   |
|            | | html      | | notif.html  | |              | | html          |
| preload.js | | preload_  | | preload_    | | preload_     | | (sem preload) |
| (electron/ | | splash.js | | update.js   | | prism.js     | |               |
|  ai/git/   | | splashAPI | | updateAPI   | | electronAPI  | |               |
|  lsp/...)  | |           | |             | | (reduzido)   | |               |
+------------+ +-----------+ +-------------+ +--------------+ +---------------+
  ctxIso=on     ctxIso=on     ctxIso=on       ctxIso=on        ctxIso=on
  sandbox=on    sandbox=on    sandbox=on      sandbox=on       sandbox=on
  nodeInt=off   nodeInt=off   nodeInt=off     nodeInt=off      nodeInt=off
\end{lstlisting}

\begin{nota}
\textbf{Invariante de segurança.} Em todas as cinco janelas vale
\lstinline|contextIsolation:true|, \lstinline|sandbox:true| e
\lstinline|nodeIntegration:false|. O renderer nunca toca em Node, \lstinline|fs|
ou \lstinline|spawn| diretamente: tudo cruza a fronteira IPC por um canal
enumerado do preload, onde o \emph{main} valida a entrada
(\lstinline|safePath|, \emph{allowlist} de binários, \emph{protected flags}) antes
de executar a operação privilegiada.
\end{nota}

% !TEX root = ../main.tex
\chapter{AURORA: o editor de código}
\label{cap:06-aurora-editor}

Este capítulo documenta o subsistema de edição de texto da \tool{AURORA}: o
motor de edição (\tool{Monaco Editor}), a camada que o instancia e gerencia
(\code{EditorManager}), o registro de modelos compartilhados entre painéis
(\code{SharedModelRegistry}), o gerenciador de abas (\code{TabManager}), o
editor dividido em múltiplos painéis (\code{SplitEditorManager}), o realce de
sintaxe (gramáticas Monarch para \cmm{} e \emph{assembly}, mais a gramática
vendorizada de Verilog), as \emph{decorations} visuais e o mecanismo de busca.

O escopo aqui é estritamente o que vive em \path{js/editor/} e \path{js/tabs/}.
Recursos de linguagem que dependem de servidores externos --- \emph{language
server} Verible, \code{clang-format}, \code{slang} e o realce semântico por
\emph{tree-sitter} --- são apenas \emph{conectados} a partir daqui (ver
\autoref{sec:editor-bootstrap}) e documentados em profundidade em outro
capítulo. Quando um assunto pertencer a outro escopo, ele é mencionado de
passagem com referência cruzada.

\section{Visão geral da arquitetura}
\label{sec:editor-overview}

O editor é construído sobre o \tool{Monaco Editor} (o mesmo motor de edição do
VS Code), carregado por meio do \emph{loader} AMD distribuído com o pacote. Em
cima desse motor a \tool{AURORA} mantém quatro módulos principais, todos em
\path{js/editor/} e \path{js/tabs/}:

\begin{tabularx}{\linewidth}{Y Y}
\toprule
\textbf{Módulo} & \textbf{Responsabilidade} \\
\midrule
\file{js/editor/monaco\_editor.js} & Classe \code{EditorManager}: bootstrap do Monaco, criação de instâncias do editor no painel principal, temas, \emph{decorations}, busca, atalhos, registro das linguagens \cmm{} e \code{asm}. \\
\file{js/editor/shared\_models.js} & \code{SharedModelRegistry}: registro de modelos de texto (\code{ITextModel}) com contagem de referências, compartilhados entre painéis. \\
\file{js/editor/split\_editor.js} & \code{SplitEditorManager} e \code{SplitPane}: editor dividido em até três painéis lado a lado, cada um com sua barra de abas. \\
\file{js/tabs/tab\_manager.js} & \code{TabManager}: abas do painel principal, ativação, marcação de modificado/salvo, contagem de instâncias, persistência de ordem, reabertura de abas fechadas. \\
\bottomrule
\end{tabularx}

\noindent
Módulos auxiliares completam o subsistema: \file{js/editor/ai\_selection\_widget.js}
(estrela \enquote{perguntar à IA} sobre uma seleção), \file{js/tabs/tab\_drag.js}
(reordenação de abas por arrastar), \file{js/tabs/tab\_viewers.js} e
\file{js/tabs/tab\_watchers.js} (visualizadores binários e observação de
arquivos), e \file{js/tabs/untitled\_docs.js} (documentos sem título).

\begin{nota}
A regra central que costura todos esses módulos é a do \textbf{modelo
compartilhado}: cada arquivo aberto corresponde a exatamente um
\code{ITextModel} do Monaco, e \emph{todas} as instâncias visuais de editor
(painel principal e painéis divididos) que mostram esse arquivo apontam para o
mesmo modelo. Edições, pilha de \emph{undo}/\emph{redo} e estado de
\enquote{modificado} ficam, portanto, sempre em sincronia, independentemente de
em qual painel o usuário digita.
\end{nota}

\section{O Monaco Editor e o \emph{loader} AMD}
\label{sec:editor-monaco}

\subsection{Versão fixada}
\label{sec:editor-pin}

A versão do \tool{Monaco Editor} é fixada \emph{exatamente} em \texttt{0.52.2}.
Isso é declarado em \file{package.json} sem nenhum operador de faixa (sem
\lstinline|^| ou \lstinline|~|):

\begin{lstlisting}[language=Java]
"monaco-editor": "0.52.2",
\end{lstlisting}

\noindent
Um \emph{semver} simples (sem modificador de faixa) sinaliza, por convenção, que
a dependência está sob verificação estrita. O script
\file{scripts/check-pinned-versions.js} percorre \code{dependencies} e
\code{devDependencies}, filtra as entradas que casam com a expressão de
\emph{semver} exato

\begin{lstlisting}[language=Java]
const EXACT_SEMVER_RE = /^\d+\.\d+\.\d+(-[0-9A-Za-z.-]+)?$/;
\end{lstlisting}

\noindent
e compara o valor declarado com a versão realmente instalada em
\path{node_modules/<pkg>/package.json}. Se divergirem, o processo encerra com
código de erro. O cabeçalho do próprio script registra o motivo histórico do
pino: a versão \texttt{0.53.0} do Monaco lança uma exceção dentro do próprio
\code{monaco.contribution.js} durante a inicialização, o que bloqueia
\code{EditorManager.initialize()}. O \code{monaco-editor} é, no momento, a única
dependência efetivamente pinada por esse mecanismo. O script roda no
\code{prestart} do \file{package.json} e na esteira de CI após \code{npm ci}.

\subsection{O \emph{loader} AMD e a configuração do caminho \texttt{vs}}
\label{sec:editor-loader}

O Monaco é distribuído como um conjunto de módulos AMD. A \tool{AURORA} carrega
primeiro o \emph{loader} AMD do pacote --- o único \emph{script} clássico (não
módulo ES) que precisa rodar antes dos módulos ES da aplicação, pois instala o
\code{require()} global do qual \file{monaco\_editor.js} depende. Em
\file{index.html}:

\begin{lstlisting}[language=Java]
<script src="node_modules/monaco-editor/min/vs/loader.js"></script>
<script>
  require.config({
    paths: { vs: new URL('node_modules/monaco-editor/min/vs',
                          document.baseURI).href }
  });
</script>
\end{lstlisting}

\noindent
O caminho-base \texttt{vs} é ancorado em uma URL \emph{absoluta} construída a
partir de \lstinline|document.baseURI|, em vez de um caminho relativo nu. O
Monaco constrói a URL de seu \emph{web worker} a partir desse caminho e o carrega
via \code{importScripts} dentro do \emph{worker}; sob o esquema \lstinline|file://|
(aplicação empacotada) um caminho relativo se resolveria contra a raiz do sistema
de arquivos e falharia. Ancorar em \lstinline|document.baseURI| funciona tanto em
desenvolvimento (\lstinline|http://localhost|) quanto na aplicação empacotada.

A inicialização efetiva dos módulos do Monaco acontece em
\file{monaco\_editor.js} por meio do \code{require} AMD:

\begin{lstlisting}[language=Java]
require(['vs/editor/editor.main'], function () {
    setupCMMLanguage();
    setupASMLanguage();
    initVerilogLSP();   // O2  -- Verible (.v/.sv): diagnostico, formatacao, outline
    initClangFormat();  // Shift+Alt+F: C / C++ / CMM via clang-format empacotado
    initSlang();        // O11 -- analise semantica slang (Verilog/SystemVerilog)
    initTreeSitter();   // O7  -- realce semantico (tree-sitter / WASM)
    monaco.editor.defineTheme('cmm-dark', { /* ... */ });
    monaco.editor.defineTheme('cmm-light', { /* ... */ });
    resolve();
});
\end{lstlisting}

\noindent
As funções \code{initVerilogLSP}, \code{initClangFormat}, \code{initSlang} e
\code{initTreeSitter} apenas \emph{conectam} provedores às linguagens já
registradas; o trabalho pesado dessas integrações pertence a outro capítulo.

\subsection{Idempotência da inicialização: \code{initMonaco}}
\label{sec:editor-initmonaco}

A função \code{initMonaco} memoiza a promessa de inicialização. Tanto
\file{renderer.js} quanto o \emph{bootstrap} de \code{DOMContentLoaded} do
próprio módulo chamam \code{initMonaco}; a memoização garante que os módulos AMD
do Monaco carreguem e que as linguagens/temas sejam registrados exatamente uma
vez, com ambos os chamadores compartilhando a mesma resolução:

\begin{lstlisting}[language=Java]
let _monacoReady = null;
function initMonaco() {
    if (_monacoReady) return _monacoReady;
    _monacoReady = new Promise((resolve) => {
        require(['vs/editor/editor.main'], function () { /* ... */ resolve(); });
    });
    return _monacoReady;
}
\end{lstlisting}

\section{Contratos de inicialização: \code{EditorManager.ready}}
\label{sec:editor-bootstrap}

A criação de qualquer instância de editor depende de duas condições: os módulos
AMD do Monaco precisam estar carregados, \emph{e} o \code{EditorManager} precisa
ter localizado e dimensionado o contêiner DOM. Como um arquivo pode ser aberto
antes de esse processo terminar (por exemplo, o usuário clica em um \path{.v}
logo após o lançamento), o módulo expõe uma promessa pública que serve de portão
de entrada:

\begin{lstlisting}[language=Java]
let _resolveEditorManagerReady;
EditorManager.ready = new Promise((resolve) => {
    _resolveEditorManagerReady = resolve;
});

document.addEventListener('DOMContentLoaded', async () => {
    try {
        await initMonaco();
        await EditorManager.initialize();
    } finally {
        _resolveEditorManagerReady();
    }
    /* ... listener de resize ... */
});
\end{lstlisting}

\noindent
\code{EditorManager.initialize()} é assíncrono. Ele primeiro aguarda
\code{ensureMonacoInitialized()} (que resolve quando \lstinline|window.monaco|
existe, com sondagem a cada 100\,ms), depois localiza o contêiner
\lstinline|#monaco-editor|, ajusta sua altura e largura para
100\%, carrega o tema salvo de \code{localStorage} e instala o
observador de redimensionamento responsivo:

\begin{lstlisting}[language=Java]
static async initialize() {
    await ensureMonacoInitialized();
    this.editorContainer = document.getElementById('monaco-editor');
    if (!this.editorContainer) { console.error('Editor container not found'); return; }
    this.editorContainer.style.height = '100%';
    this.editorContainer.style.width  = '100%';
    const savedTheme = localStorage.getItem('editorTheme');
    const isDark = savedTheme ? savedTheme === 'dark' : true;
    this.setTheme(isDark);
    this.setupResponsiveObserver();
}
\end{lstlisting}

\noindent
O contrato para os chamadores é, portanto: \textbf{sempre aguardar
\code{EditorManager.ready} antes de chamar \code{createEditorInstance}}. É
exatamente o que \code{TabManager.addTab} faz (ver \autoref{sec:editor-addtab}).

\section{O \code{EditorManager} e o ciclo de criação de instâncias}
\label{sec:editor-manager}

O \code{EditorManager} é uma classe estática: não há instâncias, apenas estado
de classe. Os campos principais são:

\begin{tabularx}{\linewidth}{Y Y}
\toprule
\textbf{Campo estático} & \textbf{Conteúdo} \\
\midrule
\code{editors} & \code{Map}: \emph{filePath} $\rightarrow$ \lstinline|{ editor, container }| (uma entrada por arquivo aberto no painel principal). \\
\code{activeEditor} & A instância de editor atualmente visível/focada no painel principal. \\
\code{editorContainer} & O elemento DOM \lstinline|#monaco-editor| que hospeda as instâncias. \\
\code{currentTheme} & Tema corrente (\code{'cmm-dark'} por padrão). \\
\code{resizeObserver} & \code{ResizeObserver} responsivo único. \\
\code{findStates} & \code{Map}: \emph{filePath} $\rightarrow$ estado do \emph{widget} de busca (aberto? termo?). \\
\code{decorationCollections} & \code{Map}: \emph{filePath} $\rightarrow$ IDs das \emph{decorations} ativas (bra-ket e barra vertical). \\
\bottomrule
\end{tabularx}

\subsection{\code{createEditorInstance}}
\label{sec:editor-create}

\code{createEditorInstance(filePath, initialContent)} é o ponto único de criação
de uma instância no painel principal. Seu fluxo:

\begin{enumerate}[leftmargin=*]
  \item \textbf{\emph{Fallback} preguiçoso do contêiner.} Se \code{editorContainer}
        ainda for nulo, tenta capturá-lo do DOM (\lstinline|#monaco-editor|). Se
        nem o contêiner nem \lstinline|window.monaco| existirem, registra um erro
        e retorna --- o editor não pode ser criado fora do contrato de
        \autoref{sec:editor-bootstrap}.
  \item \textbf{Idempotência.} Se já existir um editor para esse \emph{filePath},
        reutiliza-o; se o conteúdo do modelo estiver vazio e houver
        \code{initialContent}, semeia o modelo. Isso evita que chamadores em
        corrida (a auto-criação de \code{setActiveEditor} e a IIFE de
        \code{addTab}) empilhem dois \lstinline|div| de editor no mesmo
        contêiner.
  \item \textbf{Criação do \lstinline|div|.} Um \lstinline|div.editor-instance|
        absolutamente posicionado (\lstinline|top/left/right/bottom: 0|,
        \lstinline|display: none|) é anexado ao contêiner; seu \code{id} é
        derivado do caminho e \lstinline|dataset.filePath| recebe o caminho.
  \item \textbf{Resolução de linguagem e tema.} A linguagem vem de
        \code{getLanguageFromPath} (\autoref{sec:editor-langmap}); o tema é
        \code{cmm-dark}/\code{cmm-light} para \cmm{} e o tema corrente para o
        restante.
  \item \textbf{Aquisição do modelo compartilhado.}
        \lstinline|SharedModelRegistry.acquire(filePath, initialContent, language)|
        devolve o \code{ITextModel} (criando-o se necessário) e incrementa a
        contagem de referências.
  \item \textbf{Criação do editor Monaco} com \lstinline|model:| explícito (ver
        \autoref{sec:editor-shared} para o porquê) e o conjunto extenso de opções
        descrito em \autoref{sec:editor-options}.
  \item \textbf{Recursos extras, \emph{widget} de IA, \emph{listeners} e
        \emph{decorations}} (passos detalhados abaixo).
\end{enumerate}

\noindent
A passagem explícita de \lstinline|model:| é essencial:

\begin{lstlisting}[language=Java]
const model = SharedModelRegistry.acquire(filePath, initialContent, language);
const editor = monaco.editor.create(editorDiv, {
    theme, model, automaticLayout: true, /* ... demais opcoes ... */
});
this.setupEnhancedFeatures(editor);
attachAiSelectionWidget(editor, { getFilePath: () => filePath });
this.editors.set(filePath, { editor, container: editorDiv });
\end{lstlisting}

\subsubsection{\emph{Listeners} registrados por instância}

Após \lstinline|editors.set|, \code{createEditorInstance} conecta uma série de
ouvintes:

\begin{itemize}[leftmargin=*]
  \item \textbf{Foco/desfoco.} \code{onDidFocusEditorWidget} dispara o evento
        \code{aurora-editor-focused} (com \lstinline|paneIndex: 0|) e
        \code{aurora-editor-focusstate} (\lstinline|focused: true|);
        \code{onDidBlurEditorWidget} dispara o estado \lstinline|focused: false|.
        São eventos customizados, usados para evitar dependência circular com o
        \code{TabManager}.
  \item \textbf{Ligaduras por linguagem.} \code{syncLigatures} desliga as
        ligaduras de fonte quando a linguagem do modelo é \code{verilog} (para
        que \lstinline|<=|, atribuição não-bloqueante, não seja renderizada como
        \enquote{$\leq$}) e as mantém ligadas para as demais. É reaplicado em
        \code{onDidChangeModel} e \code{onDidChangeModelLanguage}.
  \item \textbf{Re-decoração agendada.} \code{onDidChangeModelContent},
        \code{onDidScrollChange} e \code{onDidLayoutChange} chamam
        \code{\_scheduleRedecorate}, que reexecuta as varreduras bra-ket e
        barra-vertical de forma \emph{debounced} (120\,ms).
\end{itemize}

\subsection{O conjunto de opções do editor}
\label{sec:editor-options}

A chamada a \lstinline|monaco.editor.create| em \code{createEditorInstance}
configura um conjunto amplo de opções. Os grupos mais relevantes:

\begin{tabularx}{\linewidth}{Y Y}
\toprule
\textbf{Grupo} & \textbf{Configuração efetiva} \\
\midrule
Cursor & \lstinline|cursorSmoothCaretAnimation: 'on'|, \lstinline|cursorStyle: 'line'|, \lstinline|cursorWidth: 2|, \lstinline|cursorBlinking: 'smooth'|. \\
Busca & \lstinline|addExtraSpaceOnTop|, \lstinline|seedSearchStringFromSelection: 'always'|, \lstinline|globalFindClipboard: true|, \lstinline|loop: true|. \\
Navegação & \emph{breadcrumbs} ligados (caminho + símbolos); \emph{outline} com filtros; \emph{code lens} com fonte \enquote{JetBrains Mono}. \\
Sugestões & \code{suggest} com praticamente todas as categorias habilitadas (métodos, funções, classes, \emph{snippets}, etc.). \\
Realce semântico & \lstinline|'semanticHighlighting.enabled': true| (camada tree-sitter), independentemente do tema. \\
Quebra de linha & \lstinline|wordWrap: 'bounded'| em \emph{desktop} (ou \code{'on'} abaixo de 768\,px), \lstinline|wordWrapColumn: 120|. \\
Minimapa & habilitado acima de 1024\,px; escala \lstinline|1| acima de 1200\,px, senão \lstinline|0.8|. \\
Fonte & família \enquote{JetBrains Mono}/\enquote{Fira Code}/etc., \lstinline|fontLigatures: true|, tamanho 11 (móvel) ou 12. \\
Brackets & \lstinline|bracketPairColorization.enabled: true|; guias de \emph{brackets} e indentação. \\
Edição & \lstinline|autoClosingBrackets: 'always'|, \lstinline|autoClosingQuotes: 'always'|, \lstinline|formatOnPaste: true|, \lstinline|formatOnType: true|. \\
\emph{Zoom} & \lstinline|mouseWheelZoom: false| (o \emph{zoom} de janela é a via intencional, em \file{zoom.js}). \\
\bottomrule
\end{tabularx}

\subsection{Recursos extras e atalhos: \code{setupEnhancedFeatures}}
\label{sec:editor-enhanced}

\code{setupEnhancedFeatures(editor)} registra os atalhos de teclado por meio de
\code{editor.addCommand}. A tabela resume os comandos:

\begin{tabularx}{\linewidth}{Y Y}
\toprule
\textbf{Atalho} & \textbf{Ação} \\
\midrule
\lstinline|Ctrl/Cmd + F| & \code{actions.find} (e foca o campo de busca após 50\,ms). \\
\lstinline|Ctrl/Cmd + H| & \code{editor.action.startFindReplaceAction}. \\
\lstinline|Ctrl/Cmd + Shift + F| & \code{editor.action.formatDocument}. \\
\lstinline|Ctrl/Cmd + G| & \code{editor.action.gotoLine}. \\
\lstinline|Ctrl/Cmd + P| & \code{editor.action.quickCommand}. \\
\lstinline|F12| & \code{editor.action.revealDefinition}. \\
\lstinline|Alt + F12| & \code{editor.action.peekDefinition}. \\
\lstinline|Ctrl/Cmd + F12| & \code{editor.action.goToImplementation}. \\
\lstinline|Shift + F12| & \code{editor.action.goToReferences}. \\
\bottomrule
\end{tabularx}

\noindent
O método também instala um \code{onDidChangeModelContent} que detecta o
fechamento do \emph{widget} de busca: ele primeiro consulta o estado próprio
rastreado em \code{findStates} e só então executa a consulta DOM
(\lstinline|.find-widget|), de modo que o caminho quente de digitação não pague
um \code{querySelector} por tecla quando nenhuma busca está aberta.

\subsection{Ativação, foco e \emph{layout}: \code{setActiveEditor}}
\label{sec:editor-setactive}

\code{setActiveEditor(filePath)} torna visível a instância de um arquivo. Antes
de trocar, salva o estado de busca do arquivo que está saindo (aberto? termo?).
Em seguida esconde todas as instâncias (\lstinline|display: none|), localiza a
entrada do arquivo e, se ela existir, exibe-a (\lstinline|display: block|) e a
marca como \code{activeEditor}. Não há auto-criação aqui: a criação pertence a
\code{TabManager.addTab}, então se a instância ainda não estiver no \code{Map} o
método retorna \lstinline|null| silenciosamente (a IIFE de \code{addTab}
chamará novamente quando a instância existir).

O \code{layout()} e o \code{focus()} são adiados via
\code{requestAnimationFrame}, que dispara depois de o navegador ter aplicado o
\lstinline|display: block| e computado o \emph{layout} --- exatamente quando
\code{editor.layout()} consegue medir o contêiner corretamente:

\begin{lstlisting}[language=Java]
requestAnimationFrame(() => {
    if (!this.activeEditor) return;          // fechamento rapido pode anular
    this.activeEditor.layout();
    if (!(window.SplitEditorManager && window.SplitEditorManager.focusedPane > 0)) {
        this.activeEditor.focus();           // nao rouba foco de um split ativo
    }
    /* ... restaura o widget de busca para este arquivo ... */
});
\end{lstlisting}

\noindent
A guarda \lstinline|focusedPane > 0| impede que o painel principal roube o foco
do teclado enquanto o usuário trabalha em um painel dividido.

\subsection{Fechamento: \code{closeEditor} e \code{cleanup}}
\label{sec:editor-close}

\code{closeEditor(filePath)} limpa o ponteiro \code{activeEditor} \emph{antes} de
descartar a instância (para que nenhum código chame \code{setValue}/\code{layout}
sobre uma instância já descartada), descarta a \emph{view} do editor (mas
\textbf{não} o modelo --- isso é tarefa do registro), remove o \lstinline|div| do
contêiner, apaga as entradas em \code{editors}, \code{findStates} e
\code{decorationCollections}, e por fim chama
\lstinline|SharedModelRegistry.release(filePath)|. O modelo só é efetivamente
descartado quando a contagem de referências chega a zero.

\code{cleanup()} desconecta o \code{resizeObserver}, descarta cada editor e
libera o modelo correspondente, e limpa todos os \code{Map}/\code{Set} de
estado.

\subsection{Mapeamento de extensão para linguagem}
\label{sec:editor-langmap}

\code{getLanguageFromPath(filePath)} resolve a extensão do arquivo para um
identificador de linguagem Monaco. As entradas relevantes para o domínio da
\tool{AURORA}:

\begin{tabularx}{\linewidth}{Y Y Y}
\toprule
\textbf{Extensão} & \textbf{Linguagem} & \textbf{Observação} \\
\midrule
\code{.cmm} & \code{cmm} & \cmm{}, gramática Monarch própria. \\
\code{.asm} & \code{asm} & \emph{assembly} do processor, gramática própria. \\
\code{.v}, \code{.vh} & \code{verilog} & gramática vendorizada do Monaco. \\
\code{.sv}, \code{.svh} & \code{systemverilog} & gramática vendorizada do Monaco. \\
\code{.spf} & \code{json} & arquivo de projeto (JSON) --- ganha chaves/strings/números e \emph{folding}. \\
\code{.c}, \code{.h} & \code{c} & \\
\code{.cpp}, \code{.cc}, \code{.cxx}, \code{.hpp} & \code{cpp} & \\
\bottomrule
\end{tabularx}

\noindent
Demais extensões comuns (\code{js}, \code{ts}, \code{html}, \code{css},
\code{json}, \code{md}, \code{py}) mapeiam para suas linguagens nativas; o que não
estiver no mapa cai em \code{plaintext}. O \code{SplitEditorManager} tem seu
próprio mapa equivalente em \code{\_langFromPath} (\autoref{sec:editor-split}).

\section{O \code{SharedModelRegistry}}
\label{sec:editor-shared}

O \file{js/editor/shared\_models.js} resolve um problema concreto: dividir o
editor não pode clonar o texto do arquivo em um novo modelo por painel, porque
então as edições não se propagariam e apenas a aba do painel principal receberia
a marca de \enquote{modificado}. Com modelos compartilhados, toda instância que
mostra o mesmo arquivo aponta para o mesmo \code{ITextModel}: edições,
\emph{undo}/\emph{redo} e \emph{decorations} permanecem em sincronia, e o fluxo
de modificado/salvo funciona não importa em qual painel se digita.

\subsection{Por que a \tool{AURORA} controla o ciclo de vida do modelo}

O Monaco descarta automaticamente um modelo quando o editor que o \emph{criou} é
descartado --- o que arrancaria o modelo dos demais painéis. A regra do módulo é,
portanto:

\begin{itemize}[leftmargin=*]
  \item sempre passar \lstinline|model:| explicitamente a
        \lstinline|monaco.editor.create(...)|;
  \item adquirir o modelo pelo registro (\lstinline|refCount += 1|);
  \item liberá-lo quando o editor correspondente é descartado
        (\lstinline|refCount -= 1|; descartar a zero).
\end{itemize}

\subsection{A estrutura interna}

O registro mantém um \code{Map} privado de \emph{filePath} para
\lstinline|{ model, refCount, savedAltVersionId }|. As chaves de URI são geradas
por \lstinline|monaco.Uri.file(filePath)|, que normaliza barras invertidas do
Windows e produz uma string de URI estável --- a mesma forma pela qual o Monaco
indexa modelos internamente.

\subsection{\code{acquire}}

\code{acquire(filePath, content, language)} obtém (ou cria) o modelo
compartilhado. Se o modelo já existir, \code{content} é ignorado --- o texto em
memória é a fonte da verdade, de modo que um segundo painel abrindo o arquivo vê
as edições não salvas do usuário, e não o conteúdo em disco:

\begin{lstlisting}[language=Java]
acquire(filePath, content, language) {
  let entry = entries.get(filePath);
  if (!entry) {
    const uri = uriFor(filePath);
    let model = monaco.editor.getModel(uri);
    if (!model) {
      model = monaco.editor.createModel(content || '', language, uri);
    } else if (typeof content === 'string' && model.getValue() === '') {
      model.setValue(content);   // modelo existente vazio + temos conteudo -> semeia
    }
    entry = { model, refCount: 0,
              savedAltVersionId: model.getAlternativeVersionId() };
    entries.set(filePath, entry);
  }
  entry.refCount += 1;
  return entry.model;
}
\end{lstlisting}

\noindent
No momento da aquisição, o registro tira um \emph{snapshot} do
\emph{alternative version id} do modelo. Esse é o valor contra o qual todo painel
compara para decidir se o \emph{buffer} está \enquote{sujo}. O Monaco incrementa
esse id a cada edição e o \emph{reverte} no \emph{undo}; assim, desfazer até o
estado salvo limpa corretamente a marca de modificado, exatamente como no
VS Code.

\subsection{\code{release}, \code{markSaved} e \code{isDirty}}

\code{release(filePath)} decrementa a contagem; ao chegar a zero, descarta o
modelo e remove a entrada. É seguro chamar com um caminho não rastreado (por
exemplo, arquivos binários nunca registrados).

\code{isDirty(filePath)} é uma comparação de inteiros barata e independente de
painel: retorna verdadeiro \emph{se e somente se} o
\emph{alternative version id} corrente difere do \emph{snapshot} salvo. Retorna
falso para arquivos não rastreados.

\code{markSaved(filePath)} fixa o estado corrente do \emph{buffer} como
\enquote{o estado salvo}, atualizando \code{savedAltVersionId}. Deve ser chamado
logo após uma escrita bem-sucedida em disco, de modo que \code{isDirty} retorne
falso até a próxima edição --- e o comportamento sobrevive a \emph{undo}/\emph{redo}
que cruze esse ponto.

\begin{lstlisting}[language=Java]
isDirty(filePath) {
  const entry = entries.get(filePath);
  if (!entry) return false;
  return entry.model.getAlternativeVersionId() !== entry.savedAltVersionId;
}
markSaved(filePath) {
  const entry = entries.get(filePath);
  if (entry) entry.savedAltVersionId = entry.model.getAlternativeVersionId();
}
\end{lstlisting}

\subsection{Exposição global}

Os métodos auxiliares \code{getModel}, \code{has} e \code{refCount} completam a
interface. Ao final do módulo, o registro é exposto em \lstinline|window.SharedModelRegistry|
para os caminhos de \file{split\_editor.js} e do \code{TabManager} que vivem fora
do grafo de \emph{import} do módulo de edição.

\section{O \code{TabManager}}
\label{sec:editor-tabmanager}

O \file{js/tabs/tab\_manager.js} é a classe estática \code{TabManager}, que
gerencia as abas do \emph{painel principal} (os painéis divididos têm sua própria
barra de abas, ver \autoref{sec:editor-split}). Seus campos de estado incluem:

\begin{tabularx}{\linewidth}{Y Y}
\toprule
\textbf{Campo} & \textbf{Conteúdo} \\
\midrule
\code{tabs} & \code{Map}: \emph{filePath} $\rightarrow$ conteúdo (texto) ou o marcador \lstinline|'[BINARY_FILE]'|. \\
\code{activeTab} & Caminho da aba ativa no painel principal. \\
\code{previewTab} & Caminho da aba de \emph{preview} (itálico) corrente, ou \lstinline|null|. \\
\code{unsavedChanges} & \code{Set} de caminhos com edições não salvas. \\
\code{closedTabsStack} & Pilha (máx. 10) de abas fechadas, para reabertura. \\
\code{untitledDocuments} & \code{Map} de documentos sem título. \\
\code{fileWatchers}, \code{lastModifiedTimes} & Suporte à observação de arquivos no disco. \\
\bottomrule
\end{tabularx}

\subsection{Criação de abas: \code{addTab}}
\label{sec:editor-addtab}

\code{addTab(filePath, content, options)} é o ponto de entrada para abrir um
arquivo. Seu comportamento, em ordem:

\begin{enumerate}[leftmargin=*]
  \item \textbf{Roteamento para o painel dividido focado.} Se houver um painel
        dividido focado (\lstinline|focusedPane > 0|) e a chamada não vier de
        dentro do \emph{split} (\lstinline|!options._fromSplit|), o arquivo é
        encaminhado a \code{SplitEditorManager.openInFocusedPane} e \code{addTab}
        retorna --- a regra \enquote{uma nova aba sempre cai no \emph{split}
        focado}.
  \item \textbf{Aba já existente.} Se a aba já existir, promove-a de \emph{preview}
        a permanente quando aplicável e a ativa.
  \item \textbf{\emph{Preview} substituído.} Abrir um novo \emph{preview} fecha
        silenciosamente o \emph{preview} anterior.
  \item \textbf{Elemento DOM da aba.} Cria o \lstinline|div.tab| com
        \lstinline|data-path|, ícone, nome de exibição e botão de fechar; marca
        \code{binary-file} para arquivos binários e \code{preview} para
        \emph{previews}.
  \item \textbf{\emph{Listeners} da aba.} Clique (ativa, devolve o foco ao painel
        principal), duplo-clique (promove \emph{preview}), botão de fechar, e
        clique do botão do meio (\code{auxclick}, fecha a aba).
  \item \textbf{Conteúdo.} Para binários, armazena \lstinline|'[BINARY_FILE]'| e
        ativa o visualizador; para texto, armazena o conteúdo e cria o editor.
\end{enumerate}

\noindent
A criação do editor de texto é o ponto em que o contrato de
\autoref{sec:editor-bootstrap} é honrado: uma IIFE assíncrona aguarda
\code{EditorManager.ready} antes de instanciar:

\begin{lstlisting}[language=Java]
(async () => {
    try {
        await EditorManager.ready;
        const editor = EditorManager.createEditorInstance(filePath, content || '');
        if (!editor) { this.closeTab(filePath); return; }
        this.setupContentChangeListener(filePath, editor);
        this.activateTab(filePath);
        if (options.revealPosition && typeof editor.revealLineInCenter === 'function') {
            const ln = options.revealPosition.line || 1;
            const col = options.revealPosition.column || 1;
            requestAnimationFrame(() => {
                editor.layout();
                editor.setPosition({ lineNumber: ln, column: col });
                editor.revealLineInCenter(ln);
                editor.focus();
            });
        }
    } catch (error) {
        console.error('Error creating editor:', error);
        this.closeTab(filePath);
    }
})();
\end{lstlisting}

\noindent
A opção \code{revealPosition} é usada, por exemplo, pelo \tool{PRISM} ao saltar
da definição de um módulo para o ponto correspondente no código (ver
\autoref{cap:06-aurora-editor} apenas como contexto; o \tool{PRISM} é descrito em
outro capítulo). O salto é adiado para o próximo \emph{frame} com um
\code{layout()} explícito, de modo que o editor recém-exibido tenha dimensões
reais.

\subsection{\emph{Preview tabs} (abas em itálico)}
\label{sec:editor-preview}

A \tool{AURORA} implementa o comportamento de \emph{preview} do VS Code: um
clique simples na árvore de arquivos abre o arquivo como aba \emph{preview}
(itálico), e há no máximo uma por painel. A aba é \enquote{promovida} a permanente
quando o usuário dá duplo-clique nela ou começa a editar o \emph{buffer}.
\code{promotePreviewToPermanent} remove a classe \code{preview} e zera
\code{previewTab}; \code{\_closePreviewSilently} descarta o \emph{preview}
anterior sem prompt, pois um \emph{preview} nunca está sujo (a primeira edição o
teria promovido).

\subsection{Marcação de modificado/salvo}
\label{sec:editor-dirty}

\code{markFileAsModified} adiciona o caminho a \code{unsavedChanges} e troca o
botão de fechar de \textbf{todas} as abas do arquivo (principal e \emph{splits})
por um ponto dourado (\enquote{$\bullet$}); \code{markFileAsSaved} faz o inverso,
restaurando o \enquote{$\times$}. Ambos consultam \emph{todas} as abas via
\lstinline|querySelectorAll('.tab[data-path="..."]')|, garantindo que o
indicador de modificado apareça de forma idêntica em cada instância.

A decisão de \emph{quando} marcar como modificado é tomada em
\code{setupContentChangeListener}, registrado uma única vez por editor (com
guarda de idempotência \lstinline|editor.__auroraContentDisposable|). A cada
mudança de conteúdo ele consulta \lstinline|SharedModelRegistry.isDirty| e
marca modificado ou salvo de acordo --- o que faz a marca de modificado se
limpar sozinha quando o usuário desfaz até o estado salvo:

\begin{lstlisting}[language=Java]
editor.__auroraContentDisposable = editor.onDidChangeModelContent(() => {
    if (this.isUntitledPath(filePath)) { /* ... documentos sem titulo ... */ return; }
    const dirty = window.SharedModelRegistry?.isDirty?.(filePath) ?? false;
    if (dirty) {
        this.markFileAsModified(filePath);
        if (this.previewTab === filePath) this.promotePreviewToPermanent(filePath);
    } else {
        this.markFileAsSaved(filePath);
    }
});
\end{lstlisting}

\subsection{Contagem de instâncias: \code{getInstanceCount}}
\label{sec:editor-instancecount}

\code{getInstanceCount(filePath)} conta quantas \emph{views} de editor apontam
para o arquivo: a aba do painel principal mais cada aba de painel dividido. O
fluxo de fechamento usa essa contagem para decidir se descartar uma \emph{view}
deve perguntar por alterações não salvas --- apenas a \emph{última} instância
dispara o \emph{prompt}; as anteriores apenas descartam sua \emph{view}, pois o
modelo compartilhado (e as edições) sobrevive nas demais.

\begin{lstlisting}[language=Java]
static getInstanceCount(filePath) {
    if (!filePath) return 0;
    let count = this.tabs.has(filePath) ? 1 : 0;
    const split = window.SplitEditorManager;
    if (split && Array.isArray(split.panes)) {
        for (const pane of split.panes) {
            if (pane?.tabs?.has?.(filePath)) count += 1;
        }
    }
    return count;
}
\end{lstlisting}

\subsection{Reabertura de abas fechadas}
\label{sec:editor-reopen}

Ao fechar uma aba (não sem título), o \code{TabManager} empilha
\lstinline|{ filePath, content, timestamp }| em \code{closedTabsStack},
limitada a 10 entradas. \code{reopenLastClosedTab} desempilha a última, tenta
reler o conteúdo atual do disco (caindo para o conteúdo armazenado se o arquivo
não existir mais), recria a aba via \code{addTab} e, se o conteúdo armazenado
diferir do conteúdo em disco, restaura-o e marca o arquivo como modificado.

\subsection{Ordenação e persistência de abas}
\label{sec:editor-order}

A reordenação de abas por arrastar é implementada em \file{js/tabs/tab\_drag.js}
como um \emph{mixin} instalado em \code{TabManager} via \code{Object.assign}. O
método \code{initSortableTabs} liga a API HTML5 de \emph{drag} a cada
\lstinline|.tab| dentro de \lstinline|#tabs-container|. Pontos centrais:

\begin{itemize}[leftmargin=*]
  \item \textbf{Reordenação ao vivo com FLIP.} Durante o arraste, a função
        \code{liveReorder} reposiciona a aba arrastada e anima os vizinhos da
        posição antiga para a nova com uma transição de \lstinline|transform|
        (190\,ms), de modo que deslizem para abrir espaço.
  \item \textbf{MIME customizado.} O arraste usa
        \lstinline|application/x-aurora-tab-path| para que o Monaco não trate o
        \emph{drop} como texto e cole o caminho do arquivo no \emph{buffer}; os
        alvos de \emph{drop} (painéis divididos e \emph{shell} principal) leem
        essa mesma chave.
  \item \textbf{Sinalização entre painéis.} No \code{dragstart}, sinaliza em
        \lstinline|window.SplitEditorManager._dragActive| / \code{\_dragSourcePane}
        que se trata de um arraste de aba da \tool{AURORA} originado no painel
        principal (índice 0), para que os alvos de \emph{drop} dos \emph{splits}
        o aceitem.
  \item \textbf{Auto-conexão.} Um \code{MutationObserver} conecta os
        \emph{listeners} a abas inseridas em tempo de execução.
\end{itemize}

\noindent
A persistência da ordem usa \code{localStorage}:

\begin{lstlisting}[language=Java]
getTabOrder() {
    const c = document.getElementById('tabs-container');
    return Array.from(c.querySelectorAll('.tab')).map((t) => t.getAttribute('data-path'));
}
saveTabOrder()   { localStorage.setItem('editorTabOrder', JSON.stringify(this.getTabOrder())); }
restoreTabOrder(){ /* le 'editorTabOrder' e reanexa cada .tab na ordem salva */ }
\end{lstlisting}

\noindent
\code{saveTabOrder} é chamado ao final de cada arraste; \code{restoreTabOrder}
reanexa cada \lstinline|.tab| na ordem persistida.

\section{O editor dividido (\emph{split editor})}
\label{sec:editor-split}

O \file{js/editor/split\_editor.js} implementa o editor dividido: \textbf{até três
painéis lado a lado}, cada um com sua própria barra de abas e suas próprias
instâncias Monaco. Painéis não focados recebem uma sobreposição de escurecimento
sutil, e há divisores arrastáveis (\code{SplitResizer}) entre cada par de
painéis. Os três objetos centrais são \code{SplitResizer}, \code{SplitPane} e o
gerenciador \code{SplitEditorManager}.

\subsection{Orçamento de painéis e \code{canSplit}}

O orçamento é de três painéis. O painel principal só conta quando efetivamente
tem conteúdo --- assim, depois que o usuário esvazia o painel principal (ele é
escondido), a vaga liberada permite dividir novamente:

\begin{lstlisting}[language=Java]
canSplit() {
    const mainCount = this._mainHasContent() ? 1 : 0;
    if (mainCount + this.panes.length >= 3) return false;
    if (this.focusedPane === 0) return TabManager.activeTab !== null;
    const pane = this.panes.find(p => p.paneIndex === this.focusedPane);
    return !!pane?.activeFile;
}
\end{lstlisting}

\subsection{Criação de um \emph{split}: \code{createSplit}}

\code{createSplit} resolve o arquivo de origem e seu conteúdo a partir do painel
que tem o foco. Exige um editor Monaco real (dividir um visualizador binário ou
uma aba cujo editor ainda não foi instanciado produziria um painel em branco). O
índice do novo painel é \lstinline|max(indices) + 1| (não \lstinline|panes.length + 1|,
que colidiria após o fechamento de um painel intermediário). Em seguida, o novo
painel é criado, o arquivo é aberto nele e o \emph{layout} é recalculado:

\begin{lstlisting}[language=Java]
const newIndex = this.panes.reduce((m, p) => Math.max(m, p.paneIndex), 0) + 1;
const newPane  = new SplitPane(newIndex);
this.panes.push(newPane);
this.wrapper.appendChild(newPane.element);
await newPane.openFile(filePath, content);
this.refreshLayout();
this.setFocus(newIndex);
\end{lstlisting}

\subsection{Edição ao vivo sobre o mesmo modelo}
\label{sec:editor-split-live}

O ponto-chave do \emph{split} é que ele edita o \emph{mesmo modelo} ao vivo.
Quando um painel abre um arquivo, ele adquire o modelo compartilhado pelo
\code{SharedModelRegistry} --- exatamente como o painel principal:

\begin{lstlisting}[language=Java]
const model = SharedModelRegistry.acquire(filePath, content || '', lang);
const editor = monaco.editor.create(editorDiv, {
    theme: EditorManager?.currentTheme ?? 'vs-dark',
    model, automaticLayout: true,
    fontFamily: "'JetBrains Mono', monospace", fontLigatures: true, fontSize: 12,
    minimap: { enabled: false }, scrollBeyondLastLine: false,
    cursorSmoothCaretAnimation: 'on', cursorBlinking: 'smooth',
});
\end{lstlisting}

\noindent
Como todos os painéis apontam para o mesmo \code{ITextModel}, digitar em um
painel aparece em todos os outros em tempo real, a pilha de \emph{undo}/\emph{redo}
é única e compartilhada, e a marca de modificado dispara uma vez por arquivo. O
\emph{listener} de conteúdo do painel dividido vai através do \code{TabManager}
(para que o ponto de modificado caia em todas as abas) e através de
\lstinline|SharedModelRegistry.isDirty| (para que o resultado seja o mesmo
independentemente de qual painel disparou o evento). O painel dividido também
desliga as ligaduras em Verilog e anexa o \emph{widget} de IA, espelhando o
painel principal.

\subsection{Fechamento na última instância}

\code{SplitPane.\_closeFile} usa \code{getInstanceCount}
(\autoref{sec:editor-instancecount}) para decidir se mostra o diálogo de
alterações não salvas. Se este painel for a única instância ainda segurando o
arquivo \emph{e} o \emph{buffer} estiver sujo, executa o mesmo diálogo
\enquote{salvar / não salvar / cancelar} do painel principal; com outras
instâncias vivas, o modelo sobrevive e o fechamento é silencioso. Após descartar
a \emph{view}, libera sua referência no modelo via
\lstinline|SharedModelRegistry.release|.

\subsection{Mover arquivos entre painéis}

\code{moveFileToPane(filePath, targetPaneIndex)} move um arquivo (arrastado de
sua aba) para o painel-alvo: abre no destino \emph{antes} de fechar na origem.
Como todos os painéis compartilham o modelo, abrir no destino antes mantém a
contagem de instâncias $\geq 1$ durante todo o movimento --- o modelo (e
quaisquer edições não salvas) nunca é descartado no meio da operação, e o
fechamento na origem vê que não é a última instância e não pergunta nada. O
conteúdo é sempre obtido do modelo ao vivo (\lstinline|_contentFor|), nunca do
disco.

\subsection{\emph{Layout} e a sobreposição de boas-vindas}

\code{refreshLayout} é a fonte única de verdade para o \emph{layout} do
\emph{split}: esconde painéis sem abas (inclusive o \emph{shell} principal),
reconstrói do zero os divisores entre painéis visíveis consecutivos, equaliza os
painéis visíveis com \lstinline|flex: 1|, e exibe a sobreposição de boas-vindas
apenas quando \emph{todos} os painéis estão vazios. Ao final, chama
\code{\_relayoutVisibleEditors}, que reexecuta \code{layout()} em cada editor
visível dentro de um \code{requestAnimationFrame} --- o Monaco precisa de um
\code{layout()} explícito sempre que o tamanho ou a visibilidade do contêiner
muda.

\subsection{Foco, botão flutuante e \code{getFocusedFile}}

\code{setFocus(paneIndex)} marca o painel focado, aplica o escurecimento aos
demais e dispara \code{aurora:editing-file-changed}. \code{getFocusedFile}
devolve o arquivo em edição no painel focado (do \code{TabManager.activeTab} no
painel principal, ou do \code{activeFile} do painel dividido) --- usado por
\lstinline|Ctrl+S| para que salvar funcione igual em qualquer painel.

O controle de dividir é um único botão flutuante (\code{splitFloatBtn}),
re-aninhado no painel que tem o foco a cada mudança de foco/aba via
\code{\_updateButton}, com estado \emph{enabled}/\emph{tooltip} dirigido por
\code{canSplit} e pelo sistema de \emph{tooltip} i18n.

\section{Realce de sintaxe (Monarch)}
\label{sec:editor-syntax}

O realce de sintaxe usa o mecanismo \emph{Monarch} do Monaco. A \tool{AURORA}
registra \textbf{duas linguagens próprias} --- \cmm{} e \emph{assembly} (\code{asm})
--- com gramáticas Monarch escritas à mão em \file{monaco\_editor.js}. As
linguagens \code{verilog} e \code{systemverilog} \textbf{não} têm gramática
própria neste repositório: elas vêm da gramática \emph{vendorizada} no \emph{build}
do Monaco (\path{node_modules/monaco-editor/min/vs/basic-languages/}) e são
apenas conectadas a provedores externos (Verible, slang, tree-sitter) no
\emph{bootstrap}.

\subsection{A linguagem \cmm}
\label{sec:editor-cmm}

\code{setupCMMLanguage} registra o identificador \code{cmm} e instala o
\emph{tokenizer} Monarch construído por \code{buildCMMTokenizer}. A gramática
cobre:

\begin{description}[leftmargin=*]
  \item[Palavras reservadas (\code{keywords}).] Inclui o conjunto C usual
        (\code{if}, \code{else}, \code{for}, \code{while}, \code{do},
        \code{struct}, \code{return}, \code{break}, \code{continue},
        \code{switch}, \code{case}, \code{default}, \code{goto},
        \code{sizeof}, \code{typedef}, \code{enum}, \code{union}, etc.) e os
        tipos (\code{void}, \code{int}, \code{comp}, \code{char}, \code{float},
        \code{double}, \code{bool}, \code{long}, \code{short}, \code{signed},
        \code{unsigned}, \code{const}, \code{static}). Há ainda três
        identificadores especiais do domínio: \code{Jussara}, \code{Anon} e
        \code{Chrysthofer}.
  \item[Tipos (\code{typeKeywords}).] \code{bool}, \code{int}, \code{long},
        \code{float}, \code{double}, \code{char}, \code{void}, \code{unsigned},
        \code{signed}, \code{short} --- coloridos como \code{keyword.type}.
  \item[Operadores.] O conjunto completo C, incluindo deslocamentos
        (\lstinline|<<|, \lstinline|>>|, \lstinline|>>>|) e suas variantes de
        atribuição; \lstinline|>>>| recebe um \emph{token} próprio
        (\code{operator.shift.arithmetic}).
  \item[Diretivas.] As diretivas \code{\#PRNAME}, \code{\#NUBITS},
        \code{\#NBMANT} (e demais da lista) são realçadas como
        \code{keyword.directive.cmm}.
  \item[Macros \code{\#define}.] A linha \lstinline|#define NOME corpo| colore a
        diretiva e o nome definido; cada uso posterior de \code{NOME} é capturado
        pela regra de identificador (ver \autoref{sec:editor-cmm-define}).
  \item[Funções da biblioteca padrão.] Os identificadores
        \code{in}, \code{fin}, \code{out}, \code{fout}, \code{norm}, \code{sign},
        \code{pset}, \code{abs}, \code{copy}, \code{sqrt}, \code{atan},
        \code{sin}, \code{cos}, \code{tan}, \code{exp}, \code{log}, \code{pow},
        \code{real}, \code{imag}, \code{fase}, \code{mod2}, \code{complex},
        \code{vtv} --- quando seguidos de \lstinline|(| --- recebem
        \code{keyword.function.stdlib.cmm}.
  \item[Números complexos.] O sufixo imaginário \code{im} (como em \code{8im},
        \code{9.234im}, \code{4im}) é reconhecido: a magnitude fica com cor de
        número e o \code{im} recebe \code{number.complex.imaginary.cmm}. Um
        \lstinline|\b| final impede casar \code{im} colado em um identificador
        maior.
  \item[\emph{Array} de arquivo.] O padrão \lstinline|nome[TAM] "arquivo.txt"|
        é tratado de modo que \code{TAM} saia com cor de número.
\end{description}

\subsubsection{A notação de Dirac (bra-ket)}
\label{sec:editor-dirac}

Uma característica distintiva da gramática \cmm{} é o suporte à notação de Dirac
(bra-ket) da mecânica quântica, com uma série de regras dedicadas no estado
\code{root} do \emph{tokenizer}. Os \emph{tokens} envolvidos são
\code{dirac.bracket} (para os colchetes angulares $\langle$ e $\rangle$),
\code{dirac.bar} (para a barra vertical \lstinline+|+) e
\code{keyword.special.dirac} (para símbolos especiais como \code{B}, \code{I},
\code{0}). As regras casam padrões como o produto interno
$\langle a | b \rangle$, o ket $| x \rangle$, o bra $\langle x |$, o operador
densidade $| x \rangle\langle y |$ e combinações com a função \code{in(...)}.
As regras mais específicas vêm primeiro, de modo que a
barra vertical solta caia na última regra (\code{dirac.bar}) apenas quando não
casar nenhum padrão estruturado.

\subsubsection{Conjunto dinâmico de \code{\#define}}
\label{sec:editor-cmm-define}

A gramática \cmm{} mantém um conjunto \emph{vivo} de nomes de macros
\lstinline|#define| objeto. \code{collectCMMDefineNames} varre todos os modelos
\code{cmm} abertos com a expressão âncora

\begin{lstlisting}[language=Java]
const re = /^[ \t]*#define[ \t]+([a-zA-Z_]\w*)/gm;
\end{lstlisting}

\noindent
(que só casa \lstinline|#define| no início de uma linha, evitando registrar um
falso constante escrito dentro de uma string). \code{refreshCMMDefines}
re-constrói o \emph{tokenizer} (via \code{setMonarchTokensProvider}) apenas quando
o conjunto de nomes efetivamente muda --- re-registrar re-tokeniza todos os
modelos \code{cmm}, então a atualização é \emph{debounced}
(300\,ms) e disparada por \code{onDidChangeContent},
\code{onDidCreateModel} e \code{onDidChangeModelLanguage}. Os nomes coletados
são consultados pela regra de identificador via a entrada
\lstinline|@defineConstants| do bloco de casos:

\begin{lstlisting}[language=Java]
[/[a-zA-Z_]\w*/, {
    cases: {
        '@typeKeywords':     'keyword.type',
        '@keywords':         'keyword',
        '@defineConstants':  'constant.define.cmm',
        '@default':          'identifier'
    }
}],
\end{lstlisting}

\noindent
Como a regra de identificador roda apenas no estado \code{root} (comentários e
strings têm seus próprios estados) e palavras reservadas, tipos e funções da
\emph{stdlib} são casados antes, um nome de \code{\#define} nunca pode
sobrescrever uma palavra reservada.

\subsubsection{Estados de comentário e string}

O \emph{tokenizer} \cmm{} tem estados separados para comentários
(\lstinline|/* ... */| com aninhamento via \code{@push}/\code{@pop} e
\lstinline|// ...|) e strings (\lstinline|"..."| com sequências de escape e
strings inválidas/não terminadas). O atributo \code{escapes} define a expressão
de sequências de escape reconhecidas.

\subsection{A linguagem \emph{assembly} (\code{asm})}
\label{sec:editor-asm}

\code{setupASMLanguage} registra o identificador \code{asm} com um \emph{tokenizer}
Monarch para o \emph{assembly} do processor. Os principais grupos:

\begin{description}[leftmargin=*]
  \item[Diretivas.] \code{PRNAME}, \code{NUBITS}, \code{NBMANT}, \code{NBEXPO},
        \code{NDSTAC}, \code{SDEPTH}, \code{NUIOIN}, \code{NUIOOU},
        \code{NUGAIN}, \code{FFTSIZ}, \code{array}, \code{arrays}, \code{ITRAD},
        \code{TOAQUI} --- reconhecidas com prefixo \lstinline|#| como
        \code{keyword.directive}.
  \item[Instruções.] Um conjunto extenso de mnemônicos do conjunto de instruções
        do processor (\code{LOD}, \code{LDI}, \code{SET}, \code{PSH}, \code{POP},
        \code{ADD}, \code{MLT}, \code{DIV}, \code{MOD}, \code{NEG}, \code{ABS},
        \code{NRM}, \code{I2F}, \code{F2I}, \code{AND}, \code{ORR}, \code{XOR},
        \code{INV}, \code{LES}, \code{GRE}, \code{EQU}, \code{SHL}, \code{SHR},
        \code{CAL}, \code{RET}, e muitas variantes com prefixos \code{S\_},
        \code{F\_}, \code{SF\_}, \code{P\_}, \code{PF\_}), realçados como
        \code{keyword.instruction}.
  \item[Instruções de salto.] \code{JMP} e \code{JIZ} recebem um \emph{token}
        próprio (\code{keyword.jumpInstruction}).
  \item[Rótulos.] Um identificador no início de linha seguido de \lstinline|:|
        é realçado como \code{type.identifier}.
  \item[Números.] Hexadecimal (\lstinline|0x...|), decimal e binário
        (sufixo \code{b}); strings com aspas e comentários (\lstinline|//| e
        \lstinline|;|).
\end{description}

\noindent
A regra de identificador em maiúsculas usa um bloco de casos que classifica o
\emph{token} como instrução, diretiva ou identificador comum, conforme pertença a
\lstinline|@instructions| ou \lstinline|@directives|.

\subsection{Verilog e SystemVerilog}
\label{sec:editor-verilog}

As linguagens \code{verilog} e \code{systemverilog} já estão registradas pelo
\emph{build} vendorizado do Monaco; o módulo de edição não define gramática
Monarch para elas. O que a \tool{AURORA} faz é (a) escolher o identificador
correto por extensão em \code{getLanguageFromPath}/\code{\_langFromPath}, (b)
desligar as ligaduras de fonte em Verilog (para que \lstinline|<=| não vire
\enquote{$\leq$}), (c) aplicar as \emph{decorations} de barra vertical
(\autoref{sec:editor-decorations}), e (d) conectar os provedores externos
(Verible, slang, tree-sitter) no \emph{bootstrap}. Esses provedores externos são
documentados em outro capítulo.

\section{Temas do editor}
\label{sec:editor-themes}

O módulo define quatro temas próprios via \code{monaco.editor.defineTheme}:
\code{cmm-dark}, \code{cmm-light}, \code{asm-dark} e \code{asm-light}. As cores
espelham as variáveis de tema da IDE (\path{theme_variables.css}), de modo que a
superfície do editor se misture com o restante da interface --- fundo
\lstinline|#0A0D14|, acento \lstinline|#8E83E8|. Os temas \code{cmm-*} mapeiam os
\emph{tokens} próprios da gramática \cmm{} (incluindo \code{dirac.bracket},
\code{dirac.bar}, \code{number.complex.imaginary.cmm},
\code{operator.shift.arithmetic} e \code{constant.define.cmm}) para a paleta
\enquote{Aurora}; os temas \code{asm-*} mapeiam os \emph{tokens} de
\emph{assembly}.

\code{setTheme(isDark)} aplica o tema a todas as instâncias abertas, escolhendo
por linguagem: \code{cmm-*} para \cmm{}, \code{asm-*} para \emph{assembly} e
\code{vs-dark}/\code{vs} para o restante. A preferência é persistida em
\lstinline|localStorage['editorTheme']|.

\begin{nota}
A \tool{AURORA} opera, na prática, com um tema escuro canônico único. O par
claro/escuro existe na infraestrutura de temas do editor, mas o padrão
inicializado (e a direção de projeto da interface) é o tema escuro
\code{cmm-dark}.
\end{nota}

\section{\emph{Decorations} visuais}
\label{sec:editor-decorations}

O \code{EditorManager} aplica duas \emph{decorations} inline próprias, ambas
varrendo apenas o intervalo \emph{visível} do modelo (e por isso reexecutadas em
rolagem/\emph{layout}), com os IDs de \emph{decoration} guardados por arquivo em
\code{decorationCollections}:

\begin{description}[leftmargin=*]
  \item[Bra-ket (\code{decorateBraKet}).] Procura ocorrências de $\rangle$
        nas linhas visíveis e aplica a classe inline \code{bra-ket-padding},
        ajustando o espaçamento do símbolo de \emph{ket} da notação de Dirac.
  \item[Barra vertical (\code{decorateVerticalBar}).] Procura ocorrências de
        \lstinline+|+ nas linhas visíveis e aplica a classe inline
        \code{vertical-bar-lower}. A varredura é limitada ao intervalo visível
        justamente porque a barra vertical é onipresente em Verilog (operador
        OR).
\end{description}

\noindent
Ambas usam \code{editor.getVisibleRanges()} e \code{model.findMatches(...)}, e
aplicam as \emph{decorations} via \code{editor.deltaDecorations}, reaproveitando
os IDs anteriores. As duas varreduras são reagendadas juntas por
\code{\_scheduleRedecorate}, que coalesce rajadas de digitação e \emph{flings} de
rolagem em uma passada a cada 120\,ms:

\begin{lstlisting}[language=Java]
static _scheduleRedecorate(editor) {
    if (editor.__redecorateTimer) clearTimeout(editor.__redecorateTimer);
    editor.__redecorateTimer = setTimeout(() => {
        editor.__redecorateTimer = null;
        this.decorateBraKet(editor);
        this.decorateVerticalBar(editor);
    }, 120);
}
\end{lstlisting}

\section{Busca}
\label{sec:editor-find}

A busca dentro de um arquivo usa o \emph{widget} de busca nativo do Monaco
(configurado via a opção \code{find} em \autoref{sec:editor-options}:
\code{loop}, \code{globalFindClipboard}, semear da seleção). O atalho
\lstinline|Ctrl/Cmd + F| executa a ação \code{actions.find} e foca o campo de
entrada; \lstinline|Ctrl/Cmd + H| abre buscar-e-substituir.

O \code{EditorManager} mantém o estado da busca por arquivo em \code{findStates}
(\lstinline|{ isOpen, searchTerm }|), salvo ao trocar de aba em
\code{setActiveEditor} e restaurado quando o arquivo volta a ficar ativo. A
recuperação do caminho do arquivo ativo lê \lstinline|dataset.path| da aba
ativa:

\begin{lstlisting}[language=Java]
static getActiveFilePath() {
    const activeTab = document.querySelector('.tab.active');
    return activeTab ? activeTab.dataset.path : null;
}
\end{lstlisting}

\noindent
Além da busca no arquivo, \code{searchInAllFiles(searchTerm, options)} percorre
todos os editores abertos no painel principal e devolve, por arquivo, a lista de
ocorrências (número da linha, coluna, texto da linha e \emph{range}), usando
\lstinline|model.findMatches| com suporte a \emph{regex}, sensibilidade a
maiúsculas e palavra inteira. \code{navigateToSearchResult(filePath, line, col)}
ativa o editor do arquivo, posiciona o cursor, revela a linha no centro e foca o
editor. (A busca de projeto inteiro, na árvore de arquivos, pertence a outro
escopo.)

\section{O \emph{widget} de seleção para IA}
\label{sec:editor-aiwidget}

\code{attachAiSelectionWidget(editor, { getFilePath })} anexa um botão flutuante
em forma de estrela (\enquote{$\star$}) a uma instância de editor; ambos os
\emph{factories} (painel principal e painel dividido) o chamam uma vez por
instância. O botão aparece sempre que há uma seleção não vazia, ancorado no
início da seleção (preferência ABOVE, com \emph{fallback} BELOW). Ao clicar,
abre um pequeno menu de intenções rápidas:

\begin{tabularx}{\linewidth}{Y Y}
\toprule
\textbf{Intenção} & \textbf{Comportamento} \\
\midrule
\code{explain} (Explain) & Envia imediatamente. \\
\code{fix} (Find \& fix bugs) & Envia imediatamente. \\
\code{improve} (Improve) & Envia imediatamente. \\
\code{comment} (Add comments) & Envia imediatamente. \\
\code{ask} (Ask\ldots) & Apenas semeia o compositor para o usuário digitar. \\
\bottomrule
\end{tabularx}

\noindent
A ação escolhida captura o contexto vivo da seleção (\code{code},
\code{language}, \code{filePath}, \code{lineStart}, \code{lineEnd}) e o entrega a
\lstinline|window.AuroraAPI.ai.askAboutSelection(...)|, que abre o painel de IA e
semeia o compositor. O \emph{widget} usa uma guarda de idempotência
(\lstinline|editor.__auroraAiStarAttached|) para nunca anexar-se duas vezes ao
mesmo editor reutilizado. A integração de IA (\tool{Aurora Intelligence}) é
descrita em capítulo próprio.

\section{Resumo dos contratos}
\label{sec:editor-contracts}

\begin{tabularx}{\linewidth}{Y Y}
\toprule
\textbf{Contrato} & \textbf{Garantia} \\
\midrule
\code{EditorManager.ready} & Promessa resolvida quando os módulos AMD do Monaco \emph{e} \code{initialize()} terminaram; todo chamador de \code{createEditorInstance} deve aguardá-la. \\
Modelo explícito & \lstinline|monaco.editor.create| sempre recebe \lstinline|model:| vindo do \code{SharedModelRegistry}, para que o Monaco não descarte o modelo ao fechar um painel. \\
\code{acquire}/\code{release} balanceados & Cada \code{acquire} é casado por um \code{release} no descarte da \emph{view}; o modelo só some na contagem zero. \\
\emph{Dirty} por arquivo & \code{isDirty} compara \emph{alternative version ids}: resposta única por arquivo, idêntica em qualquer painel; \emph{undo} até o estado salvo limpa a marca. \\
\emph{Prompt} na última instância & \code{getInstanceCount} garante que apenas o fechamento da última \emph{view} de um arquivo sujo pergunta por alterações não salvas. \\
\bottomrule
\end{tabularx}

% !TEX root = ../main.tex
\chapter{AURORA: projeto, processadores e árvore de arquivos}
\label{cap:07-aurora-projeto-arvore}

Este capítulo descreve o modelo de projeto da \tool{AURORA} — como um projeto
é representado no disco e em memória — e a navegação pela árvore de arquivos.
A fonte da verdade é o código-fonte do renderer (\path{js/project/},
\path{js/tree/}, \path{js/processors/}) e dos handlers do processo principal
(\path{main/ipc/project.js}, \path{main/ipc/files.js}, \path{main/recents.js}).
O fluxo de compilação que consome o que está aqui descrito (botões Verilog,
Wave, PRISM, \cmm) é assunto do \autoref{cap:08-aurora-toolchain} e não é
re-documentado aqui.

\section{Visão geral do modelo de projeto}

Um projeto da \tool{AURORA} é uma pasta no disco que contém um arquivo
\path{.spf} (\enquote{SAPHO Project File}) de mesmo nome que a pasta. O
\path{.spf} consolida toda a configuração canônica do projeto: metadados de
ciclo de vida (nome, datas, máquina, versão do app) e a \emph{estrutura}
(listas de arquivos, top-level, testbench-topo, processadores). Em torno desse
arquivo orbitam subpastas — uma por processador \tool{SAPHO} — e os arquivos
\file{.v}/\file{.sv} importados pelo usuário.

A representação em memória é dividida entre o processo principal (que faz o IO
no disco) e o renderer (que mantém o estado da UI). Três peças de software
governam o estado:

\begin{description}[leftmargin=2.2em,style=nextline]
  \item[\file{ProjectStore} (\path{js/project/project_store.js})] dono único de
    \enquote{qual projeto está aberto} no renderer — guarda o
    \lstinline|projectPath| (pasta) e o \lstinline|spfPath| (arquivo \path{.spf}).
  \item[\file{SpfStore} (\path{js/project/spf_store.ts})] escritor único e
    serializado do \path{.spf} no lado do renderer; faz leitura pura e escrita
    atômica do bloco \lstinline|structure|.
  \item[\file{state.currentOpenProjectPath} (\path{main/state.js})] no processo
    principal, é a única fonte de verdade de \enquote{qual \path{.spf} está
    aberto}; a pasta do projeto é sempre derivada por \lstinline|path.dirname|.
\end{description}

\begin{nota}
A pasta do projeto NUNCA é armazenada em paralelo ao \path{.spf} no processo
principal: ela é sempre derivada de \lstinline|path.dirname(spfPath)| sob
demanda. O comentário \enquote{A4} em \path{main/ipc/project.js} registra essa
decisão — evita o problema clássico de dois globais que divergem.
\end{nota}

\section{O arquivo de projeto \texttt{.spf}}
\label{sec:spf-formato}

O \path{.spf} é um documento JSON com dois blocos de topo: \lstinline|metadata|
e \lstinline|structure|. O esqueleto é produzido pela classe
\lstinline|ProjectFile| em \path{main/ipc/project.js} no momento da criação.

\subsection{Estrutura completa}

A tabela seguinte lista todos os campos definidos no código. As interfaces
\lstinline|SpfDocument|, \lstinline|SpfStructure| e \lstinline|SpfMetadata| em
\path{spf_store.ts} são a referência do lado do renderer; o construtor de
\lstinline|ProjectFile| é a referência do lado do processo principal.

\begin{longtable}{p{0.30\linewidth} p{0.16\linewidth} p{0.46\linewidth}}
  \toprule
  \textbf{Campo} & \textbf{Bloco} & \textbf{Significado} \\
  \midrule
  \endhead
  \lstinline|projectName|   & metadata  & Nome do projeto (basename da pasta). Exibido na barra de título como \lstinline|<nome>.spf|. \\
  \lstinline|createdAt|     & metadata  & Timestamp ISO de criação. \\
  \lstinline|lastModified|  & metadata  & Timestamp ISO carimbado a cada escrita do \path{.spf}. \\
  \lstinline|lastOpened|    & metadata  & Timestamp ISO da última abertura (gravado em \lstinline|project:open|). \\
  \lstinline|computerName|  & metadata  & \lstinline|COMPUTERNAME| ou \lstinline|os.hostname()|. \\
  \lstinline|appVersion|    & metadata  & Versão do app (\lstinline|app.getVersion()|). \\
  \lstinline|projectPath|   & metadata  & Caminho absoluto da pasta do projeto. \\
  \midrule
  \lstinline|basePath|          & structure & Caminho absoluto da pasta do projeto; reconciliado na abertura. \\
  \lstinline|folders|           & structure & Subpastas registradas pela UI (criadas via \enquote{New Folder}). \\
  \lstinline|processors|        & structure & Lista de processadores \tool{SAPHO}; cada entrada é \lstinline|{ name }|. \\
  \lstinline|topLevelFile|      & structure & Top do design (modo \emph{check} do iverilog, flag \lstinline|-s|). \\
  \lstinline|testbenchFile|     & structure & Testbench-topo (modo \emph{simulação}, flag \lstinline|-s| do botão Wave). \\
  \lstinline|synthesizableFiles|& structure & Arquivos \file{.v} sintetizáveis; cada entrada tem \lstinline|path| (e \lstinline|name|). \\
  \lstinline|testbenchFiles|    & structure & Arquivos \file{.v} de testbench; mesma forma de entrada. \\
  \bottomrule
\end{longtable}

As constantes de \lstinline|STRUCTURE_DEFAULTS| em \path{spf_store.ts}
congelam (\lstinline|Object.freeze|) os valores-padrão de \lstinline|structure|:
\lstinline|basePath: ''|, \lstinline|folders: []|, \lstinline|processors: []|,
\lstinline|topLevelFile: ''|, \lstinline|testbenchFile: ''|,
\lstinline|synthesizableFiles: []| e \lstinline|testbenchFiles: []|. Na
leitura, os defaults são aplicados primeiro e os valores do disco depois
(\code{\{ ...STRUCTURE\_DEFAULTS, ...parsed.structure \}}), de modo que chaves
desconhecidas sobrevivem ao \emph{round-trip} (\enquote{unknown keys survive}).

\subsection{Exemplo de um \texttt{.spf}}

\begin{lstlisting}
{
  "metadata": {
    "projectName": "demo",
    "createdAt": "2026-06-21T12:00:00.000Z",
    "lastModified": "2026-06-21T12:34:56.000Z",
    "computerName": "NIPS-CERN-PC",
    "appVersion": "1.0.0",
    "projectPath": "C:\\Users\\dev\\Documents\\demo"
  },
  "structure": {
    "basePath": "C:\\Users\\dev\\Documents\\demo",
    "folders": [],
    "processors": [ { "name": "cpu" } ],
    "topLevelFile": "cpu\\Hardware\\cpu.v",
    "testbenchFile": "cpu\\Simulation\\cpu_tb.v",
    "synthesizableFiles": [
      { "name": "cpu.v", "path": "cpu\\Hardware\\cpu.v" }
    ],
    "testbenchFiles": [
      { "name": "cpu_tb.v", "path": "cpu\\Simulation\\cpu_tb.v" }
    ]
  }
}
\end{lstlisting}

\subsection{Normalização de caminhos: relativo no disco, absoluto em memória}
\label{sec:spf-paths}

O \file{SpfStore} grava os caminhos de arquivo \emph{relativos} ao
\lstinline|basePath| quando o arquivo está dentro da pasta do projeto, e
\emph{absolutos} quando está fora. Tudo que é exposto fora do \file{SpfStore}
(via \lstinline|read|/\lstinline|update|) é \emph{sempre} absoluto — o chamador
não precisa conhecer essa normalização. As funções relevantes em
\path{spf_store.ts} são:

\begin{itemize}[leftmargin=1.6em]
  \item \lstinline|isAbsolutePath| — reconhece caminhos Windows
    (\lstinline|C:\|, \lstinline|\\server\|) e POSIX (\lstinline|/|).
  \item \lstinline|toRelativeIfInside| — converte para relativo se o absoluto
    estiver sob o \lstinline|baseDir| (comparação \emph{case-insensitive},
    idioma Windows); caso contrário retorna o absoluto intacto.
  \item \lstinline|toAbsoluteFromBase| — junta um relativo com o
    \lstinline|baseDir| na leitura.
  \item \lstinline|expandStructurePaths| / \lstinline|contractStructurePaths| —
    aplicam a conversão sobre os campos de caminho (\lstinline|topLevelFile|,
    \lstinline|testbenchFile|, \lstinline|synthesizableFiles[].path|,
    \lstinline|testbenchFiles[].path|).
\end{itemize}

\begin{nota}
\path{.spf} antigos com caminhos absolutos continuam sendo lidos
(\lstinline|isAbsolutePath| faz \emph{bypass} do \lstinline|resolve|); o
primeiro \lstinline|save| migra automaticamente para a forma relativa. Isso
permite copiar/mover a pasta do projeto sem quebrar as referências internas.
\end{nota}

\subsection{Parsing tolerante}

A leitura do \path{.spf} usa \lstinline|parseSpfTolerant| (presente tanto em
\path{spf_store.ts} quanto em \path{main/ipc/project.js}): tenta
\lstinline|JSON.parse| estrito primeiro e, em caso de falha, faz uma única
passagem leniente que remove BOM, comentários (\lstinline|//| e
\lstinline|/* */| fora de strings) e vírgulas finais (\emph{trailing commas}).
Um arquivo recuperável não é silenciosamente resetado para os defaults; só um
arquivo genuinamente ilegível cai no fallback (defaults), com aviso no console.

\section{O \texttt{SpfStore}: escritor único e serializado}
\label{sec:spfstore}

O \file{SpfStore} (\path{js/project/spf_store.ts}, compilado por
\lstinline|tsc| em \path{spf_store.js}) é o único ponto de escrita do
\path{.spf} no renderer. Sua API espelha um antigo \lstinline|ProjectConfigStore|:

\begin{description}[leftmargin=2em,style=nextline]
  \item[\lstinline|read(spfPath)|] devolve só o bloco \lstinline|structure|
    (com defaults aplicados, caminhos expandidos para absolutos).
  \item[\lstinline|readFull(spfPath)|] devolve \lstinline|metadata| +
    \lstinline|structure|.
  \item[\lstinline|update(spfPath, mutator)|] leitura-mutação-escrita atômica;
    o \lstinline|mutator| recebe o \lstinline|structure| e o \lstinline|metadata|
    é preservado com \lstinline|lastModified| recarimbado.
\end{description}

\subsection{Serialização das escritas}

Cada \path{.spf} tem uma cadeia de promises própria
(\lstinline|writeChainByPath: Map|): chamadas de \lstinline|update| para o mesmo
caminho serializam em ordem de chegada (\lstinline|prev.then(...)|), de modo que
duas mutações no mesmo \emph{tick} não operem sobre o mesmo \emph{snapshot} pré-escrita.

\subsection{Coalescência de leituras}

Para evitar \emph{fanout} de \lstinline|readFile|/\lstinline|JSON.parse|, o
\file{SpfStore} coalesce leituras em voo (\lstinline|inFlightReads: Map|): se
vários consumidores chamam \lstinline|read| no mesmo tick, todos compartilham a
mesma promise. O mapa é limpo quando a promise é resolvida — não há cache
atravessando ticks (sem \emph{mtime} obsoleto). O \lstinline|update|, por outro
lado, sempre lê do disco com \lstinline|readRawUncached| (o \lstinline|writeChain|
já serializa, então pular o coalescing aqui não introduz \emph{race}).

\subsection{Evento \texttt{aurora:spf-changed}}

Após cada escrita, o \file{SpfStore} compara o \lstinline|structure| antes e
depois do \lstinline|mutator| (via \lstinline|JSON.stringify|). Somente quando o
\lstinline|structure| efetivamente mudou é disparado o evento de janela
\lstinline|aurora:spf-changed| (com \lstinline|detail: { spfPath, source }|).
Mutações \emph{no-op} não geram refresh — isso elimina a classe de laços
\enquote{save $\rightarrow$ event $\rightarrow$ refresh $\rightarrow$ load
$\rightarrow$ save (no-op) $\rightarrow$ \ldots}. Os assinantes desse evento
incluem a barra de status, o painel de configuração do processador e o
\file{file\_mode}.

\section{O \texttt{ProjectStore}: dono do projeto aberto}

O \file{ProjectStore} (\path{js/project/project_store.js}) mantém em variáveis
de módulo o \lstinline|projectPath| e o \lstinline|spfPath| do projeto aberto.
A API é \lstinline|setProject(spfPath, projectPath)|, \lstinline|clearProject()|,
\lstinline|getProjectPath()|, \lstinline|getSpfPath()|, \lstinline|hasProject()|
e \lstinline|subscribe(fn)|.

Por compatibilidade com dezenas de pontos de leitura legados, o store
\emph{espelha} seus valores para \lstinline|window.currentProjectPath| e
\lstinline|window.currentSpfPath| a cada mudança (\lstinline|mirrorToWindow|).
Código novo deve preferir os getters. Assinantes (\lstinline|subscribe|) são
notificados com um \emph{snapshot} \lstinline|{ projectPath, spfPath }| — o
controlador de views da árvore, por exemplo, usa isso para habilitar/desabilitar
a view \enquote{standard} conforme um projeto abre ou fecha.

\section{Ciclo de vida: criar, abrir e fechar projeto}

\subsection{Criar projeto}

O fluxo começa na UI: \path{js/ui/new_project_modal.js} valida o nome
(\lstinline|/^[a-zA-Z0-9_-]+$/|) e o caminho, monta
\lstinline|projectPath = <local>\<nome>| e
\lstinline|spfPath = <projectPath>\<nome>.spf|, e chama
\lstinline|electronAPI.createProjectStructure|. No processo principal, o handler
\lstinline|project:createStructure| (\path{main/ipc/project.js}):

\begin{enumerate}[leftmargin=1.8em]
  \item cria a pasta do projeto (\lstinline|fse.mkdir| recursivo);
  \item instancia \lstinline|new ProjectFile(projectPath)| e grava o
    \path{.spf} com \lstinline|metadata| + \lstinline|structure| iniciais;
  \item confirma a existência da pasta e do \path{.spf};
  \item lista o conteúdo da pasta para devolver ao renderer;
  \item registra o novo \path{.spf} nos recentes e reconstrói a \emph{jumplist}
    do Windows.
\end{enumerate}

A localização escolhida é lembrada em
\lstinline|aurora-last-new-project-location| (localStorage) para que a próxima
caixa de diálogo \enquote{New Project} abra na mesma pasta-pai. Em seguida o
modal chama \lstinline|projectManager.loadProject(spfPath)|.

\subsection{Abrir projeto}

Há vários gatilhos de abertura, todos convergindo para
\lstinline|loadProject| em \path{js/project/project_manager.js}:

\begin{itemize}[leftmargin=1.6em]
  \item botões \enquote{Open Project} (toolbar e tela de boas-vindas), que
    chamam \lstinline|electronAPI.showOpenDialog|;
  \item clique numa linha de projeto recente (tela de boas-vindas);
  \item abertura via \enquote{File > Open} / atalhos
    (\lstinline|onSimulateOpenProject|).
\end{itemize}

O handler do processo principal \lstinline|project:open|
(\path{main/ipc/project.js}) faz o trabalho pesado:

\begin{enumerate}[leftmargin=1.8em]
  \item Corrige o caminho se o \path{.spf} não existir no formato esperado
    (tenta \lstinline|<root>/<nome>/<nome>.spf|).
  \item Define \lstinline|state.currentOpenProjectPath = spfPath| (fonte de
    verdade no main).
  \item Registra nos recentes (\path{main/recents.js}) e reconstrói a
    \emph{jumplist}; chama também \lstinline|app.addRecentDocument|.
  \item Lê e faz parse tolerante do \path{.spf}, carimba
    \lstinline|metadata.lastOpened|.
  \item \textbf{Reconcilia o \lstinline|basePath|}: alinha sempre com
    \lstinline|path.dirname(spfPath)|. Se diferir, faz \lstinline|deepRemapPaths|
    de todo caminho absoluto que ainda apontava para a raiz antiga (caso de
    projeto copiado/movido) — ver \autoref{sec:rename}.
  \item Para cada processador, anexa \lstinline|exists| (existência da pasta no
    disco).
  \item Reescreve o \path{.spf} (já reconciliado) e devolve
    \lstinline|{ projectData, files, spfPath }|, enviando também os eventos IPC
    \lstinline|project:processorHubState| e \lstinline|project:processors| para
    a janela que originou a requisição.
\end{enumerate}

De volta no renderer, \lstinline|loadProject| orquestra a UI:

\begin{enumerate}[leftmargin=1.8em]
  \item resolve o \lstinline|basePath| de forma tolerante (várias formas de
    payload) e chama \lstinline|ProjectStore.setProject(spfPath, basePath)|;
  \item \lstinline|projectTreeManager.reset()| — \emph{clean slate} antes da
    nova árvore (uma troca direta projeto$\rightarrow$projeto não passa pelo
    fluxo de fechar);
  \item \lstinline|setAvailableProcessors(structure.processors)| — semeia a
    lista global de processadores ANTES de renderizar;
  \item atualiza o nome na UI, fecha todas as abas
    (\lstinline|TabManager.closeAllTabs|);
  \item \lstinline|projectTreeManager.activateTree()| popula e renderiza a
    árvore;
  \item inicia o \emph{watcher} de arquivos, registra o projeto nos recentes da
    tela de boas-vindas e habilita os botões de compilação;
  \item sintetiza \lstinline|aurora:spf-changed| (\lstinline|source: 'project-loaded'|)
    para forçar a reaplicação do estado derivado do \path{.spf} (testbench,
    top-level, processadores) em todos os assinantes;
  \item se o \path{.spf} referencia arquivos que sumiram do disco, gera o
    relatório \path{.aurora-missing-files.log} (aberto no Monaco como preview) e
    uma notificação; caso contrário, remove um log obsoleto de corrida anterior.
\end{enumerate}

\subsection{Fechar projeto}

\path{js/project/close_project.js} trata o botão de fechar: pede confirmação
(\lstinline|showDialog|), chama \lstinline|electronAPI.closeProject| (handler
\lstinline|project:close| no main, que zera \lstinline|currentOpenProjectPath| e
emite \lstinline|project:processors|/\lstinline|project:fileTree|/\lstinline|project:closed|
vazios), fecha todas as abas, limpa a interface
(\lstinline|clearProjectInterface| $\rightarrow$ \lstinline|treeView.clearAll()|
+ card de \enquote{no project}), chama \lstinline|ProjectStore.clearProject()|,
esquece o último projeto aberto (\lstinline|appInitializer.clearLastProject|) e
faz \lstinline|projectTreeManager.reset()| para que a próxima abertura dispare
um \emph{re-activate} completo.

\section{Processadores}
\label{sec:processadores}

Um processador \tool{SAPHO} corresponde a uma subpasta da raiz do projeto, com
três subpastas reconhecidas: \path{Software/}, \path{Hardware/} e
\path{Simulation/}. A lista canônica de processadores do projeto é
\lstinline|structure.processors| no \path{.spf} (cada entrada é
\lstinline|{ name }|); em memória, no renderer, a lista de nomes vive em
\lstinline|window.availableProcessors|.

\subsection{A lista global de processadores}

O módulo \file{processor\_list} (\path{js/project/processor_list.ts}, compilado
em \path{.js}) é o dono único de \lstinline|window.availableProcessors| e expõe:

\begin{itemize}[leftmargin=1.6em]
  \item \lstinline|setAvailableProcessors(list)| — substitui a lista inteira;
    aceita arrays de strings ou de \lstinline|{name}| e faz \emph{dedup}
    \emph{case-insensitive} (pastas no Windows são \emph{case-insensitive});
  \item \lstinline|getAvailableProcessors()| — sempre retorna array;
  \item \lstinline|addAvailableProcessor(name)| — acrescenta um nome se ainda
    não presente (match \emph{case-insensitive}).
\end{itemize}

\subsection{Criar processador}

A criação é feita pelo \emph{Processor Hub} (\path{js/processors/processor_hub.js}):
um formulário modal com validação ao vivo. As regras de validação são:

\begin{longtable}{p{0.30\linewidth} p{0.62\linewidth}}
  \toprule
  \textbf{Campo} & \textbf{Regra} \\
  \midrule
  \endhead
  \lstinline|name|         & \lstinline|/^[a-zA-Z0-9_-]+$/| (letras, números, hífen, sublinhado). \\
  \lstinline|nBits|        & Inteiro positivo; deve satisfazer a consistência abaixo. \\
  \lstinline|nbMantissa|   & Inteiro positivo. \\
  \lstinline|nbExponent|   & Inteiro positivo. \\
  \lstinline|gain|         & Potência de 2 (\lstinline|(num & (num-1)) === 0|). \\
  \lstinline|instructionStackSize| & Inteiro positivo. \\
  \lstinline|dataStackSize|        & Inteiro positivo. \\
  \lstinline|inputPorts|   & Inteiro positivo. \\
  \lstinline|outputPorts|  & Inteiro positivo. \\
  \bottomrule
\end{longtable}

A consistência de bits exige \lstinline|nBits === nbMantissa + nbExponent + 1|
(\lstinline|checkBitConsistency|). O botão \enquote{Generate} só fica habilitado
quando todos os campos passam (\lstinline|validateAll|).

No envio, o renderer chama \lstinline|electronAPI.createProcessorProject|; o
handler \lstinline|create-processor-project| (\path{main/ipc/project.js}):

\begin{enumerate}[leftmargin=1.8em]
  \item re-valida o nome (defesa contra \emph{path traversal}: o nome vai direto
    para \lstinline|path.join|, então é confinado a um único segmento);
  \item cria \path{<proc>/Software}, \path{<proc>/Hardware} e
    \path{<proc>/Simulation};
  \item escreve \path{Software/<proc>.cmm} com um cabeçalho de diretivas
    \tool{SAPHO} preenchido a partir do formulário;
  \item acrescenta \lstinline|{ name }| em \lstinline|structure.processors| do
    \path{.spf} (com \emph{dedup case-insensitive}) e reescreve o arquivo;
  \item emite o evento IPC \lstinline|processor:created| para o renderer.
\end{enumerate}

O cabeçalho \file{.cmm} gerado tem a forma:

\begin{lstlisting}[style=cmm]
#PRNAME <nome>
#NUBITS <nBits>
#NDSTAC <dataStackSize>
#SDEPTH <instructionStackSize>
#NUIOIN <inputPorts>
#NUIOOU <outputPorts>
#NBMANT <nbMantissa>
#NBEXPO <nbExponent>
#NUGAIN <gain>

void main()
{
    // ...
}
\end{lstlisting}

\begin{nota}
As diretivas \lstinline|#NUBITS|, \lstinline|#NBMANT| etc. no \file{.cmm} são a
fonte \emph{canônica} dos parâmetros do processador para o pipeline de
compilação. O \path{.spf} guarda apenas \lstinline|{ name }|; os números vivem
no \file{.cmm}. A semântica das diretivas é detalhada nos capítulos do
\tool{YANC}.
\end{nota}

\subsection{Renomear processador}
\label{sec:rename}

O handler \lstinline|rename-processor| (\path{main/ipc/project.js}) renomeia o
processador em todas as superfícies internas do \tool{SAPHO}, de forma atômica:

\begin{enumerate}[leftmargin=1.8em]
  \item libera os \emph{watchers} sob a raiz do projeto
    (\lstinline|releaseWatchersUnder|) — sem isso o \tool{chokidar} mantém um
    \emph{handle} e a renomeação da pasta falharia com EPERM/EBUSY no Windows;
  \item move a pasta \path{<root>/<old>} $\rightarrow$ \path{<root>/<new>}
    (\lstinline|moveWithRetry|; renomeação só de \emph{case} usa um salto por
    pasta temporária);
  \item renomeia os artefatos nomeados pelo processador:
    \path{Software/<old>.cmm}, \path{Software/<old>.asm},
    \path{Hardware/<old>.v}, \path{Simulation/<old>_tb.v};
  \item corrige a diretiva \lstinline|#PRNAME| no \file{.cmm} — \emph{apenas} a
    linha da diretiva; comentários e código do usuário não são tocados;
  \item atualiza a entrada em \lstinline|processors[]| preservando a
    configuração por processador;
  \item faz \lstinline|remapProcessorPath| de qualquer referência do \path{.spf}
    (\lstinline|topLevelFile|, \lstinline|testbenchFile|,
    \lstinline|synthesizableFiles|, \lstinline|testbenchFiles|) que apontava para
    dentro da pasta antiga;
  \item emite \lstinline|project:processors| e \lstinline|processor:renamed|.
\end{enumerate}

\subsection{Deletar processador}

O handler \lstinline|delete-processor| remove a pasta do processador do disco
(\lstinline|fse.remove|), filtra a entrada de \lstinline|structure.processors| no
\path{.spf} e emite \lstinline|project:processors| atualizado. No renderer, a
opção \enquote{Delete processor} aparece no menu de contexto do separador de
processador na árvore (\lstinline|_deleteProcessorByName| em
\path{project_tree_actions.js}, que chama \lstinline|electronAPI.deleteProcessor|).

\subsection{Agrupamento e identificação na árvore}

Na view de arquivos (\enquote{verilog picker}), os arquivos são agrupados por
processador. A função \lstinline|_getProcessorForFile| em \path{file_mode.js}
decide o \enquote{dono} de um arquivo: o caminho precisa cair em
\path{<projeto>/<proc>/{Hardware|Software|Simulation}/<arquivo>} e
\lstinline|<proc>| precisa estar em \lstinline|window.availableProcessors|
(comparação \emph{case-insensitive}). Cada grupo é precedido por um separador
horizontal rotulado com o nome do processador
(\lstinline|_createProcessorSeparator| em \path{project_tree_render.js}); o
grupo de arquivos importados que não pertencem a nenhum processador aparece sob
um separador \enquote{Imported} (apenas quando há também grupos de processador).

\section{As três views da árvore de arquivos}
\label{sec:views}

A árvore de arquivos da \tool{AURORA} pode se apresentar de três formas, todas
dentro do mesmo elemento \lstinline|#file-tree|:

\begin{description}[leftmargin=2em,style=nextline]
  \item[\enquote{verilog}] o \emph{verilog picker} — a view de arquivos do
    projeto, agrupada por processador, com indicadores de synth/testbench. É a
    view padrão e também hospeda os cartões de estado vazio.
  \item[\enquote{hierarchy}] a árvore de instâncias de módulo produzida pelo
    \tool{Yosys} (resultado de compilação).
  \item[\enquote{standard}] uma árvore simples de pastas/arquivos enraizada na
    pasta do projeto, lendo filhos sob demanda (\emph{lazy}).
\end{description}

\subsection{Subcontêineres físicos e o \texttt{tree\_view}}

O controlador de baixo nível \file{tree\_view} (\path{js/tree/tree_view.js})
cria três subcontêineres DOM fisicamente separados dentro de
\lstinline|#file-tree| — um por view — e governa a visibilidade por um único
atributo \lstinline|data-active-view| (a CSS em \path{css/tree/file_tree.css}
mostra/oculta com base nele). Cada \emph{renderer} escreve apenas no \emph{seu}
subcontêiner — eles não podem colidir.

\begin{lstlisting}
<div id="file-tree" data-active-view="verilog">
  <div class="tree-view tree-view-verilog">...</div>
  <div class="tree-view tree-view-hierarchy">...</div>
  <div class="tree-view tree-view-standard">...</div>
</div>
\end{lstlisting}

A API é \lstinline|getContainer(name)| (devolve o subcontêiner), \lstinline|setActive(name)|
(troca a view visível — uma única mudança de atributo, sem mutação de DOM),
\lstinline|getActive()|, \lstinline|clear(name)| e \lstinline|clearAll()|. O
array \lstinline|VIEW_NAMES| congela os três nomes.

\subsection{O controlador único de views}

O \file{file\_tree\_view\_controller} (\path{js/tree/file_tree_view_controller.js})
é o dono único do par \enquote{qual view está ativa} + o \emph{listener} do
botão de alternância (exatamente um, anexado uma vez). Ele:

\begin{itemize}[leftmargin=1.6em]
  \item mantém o nome da view ativa (\lstinline|_activeView|) e os dados de
    hierarquia (\lstinline|_hierarchyData|);
  \item registra um \emph{renderer} por view (\lstinline|registerRenderer|); ao
    ativar a view, invoca o renderer correspondente;
  \item expõe \lstinline|showFileMode()|, \lstinline|showHierarchyMode()| e
    \lstinline|showStandardMode()|, além dos predicados
    \lstinline|isShowingFileMode/Hierarchy/Standard|;
  \item assina o \file{ProjectStore} (\lstinline|_wireProjectStore|): a view
    \enquote{standard} só entra no ciclo com um projeto aberto, então o botão
    habilita ao abrir e desabilita ao fechar (deixando \enquote{standard} para
    não ficar numa árvore de pastas vazia).
\end{itemize}

O conjunto de views pelas quais o botão rotaciona é dinâmico
(\lstinline|_cyclableViews|): \enquote{verilog} sempre presente, \enquote{hierarchy}
apenas com dados compilados, \enquote{standard} apenas com projeto aberto. O
botão fica habilitado quando há mais de uma view; um clique avança para a
próxima view do ciclo (\lstinline|_nextView|, com \emph{wrap}). Os renderers
registrados são delegadores que buscam o gerente correto em tempo de chamada:
\enquote{verilog} chama \lstinline|window.projectTreeManager.renderTree()|,
\enquote{hierarchy} chama \lstinline|renderHierarchicalTree| do
\lstinline|CompilationModule| mais recente e \enquote{standard} chama
\lstinline|window.standardTreeRenderer.render()|.

\subsection{Renderer da view \enquote{verilog}}

A renderização do \emph{verilog picker} fica no \file{RenderMixin}
(\path{js/project/project_tree_render.js}), misturado em
\lstinline|ProjectTreeManager.prototype| por \path{file_mode.js}. O
\lstinline|renderTree()| é um \emph{reconciler keyed-por-path}: nunca faz
\lstinline|innerHTML = ''|; compara \lstinline|this.verilogFiles| (chaveado por
caminho) contra as linhas existentes — linhas cujo caminho sumiu são removidas,
linhas ainda desejadas são atualizadas \emph{in-place}
(\lstinline|_updateFileItem|), e caminhos sem linha são criados
(\lstinline|_createFileItem|) e posicionados via \lstinline|insertBefore|. É
idempotente: o mesmo input duas vezes não causa mutações na segunda chamada.

As linhas software (\file{.cmm} em \path{Software/}) não recebem botão de
remoção (são geridas pelo processador). Arquivos importados que sumiram do disco
são exibidos num cartão de aviso (\lstinline|_renderMissingFilesNotice|).

\subsection{Renderer da view \enquote{standard}}

O \file{standard\_tree\_render} (\path{js/tree/standard_tree_render.js})
desenha a árvore de pastas enraizada em \lstinline|window.currentProjectPath|,
ao modo de um explorador de arquivos genérico. As leituras são \emph{lazy}: só
a raiz é listada de início; os filhos de uma pasta são buscados
(\lstinline|electronAPI.getFolderFiles|) na primeira expansão e mantidos no DOM
para reexpansão instantânea. Caminhos expandidos são lembrados num
\lstinline|Set| (\lstinline|_expanded|) para sobreviver a um re-render. A
ordenação é \enquote{pastas primeiro, depois alfabética} (\lstinline|sortEntries|).
Arquivos ignorados (\lstinline|isIgnored|): \emph{dotfiles}, o próprio
\path{.spf} e blobs de configuração legados
(\lstinline|projectOriented.json|, \lstinline|processorConfig.json|,
\lstinline|fileOriented.json|). O método \lstinline|revealFolder| expande os
ancestrais até uma pasta-alvo e a destaca; \lstinline|expandAll| /
\lstinline|collapseAll| operam sobre toda a árvore.

\begin{nota}
Há um único caminho de abertura de arquivo, idêntico ao da view \enquote{verilog}
(\lstinline|readFile| $\rightarrow$ painel dividido em foco ou \emph{preview} do
\file{TabManager}; o duplo-clique promove a aba a permanente). Não há um segundo
\emph{handler} de abertura — a classe de bug que motivou a remoção da antiga
view \enquote{standard} não retorna.
\end{nota}

\section{Marcação de top-level e testbench}

A categoria \emph{synth} vs \emph{testbench} de cada arquivo Verilog é
\emph{auto-detectada} do conteúdo pelo classificador
(\path{js/project/verilog_classifier.ts}), não por uma marca manual. O
classificador é um \emph{scorer} com limiar \lstinline|TB_THRESHOLD = 3|: cada
sinal definitivo (dump de VCD via \lstinline|$dump*|, terminação de simulação
\lstinline|$finish|/\lstinline|$stop|, módulo sem portas) soma \lstinline|+3| e
já passa o limiar sozinho; o bloco \lstinline|initial| soma \lstinline|+2|
(forte, mas não classifica isolado, pois código sintetizável pode tê-lo); e os
sinais fracos (\lstinline|$display| e afins, atrasos procedurais como
\lstinline|#10|, hint de nome \lstinline|tb|/\lstinline|test|/\lstinline|testbench|)
somam \lstinline|+1| ou \lstinline|+2| cada. Em dúvida, o default seguro é
\emph{synthesizable}. Arquivos \file{.py} (cocotb) são sempre testbench
(decisão em \lstinline|_classifyAll| de \path{file_mode.js}, antes do
\emph{scorer}).

A escolha de \emph{top-level} e \emph{testbench top}, porém, é explícita do
usuário (menu de contexto na linha do arquivo, em \path{project_tree_actions.js}).
A marca \lstinline|isTopLevel| é relativa à categoria atual do arquivo: um synth
com \lstinline|isTopLevel=true| é o \emph{synth top}; um testbench com
\lstinline|isTopLevel=true| é o \emph{testbench top}. As marcações são
persistidas no \path{.spf} via \lstinline|SpfStore.update| (campos
\lstinline|topLevelFile| / \lstinline|testbenchFile| e as listas
\lstinline|synthesizableFiles| / \lstinline|testbenchFiles|), e a remoção de um
arquivo limpa essas referências quando apontavam para o caminho removido.

\section{Tela de boas-vindas e projetos recentes}
\label{sec:recentes}

Há \emph{duas} listas de recentes, com propósitos distintos:

\begin{description}[leftmargin=2em,style=nextline]
  \item[Tela de boas-vindas (renderer)] gerida por
    \file{RecentProjectsManager} (\path{js/project/recent_projects.js}),
    persistida em \lstinline|localStorage| sob a chave
    \lstinline|aurora-recent-projects| (até 10 entradas, cada uma com
    \lstinline|name|, \lstinline|path|, \lstinline|lastOpened|). Alimenta o
    componente \lstinline|<aurora-welcome>| (Lit), que renderiza a etapa
    Start/Recent, anima as linhas e mostra um \emph{preview} dos processadores
    de cada projeto (lido sob demanda de cada \path{.spf} por
    \lstinline|_readProcessors|, cacheado em \lstinline|p._procs|).
  \item[\emph{Jumplist} do Windows (processo principal)] gerida por
    \path{main/recents.js}, persistida num único JSON em \lstinline|userData|
    (\path{aurora-recents.json}), também limitada a 10 (\lstinline|MAX_RECENTS|).
    A API é \lstinline|read|, \lstinline|push| (insere no topo,
    deduplicado/normalizado), \lstinline|prune| (remove entradas cujo \path{.spf}
    não existe mais) e a constante \lstinline|MAX_RECENTS|.
\end{description}

A separação existe porque o processo principal não alcança o
\lstinline|localStorage|, e a lista no main permite limpar os resíduos
\enquote{Electron} de execuções antigas da \emph{jumplist} gerida pela shell do
Windows. Os recentes do main são atualizados em toda abertura/criação de
projeto, junto da reconstrução da \emph{jumplist} (\lstinline|rebuildJumpList|).
O handler \lstinline|list-recent-projects| devolve a lista do main após
\lstinline|prune|.

A \file{RecentProjectsManager} valida a existência do \path{.spf} antes de abrir
(\lstinline|checkProjectExists|); se sumiu, remove da lista e avisa o usuário.
Ao clicar numa linha, fecha as abas abertas e delega para o callback de abertura
injetado (\lstinline|loadProject|).

\section{Backups (zip sob demanda)}

A \tool{AURORA} produz backups \file{.zip} sob demanda (botão
\lstinline|backupFolderBtn| na toolbar, habilitado com projeto aberto). O
handler \lstinline|create-backup| (\path{main/ipc/files.js}):

\begin{enumerate}[leftmargin=1.8em]
  \item valida o \lstinline|folderPath| (\lstinline|safePath|);
  \item garante a pasta \path{Backup/} dentro do projeto e uma pasta de
    \emph{staging} temporária \path{backup_<timestamp>/};
  \item copia todo o conteúdo do projeto para a \emph{staging} (exceto a própria
    \path{Backup/} e a \emph{staging});
  \item gera \path{Backup/<nome>_<timestamp>.zip} via
    \lstinline|Compress-Archive| do PowerShell (nativo do Windows, sem depender
    de 7-Zip no PATH);
  \item remove a pasta de \emph{staging} em qualquer caminho (sucesso ou
    falha), devolvendo \lstinline|{ success, message }|.
\end{enumerate}

O \emph{quoting} para o PowerShell escapa aspas simples (\lstinline|'| para
\lstinline|''|), e o comando roda com \lstinline|-NoProfile -NonInteractive|.

\section{Estrutura de um projeto no disco}

O diagrama abaixo resume a estrutura típica de um projeto com um processador
\lstinline|cpu| e a pasta de backups:

\begin{lstlisting}
demo/                          <- pasta do projeto (== nome do .spf)
|-- demo.spf                   <- arquivo de projeto (JSON: metadata + structure)
|-- cpu/                       <- processador SAPHO
|   |-- Software/
|   |   |-- cpu.cmm            <- fonte C+-  (#PRNAME, #NUBITS, ...)
|   |   `-- cpu.asm            <- gerado (nao aparece na arvore)
|   |-- Hardware/
|   |   `-- cpu.v              <- Verilog sintetizavel (synth top)
|   `-- Simulation/
|       `-- cpu_tb.v           <- testbench (testbench top)
|-- (arquivos .v importados)   <- grupo "Imported" na arvore
`-- Backup/
    `-- demo_2026-06-21.zip    <- backup sob demanda
\end{lstlisting}

\begin{nota}
A categoria de cada \file{.v} (synth ou testbench) é decidida pelo conteúdo, não
pela subpasta: \lstinline|_discoverProcessorFiles| em \path{file_mode.js} varre
as três subpastas (\path{Hardware}, \path{Software}, \path{Simulation}) e aceita
\file{.v}/\file{.sv}/\file{.vh} como sintetizáveis e \file{.cmm} como software,
independentemente da pasta; o \lstinline|_classifyAll| posteriormente reclassifica
cada Verilog. O \file{.asm} é gerado e re-descoberto a cada \emph{load}, ficando
fora da árvore para não gerar laço com o \emph{watcher}.
\end{nota}

% !TEX root = ../main.tex
\chapter{AURORA: o subsistema de compilação}
\label{cap:08-aurora-compilacao}

Este capítulo documenta como a IDE \tool{AURORA} orquestra a compilação e a
simulação de soft-processors. O foco é o caminho que vai do clique num botão da
barra de ferramentas até o processo-filho real disparado pela toolchain
(\tool{cmmcomp}, \tool{asmcomp}, \tool{iverilog}, \tool{vvp}, \tool{verilator},
\tool{yosys}, \tool{gtkwave}, etc.). A documentação que segue foi verificada
diretamente contra o código-fonte; quando o texto descreve um comportamento, ele
corresponde ao que os arquivos \path{js/compilation/*}, \path{main/compile/*} e
\path{main/ipc/compile.js} efetivamente fazem.

A arquitetura tem duas metades bem separadas, ligadas por IPC:

\begin{description}[leftmargin=1.4em]
  \item[Renderer (processo da janela).] Decide \emph{o que} rodar. Os botões
    expandem para sequências de etapas; cada etapa é montada por um
    \emph{builder} puro que devolve um \code{CommandSpec} estruturado
    (\lstinline|{ binary, args[], cwd, env, prependPath }|). Vive sob
    \path{js/compilation/}.
  \item[Main (processo principal do Electron).] Decide \emph{como} rodar. O
    \emph{executor} (\path{main/compile/executor.js}) é o único ponto que
    efetivamente chama \lstinline|spawn(binary, args, { shell:false })|, depois de
    validar o spec contra uma \emph{allowlist} de binários e contra as
    \emph{protected flags}. Vive sob \path{main/compile/} e
    \path{main/ipc/compile.js}.
\end{description}

A fronteira de confiança é nítida: o renderer nunca executa nada por conta
própria; ele compõe um spec e o envia por IPC. O main re-valida tudo antes de
gerar o processo-filho.

\begin{nota}
A modelagem em \code{CommandSpec} substituiu o antigo idioma de concatenar uma
\emph{string de shell} no renderer. Como o spawn usa \lstinline|shell:false|,
cada argumento chega ao processo-filho exatamente como está no array
\code{args}; não há parsing de shell, portanto não há bug de aspas nem vetor de
injeção. É justamente isso que permite que o subsistema de \emph{overrides}
dirigido pela Aurora Intelligence (\autoref{cap:11-aurora-intelligence}) exista
com segurança.
\end{nota}

\section{Visão geral do fluxo}
\label{sec:08-visao-geral}

O ciclo de vida de uma etapa de compilação percorre sempre o mesmo encadeamento
de módulos. O diagrama a seguir resume o caminho de uma invocação, do clique ao
processo-filho e de volta ao terminal.

\begin{lstlisting}[basicstyle=\ttfamily\footnotesize]
 [botao toolbar]  -> window.AuroraAPI.compile.compileStep('wave')
        |
        v
 CompilationFlowManager.runSingleStep('wave')      (compilation_flow.js)
        |  expande o botao numa sequencia de etapas
        v
 CompilationModule.<fase>()                         (compilation_module.js)
        |  monta o CommandSpec via um builder puro
        v
 builders/<tool>.js  ->  { step, binary, args[], cwd, env?, prependPath? }
        |
        v
 spec_runner.runSpec / runSpecStreamed             (spec_runner.js)
        |  aplica overrides (command_overrides.js), loga o diff
        v
 electronAPI.execSpec / execSpecStreamed   ==IPC==> main
        |
        v
 executor.js  validateSpecForExec:                  (main/compile/executor.js)
        |   1) binario na allowlist  (binary_allowlist.js)
        |   2) protected flags ok    (protected_flags.js)
        |   3) shape do spec valido
        v
 spawnTracked(binary, args, { cwd, env, shell:false })
        |
        |  stdout/stderr  ==IPC (one-shot ou stream)==>  terminal (renderer)
        v
 { code, stdout, stderr, pid }
\end{lstlisting}

\subsection{Os módulos e seus papéis}

\begin{longtable}{p{0.34\linewidth} p{0.58\linewidth}}
  \toprule
  \textbf{Arquivo} & \textbf{Responsabilidade} \\
  \midrule
  \endhead
  \path{js/compilation/compilation_flow.js} &
    Handlers dos botões. Expande cada botão na sequência de etapas, gerencia
    cancelamento, troca de terminal e o \emph{gating} (habilitar/desabilitar)
    dos botões conforme o estado do \path{.spf}. \\
  \path{js/compilation/compilation_module.js} &
    A classe \code{CompilationModule} — o ``backend'' dos botões. Cada método
    público é uma fase (cmm, asm, verilog, wave, etc.). Monta os specs e os
    despacha. \\
  \path{js/compilation/processor_compiler.js} &
    As fases \code{cmmCompilation} e \code{asmCompilation} (impuras: rodam
    \tool{cmmcomp}/\tool{appcomp}/\tool{asmcomp}, salvam buffers, copiam
    artefatos). \\
  \path{js/compilation/builders/*} &
    Builders puros: dado um contexto, devolvem um \code{CommandSpec} sem efeito
    colateral nem I/O. \\
  \path{js/compilation/command_spec.js} &
    O contrato \code{CommandSpec}: IDs de etapa, \code{applyOverride},
    \code{formatSpec}, \code{diffSpecs}, \code{validateShape}. \\
  \path{js/compilation/command_overrides.js} &
    Registro de overrides (efêmeros e persistidos) que a Aurora Intelligence
    aplica sobre os specs. \\
  \path{js/compilation/spec_factory.js} &
    Constrói o spec-base de qualquer etapa \emph{sem} executá-lo (para
    pré-visualização pela IA). \\
  \path{js/compilation/spec_runner.js} &
    O ponto único renderer-side que aplica overrides e chama o executor via
    IPC (\code{runSpec} / \code{runSpecStreamed}). \\
  \path{main/compile/executor.js} &
    O runner real: \lstinline|spawn(..., shell:false)|, monta o ambiente do
    filho, valida e faz streaming. \\
  \path{main/compile/binary_allowlist.js} &
    Conjunto fechado de binários que o executor aceita gerar. \\
  \path{main/compile/protected_flags.js} &
    Flags invariantes por etapa que um override não pode remover. \\
  \path{main/compile/python_locator.js} &
    Localização do interpretador Python para o fluxo cocotb. \\
  \path{main/ipc/compile.js} &
    Handlers IPC auxiliares: lançar GTKWave/Surfer (detached), cancelar
    processos, decodificar números complexos. \\
  \bottomrule
\end{longtable}

\section{Os botões da toolbar e suas sequências}
\label{sec:08-botoes}

Todos os botões da toolbar passam por \code{window.AuroraAPI.compile} —
intencionalmente, para que o clique do usuário percorra exatamente o mesmo
caminho que a Aurora Intelligence usa via \emph{function-calling}. Os listeners
são registrados em \code{CompilationFlowManager.initialize}:

\begin{lstlisting}[language=Java]
document.getElementById('cmmcomp')?.addEventListener('click',  () => window.AuroraAPI?.compile.compileStep('cmm'));
document.getElementById('vericomp')?.addEventListener('click', () => window.AuroraAPI?.compile.compileStep('verilog'));
document.getElementById('wavecomp')?.addEventListener('click', () => window.AuroraAPI?.compile.compileStep('wave'));
document.getElementById('prismcomp')?.addEventListener('click',() => window.AuroraAPI?.compile.compileStep('prism'));
document.getElementById('verilatorproc')?.addEventListener('click',() => window.AuroraAPI?.compile.compileStep('verilator-proc'));
document.getElementById('fastsim')?.addEventListener('click',() => window.AuroraAPI?.compile.compileStep('verilator-fast'));
document.getElementById('allcomp')?.addEventListener('click',  () => window.AuroraAPI?.compile.compileAll());
document.getElementById('cancel-everything')?.addEventListener('click', () => window.AuroraAPI?.compile.cancel());
\end{lstlisting}

\subsection{Filosofia: cada botão é \emph{self-contained}}

O cabeçalho de \path{compilation_flow.js} descreve a filosofia adotada como
\enquote{lazy do ponto de vista do usuário}: cada botão é auto-suficiente, e o
usuário não precisa lembrar a ordem dos cliques nem rodar uma compilação prévia
antes de pedir \tool{Wave} ou \tool{PRISM}. Aceita-se re-trabalho (recompilar a
cada clique) em troca de previsibilidade. Cada botão expande para a mesma
sequência-base:

\begin{lstlisting}[basicstyle=\ttfamily\footnotesize]
 Verilog =        cmm + asm + iverilog(-tnull, top-level)
 PRISM   = Verilog                                        + yosys (via IPC PRISM)
 Wave    =        cmm + asm + iverilog(-o vvp, testbench) + vvp + gtkwave
\end{lstlisting}

Os laços que rodam \code{cmm}\,+\,\code{asm} iteram sobre
\code{window.availableProcessors}. Em um projeto de Verilog puro (sem
processadores), o array é vazio e o laço é um \emph{no-op} natural — não há
\lstinline|if (hasProcessor)| no pipeline. Esse princípio (laço sobre array
vazio em vez de ramificação estrutural) é recorrente em todo o subsistema.

\subsection{Mapa de botões}

\begin{longtable}{p{0.16\linewidth} p{0.20\linewidth} p{0.56\linewidth}}
  \toprule
  \textbf{Etapa} & \textbf{Handler} & \textbf{Sequência expandida} \\
  \midrule
  \endhead
  \code{cmm} & \code{handleCmmStep} &
    \tool{cmmcomp} ($\to$ \path{Software/<base>.asm} + \path{cmm_log.txt}),
    depois \tool{asmcomp} ($\to$ \path{Hardware/<proc>.v} + memória +
    \path{Simulation/<proc>_tb.v}). Exige um \path{.cmm} em foco (ou inferível). \\
  \code{asm} & \code{handleAsmStep} &
    \tool{asmcomp} (\emph{sem} \tool{cmmcomp}) + \tool{iverilog -tnull}. Etapa
    sem botão na toolbar; existe para a Aurora Intelligence testar um
    \path{.asm} otimizado à mão sem regenerá-lo a partir do \path{.cmm}. \\
  \code{verilog} & \code{handleVerilogStep} &
    Pré-compila cmm+asm de todos os processadores, depois
    \tool{iverilog -tnull} com o top-level (sem testbench, sem \path{.vvp}). \\
  \code{wave} & \code{handleWaveStep} &
    Pré-compila cmm+asm, depois \code{runGtkWave} (que faz
    \tool{iverilog -o vvp} com testbench, roda \tool{vvp} e abre o visualizador). \\
  \code{prism} & \code{handlePrismStep} &
    Tudo que \code{verilog} faz e em seguida \tool{yosys} via IPC PRISM. PRISM é
    superset de Verilog. \\
  \code{verilator-proc} & \code{handleVerilatorProcStep} &
    Pré-compila cmm+asm e roda o top-level \path{<proc>.v} gerado sob
    \tool{Verilator} com a fiação previsível do processador SAPHO
    (req\_in/out\_en one-hot, I/O por arquivo). \\
  \code{verilator-fast} & \code{handleFastSimStep} &
    Roda o testbench via \tool{Verilator} binário \emph{sem} gerar onda nem
    abrir o visualizador (\code{runFastSim}). \\
  \code{all} & \code{runAll} &
    Full Build: \code{precompileAllProcessors} + \code{runGtkWave}. \\
  \bottomrule
\end{longtable}

\subsection{Gating dos botões pelo estado do design}
\label{sec:08-gating}

\code{syncToolbarEnabledState} habilita/desabilita os botões lendo a mesma fonte
de verdade da barra de status (o \code{SpfStore}). As regras observadas no
código:

\begin{itemize}[leftmargin=1.4em]
  \item \textbf{top-level definido} $\to$ habilita \code{vericomp},
    \code{prismcomp};
  \item \textbf{processador ativo} (\path{.cmm} em foco) $\to$ habilita
    \code{verilatorproc};
  \item \textbf{testbench definido} $\to$ habilita \code{wavecomp},
    \code{waveConfigBtn}, \code{cancel-everything};
  \item \textbf{Fast Sim} (\code{fastsim}) habilita quando há testbench
    \emph{e} ou ele é \path{.py} (cocotb roda em qualquer engine) ou o simulador
    selecionado é Verilator (\lstinline|getSimulator() === 'verilator'|).
\end{itemize}

\noindent O botão C± (\code{cmmcomp}) tem regra própria
(\code{syncCmmcompEnabled}): só fica habilitado quando o arquivo em foco no
Monaco termina em \path{.cmm}. A re-sincronização é disparada por eventos
(\lstinline|aurora:editing-file-changed|, \lstinline|aurora:spf-changed|,
\lstinline|aurora:wave-simulator-changed|, criação/remoção de processador).

\section{Os builders}
\label{sec:08-builders}

Cada builder em \path{js/compilation/builders/} é uma função pura: recebe um
objeto de contexto (caminhos, configuração, valores derivados) e devolve um
\code{CommandSpec}, sem I/O, sem \code{console} e sem efeito colateral. Isso os
torna triviais de testar e fáceis de alimentar às ferramentas
\code{inspect_compile_command} / \code{preview_compile_command} da IA. O
\path{builders/index.js} apenas re-exporta todos.

\begin{nota}
Os arquivos carregados em runtime são os \path{.js}; os \path{.ts} ao lado são a
fonte TypeScript compilada por \lstinline|npm run build:ts|. O comportamento
descrito aqui foi lido dos \path{.js} efetivamente executados.
\end{nota}

\subsection{\file{cmm.js} — \texttt{buildCmmSpec}}

Espelha o call site \code{cmmCompilation}. Monta os argumentos nomeados do
compilador C± (\tool{cmmcomp.exe}). A flag de idioma (\code{-pt}/\code{-en}) vai
\emph{primeiro}, pois o yanc a consome em \code{parse_lang_flag} antes do
parsing das demais opções. \code{-A} (\code{--array}) liga o dump de arrays do
waveform quando \code{showArrays} está ativo.

\begin{lstlisting}[language=Java]
export function buildCmmSpec(ctx) {
  const args = [
    `-${ctx.lang}`,
    '-i', ctx.inputFile,   // .cmm de entrada
    '-n', ctx.baseName,    // nome base do processador
    '-p', ctx.projectPath, // diretorio do processador
    '-m', ctx.macrosPath,  // components/Macros
    '-t', ctx.tempPath,    // components/Temp/<proc>
  ];
  if (ctx.showArrays) args.push('-A');
  return { step: 'cmm', binary: ctx.cmmCompPath, args, cwd: ctx.projectPath,
           processorName: ctx.processorName, label: `cmm: ${ctx.processorName}` };
}
\end{lstlisting}

\subsection{\file{asm.js} — \texttt{buildAsmPreSpec} e \texttt{buildAsmSpec}}

Dois builders para os dois passos do estágio ASM:

\begin{itemize}[leftmargin=1.4em]
  \item \code{buildAsmPreSpec} — o pré-processador de macros (\tool{appcomp.exe}),
    com \code{-i <asm>} e \code{-t <temp>};
  \item \code{buildAsmSpec} — o compilador final (\tool{asmcomp.exe}), com
    \code{-i -p -d -m -t -f -c}. \code{-f} (frequência/clk) e \code{-c}
    (\emph{clocks}) devem ser inteiros — o yanc rejeita valor não-numérico.
\end{itemize}

\begin{lstlisting}[language=Java]
export function buildAsmSpec(ctx) {
  const args = [
    `-${ctx.lang}`,
    '-i', ctx.asmFile,
    '-p', ctx.projectPath,
    '-d', ctx.hdlPath,     // components/HDL
    '-m', ctx.macrosPath,
    '-t', ctx.tempPath,
    '-f', String(ctx.freq),
    '-c', String(ctx.clocks),
  ];
  return { step: 'asm', binary: ctx.asmCompPath, args, cwd: ctx.tempPath,
           processorName: ctx.processorName, label: `asm: ${ctx.processorName}` };
}
\end{lstlisting}

\subsection{\file{iverilog.js} — \texttt{buildIverilogCheckSpec} e \texttt{buildIverilogBuildSpec}}

Dois usos do \tool{iverilog}:

\begin{itemize}[leftmargin=1.4em]
  \item \code{iverilog-check} — \code{-tnull} (verificação de sintaxe e
    elaboração, sem emitir \path{.vvp}); usado pelo botão Verilog e pelo
    \emph{gate} do Wave-Config;
  \item \code{iverilog-build} — \code{-o <vvp>} (build completo, produz o
    \path{.vvp}); usado pelo botão Wave antes de o \tool{vvp} rodar.
\end{itemize}

\noindent Ambos sempre incluem \code{-y <components/HDL>} para que os módulos da
biblioteca SAPHO (\path{processor.v}, \path{ula.v}, \path{myFIFO.v}, etc.)
resolvam sem o usuário precisar listá-los. Ambos também retornam
\code{prependPath} com o diretório do binário, porque o \tool{iverilog} e os
módulos que ele faz \emph{dlopen} são MinGW-linked e dependem das DLLs de runtime
em \path{mingw64/bin}.

\begin{lstlisting}[language=Java]
export function buildIverilogBuildSpec(ctx) {
  const args = [
    '-y', ctx.hdlPath,
    '-s', ctx.simTopModule,   // top-level da simulacao
    '-o', ctx.outputFile,     // <temp>/<simTop>.vvp
    ...ctx.sourceFiles,
  ];
  return { step: 'iverilog-build', binary: ctx.iveriCompPath, args, cwd: ctx.cwd,
           ...prependBinDir(ctx.iveriCompPath), label: ... };
}
\end{lstlisting}

\subsection{\file{vvp.js} — \texttt{buildVvpRunSpec}}

Uma única simulação completa com saída FST: \lstinline|<vvpBin> <vvpFile> -fst|.
O cabeçalho VCD não é mais capturado por um passo separado \emph{header-only}; é
extraído diretamente do FST finalizado (\code{_extractFstHeaderVcd}), de modo
que uma execução produz tudo que o seletor do Wave-Config e o \path{.gtkw}
automático precisam. O \code{cwd} resolve os caminhos relativos de
\code{$readmemb} / \code{$fopen}, dispensando o antigo idioma
\lstinline|cd && vvp|. O builder prepara o diretório do \tool{vvp} no
\code{prependPath}, porque o \tool{vvp} carrega \path{system.vpi} com
\code{LOAD_WITH_ALTERED_SEARCH_PATH} e precisa das DLLs MinGW vizinhas.

\subsection{\file{cocotb.js} — \texttt{buildCocotbRunSpec}}

Para o fluxo cocotb (testbench em Python), a AURORA é dona do \emph{runner
script} e passa todos os dados específicos do projeto por variáveis de ambiente.
O arquivo Python do usuário contém apenas testes cocotb; nenhum \tool{Makefile}
ou boilerplate de runner é exigido no projeto.

\begin{lstlisting}[language=Java]
export function buildCocotbRunSpec(ctx) {
  return {
    step: 'cocotb-run',
    binary: ctx.pythonPath,             // python.exe do bundle
    args: [ctx.runnerScript],           // aurora_cocotb_runner.py
    cwd: ctx.cwd,
    env: ctx.env,                       // AURORA_COCOTB_*, SIM, PYTHONHOME, ...
    prependPath: ctx.prependPath || [], // mingw64/bin (iverilog) + gtkwave-nipscern
    label: 'python aurora_cocotb_runner.py',
  };
}
\end{lstlisting}

\subsection{\file{verilator.js} — pipeline do Verilator}
\label{sec:08-verilator-builders}

O arquivo \path{builders/verilator.js} concentra cinco builders e os ``botões''
de performance do Verilator. O script \code{verilator} é Perl, então é sempre
invocado como \lstinline|perl.exe verilator <args>|, com \path{mingw64/bin} e
\path{usr/bin} prefixados no PATH para que \tool{perl}, \tool{g++}, \tool{make} e
os utilitários do \path{verilated.mk} resolvam. O \path{V<top>.exe} gerado é
aceito pela regra de prefixo da allowlist (vive sob \path{components/Temp/obj_dir_*}).

\paragraph{Botões de performance (constantes do módulo).}

\begin{tabularx}{\linewidth}{Y Y}
  \toprule
  \textbf{Constante} & \textbf{Valor / efeito} \\
  \midrule
  \code{VERILATOR_BUILD_JOBS} & \code{'0'} — um job por núcleo (\code{-j 0}). \\
  \code{VERILATOR_OPT_LEVEL}  & \code{'-O3'} — otimização do g++ do modelo. \\
  \code{VERILATOR_SIM_THREADS}& \code{0} — simulação single-thread por padrão. \\
  \bottomrule
\end{tabularx}

\paragraph{\texttt{buildVerilatorBuildSpec} (fluxo Wave, \texttt{verilator-build}).}
Gera o \path{V<top>.exe} a partir das fontes com \code{--binary --main}. Inclui
\code{--trace-fst} por padrão (o Fast Sim passa \code{trace:false} para suprimir
o dump), \code{--timing} (o testbench dirige o clock por \code{#delay}),
\code{+define+YANC_TRACE} (liga a visibilidade de simulação do harness gerado
pelo YANC), e uma sequência de \code{-CFLAGS} (\code{-O3}, \code{-march=native},
\code{-fstrict-aliasing}, \code{-pipe}, \code{-Wno-attributes}). Para evitar uma
falha do \tool{ccache} empacotado (que pode abortar o build com erro 127), passa
\code{-MAKEFLAGS OBJCACHE=}, forçando o Makefile gerado a chamar \tool{g++}
diretamente. O ambiente carrega \lstinline|LC_ALL=C| para silenciar o aviso de
locale do perl MSYS2.

\paragraph{\texttt{buildVerilatorRunSpec} (\texttt{verilator-run}).}
Roda o \path{V<top>.exe} sem argumentos; o cabeçalho VCD é extraído do FST
finalizado, igual ao fluxo do \tool{vvp}.

\paragraph{Pipeline top-level em três passos (botão Verilator do processador).}
Diferente do fluxo Wave, o harness para o processador CMM roda o Verilator em
três etapas distintas:

\begin{enumerate}[leftmargin=1.6em]
  \item \code{buildVerilatorJsonSpec} (\code{verilator-json}) — \code{--json-only}
    dumpa \path{V<top>.tree.json} (as portas do top-level); nenhum C++ é gerado.
  \item \code{buildVerilatorTbBuildSpec} (\code{verilator-tb-build}) —
    \code{--cc --exe --build} compila o testbench C++ manual (gerado em
    \path{verilator_tb.js}) junto com as fontes. Aqui \emph{não} há
    \code{--timing}: o clock é dirigido à mão pelo harness C++, e o RTL é 100\%
    sintetizável.
  \item \code{buildVerilatorTbRunSpec} (\code{verilator-tb-run}) — roda o
    \path{V<top>.exe}, que lê \path{input_<N>.txt} e escreve
    \path{output_<N>.txt}. \code{+cycles=N} é apenas o teto.
\end{enumerate}

\subsection{\file{yosys.js} — \texttt{buildYosysHierarchySpec} e \texttt{buildPrismYosysSpec}}

Dois usos do \tool{yosys}, ambos como \lstinline|yosys.exe -s <script>|; a única
diferença é qual script o chamador grava em disco antes. \code{yosys-hierarchy}
emite o JSON de hierarquia para o painel do renderer; \code{prism-yosys} faz a
síntese para o esquemático do PRISM (ver \autoref{cap:10-aurora-prism}). O
builder desconhece o conteúdo do script — só o caminho.

\subsection{\file{wave_tools.js} — \texttt{buildFst2VcdSpec} e \texttt{buildGtkwaveSpec}}

\code{buildFst2VcdSpec} monta \lstinline|fst2vcd -f <in> -o <out>|, usado para
extrair um cabeçalho VCD em texto a partir do FST. \code{buildGtkwaveSpec} monta
a linha do \tool{gtkwave} do fork \repo{gtkwave-nipscern}:
\code{--dark --zoom-fit --left-justify <vcd>} (com \code{-a <save>} opcional). O
GTKWave é caso especial: é lançado \emph{detached} e monitorado, não
\emph{exec'd-and-awaited} — o renderer alimenta o spec ao IPC
\code{launch-gtkwave-only} depois de renderizar os args; o pipeline de overrides
ainda roda, só o spawn real é diferente.

\subsection{Lista completa de builders}

\begin{longtable}{p{0.30\linewidth} p{0.16\linewidth} p{0.46\linewidth}}
  \toprule
  \textbf{Builder} & \textbf{Etapa} & \textbf{Binário / efeito} \\
  \midrule
  \endhead
  \code{buildCmmSpec} & \code{cmm} & \tool{cmmcomp.exe}: \path{.cmm} $\to$ \path{.asm}. \\
  \code{buildAsmPreSpec} & \code{asm-pre} & \tool{appcomp.exe}: expande macros. \\
  \code{buildAsmSpec} & \code{asm} & \tool{asmcomp.exe}: \path{.asm} $\to$ \path{.v} + memória + \path{_tb.v}. \\
  \code{buildIverilogCheckSpec} & \code{iverilog-check} & \tool{iverilog -tnull}: sintaxe/elaboração. \\
  \code{buildIverilogBuildSpec} & \code{iverilog-build} & \tool{iverilog -o}: produz \path{.vvp}. \\
  \code{buildVvpRunSpec} & \code{vvp-run} & \tool{vvp -fst}: simulação $\to$ FST. \\
  \code{buildCocotbRunSpec} & \code{cocotb-run} & Python runner cocotb. \\
  \code{buildVerilatorBuildSpec} & \code{verilator-build} & \tool{verilator --binary --main} $\to$ \path{.exe}. \\
  \code{buildVerilatorRunSpec} & \code{verilator-run} & roda \path{V<top>.exe}. \\
  \code{buildVerilatorJsonSpec} & \code{verilator-json} & \tool{verilator --json-only}: AST de portas. \\
  \code{buildVerilatorTbBuildSpec} & \code{verilator-tb-build} & \tool{verilator --cc --exe --build} + harness C++. \\
  \code{buildVerilatorTbRunSpec} & \code{verilator-tb-run} & roda \path{V<top>.exe} (harness I/O por arquivo). \\
  \code{buildFst2VcdSpec} & \code{fst2vcd} & \tool{fst2vcd}: FST $\to$ VCD texto. \\
  \code{buildGtkwaveSpec} & \code{gtkwave} & \tool{gtkwave} (lançado detached). \\
  \code{buildYosysHierarchySpec} & \code{yosys-hierarchy} & \tool{yosys -s}: hierarquia (JSON). \\
  \code{buildPrismYosysSpec} & \code{prism-yosys} & \tool{yosys -s}: síntese PRISM. \\
  \bottomrule
\end{longtable}

\section{O contrato \texttt{CommandSpec}}
\label{sec:08-commandspec}

Um \code{CommandSpec} descreve \emph{o que} rodar; o executor decide \emph{como}.
É a forma sobre a qual os overrides da Aurora Intelligence se conectam. Suas
invariantes (documentadas em \path{command_spec.js}):

\begin{itemize}[leftmargin=1.4em]
  \item \code{binary} é um caminho absoluto (validado pelo main contra a
    allowlist; o renderer nunca executa nada);
  \item \code{args} é um array de tokens \emph{posicionais}, sem aspas de shell —
    um caminho com espaços é \emph{um} elemento, não \lstinline|"path with spaces"|;
  \item \code{cwd} e \code{env} são opcionais e, quando presentes, estruturados
    (nada de idiomas \lstinline|cd "..." && cmd|).
\end{itemize}

\noindent O módulo exporta \code{STEP_IDS} (a lista congelada dos 16 IDs de
etapa), \code{STEP_DESCRIPTIONS} (descrição de cada uma) e quatro funções puras:

\begin{description}[leftmargin=1.4em]
  \item[\code{applyOverride(spec, ov)}] devolve uma cópia do spec com o override
    aplicado, sem mutar a entrada. A ordem é: filtra \code{removeArgs},
    prepende \code{prependArgs}, anexa \code{appendArgs}, mescla \code{envSet},
    remove \code{envUnset}.
  \item[\code{formatSpec(spec)}] renderiza o spec numa linha de comando legível
    (somente exibição — nunca alimentada a um shell); tokens com espaço são
    citados.
  \item[\code{diffSpecs(before, after)}] reporta o que mudou (args adicionados/
    removidos, env adicionado/removido/alterado). Alimenta o modal de
    confirmação e o log de auditoria.
  \item[\code{validateShape(spec)}] validação rasa de formato (etapa conhecida,
    \code{binary} string, \code{args} array de strings, \code{cwd} string).
\end{description}

\section{Overrides dirigidos pela Aurora Intelligence}
\label{sec:08-overrides}

O subsistema de overrides (\path{command_overrides.js}) permite que a IA
acrescente, remova ou ajuste flags e variáveis de ambiente de qualquer etapa,
sem nunca trocar o binário. Há duas camadas com dois tempos de vida:

\begin{description}[leftmargin=1.4em]
  \item[Efêmera (\code{Map} em memória).] Vale só até a próxima invocação da
    etapa correspondente; o wrapper do executor a consome e limpa o slot. Bom
    para \enquote{tente esta flag no próximo compile}.
  \item[Persistida (\path{.spf} \code{structure.commandOverrides}).] Sobrevive
    entre sessões, viaja com o projeto e fica sob a efêmera quando ambas
    existem (a efêmera vence naquela execução).
\end{description}

\noindent As chaves são \code{<step>} (global) ou \code{<step>:<processorName>}
(por processador); a mais específica vence. \code{resolveOverride} combina as
camadas na ordem persistida-global, persistida-específica, efêmera-global,
efêmera-específica, de modo que a camada mais específica e mais recente
prevaleça. \code{applyResolved} aplica o override resolvido sobre o spec-base e,
com \lstinline|consumeEphemeral:true|, remove o slot efêmero após a leitura — é o
que o wrapper do executor faz a cada invocação.

\section{O ponto único de execução: \texttt{spec\_runner.js}}
\label{sec:08-spec-runner}

\code{spec_runner.js} é o ``chokepoint'' renderer-side que transforma um
\code{CommandSpec} base numa chamada de exec real. Qualquer caminho que
\emph{burle} o spec\_runner não honra overrides — intencional, pois esses são
caminhos não-toolchain ou gerenciados pelo main. Há duas funções de execução:

\begin{itemize}[leftmargin=1.4em]
  \item \code{runSpec(baseSpec, options)} — variante one-shot; devolve
    \lstinline|{ code, stdout, stderr, pid }|;
  \item \code{runSpecStreamed(baseSpec, options)} — variante com streaming;
    stdout/stderr disparam eventos \code{exec-spec-stream}.
\end{itemize}

\noindent Ambas validam o formato (\code{validateShape}), resolvem e aplicam
overrides (\code{applyResolved}), logam o diff no terminal certo (mapeado por
\code{terminalForStep}) e despacham por IPC. Crucialmente, ambas enviam tanto o
\code{appliedSpec} (com override) quanto o \code{baseSpec} (sem override): o main
re-roda a verificação de protected flags contra a base.

\begin{lstlisting}[language=Java]
export async function runSpec(baseSpec, options = {}) {
  const shape = CommandSpec.validateShape(baseSpec);
  if (!shape.ok) throw new Error(`runSpec: invalid base spec: ${shape.error}`);
  const { appliedSpec, override, sources } =
    await applyResolved(baseSpec, { consumeEphemeral: !!options.consumeEphemeral });
  logOverride(appliedSpec, baseSpec, override, sources, terminalForStep(baseSpec.step));
  return electronAPI.execSpec({ spec: appliedSpec, baseSpec });
}
\end{lstlisting}

\noindent O mapeamento etapa $\to$ terminal (em \code{terminalForStep}) define
onde a saída e as dicas de override aparecem:

\begin{tabularx}{\linewidth}{Y Y}
  \toprule
  \textbf{Etapa(s)} & \textbf{Terminal} \\
  \midrule
  \code{cmm} & \code{tcmm} \\
  \code{asm-pre}, \code{asm} & \code{tasm} \\
  \code{iverilog*}, \code{prism-yosys} & \code{tveri} \\
  \code{vvp*}, \code{cocotb*}, \code{verilator*}, \code{fst2vcd}, \code{gtkwave}, \code{yosys-hierarchy} & \code{twave} \\
  \bottomrule
\end{tabularx}

\section{O executor estruturado (main)}
\label{sec:08-executor}

\path{main/compile/executor.js} é o único módulo que efetivamente chama
\lstinline|spawn(binary, args, { cwd, env, shell:false })|. Ele registra dois
handlers IPC principais e dois de diagnóstico.

\subsection{Validação antes do spawn}

Tanto a chamada one-shot quanto a com streaming passam por
\code{validateSpecForExec}, que enforça, em ordem:

\begin{enumerate}[leftmargin=1.6em]
  \item formato: \code{spec.binary} string não-vazia, \code{spec.args} array de
    strings, \code{spec.cwd} string não-vazia;
  \item \textbf{allowlist}: o binário precisa estar em \code{binary_allowlist.js}
    (\autoref{sec:08-allowlist});
  \item \textbf{protected flags}: quando o \code{baseSpec} foi fornecido (i.e.,
    houve override), as flags protegidas da base precisam ter sido preservadas
    (\autoref{sec:08-protected}).
\end{enumerate}

\subsection{Os handlers IPC}

\begin{description}[leftmargin=1.4em]
  \item[\code{exec-spec}] one-shot: gera o binário, acumula stdout/stderr em
    strings e resolve com \lstinline|{ code, stdout, stderr, pid }| no
    fechamento.
  \item[\code{exec-spec-streamed}] streaming: cada chunk de stdout/stderr dispara
    um evento \code{exec-spec-stream} para o renderer; resolve com
    \lstinline|{ code }| no fechamento. Usado pelas simulações de pass-2
    (\tool{vvp}/\tool{verilator}) para que linhas de \code{$display} apareçam ao
    vivo.
  \item[\code{exec-spec-allowed-binaries}] diagnóstico: lista os binários
    permitidos (usado pela ferramenta \code{list_allowed_binaries} da IA).
  \item[\code{exec-spec-protected-flags}] diagnóstico: tabela de flags protegidas
    por etapa (usado por \code{list_allowed_flags}).
\end{description}

\subsection{Spawn rastreado e cancelamento}

O filho é gerado por \code{spawnTracked}, que o registra no \emph{registry}
central da toolchain — assim, fechar a janela principal (ou sair) faz um
\emph{tree-kill} da etapa \emph{e} de qualquer ferramenta que ela fan-out
(Verilator $\to$ make/g++/ccache, cocotb $\to$ python), não só do PID direto. O
filho também é estacionado em \code{state.currentVvpProcess} para que
\code{cancel-vvp-process} possa matá-lo — o slot serve para qualquer etapa longa,
não só \tool{vvp} (o nome é histórico). O pipeline roda as etapas
sequencialmente (apenas um filho vivo por vez), então um único slot basta.

\subsection{Prioridade de processo}

Os filhos da toolchain recebem mais CPU que o trabalho normal do desktop sem
inanir a UI. O padrão é \code{ABOVE_NORMAL} (passo seguro: \code{HIGH}/
\code{REALTIME} podem engasgar o desktop e o próprio renderer da AURORA). É
\emph{best-effort} — \code{os.setPriority} pode lançar \code{EPERM} em sistemas
restritos, nunca fatal. A variável de ambiente
\lstinline|AURORA_TOOLCHAIN_PRIORITY| (\code{normal} | \code{above} | \code{high}
| \code{realtime}) permite optar por outro nível em uma máquina dedicada a build.

\subsection{Streaming até o terminal}

No fluxo do \tool{vvp}, por exemplo, o renderer assina \code{onExecSpecStream}
\emph{antes} de invocar \code{runSpecStreamed}, classifica cada linha e a escreve
em \code{twave}. Linhas de bookkeeping do simulador (\code{FST info:},
\code{$finish called at}, etc.) são descartadas; as linhas de \code{$display}
do usuário recebem a tag \code{raw} (sempre visíveis). A assinatura é removida no
\code{finally}, garantindo limpeza mesmo em erro.

\begin{lstlisting}[language=Java]
let unsubscribe = null;
if (typeof electronAPI.onExecSpecStream === 'function') {
  unsubscribe = electronAPI.onExecSpecStream((payload) => {
    if (!payload || !payload.data) return;
    for (const line of payload.data.split(/\r?\n/)) {
      if (!line.trim()) continue;
      if (isVvpNoise(line)) continue;
      this.terminalManager.appendToTerminal('twave', line, 'raw');
    }
  });
}
try {
  const r = await runSpecStreamed(buildVvpRunSpec({ ... }), { consumeEphemeral: true });
  code = r.code;
} finally { if (unsubscribe) unsubscribe(); }
\end{lstlisting}

\noindent A ponte IPC fica em \path{js/app/preload.js}:

\begin{lstlisting}[language=Java]
execSpec:           (payload)  => ipcRenderer.invoke('exec-spec', payload),
execSpecStreamed:   (payload)  => ipcRenderer.invoke('exec-spec-streamed', payload),
onExecSpecStream:   (callback) => { /* on('exec-spec-stream'); returns unsubscribe */ },
killCurrentSpecProcess: ()     => ipcRenderer.invoke('kill-current-spec-process'),
\end{lstlisting}

\section{A allowlist de binários}
\label{sec:08-allowlist}

\path{main/compile/binary_allowlist.js} é o conjunto fechado de executáveis que o
executor aceita gerar — a fronteira de confiança. A superfície de edição de
linha de comando exposta à IA \emph{não} permite trocar o binário, mas um
renderer com defeito ou comprometido poderia tentar enviar qualquer coisa; a
allowlist rejeita qualquer spec cujo \code{binary} não esteja na tabela, com erro
claro, independentemente do que os args digam.

A tabela é indexada pelo \emph{basename} do binário; cada entrada lista os
diretórios legais (relativos a \path{components/<subdir>}). Um spec precisa
terminar com um dos nomes listados \emph{e} viver em um dos diretórios listados.

\subsection{Os binários permitidos}

\begin{longtable}{p{0.30\linewidth} p{0.42\linewidth} p{0.18\linewidth}}
  \toprule
  \textbf{Binário} & \textbf{Diretório permitido (sob \path{components})} & \textbf{Papel} \\
  \midrule
  \endhead
  \path{cmmcomp.exe} & \path{bin} & compilador C± \\
  \path{appcomp.exe} & \path{bin} & pré-proc. de macros ASM \\
  \path{asmcomp.exe} & \path{bin} & compilador ASM \\
  \path{iverilog.exe} & \path{Packages/msys/mingw64/bin} & Icarus Verilog \\
  \path{vvp.exe} & \path{Packages/msys/mingw64/bin} & runtime do Icarus \\
  \path{verilator} / \path{verilator.exe} & \path{Packages/msys/mingw64/bin} & Verilator \\
  \path{perl.exe} & \path{Packages/msys/mingw64/bin} & invoca o script verilator \\
  \path{g++.exe} & \path{Packages/msys/mingw64/bin} & compila o modelo C++ \\
  \path{make.exe} & \path{Packages/msys/mingw64/bin} & build do Verilator \\
  \path{yosys.exe} & \path{Packages/msys/mingw64/bin} & síntese (hierarquia/PRISM) \\
  \path{gtkwave.exe} & \path{Packages/gtkwave-nipscern} & visualizador de onda \\
  \path{fst2vcd.exe} & \path{Packages/gtkwave-nipscern} & FST $\to$ VCD \\
  \path{surfer.exe} & \path{Packages/surfer} & visualizador opt-in \\
  \path{verible-verilog-ls.exe} & \path{Packages/verible/bin} & LSP Verilog \\
  \path{clang-format.exe} & \path{Packages/clang-format/bin} & formatador C/C++/CMM \\
  \path{slang-server.exe} & \path{Packages/slang-server/bin} & LSP SystemVerilog \\
  \path{comp2gtkw.exe} & \path{bin} & decode de números complexos \\
  \bottomrule
\end{longtable}

\begin{nota}
A tabela estática (\code{RAW_ALLOWLIST}) também é a documentação canônica de
\enquote{onde cada ferramenta mora}. Adicionar um binário novo à toolchain é
literalmente adicionar uma linha aqui — esse é o portão inteiro.
\end{nota}

\subsection{Os dois casos especiais}

Além da tabela estática, \code{isAllowed} trata dois casos por regra:

\begin{description}[leftmargin=1.4em]
  \item[Python.] Qualquer \path{python}/\path{python3}/\path{py} (com ou sem
    \code{.exe}) só é aceito se for exatamente o Python empacotado, verificado
    por \code{isBundledPythonPath} (\autoref{sec:08-python}). O único Python do
    bundle serve a ambos os fluxos cocotb (Icarus e Verilator).
  \item[Executável gerado pelo Verilator.] O \path{V<top>.exe} não é empacotado;
    ele nasce em \path{components/Temp/obj_dir_<simTop>/}. A allowlist o aceita
    por casamento de prefixo: o caminho precisa começar com
    \path{<components>/Temp/} e bater com o padrão
    \lstinline|/obj_dir*/V<algo>(.exe)?$|.
\end{description}

\noindent A comparação de diretório é case-insensitive no Windows (o
casing da letra de unidade pode variar conforme o app foi lançado), e o caminho
do binário precisa ser absoluto.

\section{As protected flags}
\label{sec:08-protected}

\path{main/compile/protected_flags.js} mantém, por etapa, as flags que as camadas
de override \emph{não} podem remover nem substituir. A superfície da IA permite
\emph{adicionar} flags — esse é o ponto de ter poder de override —, mas algumas
flags são invariantes do pipeline: remover \code{-o <vvpFile>} do
\code{iverilog-build}, por exemplo, quebraria o fluxo Wave, pois nenhum
\path{.vvp} cairia no caminho que a etapa seguinte lê.

A verificação dispara \emph{depois} de o override ser aplicado: compara os args
resultantes contra a lista protegida. Se um token protegido da base sumiu,
rejeita; se uma env var protegida foi removida ou re-vinculada a outro valor,
rejeita. O casamento é por token exato. Para pares flag-e-valor (\code{-o file},
\code{-y dir}, \code{-s tb}) verifica-se apenas o \emph{nome} da flag (o valor
pode mudar — a AURORA o regenera a cada run); para tokens isolados
(\code{--binary}, \code{--main}, \code{-tnull}) verifica-se o literal.

\subsection{Tabela de regras (\texttt{RULES})}

\begin{longtable}{p{0.24\linewidth} p{0.30\linewidth} p{0.36\linewidth}}
  \toprule
  \textbf{Etapa} & \textbf{Flags com valor (nome preservado)} & \textbf{Literais / env protegidos} \\
  \midrule
  \endhead
  \code{cmm} & \code{-i -n -p -m -t} & --- \\
  \code{asm-pre} & \code{-i -t} & --- \\
  \code{asm} & \code{-i -p -d -m -t -f -c} & --- \\
  \code{iverilog-check} & \code{-s -y} & \code{-tnull} \\
  \code{iverilog-build} & \code{-s -y -o} & --- \\
  \code{vvp-run} & --- & \code{-fst} \\
  \code{cocotb-run} & --- & env: \code{AURORA_COCOTB_SOURCES_JSON}, \code{_TOP}, \code{_TEST_MODULE}, \code{_BUILD_DIR} \\
  \code{verilator-build} & \code{-Mdir --top-module -y} & \code{--binary --main --trace-fst} \\
  \code{verilator-json} & \code{--top-module -Mdir -y} & \code{--json-only} \\
  \code{verilator-tb-build} & \code{--top-module -Mdir -y} & \code{--cc --exe --build} \\
  \code{fst2vcd} & \code{-f -o} & --- \\
  \code{yosys-hierarchy} & \code{-s} & --- \\
  \code{prism-yosys} & \code{-s} & --- \\
  \bottomrule
\end{longtable}

\noindent As etapas \code{verilator-run}, \code{verilator-tb-run} e
\code{gtkwave} têm regras vazias (nada a proteger além do binário). A função
\code{describe(step)} expõe a tabela para a ferramenta de introspecção
\code{list_allowed_flags} da IA.

\section{Localização do Python}
\label{sec:08-python}

\path{main/compile/python_locator.js} localiza um interpretador Python para o
fluxo cocotb. A função \code{getBundledPythonPath} aponta para o único Python do
bundle unificado:
\path{components/Packages/msys/mingw64/bin/python.exe} (ou \path{python3} fora do
Windows). Esse Python carrega ambos os VPIs do cocotb
(\path{libcocotbvpi_icarus.vpl} e o estático \path{libcocotbvpi_verilator.a}),
construídos pelo repositório \repo{aurora-toolchain}.

\subsection{Ordem de preferência da descoberta}

\code{listPythonCandidatesSync} monta a lista de candidatos nesta ordem:

\begin{enumerate}[leftmargin=1.6em]
  \item o Python empacotado (\code{getBundledPythonPath});
  \item a variável de ambiente \lstinline|AURORA_PYTHON|, depois \lstinline|PYTHON|;
  \item o prefixo Conda ativo (\lstinline|CONDA_PREFIX|);
  \item candidatos comuns do Windows (miniconda3/anaconda3/mambaforge/miniforge3
    sob o perfil do usuário e Program Files; instalações em
    \path{LOCALAPPDATA/Programs/Python}; o \tool{py.exe});
  \item Python encontrado no \code{PATH} (e via \tool{where.exe}/\tool{which}).
\end{enumerate}

\subsection{Sonda de capacidade}

\code{getPythonStatus} prioriza o candidato empacotado; se presente, roda nele
uma sonda (\code{runPythonProbe}) que executa um script curto reportando
\code{executable}, versão do Python, se tem cocotb e qual a versão do cocotb. O
status enriquecido (\code{enrichStatus}) marca \code{isBundled} e compara a versão
do cocotb com a esperada (\code{EXPECTED_COCOTB_VERSION = '2.0.1'}), expondo
\code{isPinnedCocotb}. A fábrica de specs (\autoref{sec:08-spec-factory}) usa esse
status: o \code{cocotb-run} exige que o Python seja o do bundle, tenha cocotb e
que a versão case com a esperada — caso contrário, lança erro descritivo antes de
qualquer spawn.

\section{O ambiente do processo-filho}
\label{sec:08-env}

\code{buildChildEnv} (em \path{executor.js}) monta o ambiente do filho em
camadas: \code{process.env} $\to$ defaults de performance da toolchain $\to$
\code{spec.env} (overrides vencem). Os defaults injetados são
\lstinline|OMP_NUM_THREADS| e \lstinline|OMP_THREAD_LIMIT|, ambos com a contagem
de CPUs.

\subsection{Sanitização do \texttt{spec.env}}

Antes de mesclar, \code{spec.env} é \emph{sanitizado} como defesa contra um spec
adulterado: só passam chaves que casam \lstinline|^[A-Za-z_][A-Za-z0-9_]*$| com
valor string sem \emph{null byte}. Specs legítimos (\code{OMP_*},
\code{MAKEFLAGS}, \code{LC_ALL}, \code{AURORA_COCOTB_*}, etc.) passam intactos.

\subsection{Composição do PATH (\texttt{prependPath})}

Quando o spec tem \code{prependPath}, esses diretórios são prefixados ao
\code{PATH} \emph{depois} da mescla, para que sempre tenham precedência. Apenas
diretórios string, não-vazios, sem \emph{null byte} e \emph{sem} o separador de
PATH (uma entrada com \code{;}/\code{:} contrabandearia várias) entram no
prefixo. É por meio do \code{prependPath} que o \tool{iverilog}/\tool{vvp} acham
suas DLLs MinGW e que o Verilator acha \tool{perl}/\tool{g++}/\tool{make} —
substituindo o antigo \lstinline|cmd.exe && set "PATH=..."|.

\begin{lstlisting}[language=Java]
function buildChildEnv(spec) {
  const cpuCount = getCPUCount();
  const sep = process.platform === 'win32' ? ';' : ':';
  const safeSpecEnv = {};
  for (const [k, v] of Object.entries(spec.env || {})) {
    if (/^[A-Za-z_][A-Za-z0-9_]*$/.test(k) && typeof v === 'string' && !v.includes('\0'))
      safeSpecEnv[k] = v;
  }
  const env = { ...process.env,
    OMP_NUM_THREADS: cpuCount.toString(),
    OMP_THREAD_LIMIT: cpuCount.toString(),
    ...safeSpecEnv };
  if (Array.isArray(spec.prependPath) && spec.prependPath.length) {
    const safeDirs = spec.prependPath.filter((p) =>
      typeof p === 'string' && p.length > 0 && !p.includes('\0') && !p.includes(sep));
    if (safeDirs.length) env.PATH = safeDirs.join(sep) + sep + (env.PATH || '');
  }
  return env;
}
\end{lstlisting}

\section{Modos de pipeline: simulação completa vs.\ Verilog-only}
\label{sec:08-modos}

O pipeline não tem ramos do tipo \lstinline|if (hasProcessor)|; em vez disso, os
laços de cmm+asm iteram sobre \code{availableProcessors} e o conjunto de fontes
(\code{synthesizableFiles}) já inclui os \path{.v} dos processadores. O que
\emph{varia} entre os botões é \emph{quanto} do pipeline executa, decidido pelo
botão (e pelo estado do design lido do \path{.spf}), não por uma ramificação
condicional dentro de uma etapa.

\subsection{Pré-compilação dos processadores}

\code{precompileAllProcessors} é o pré-flight comum a Verilog/Wave/PRISM: para
cada processador conhecido, roda cmm+asm. Pula em silêncio aqueles cuja
convenção \path{<proj>/<proc>/Software/<proc>.cmm} não existe em disco (sem
\path{.cmm}, o \tool{cmmcomp} falharia). Como itera sobre o array de
processadores, é \emph{no-op} natural num projeto de Verilog puro.

\subsection{Verilog-only}

O botão \code{verilog} (e o \code{prism}, superset dele) termina em
\tool{iverilog -tnull} sobre o top-level. É uma verificação de sintaxe e
elaboração: nenhum \path{.vvp} é produzido, nenhuma simulação roda. O top-level
da \enquote{sim} usado nesse modo é o módulo do top-level do projeto (e não o do
testbench), pois não há testbench em jogo.

\subsection{Simulação completa (Wave)}

O botão \code{wave} acrescenta \tool{iverilog -o vvp} \emph{com} o testbench,
roda \tool{vvp -fst} (com streaming ao vivo) e abre o visualizador. A escolha do
top-level da simulação reflete a presença de um testbench Verilog:

\begin{lstlisting}[language=Java]
const hasVerilogTestbench = config.testbenchFile && !isPythonFile(config.testbenchFile);
const simTopModule = hasVerilogTestbench
  ? moduleStemFromPath(config.testbenchFile)  // tem testbench -> top = testbench
  : topLevelModuleName;                        // sem tb -> top = top-level do projeto
\end{lstlisting}

\subsection{Auto-decisão de engine e fluxo cocotb}

A fábrica de specs (\autoref{sec:08-spec-factory}) deriva o engine a partir do
tipo do testbench e da preferência de simulador:

\begin{itemize}[leftmargin=1.4em]
  \item testbench \path{.py} $\to$ fluxo \code{cocotb-run} (roda em qualquer
    engine, headless ou com onda); exige o Python do bundle com cocotb na versão
    correta;
  \item testbench \path{.v} $\to$ fluxo \tool{iverilog}+\tool{vvp} (padrão) ou
    \tool{Verilator} binário, conforme \code{getSimulator}; o
    \code{iverilog-build} rejeita explicitamente um testbench Python
    (\enquote{use cocotb-run}).
\end{itemize}

\noindent O Fast Sim (\code{verilator-fast}) é o caminho headless: roda o
testbench via Verilator binário (ou cocotb headless, se \path{.py}) sem gerar
onda nem abrir o visualizador — só a velocidade.

\section{A fábrica de specs (pré-visualização)}
\label{sec:08-spec-factory}

\path{js/compilation/spec_factory.js} constrói o spec-base de qualquer etapa
\emph{sem} executá-lo. É o que alimenta as ferramentas
\code{inspect_compile_command} / \code{preview_compile_command} da IA: o
modelo pergunta \enquote{o que você \emph{rodaria} para a etapa X?}, a fábrica
monta o spec usando os \emph{mesmos} builders do pipeline real, e as camadas de
override são depois sobrepostas por \code{resolveSpec} (em \code{spec_runner.js})
para exibição.

Para etapas por-processador (\code{cmm}, \code{asm-pre}, \code{asm}),
\code{processorName} é exigido — senão, a fábrica escolhe o primeiro processador
do projeto ou falha com mensagem útil. Para etapas do pipeline inteiro
(\code{iverilog-*}, \code{vvp-*}, \code{verilator-*}, \code{gtkwave},
\code{prism-yosys}), o processador é ignorado. Algumas etapas lançam erro quando
precisam de informação de runtime que a fábrica não consegue reconstruir — por
exemplo, \code{verilator-run} precisa do \path{V<top>.exe} por-build, que só
existe depois de o build rodar.

\section{Os handlers IPC auxiliares de compilação}
\label{sec:08-ipc-auxiliares}

\path{main/ipc/compile.js} concentra a execução de processos que \emph{não}
passam pelo executor estruturado, sempre mantendo o cancelamento numa única
fonte de verdade. O handler legado \code{exec-command} (string de shell crua) foi
removido — toda a execução da toolchain passa agora pelo executor estruturado.
Os handlers presentes:

\begin{description}[leftmargin=1.4em]
  \item[\code{launch-gtkwave-only}] lança o GTKWave \emph{detached} (sobrevive ao
    run; rastreado por \code{spawnTracked} para ser derrubado junto com a IDE). O
    binário fornecido pelo renderer é passado pela mesma allowlist
    (\code{isAllowed}) antes de qualquer spawn.
  \item[\code{launch-surfer}] lança o visualizador opt-in Surfer, com a mesma
    porta de allowlist e um \code{existsSync} adicional (Surfer não é
    empacotado por padrão). Reaproveita uma única janela fechando a instância
    anterior antes de abrir outra.
  \item[\code{write-surfer-mappings}] grava os \emph{translators} de mapeamento
    (decode de Assembly/linha-fonte) que o \path{.surf.ron} referencia.
  \item[\code{decode-complex}] decodifica padrões de bits de números complexos
    via o canônico \tool{comp2gtkw.exe} (também passado pela allowlist),
    alimentando os valores por stdin.
  \item[\code{cancel-vvp-process}] cancelamento amplo: mata o filho estacionado
    \emph{e} varre \path{vvp.exe}/\path{gtkwave.exe} por nome.
  \item[\code{kill-current-spec-process}] parada cirúrgica: mata \emph{apenas} o
    filho com streaming estacionado (sem a varredura por nome). Usado por
    \code{_extractFstHeaderVcd} para parar o \tool{fst2vcd} assim que o cabeçalho
    é capturado.
\end{description}

\section{Resumo}
\label{sec:08-resumo}

O subsistema de compilação da AURORA é uma esteira de duas metades. No renderer,
os botões expandem para sequências de etapas; cada etapa vira um
\code{CommandSpec} produzido por um \emph{builder} puro, eventualmente alterado
por \emph{overrides} da Aurora Intelligence, e despachado por um ponto único
(\code{spec_runner.js}). No main, um \emph{executor} valida o spec contra uma
\emph{allowlist} fechada de \textasciitilde17 binários e contra \emph{protected
flags} por etapa, monta o ambiente do filho (com sanitização e
\code{prependPath}) e o gera com \lstinline|shell:false| — eliminando, por
construção, qualquer parsing de shell. A distinção entre Verilog-only e simulação
completa não é uma ramificação dentro das etapas, e sim o quanto da esteira cada
botão executa, com o engine (Icarus, Verilator ou cocotb) auto-decidido pelo tipo
do testbench e pela preferência de simulador.

% !TEX root = ../main.tex
\chapter{AURORA: simulação e o fluxo de ondas}
\label{cap:09-aurora-simulacao-ondas}

Este capítulo documenta o subsistema da \tool{AURORA} que vai dos arquivos
Verilog de um projeto até um arquivo de ondas (\file{.fst} ou \file{.vcd})
pronto para ser visualizado. O escopo cobre a orquestração da simulação, a
escolha do simulador e do visualizador, a preparação do que entra na onda
(o \lstinline|$dumpvars|) e a geração automática do layout \file{.gtkw}. O
visualizador em si (a janela do GTKWave/fork \tool{nipscern} e a integração com
o \tool{Surfer}) é assunto do \autoref{cap:10-aurora-visualizacao}; aqui o fluxo
termina no instante em que o arquivo de ondas existe e o processo do
visualizador é disparado.

A fonte de verdade deste capítulo é o código em \path{js/wave/*} e as fases
\lstinline|_wave*| do método \lstinline|runGtkWave| em
\path{js/compilation/compilation_module.js}, além do resolvedor de seleção em
\path{js/compilation/wave_signal_validator.js}.

\section{Visão geral: o botão Wave}

O ponto de entrada é \lstinline|CompilationModule.runGtkWave()|. O nome do
método é histórico: ele orquestra o pipeline completo independentemente do
simulador e do visualizador escolhidos. A primeira linha do método valida a
configuração com \lstinline|validateForWave()|, que exige um testbench --- sem
testbench não há o que simular. Os arquivos sintetizáveis e o \emph{top-level}
são opcionais: um testbench autônomo pode definir o DUT internamente.

\subsection{Os quatro caminhos de simulação}

A \tool{AURORA} oferece quatro caminhos para produzir uma onda. Eles resultam da
combinação de dois eixos independentes: a \emph{linguagem do testbench}
(Verilog puro ou Python/cocotb) e o \emph{motor de simulação} (Icarus Verilog ou
Verilator). A \autoref{tab:09-quatro-caminhos} resume as combinações.

\begin{table}[h]
\centering
\begin{tabularx}{\linewidth}{Y Y Y}
\toprule
\textbf{Testbench} & \textbf{Simulador} & \textbf{Cadeia de ferramentas} \\
\midrule
Verilog (\file{.v}) & Icarus & \tool{iverilog} $\rightarrow$ \file{.vvp} $\rightarrow$ \tool{vvp} \\
Verilog (\file{.v}) & Verilator & \tool{verilator} $\rightarrow$ C++ $\rightarrow$ \tool{g++} $\rightarrow$ \file{.exe} \\
Python (\file{.py}, cocotb) & Icarus & runner cocotb com \lstinline|SIM=icarus| \\
Python (\file{.py}, cocotb) & Verilator & runner cocotb com \lstinline|SIM=verilator| \\
\bottomrule
\end{tabularx}
\caption{Os quatro caminhos de simulação e suas cadeias de ferramentas.}
\label{tab:09-quatro-caminhos}
\end{table}

Os quatro caminhos convergem: cada um produz um arquivo de ondas que o restante
do pipeline trata de modo uniforme. A simulação roda \emph{exatamente uma vez}
em todos os caminhos. O cabeçalho do VCD (a hierarquia de \lstinline|$scope| e
\lstinline|$var| usada pelo \emph{picker} de sinais e pelo gerador automático de
\file{.gtkw}) é extraído diretamente do \file{.fst} já finalizado por
\lstinline|_extractFstHeaderVcd| --- não há mais um passe separado só para
gerar o cabeçalho.

\subsection{O diagrama de fluxo}

\begin{lstlisting}[caption={Os quatro caminhos convergindo no arquivo de ondas e no visualizador.},captionpos=b]
                         runGtkWave()
                              |
                    validateForWave()  (exige testbench)
                              |
           resolveWaveToolchain(componentsPath)
                              |
              _waveDeriveSimTopModule(config)
                              |
        +---------------------+---------------------+
        |                                           |
   testbench .py?                            testbench .v
   (isPythonFile)                                   |
        |                                  getSimulator()
        | sim escolhido:                  +---------+---------+
        |  icarus | verilator             |                   |
        v                            'verilator'          'iverilog'
 _waveValidateCocotbConfig                |                   |
        |                       _waveBuildVerilator   _waveBuildAndVerifyVvp
 _waveRunCocotbSimulation               |                   |
   (runner cocotb)       _waveRunVerilatorSimulation   _waveRunVvpSimulation
        |                                |                   |
        |                                +---------+---------+
        |                                          |
        |                              _waveResolveVcdFile()
        |                                          |
        +----------------------+-------------------+
                               |
                  (CONVERGENCIA: .fst / .vcd)
                               |
                  _extractFstHeaderVcd()  -> <top>.header.vcd
                               |
                       getViewer()
              +----------------+----------------+
              |                                 |
          'surfer'                          'gtkwave'
              |                                 |
 _waveResolveSurferSaveFile      _waveResolveGtkwSaveFile (auto-gtkw)
 _waveLaunchSurfer               _waveLaunchGtkwave
\end{lstlisting}

\section{Escolha do simulador e do visualizador}

Dois \emph{toggles} globais, persistidos em \lstinline|localStorage|, controlam
qual simulador e qual visualizador o botão Wave usa. Ambas as escolhas são do
usuário, valem para o aplicativo inteiro (não por projeto nem por testbench) e
são lidas \enquote{frescas} a cada clique no Wave.

\subsection{Preferência de simulador}

O módulo \file{js/wave/simulator_preference.js} expõe \lstinline|getSimulator()|
e \lstinline|setSimulator(value)|. A chave de storage é
\lstinline|aurora.waveSimulator| e os valores válidos são \lstinline|iverilog|
(padrão) e \lstinline|verilator|.

\begin{lstlisting}[language=Java,caption={Leitura defensiva da preferência de simulador.}]
const STORAGE_KEY = 'aurora.waveSimulator';
const VALID = new Set(['iverilog', 'verilator']);

export function getSimulator() {
    try {
        const v = (typeof localStorage !== 'undefined')
            ? localStorage.getItem(STORAGE_KEY) : null;
        return VALID.has(v) ? v : 'iverilog';
    } catch (_e) {
        return 'iverilog';
    }
}
\end{lstlisting}

A função nunca lança exceção: ela roda em \emph{hot path} (a cada clique no
Wave) e é defensiva quanto à corrupção do storage --- qualquer valor inválido
ou ausente cai no padrão \lstinline|iverilog|. A \tool{UI} do toggle vive em
\file{js/wave/simulator_toggle.js} (um \emph{segmented switch} na barra de
ferramentas) e dispara o evento \lstinline|aurora:wave-simulator-changed| para
que outras superfícies sigam a mudança.

\begin{nota}
Testbenches em Python/cocotb também respeitam \lstinline|getSimulator()|: o
valor é mapeado para o campo \lstinline|SIM| do runner cocotb
(\lstinline|verilator| ou \lstinline|icarus|). Não há, portanto, um quinto
caminho --- o eixo do simulador é o mesmo para Verilog e para Python.
\end{nota}

\subsection{Preferência de visualizador}

O módulo \file{js/wave/viewer_preference.js} é o espelho exato do anterior, com
\lstinline|getViewer()| / \lstinline|setViewer(value)|, chave
\lstinline|aurora.waveViewer| e valores \lstinline|gtkwave| (padrão) e
\lstinline|surfer|. A escolha do visualizador só importa no fim do pipeline,
após a onda existir: \lstinline|runGtkWave| ramifica em
\lstinline|getViewer() === 'surfer'| para decidir entre o caminho do Surfer
(layout \file{.surf.ron} / \file{.sucl}) e o caminho do GTKWave (save-file
\file{.gtkw}).

\section{As fases do orquestrador}

O corpo de \lstinline|runGtkWave| é curto de propósito: ele expressa a
\emph{ordem das fases} e delega cada preocupação a um método privado
\lstinline|_wave*| com contrato próprio. Esta seção percorre as fases na ordem
em que executam.

\subsection{\texttt{resolveWaveToolchain}}

Resolve os caminhos e binários necessários à simulação a partir de
\lstinline|this.componentsPath|: \lstinline|tempBaseDir| (a pasta
\file{components/Temp/}, criada se não existir), o binário do GTKWave, o
\tool{vvp}, o \tool{fst2vcd} (usado na extração de cabeçalho), entre outros. O
diretório \lstinline|tempBaseDir| é o \emph{working directory} de todas as
simulações --- caminhos relativos em \lstinline|$readmemb| e
\lstinline|$fopen| são resolvidos a partir dele.

\subsection{\texttt{\_waveDeriveSimTopModule}}

Deriva o nome do módulo \emph{top} de simulação --- o valor passado como
\lstinline|-s| ao \tool{iverilog} e o radical (\emph{stem}) que a \tool{AURORA}
espera para a tríade \file{.vvp} / \file{.vcd} / \file{.gtkw}. A regra é
direta: o módulo do testbench é o top quando há testbench; o fallback para o
top sintetizável só existe para o fluxo raro de \enquote{compilar um único
\file{.v} sem testbench}, que o botão Wave rejeita a montante.

\begin{lstlisting}[language=Java,caption={Derivação do módulo top de simulação.}]
_waveDeriveSimTopModule(config) {
    if (config.testbenchFile) {
        return moduleStemFromPath(config.testbenchFile);
    }
    if (!config.topLevelFile) return null;
    return moduleStemFromPath(config.topLevelFile);
}
\end{lstlisting}

No caminho cocotb, o módulo top é \emph{re-derivado}: ele vem da diretiva
\lstinline|# aurora-toplevel: <module>| do \file{.py} (se presente) ou do
top-level do \file{.spf} (com aviso), via
\lstinline|_waveValidateCocotbConfig|.

\subsection{Construção e simulação por caminho}

\subsubsection{Caminho Icarus (\texttt{iverilog} + \texttt{vvp})}

Duas fases: \lstinline|_waveBuildAndVerifyVvp(simTopModule, tempBaseDir)| e
\lstinline|_waveRunVvpSimulation(simTopModule, tools)|.

A primeira chama \lstinline|waveBuildVvp()|, que compila o conjunto de fontes
(sintetizáveis $+$ testbench instrumentado $+$ \lstinline|-y components/HDL|)
com \tool{iverilog} produzindo \file{<top>.vvp}, e depois confirma que o
arquivo foi escrito em disco. O \file{.vvp} é \emph{sempre} reconstruído ---
o testbench instrumentado \enquote{assa} a seleção de \lstinline|$dumpvars| do
usuário no momento da compilação, então um \file{.vvp} antigo travaria uma
seleção antiga.

A segunda executa \lstinline|vvp -fst| sobre o \file{.vvp}, com \emph{cwd} em
\lstinline|tempBaseDir|. Antes de rodar, ela faz \emph{staging} de arquivos: os
\file{pc_*_mem.txt} gerados pelo \tool{cmmcomp} em subpastas
\file{Temp/<proc>/} são copiados para \file{Temp/} direto (para que os
\lstinline|$readmemb| relativos resolvam), e os arquivos de dado referenciados
pelo testbench (\lstinline|$fopen|, \lstinline|$readmemh|) são copiados da
pasta do testbench. A saída é \file{<top>.vcd} com conteúdo binário FST escrito
sob o nome do \lstinline|$dumpfile| --- o \tool{vvp} honra o nome sem trocar a
extensão.

\subsubsection{Caminho Verilator (testbench Verilog)}

Duas fases: \lstinline|_waveBuildVerilator(...)| e
\lstinline|_waveRunVerilatorSimulation(...)|. O Verilator é \emph{opt-in} via o
toggle de simulador. O build reusa \lstinline|_prepareWaveBuildInputs| (a mesma
resolução de seleção e instrumentação de testbench do caminho Icarus) e invoca
o Verilator com, entre outras, as flags da \autoref{tab:09-verilator-flags}.

\begin{table}[h]
\centering
\begin{tabularx}{\linewidth}{Y Z}
\toprule
\textbf{Flag} & \textbf{Função} \\
\midrule
\lstinline|--binary| & gera o \lstinline|main| e invoca o \tool{make} para produzir o \file{.exe} \\
\lstinline|--trace-fst| & instrumenta o runtime para escrever FST (honra \lstinline|$dumpvars|/\lstinline|$dumpfile|) \\
\lstinline|--timing| & habilita \lstinline|#| delays (testbenches SAPHO geram clock por delay) \\
\lstinline|+define+YANC_TRACE| & ativa o bloco de sim-visibility no \file{<proc>.v} gerado \\
\lstinline|-Mdir <objDir>| & diretório de saída dos arquivos C++ e do Makefile \\
\lstinline|--top-module| & módulo top da simulação (nome do testbench) \\
\lstinline|-y <hdl>| & biblioteca de módulos (mesma lógica do \tool{iverilog}) \\
\lstinline|-Wno-fatal| & avisos não abortam o build (Verilator é mais estrito que Icarus) \\
\bottomrule
\end{tabularx}
\caption{Principais flags do build Verilator no fluxo Wave.}
\label{tab:09-verilator-flags}
\end{table}

O executável resultante é \file{obj\_dir\_<top>/V<top>.exe} (com fallback para
\file{V<top>} sem extensão). A fase de simulação roda esse \file{.exe} uma
única vez, com \emph{cwd} em \lstinline|tempBaseDir| e com o \lstinline|PATH|
estendido com \lstinline|mingw64\bin| e \lstinline|usr\bin| do bundle (o
\file{.exe} linka dinamicamente contra \file{libstdc++-6.dll} e afins).

\begin{nota}
\textbf{Visibilidade de sinais sob Verilator (YANC v4.3).} O harness compila
sob Verilator via \lstinline|+define+YANC_TRACE|, e o \file{<proc>.v} gerado
espelha cada variável/array do usuário, a tabela de linha PC$\rightarrow$\cmm{}
e os ports de I/O em sinais marcados \lstinline|/* verilator public_flat */|.
Esses sinais curados aparecem no FST igual ao Icarus. Permanecem
\emph{não-visíveis} sob Verilator os internos do CPU cercados por
\lstinline|/* verilator tracing_off */| (por exemplo os sinais de Stack/ULA
dentro do \lstinline|.core|): mesmo que o \lstinline|$dumpvars| os nomeie, a
\emph{fence} vence. Sob Icarus, tudo é \emph{dumpável}.
\end{nota}

\subsubsection{Caminhos cocotb (testbench Python em Icarus ou Verilator)}

Quando \lstinline|isPythonFile(config.testbenchFile)| é verdadeiro, o
orquestrador toma o ramo cocotb. \lstinline|_waveValidateCocotbConfig| resolve o
DUT (módulo que o cocotb elabora e liga a \lstinline|dut|), preferindo a
diretiva \lstinline|# aurora-toplevel:| no \file{.py} e, na ausência dela, o
top-level do \file{.spf} (com aviso). \lstinline|_waveRunCocotbSimulation|
então:

\begin{enumerate}[leftmargin=*]
  \item resolve o perfil de simulador com \lstinline|_resolveCocotbSimProfile|.
    Ambos os perfis rodam sobre o \emph{único} Python do bundle unificado
    (\path{components/Packages/msys/mingw64/bin/python.exe}), cujo cocotb carrega
    os dois VPIs (Icarus e Verilator). A única diferença por simulador é o campo
    \lstinline|SIM| e os \emph{build args} (\lstinline|-g2012| é exclusivo do
    Icarus; o Verilator recebe o conjunto \lstinline|-Wno-*|,
    \lstinline|+define+YANC_TRACE| e \lstinline|--trace-fst|);
  \item escreve o script runner (\lstinline|_writeCocotbRunnerScript|), que usa
    \lstinline|cocotb_tools.runner.get_runner| para \lstinline|build| e
    \lstinline|test|;
  \item coleta as fontes (\lstinline|_collectCocotbSources|: sintetizáveis $+$
    top-level $+$ HDL do bundle), faz \emph{staging} dos \file{pc_*_mem.txt}
    para o \emph{build dir} e monta o ambiente (\lstinline|SIM|,
    \lstinline|TOPLEVEL_LANG=verilog|, \lstinline|WAVES|, \lstinline|PYTHONUTF8|,
    etc.);
  \item após a execução, \emph{adota} a onda com \lstinline|_adoptCocotbWaveform|:
    procura o dump no \emph{build dir} e na pasta do testbench, normaliza-o para
    \file{Temp/<top>.fst} (ou \file{.vcd}) e devolve o caminho.
\end{enumerate}

O cocotb usa \lstinline|--trace-fst| (FST é cerca de 10$\times$ menor que o VCD
bruto, barateando o I/O da simulação longa). O cabeçalho para o auto-gtkw é
puxado desse FST pelo mesmo \lstinline|_extractFstHeaderVcd| dos demais
caminhos.

\subsection{\texttt{\_waveResolveVcdFile}}

Após a simulação (nos caminhos não-cocotb), descobre o caminho absoluto do
arquivo de ondas. A ordem de busca em \lstinline|tempBaseDir| é:
\file{<top>.fst} (esperado) $\rightarrow$ \file{<top>.vcd} (fallback legado)
$\rightarrow$ varredura por um único \file{.fst}/\file{.vcd} restante. Quando há
exatamente um candidato com nome inesperado, ele é adotado com um aviso de
\enquote{dumpfile mismatch} (o \lstinline|$dumpfile()| do usuário escolheu outro
nome). Zero candidatos ou múltiplos candidatos resultam em erro com detalhe.

\subsection{\texttt{\_extractFstHeaderVcd}}

Esta é a fase de \emph{convergência} dos quatro caminhos. Ela extrai
\emph{apenas} o cabeçalho VCD (a hierarquia \lstinline|$scope|/\lstinline|$var|
até \lstinline|$enddefinitions|) de um FST, \emph{sem} materializar o VCD-texto
completo. A técnica:

\begin{itemize}[leftmargin=*]
  \item executa o \tool{fst2vcd} em modo \emph{stream} (sem \lstinline|-o|, ele
    emite o VCD para \lstinline|stdout|, cabeçalho primeiro);
  \item acumula \lstinline|stdout| e, no instante em que vê
    \lstinline|$enddefinitions $end|, mata o \tool{fst2vcd}
    (\lstinline|killCurrentSpecProcess|) --- assim ele percorre apenas a
    geometria do FST mais o primeiro bloco bufferizado, nunca o corpo de
    centenas de MB;
  \item grava o cabeçalho num arquivo irmão \file{<top>.header.vcd}.
\end{itemize}

Há um \emph{fallback} de conversão completa (\tool{fst2vcd} normal) para o caso
de o streaming não estar disponível ou o cabeçalho não surgir. Uma entrada que
não é FST (por exemplo um dump que já é VCD-texto) falha a verificação de
\emph{magic} do \tool{fst2vcd}; nesse caso a função retorna \lstinline|false| e
o VCD-texto é parseado diretamente a jusante. Um arquivo de cabeçalho vazio
(FST corrompido, ou exit limpo sem conteúdo) é tratado como falha de captura.

\subsection{\texttt{\_waveResolveGtkwSaveFile}}

Decide qual \file{.gtkw} o GTKWave abre. Duas fontes, em ordem de prioridade:

\begin{enumerate}[leftmargin=*]
  \item \textbf{\file{.gtkw} curado pelo usuário}: a entrada com
    \lstinline|isActive === true| na lista per-testbench do \tool{WaveStore}
    (marcada pelo dropdown do \emph{gtkw picker} na barra). É submetida a um
    cross-check contra o VCD (\lstinline|_waveValidateUserGtkwAgainstVcd|) e
    devolvida intacta;
  \item \textbf{auto-gerado por \lstinline|buildAuroraGtkw|}: parseia o
    cabeçalho (preferindo \file{<top>.header.vcd}), constrói o layout
    (seção \enquote{Top-level} $+$ uma seção SAPHO por processador detectado),
    filtra por \lstinline|_validatedWaveSelection| quando há seleção e escreve
    \file{<top>.gtkw}.
\end{enumerate}

Quando o arquivo de ondas é \file{.fst} e não há \file{.header.vcd}
parseável, o método retorna \lstinline|null| e o GTKWave abre sem save-file.
Antes de escrever o layout, sinais selecionados que não chegaram ao VCD são
reportados; sob Verilator, sinais ausentes que vivem dentro de um
\lstinline|.core| de processador recebem uma nota amigável por processador
(limitação conhecida, não erro) em vez do aviso genérico de \enquote{stale}.

\subsection{\texttt{\_waveLaunchGtkwave}}

Monta a linha de comando do GTKWave (fork \tool{nipscern}, com \lstinline|--dark|,
\lstinline|--zoom-fit|, \lstinline|--left-justify| e \lstinline|-a <gtkw>| quando
aplicável), aplica eventuais \emph{overrides} (de
\path{js/compilation/command_overrides.js}, usados pela \tool{Aurora Intelligence})
e dispara o processo via \lstinline|launchGtkwaveOnly| (spawn detached,
monitorado por poll). O lançamento do visualizador é detalhado no
\autoref{cap:10-aurora-visualizacao}; a fase paralela para o Surfer é
\lstinline|_waveLaunchSurfer|.

\section{O WaveStore: estado por testbench}

O módulo \file{js/wave/wave_state_store.ts} (compilado para
\file{wave\_state\_store.js}) mantém o estado do fluxo de ondas isolado
\emph{por testbench}. Cada testbench tem seu próprio escopo, persistido como um
JSON em \file{<project>/testbench/<tbKey>.json}, onde \lstinline|tbKey| é o nome
do \file{.v} sem extensão.

\subsection{O formato do estado}

\begin{table}[h]
\centering
\begin{tabularx}{\linewidth}{Y Z}
\toprule
\textbf{Campo} & \textbf{Significado} \\
\midrule
\lstinline|tbPath| & caminho absoluto do \file{.v} do testbench \\
\lstinline|tbModule| & nome do módulo Verilog ($=$ \lstinline|tbKey|) \\
\lstinline|hadOriginalDumpvars| & na 1ª visita, o testbench já tinha \lstinline|$dumpfile|/\lstinline|$dumpvars| escrito à mão? \\
\lstinline|gtkwFiles| & \file{.gtkw} registrados via dropdown; um com \lstinline|isActive: true| \\
\lstinline|surferFiles| & layouts do Surfer (\file{.surf.ron}/\file{.sucl}) registrados \\
\lstinline|waveSignals| & seleção da Wave Configuration (caminhos pontilhados, ex.: \lstinline|tb.dut.q|) \\
\lstinline|wcInitialized| & o modal Wave Configuration já foi aberto pelo menos uma vez? \\
\lstinline|wcCustomized| & a seleção salva difere do default (o usuário alterou)? \\
\bottomrule
\end{tabularx}
\caption{Campos persistidos do \tool{WaveStore} (\lstinline|WaveState|).}
\label{tab:09-wavestore}
\end{table}

O campo \lstinline|hadOriginalDumpvars| é um \emph{snapshot} da 1ª visita: ele
NÃO muda depois do registro inicial. Mesmo que o usuário edite o testbench, a
re-instrumentação do botão Wave continua baseada nessa decisão original.

\subsection{A API e a serialização das escritas}

A API espelha a do \lstinline|SpfStore|: \lstinline|read| puro e
\lstinline|update| atômico via cadeia de promessas serializada por
\lstinline|(projectPath, tbKey)|.

\begin{description}[leftmargin=*]
  \item[\lstinline|get(projectPath, tbKey)|] retorna o estado salvo ou
    \lstinline|null| se não registrado.
  \item[\lstinline|read(projectPath, tbKey)|] retorna o estado aplicando os
    \lstinline|DEFAULTS| (não persiste).
  \item[\lstinline|ensureRegistered(...)|] registra o testbench se ausente;
    idempotente --- re-registrar não sobrescreve \lstinline|hadOriginalDumpvars|.
  \item[\lstinline|update(...)|] read-mutate-write atômico. Updates para o mesmo
    \lstinline|(projectPath, tbKey)| serializam; updates para testbenches
    distintos correm em paralelo.
  \item[\lstinline|list(projectPath)|] lista todos os testbenches registrados.
\end{description}

A chave de chaveamento usa \lstinline|safeKey|, que substitui
separadores de caminho e \lstinline|..| por \lstinline|_| (defesa contra
\emph{path traversal}). Quem escreve as flags do WaveStore: o \emph{gtkw
picker} (\lstinline|gtkwFiles|/\lstinline|surferFiles|, \lstinline|isActive|) e
o modal Wave Configuration (\lstinline|waveSignals|, \lstinline|wcInitialized|,
\lstinline|wcCustomized|).

\section{O parser de sinais}

O módulo \file{js/wave/signal_parser.ts} é um parser Verilog leve, baseado em
expressões regulares, usado pelo \emph{picker} da Wave Configuration e pela
resolução de seleção. Ele não usa o \tool{Yosys} porque (a) a geração de
hierarquia do \tool{Yosys} alimenta apenas os arquivos sintetizáveis --- o
testbench é excluído, então sinais declarados nele nunca apareceriam; (b) o
\tool{Yosys} otimiza fora redes não usadas; (c) o picker precisa da lista de
sinais \emph{antes} de qualquer simulação rodar, então também não dá para
extrair os nomes de um VCD.

\subsection{O que o parser captura}

O parser extrai apenas \emph{nomes} de sinais e instanciações de módulos, não
tipos, expressões ou timing. As estruturas principais:

\begin{itemize}[leftmargin=*]
  \item \lstinline|VerilogSignal|: \lstinline|name|, \lstinline|kind| (palavra
    primária por prioridade: direção $>$ tipo de rede $>$ \lstinline|signed|),
    \lstinline|isSigned| e \lstinline|range| (texto da faixa de bits sem
    colchetes, ou \lstinline|null| se escalar);
  \item \lstinline|ModuleInstance|: \lstinline|instanceName| e
    \lstinline|moduleType|;
  \item \lstinline|ModuleInfo|: arquivo, sinais e instâncias de um módulo;
  \item \lstinline|HierarchyNode|: um nó da árvore de hierarquia, com
    \lstinline|scopePath| (caminho pontilhado) e filhos.
\end{itemize}

Os tokens de tipo reconhecidos incluem \lstinline|input|, \lstinline|output|,
\lstinline|inout|, \lstinline|wire|, \lstinline|reg|, \lstinline|logic|,
\lstinline|signed| e --- relevante para o SAPHO --- \lstinline|real|,
\lstinline|integer| e \lstinline|time|. Esses três últimos são tipos não
sintetizáveis que aparecem no source gerado pelo \tool{asmcomp} para variáveis
\cmm{} (por exemplo \lstinline|real| para \emph{float}). Sem eles, o picker não
mostraria variáveis \lstinline|real|, e o \lstinline|$dumpvars| ficaria sem.

\subsection{Robustez do parser}

O parser faz várias passagens de limpeza antes de aplicar os regex de extração:

\begin{itemize}[leftmargin=*]
  \item \lstinline|stripComments|: remove comentários de bloco e de linha;
  \item \lstinline|expandPortHeader|: em portas ANSI, \lstinline|input a, b, c|
    declara três portas do mesmo tipo --- a vírgula separa nomes, não
    declarações. A função propaga o prefixo (tipo $+$ faixa) para as entradas
    seguintes;
  \item \lstinline|stripDirectives|: remove diretivas de pré-processador
    (\lstinline|`ifdef|, \lstinline|`define|, etc.) mantendo ambos os ramos;
  \item \lstinline|stripConditionalGenerates|: remove blocos
    \lstinline|generate if/case| --- como o parser não avalia parâmetros,
    capturá-los faria o \tool{iverilog} falhar com \enquote{Unable to bind};
  \item \lstinline|stripParamLists|: substitui blocos \lstinline|#(...)| por
    \lstinline|#()| (com balanceamento de parênteses e ciência de literais de
    string), evitando que o regex não-guloso confunda instâncias adjacentes.
\end{itemize}

A construção da hierarquia (\lstinline|buildHierarchyTree|) parte do módulo top
e desce por cada instanciação, quebrando ciclos com um conjunto de
ancestrais.

\section{A instrumentação do testbench}

O módulo \file{js/wave/testbench_instrumenter.ts} decide \emph{se} e
\emph{como} injetar \lstinline|$dumpfile| $+$ \lstinline|$dumpvars| num
testbench Verilog, e constrói o texto instrumentado. É uma função pura: recebe o
conteúdo original mais a seleção e retorna o novo texto (ou indica que nenhuma
escrita é necessária); a leitura/escrita de arquivos fica no fluxo de
compilação (\lstinline|instrumentTestbench|, que escreve
\file{Temp/instr\_<basename>}).

\subsection{A intenção e os argumentos do dump}

A intenção é dupla: testbenches que já têm \lstinline|$dumpfile| ou
\lstinline|$dumpvars| escrito à mão são deixados em paz (o usuário sabe o que
quer); caso contrário, a \tool{AURORA} injeta um bloco \lstinline|initial| antes
do último \lstinline|endmodule|. A lista de argumentos do \lstinline|$dumpvars|
vem de \lstinline|selectedSignals|:

\begin{itemize}[leftmargin=*]
  \item não-vazia $\rightarrow$ \lstinline|$dumpvars(0, sig1, sig2, ...)| ---
    exatamente os sinais escolhidos no modal Wave Configuration;
  \item vazia (default) $\rightarrow$ \lstinline|$dumpvars(1, <tb>)| --- sinais
    apenas no escopo do módulo do testbench.
\end{itemize}

\begin{lstlisting}[language=Verilog,caption={Forma do bloco injetado pela auto-instrumentação.}]
// --- AURORA AUTO-INSTRUMENTATION ---
// <comentario de cabecalho> <nota sobre a fonte da selecao>
initial begin
    $dumpfile("<tbModule>.vcd");
    $dumpvars(<args>);
end
// --------------------------------------------------
\end{lstlisting}

\subsection{Funções auxiliares}

\begin{description}[leftmargin=*]
  \item[\lstinline|hasUserDumpCalls(src)|] retorna verdadeiro se o source tem
    \lstinline|$dumpfile| ou \lstinline|$dumpvars| \emph{fora} de comentário
    (usa \lstinline|stripVerilogComments|, que é ciente de literais de string);
  \item[\lstinline|commentOutDumpCalls(src)|] substitui chamadas a
    \lstinline|$dumpfile|/\lstinline|$dumpvars| por comentários, neutralizando o
    efeito (usado no caso de \emph{override} e no Fast Sim);
  \item[\lstinline|lastEndmoduleIndex(src)|] acha o índice do último
    \lstinline|endmodule| que é um \emph{token} de verdade (fora de comentário,
    fora de string, delimitado por não-identificadores) --- garantindo que a
    injeção entre dentro do módulo.
\end{description}

O parâmetro \lstinline|overrideUserDumpvars| ativa o caso especial: quando o
testbench tem dump escrito à mão \emph{e} o usuário customizou a Wave
Configuration, a função comenta as chamadas originais e injeta o
\lstinline|$dumpvars| da \tool{AURORA}. O resultado carrega um campo
\lstinline|reason| diagnóstico: \lstinline|user-defined|, \lstinline|malformed|,
\lstinline|auto|, \lstinline|auto-selection| ou \lstinline|override-user|.

\begin{nota}
A simulação roda exatamente uma vez para todos os caminhos. Não há mais um
gate \lstinline|+AURORA_HEADER_ONLY| no testbench: o cabeçalho do VCD é puxado
do FST finalizado por \lstinline|_extractFstHeaderVcd|. O dump corre livre até o
\lstinline|$finish|, sem \emph{flush} periódico --- os \lstinline|$display| do
usuário saem em bloco quando a simulação termina.
\end{nota}

\section{A precedência de \texttt{\$dumpvars}: o que dumpar}

A decisão central do fluxo de ondas --- \emph{quais sinais entram no VCD} --- é
tomada por \lstinline|resolveWaveSelection| em
\file{js/compilation/wave\_signal\_validator.js}. Três eixos agem em conjunto,
numa ordem fixa de precedência (do maior para o menor). A
\autoref{tab:09-precedencia} resume.

\begin{table}[h]
\centering
\begin{tabularx}{\linewidth}{Y Y Z}
\toprule
\textbf{Prioridade} & \textbf{Fonte (\lstinline|source|)} & \textbf{Condição e efeito} \\
\midrule
1 (maior) & \lstinline|gtkw| & há \file{.gtkw} ativo (\lstinline|gtkwFiles[].isActive|); extrai refs do arquivo, valida e usa. \lstinline|override = true| \\
2 & \lstinline|wc| & \lstinline|wcCustomized| e \lstinline|waveSignals| não-vazio; o WC dita o dump. \lstinline|override = true| \\
3 & \lstinline|tb| & \lstinline|hadOriginalDumpvars|; nada é injetado, o testbench domina. \lstinline|override = false| \\
4 (menor) & \lstinline|default| & \lstinline|$dumpvars(1, tbModule)|; sinais só do escopo do testbench. \lstinline|override = false| \\
\bottomrule
\end{tabularx}
\caption{Precedência da resolução de \lstinline|$dumpvars|.}
\label{tab:09-precedencia}
\end{table}

\subsection{O algoritmo de resolução}

\lstinline|resolveWaveSelection| executa, em ordem:

\begin{enumerate}[leftmargin=*]
  \item lê o testbench e, na 1ª visita, registra o estado no \tool{WaveStore}
    com o snapshot \lstinline|hadOriginalDumpvars = hasUserDumpCalls(tbContent)|;
  \item \textbf{(.gtkw ativo)}: se há entrada \lstinline|isActive|, lê o arquivo,
    extrai os caminhos com \lstinline|extractSignalRefs|, valida-os contra a
    hierarquia parseada (\lstinline|validateSelection|) e retorna o conjunto
    válido com \lstinline|source: 'gtkw'|. Refs que sumiram do source geram aviso
    em \lstinline|twave| $+$ \emph{toast}; o build segue sem elas;
  \item \textbf{(Wave Config customizada)}: se \lstinline|wcCustomized| e
    \lstinline|waveSignals| não-vazio, valida a seleção salva (com auto-prune) e
    retorna com \lstinline|source: 'wc'|;
  \item \textbf{(testbench com dump à mão)}: se \lstinline|hadOriginalDumpvars|,
    retorna \lstinline|signalsToDump: []| e \lstinline|source: 'tb'| --- nada
    é instrumentado;
  \item \textbf{(default)}: retorna \lstinline|signalsToDump: []| e
    \lstinline|source: 'default'|.
\end{enumerate}

O parse do source é \emph{on-demand}: só ocorre se for preciso validar um
conjunto de sinais (vindo do \file{.gtkw} ou do WC), e o resultado é cacheado.
A fonte vencedora é registrada em \lstinline|twave| (\enquote{Wave source:
...}), o que ajuda a depurar quando o usuário esperava outra seleção.

\subsection{Como cada flag chega ao WaveStore}

A precedência depende de flags escritas por superfícies de UI distintas:

\begin{itemize}[leftmargin=*]
  \item \lstinline|isActive| em \lstinline|gtkwFiles| é marcada pelo \emph{gtkw
    picker} (\file{js/wave/gtkw\_picker.js}) ao escolher um \file{.gtkw} ativo no
    dropdown da barra;
  \item \lstinline|wcCustomized| é ligada pelo modal Wave Configuration
    (\file{js/wave/wave\_config\_manager.js}) quando a seleção salva difere da
    seleção inicial. A comparação é por conjunto (tamanhos iguais \emph{e} todos
    os elementos em comum); a flag só liga --- uma vez customizado, fica;
  \item \lstinline|hadOriginalDumpvars| é o snapshot da 1ª visita, escrito por
    \lstinline|ensureRegistered|.
\end{itemize}

\section{A validação de seleção}

O módulo \file{js/wave/selection\_validator.ts} filtra uma seleção salva contra
a hierarquia Verilog parseada atual. O picker armazena seleções como
\emph{strings} de caminho pontilhado (ex.: \lstinline|tb.dut.q_next|); quando o
usuário renomeia uma instância, remove um sinal ou edita o testbench, as
entradas salvas podem ficar pendentes (\emph{dangling}). Alimentar caminhos
pendentes a \lstinline|$dumpvars(0, ...)| faz o \tool{iverilog} falhar com erro
de elaboração críptico.

\lstinline|validateSelection(selectedSignals, hierarchyTree)| é pura e retorna
um \lstinline|SelectionSplit| com dois conjuntos: \lstinline|valid| (caminhos
ainda presentes na hierarquia) e \lstinline|dropped| (não resolvem mais, para
que o chamador avise o usuário). Detalhes de borda:

\begin{itemize}[leftmargin=*]
  \item seleção vazia $\rightarrow$ ambos os conjuntos vazios;
  \item \lstinline|hierarchyTree === null| (parse falhou, top ausente)
    $\rightarrow$ mantém a seleção como está --- é melhor deixar o
    \tool{iverilog} falar do que apagar silenciosamente as escolhas do usuário;
  \item caso normal: percorre a árvore coletando todo
    \lstinline|{scopePath}.{signal}| legal e separa a seleção contra esse
    conjunto.
\end{itemize}

\section{Os escritores de \texttt{.gtkw}}

A geração do layout \file{.gtkw} é dividida em dois módulos.

\subsection{\texttt{gtkw\_writer.ts}: inspeção}

\file{js/wave/gtkw\_writer.ts} contém \lstinline|extractSignalRefs|, que extrai
os caminhos pontilhados referenciados por um \file{.gtkw} existente. É usado
para (a) decodificar o que dumpar a partir de um \file{.gtkw} ativo
(precedência 1) e (b) cross-checar um \file{.gtkw} curado contra o VCD da
simulação atual. O \file{.gtkw} é um formato orientado a linha; a função pula
decorações (\lstinline|[*]|, \lstinline|@<hex>|, \lstinline|-<nome>|, etc.),
trata o prefixo de alias \lstinline|+{alias} <path>|, remove uma faixa final
\lstinline|[a:b]| e aceita como sinal apenas linhas com forma de identificador
pontilhado (começa com letra/underscore e contém pelo menos um ponto).

\subsection{\texttt{gtkw\_proc\_writer.ts}: o construtor}

\file{js/wave/gtkw\_proc\_writer.ts} contém \lstinline|buildAuroraGtkw|, o
construtor principal do layout. O arquivo resultante tem duas zonas:

\begin{enumerate}[leftmargin=*]
  \item \textbf{Top-level}: todos os sinais do VCD que NÃO estão dentro de uma
    instância de processador SAPHO --- a \enquote{glue logic} do design. Lista
    plana com formato escolhido por sinal: \lstinline|FMT_BIN| para escalares,
    \lstinline|FMT_SIGNED_DEC| para barramentos \lstinline|signed| (detectados
    pelo source via \lstinline|buildSignedSet|) e \lstinline|FMT_DEC| para os
    demais;
  \item \textbf{Seções por processador}: para cada instância de processador
    SAPHO detectada, uma seção completa com cores, aliases e grupos.
\end{enumerate}

\subsubsection{Detecção de processador}

Um escopo conta como instância de processador se contém \emph{ambos} os
registradores \lstinline|valr2| e \lstinline|linetabs| (que o \tool{asmcomp}
emite). O \lstinline|procType| (nome do módulo $=$ nome da pasta
\file{Temp/<procType>/} com os arquivos de tradução) é resolvido por
preferência: (1) \lstinline|moduleType| resolvido pelo source via
\lstinline|resolveScopeModules|; (2) padrão \lstinline|<inst>.p_<X>.core| no
VCD; (3) heurística do pai \lstinline|<X>_inst|; (4) fallback
\lstinline|<X>_tb|; (5) último recurso, o \lstinline|instanceName|. O
\lstinline|corePath| (onde vivem \lstinline|clk|/\lstinline|rst|/\lstinline|itr|
e Stack/ULA) prefere o sub-escopo \lstinline|.p_<X>.core| se existir.

\subsubsection{O layout por processador}

\begin{table}[h]
\centering
\begin{tabularx}{\linewidth}{Y Y Z}
\toprule
\textbf{Seção} & \textbf{Sinais} & \textbf{Formato / decoração} \\
\midrule
Banner & \lstinline|clk|, \lstinline|rst|, \lstinline|itr| & \lstinline|FMT_BIN|, do core \\
I/O & pares \lstinline|req_in_sim_N|/\lstinline|in_sim_N| e \lstinline|out_en_sim_N|/\lstinline|out_sig_N| & amarelo, alias \enquote{req\_in N} / \enquote{input N} etc. \\
Instructions & \lstinline|valr2| / \lstinline|linetabs| & indigo (\enquote{Assembly}, filtro \file{trad\_opcode.txt}) / violeta (\enquote{\cmm{}}, filtro \file{trad\_cmm.txt}) \\
Variables & \lstinline|me1_*| (int), \lstinline|me2_*| (float), \lstinline|comp_me3_*| (comp) e arrays & laranja; \lstinline|FMT_SIGNED_DEC|/\lstinline|FMT_REAL|/\lstinline|FMT_BIN| \\
Flags & grupo Stack (\lstinline|sp|/\lstinline|isp|) e grupo ULA (\lstinline|delta_int|/\lstinline|delta_float|) & \emph{analog step} $+$ signed/hex \\
\bottomrule
\end{tabularx}
\caption{Seções do layout per-processador gerado por \lstinline|buildAuroraGtkw|.}
\label{tab:09-gtkw-secoes}
\end{table}

Os formatos foram calibrados a partir de um \file{.gtkw} salvo pelo próprio
GTKWave (não por adivinhação). Cada linha de sinal é emitida por
\lstinline|emitSignal|, que escreve a flag \lstinline|@<hex>|, a cor opcional
\lstinline|[color] N|, os filtros inline e a referência (com alias
\lstinline|+{alias} <path>| quando há). Há dois tipos de filtro de tradução:

\begin{itemize}[leftmargin=*]
  \item \textbf{file filter} (tabela \file{.txt} como \file{trad\_opcode.txt},
    \file{trad\_cmm.txt}): sintaxe \lstinline|^N <path>|, acende
    \lstinline|TR_FTRANSLATED|;
  \item \textbf{process filter} (executável \tool{comp2gtkw.exe}, traduz binário
    bruto para o literal \cmm{}): sintaxe \lstinline|^>N <path>|, acende
    \lstinline|TR_PTRANSLATED|.
\end{itemize}

Quando \lstinline|selectedSignals| é fornecido, o layout completo é percorrido
(comentários, grupos, cores, tradutores) mas apenas sinais cujo caminho bate são
emitidos --- os aliases continuam refletindo a hierarquia inteira (posições
estáveis). \lstinline|buildAliasMap| reusa as mesmas regras para que o modal
Wave Configuration mostre os mesmos rótulos amigáveis.

\section{A decodificação de complexos}

O módulo \file{js/wave/complex\_decode.ts} é o pré-passe de decodificação dos
números complexos do SAPHO \emph{para o Surfer}. Diferença em relação ao
GTKWave: no GTKWave, os complexos são decodificados \emph{ao vivo} pelo
\emph{process filter} \tool{comp2gtkw.exe} (sintaxe \lstinline|^>N <exe>|, que
roda sobre os bits crus em tempo de render). O Surfer não tem esse filtro
externo e seu tradutor é uma tabela estática valor$\rightarrow$string. Portanto,
decodificar complexo no Surfer é um \emph{pré-passe}.

\subsection{O pré-passe}

Os prefixos dos sinais complexos que o \tool{asmcomp} emite são
\lstinline|comp_me3_| (escalar) e \lstinline|comp_arr_me3_| (array). O pré-passe:

\begin{enumerate}[leftmargin=*]
  \item usa \lstinline|ComplexVcdScanner|, um scanner \emph{incremental} do stream
    de texto VCD do \tool{fst2vcd}, para coletar o conjunto de valores
    \emph{distintos} (zero-extended à largura do sinal) de todos os sinais
    complexos. A memória é limitada (buffer de linha $+$ o \lstinline|Set| de
    distintos, com cap defensivo via \lstinline|maxValues|);
  \item decodifica cada valor com o próprio \tool{comp2gtkw.exe} (o binário
    canônico do GTKWave);
  \item monta um mapping \lstinline|bitpattern -> "re imi"| com
    \lstinline|buildComplexMapping| que o Surfer mostra direto.
\end{enumerate}

O módulo é \emph{puro} (sem \lstinline|fs|/\lstinline|spawn|): a orquestração (o
stream do \tool{fst2vcd} $+$ o IPC do \tool{comp2gtkw}) vive no
\lstinline|compilation_module|. Sobre o \emph{encoding} (de \file{comp2gtkw.c}):
os 16 primeiros bits são auto-descritivos --- \lstinline|bits[0:8]| é a largura
da mantissa (\lstinline|nbm|) e \lstinline|bits[8:16]| a largura do expoente
(\lstinline|nbe|); em seguida vêm \lstinline|re| e \lstinline|im|, cada um com
\lstinline|nbm + nbe + 1| bits. O módulo apenas \emph{extrai} o bitpattern; a
aritmética do decode roda no \tool{comp2gtkw.exe}.

\begin{nota}
O mapping do Surfer usa cabeçalho \lstinline|Name =| (sem \lstinline|Bits =|, a
largura varia por formato de processador e o match é por valor) seguido de uma
linha \lstinline|0b<bits> <re imi>| por valor decodificado. Sem entradas
decodificadas, \lstinline|buildComplexMapping| retorna \lstinline|null| --- não
se cria um mapping vazio, que o Surfer rejeitaria.
\end{nota}

\section{Resumo do fluxo de dados}

A \autoref{tab:09-resumo-modulos} mapeia cada módulo à sua responsabilidade no
fluxo de ondas, fechando o capítulo.

\begin{table}[h]
\centering
\begin{tabularx}{\linewidth}{Y Z}
\toprule
\textbf{Módulo / símbolo} & \textbf{Responsabilidade} \\
\midrule
\lstinline|runGtkWave| (\lstinline|compilation_module.js|) & orquestrador dos quatro caminhos \\
\lstinline|simulator_preference.js| & escolha Icarus vs Verilator (localStorage) \\
\lstinline|viewer_preference.js| & escolha GTKWave vs Surfer (localStorage) \\
\lstinline|wave_state_store.ts| & estado por testbench (\lstinline|WaveState|, JSON) \\
\lstinline|wave_signal_validator.js| & precedência do \lstinline|$dumpvars| (\lstinline|resolveWaveSelection|) \\
\lstinline|signal_parser.ts| & parser Verilog leve $\rightarrow$ hierarquia \\
\lstinline|selection_validator.ts| & filtra seleção salva contra a hierarquia \\
\lstinline|testbench_instrumenter.ts| & injeta/neutraliza \lstinline|$dumpvars| \\
\lstinline|gtkw_writer.ts| & \lstinline|extractSignalRefs| (inspeção de \file{.gtkw}) \\
\lstinline|gtkw_proc_writer.ts| & \lstinline|buildAuroraGtkw| (layout SAPHO) \\
\lstinline|complex_decode.ts| & pré-passe de complexos para o Surfer \\
\lstinline|_extractFstHeaderVcd| & convergência: cabeçalho VCD a partir do FST \\
\bottomrule
\end{tabularx}
\caption{Módulos do fluxo de ondas e suas responsabilidades.}
\label{tab:09-resumo-modulos}
\end{table}

% !TEX root = ../main.tex
\chapter{AURORA: visualização de ondas (GTKWave e Surfer)}
\label{cap:10-aurora-visualizacao-gtkwave-surfer}

Depois que a simulação produz um dump de formas de onda (\path{.vcd} ou \path{.fst}),
a AURORA oferece dois visualizadores externos para inspecioná-lo: o \tool{GTKWave}
(em um \emph{fork} próprio do laboratório, \repo{gtkwave-nipscern}) e o \tool{Surfer}
(viewer escrito em Rust, integrado de forma \emph{opt-in}). O \tool{GTKWave} é o
visualizador padrão e vem empacotado com a IDE; o \tool{Surfer} é uma alternativa
opcional que o usuário liga por um \emph{toggle}, e que cai de volta para o
\tool{GTKWave} caso não esteja instalado.

Este capítulo descreve os dois visualizadores a partir do código-fonte:
\textbf{(A)} as customizações de UI/comportamento do \emph{fork} \tool{GTKWave}
\repo{gtkwave-nipscern} frente ao \emph{upstream}, como ele é construído (Meson +
MSYS2) e como a AURORA o invoca; e \textbf{(B)} o \tool{Surfer} como viewer
\emph{opt-in}, com ênfase no \emph{layout} declarativo \path{.surf.ron} gerado por
\lstinline|buildSurferLayout| (grupos colapsáveis por processador, cores, aliases,
tracks de Assembly/\cmm{} via \emph{mapping translators}, decode de números
complexos via pré-passo). A geração do arquivo de \emph{save} \path{.gtkw} para o
\tool{GTKWave} é o tema do \autoref{cap:09-aurora-simulacao-ondas}; aqui ela é
mencionada apenas de passagem, no que toca aos dois visualizadores.

\section{Visão geral: dois visualizadores, um pipeline}
\label{sec:cap10-visao-geral}

A escolha de visualizador é uma preferência global do usuário, persistida em
\lstinline|localStorage| sob a chave \lstinline|aurora.waveViewer|, com dois valores
válidos: \lstinline|'gtkwave'| (default) e \lstinline|'surfer'|. O módulo
\file{js/wave/viewer\_preference.js} expõe \lstinline|getViewer()| e
\lstinline|setViewer(value)|; valores inválidos normalizam para \lstinline|'gtkwave'|,
e \lstinline|getViewer()| nunca lança (é chamado no \emph{hot path}, a cada clique no
botão Wave).

No passo \enquote{Wave} do fluxo de compilação (\file{js/compilation/compilation\_module.js})
o ramo é simples: depois de localizar o dump e extrair o cabeçalho VCD, se
\lstinline|getViewer() === 'surfer'| a AURORA resolve o \emph{layout} do Surfer e o
lança; caso contrário segue o caminho do \tool{GTKWave}.

\begin{lstlisting}[language=Java,caption={Ramo do visualizador no passo Wave (compilation\_module.js).}]
await this._extractFstHeaderVcd(vcdFile, headerVcd, tools.fst2vcdBin, tools.tempBaseDir);
if (getViewer() === 'surfer') {
    const surferLayout =
        await this._waveResolveSurferSaveFile(simTopModule, vcdFile, tools.tempBaseDir);
    await this._waveLaunchSurfer(vcdFile, surferLayout, tools);
}
// senao: caminho GTKWave (auto-.gtkw + _waveLaunchGtkwave)
\end{lstlisting}

\begin{nota}
Ambos os viewers leem o \emph{mesmo} dump de simulação e ambos são processos
\emph{externos} lançados de forma \emph{detached} pelo processo \emph{main} do
Electron (não há embutimento na janela da IDE). A relação entre o ramo
\enquote{surfer} e o \enquote{gtkwave} é de \emph{fallback}: se o \path{surfer.exe}
não existe, \lstinline|_waveLaunchSurfer| degrada para o \tool{GTKWave}, de modo que
o botão Wave \emph{sempre} produz um visualizador.
\end{nota}

A \autoref{tab:cap10-viewers} resume o contraste entre os dois viewers conforme o
código.

\begin{table}[h]
\centering
\caption{\tool{GTKWave} (fork nipscern) vs.\ \tool{Surfer} na AURORA.}
\label{tab:cap10-viewers}
\begin{tabularx}{\linewidth}{l Y Y}
\toprule
 & \textbf{GTKWave (nipscern)} & \textbf{Surfer} \\
\midrule
Status     & padrão, empacotado          & \emph{opt-in}, opcional      \\
Origem     & \emph{fork} C/GTK3 do lab   & binário Rust de \emph{upstream} \\
Curadoria  & arquivo \path{.gtkw}        & arquivo \path{.surf.ron}     \\
Decode ASM/\cmm & \emph{file filter} (\path{trad\_*.txt}) & \emph{mapping translator} \\
Decode complexo & \emph{process filter} (\tool{comp2gtkw}) ao vivo & pré-passo (\tool{comp2gtkw}) assado em mapping \\
Local instalação & \path{components/Packages/gtkwave-nipscern/} & \path{components/Packages/surfer/} \\
Lançamento & \lstinline|launch-gtkwave-only| & \lstinline|launch-surfer| \\
\bottomrule
\end{tabularx}
\end{table}

\section{O \emph{fork} GTKWave do nipscern}
\label{sec:cap10-gtkwave-fork}

O repositório \repo{gtkwave-nipscern} é um \emph{fork} do \tool{GTKWave} de
\emph{upstream} (que parte da base v3.3.116, já migrada de autotools para Meson e
de GTK2 para GTK3). Sobre essa base, o laboratório aplicou um conjunto de
customizações de interface e comportamento, além de tornar a build no Windows
viável sem WSL. As customizações foram introduzidas em quatro \emph{commits}
sequenciais (verificados no histórico do \emph{fork}):

\begin{enumerate}[label=\arabic*.]
  \item \lstinline|b83a5ee| — \enquote{Add nipscern customizations}: \lstinline|--zoom-fit|,
        painel de sinais escuro, remoção do SST, \lstinline|--left-justify|;
  \item \lstinline|6aa80b2| — link do \path{gtkwave.exe} com o \emph{subsystem}
        Windows GUI;
  \item \lstinline|f17e0b5| — filhos lançados com \lstinline|CREATE_NO_WINDOW| no
        Windows;
  \item \lstinline|4491579| — \enquote{Surfer dark theme, Phosphor icons, custom
        title bar, animated zoom + Windows build}: tema escuro padrão alinhado ao
        Surfer, ícones Phosphor, barra de título customizada e \emph{zoom} animado.
\end{enumerate}

\subsection{Flags de linha de comando adicionadas}
\label{sec:cap10-gtkwave-flags}

Duas opções novas de CLI foram registradas em \file{src/main.c}, na tabela
\lstinline|long_options| consumida por \lstinline|getopt_long|.

\subsubsection{\lstinline|-Z| / \lstinline|--zoom-fit| — zoom-fit inicial}

Força um \emph{Zoom Fit} (o mesmo do botão da \emph{toolbar}) assim que o
\emph{main loop} entra, de modo que a forma de onda inteira fique visível no
\emph{startup}. No código: o \lstinline|case 'Z'| seta
\lstinline|GLOBALS->do_initial_zoom_fit = 1|, e logo antes de \lstinline|gtk_main()|
um \lstinline|g_idle_add(initial_zoom_fit_idle_cb, NULL)| agenda a aplicação. O
\emph{callback} ocioso chama \lstinline|service_zoom_fit(NULL, NULL)| e marca
\lstinline|do_initial_zoom_fit_used = 1| para não repetir.

\begin{lstlisting}[language=C,caption={Zoom-fit inicial em src/main.c (verificado).}]
static gboolean initial_zoom_fit_idle_cb(gpointer data)
{
    service_zoom_fit(NULL, NULL);
    GLOBALS->do_initial_zoom_fit_used = 1;
    /* ... */
}
/* ... no parser de opcoes ... */
case 'Z':
    GLOBALS->do_initial_zoom_fit = 1;
    break;
/* ... antes de gtk_main(): */
if (GLOBALS->do_initial_zoom_fit) {
    g_idle_add(initial_zoom_fit_idle_cb, NULL);
}
\end{lstlisting}

\subsubsection{\lstinline|-L| / \lstinline|--left-justify| — nomes à esquerda}

Alinha os nomes de sinal à \emph{esquerda} no painel (em vez do alinhamento padrão
à direita). Expõe na CLI o \emph{rcvar} \lstinline|left_justify_sigs| já existente
(também acessível pelo menu \textbf{Edit $\rightarrow$ Set Left Justify Signals}). No
\lstinline|case 'L'|: \lstinline|GLOBALS->left_justify_sigs = ~0|.

\subsection{Painel de sinais escuro}
\label{sec:cap10-gtkwave-darkpanel}

A flag \lstinline|-6| / \lstinline|--dark| de \emph{upstream} apenas setava
\lstinline|gtk-application-prefer-dark-theme|, o que não afetava a lista de sinais
desenhada via Cairo. O \emph{fork} introduziu uma paleta escura própria para esse
painel. Posteriormente (\emph{commit} \lstinline|4491579|) o tema escuro foi
promovido a \textbf{padrão}: em \file{src/globals.c}, o campo \lstinline|use_dark|
é inicializado como \lstinline|TRUE| (comentário: \enquote{nipscern: dark theme is
the default}). A paleta do canvas e do painel de nomes foi recolorida estritamente
para a paleta \emph{default-dark} do Surfer.

\subsection{Remoção do painel SST (Signal Search Tree)}
\label{sec:cap10-gtkwave-sst}

A árvore de hierarquia do lado esquerdo (SST) deixou de ser construída: a janela
principal contém apenas a lista de sinais e a área de formas de onda. Em
\file{src/main.c} o bloco que montava \lstinline|toppanedwindow| /
\lstinline|sstpane| / \lstinline|expanderwindow| foi substituído por atribuições
\lstinline|NULL| (verificado no código atual):

\begin{lstlisting}[language=C,caption={SST desativado em src/main.c.}]
/* SST panel disabled: do not create toppanedwindow/sstpane/expanderwindow */
GLOBALS->toppanedwindow = NULL;
GLOBALS->sstpane = NULL;
GLOBALS->expanderwindow = NULL;
\end{lstlisting}

A página do \emph{notebook} já trazia um \emph{fallback}
\lstinline|toppanedwindow ? toppanedwindow : panedwindow|, então passa a usar o
\lstinline|panedwindow| diretamente. A reconstrução da SST em tempo de
\emph{reload} (em \file{src/globals.c}) fica protegida por
\lstinline|if (GLOBALS->expanderwindow)|, tornando-se um \emph{no-op}.

\subsection{Tema visual: paleta Surfer, ícones Phosphor, barra de título e zoom animado}
\label{sec:cap10-gtkwave-theme}

O \emph{commit} \lstinline|4491579| fez uma reformulação visual da GUI, alinhando-a
à paleta padrão (escura) do Surfer para dar consistência visual entre os dois
viewers do SAPHO.

\paragraph{Paleta do canvas e do painel.} Em
\file{lib/libgtkwave/src/gw-color-theme.c}, o canvas e o painel de nomes são
recoloridos com os valores da paleta \emph{default-dark} do Surfer. Alguns valores
verificados no código: fundo do canvas \lstinline|#0b151d|, \emph{foreground}
\lstinline|#d4d4d4|, traços de variável em verde \lstinline|#56c126|, cursor
\lstinline|#b63935|, \emph{undef} \lstinline|#f44747|, \emph{high-impedance}
\lstinline|#c9e124|, \emph{don't-care} \lstinline|#4040ff|, \emph{weak}
\lstinline|#808080|.

\paragraph{Folha de estilo da GUI.} O arquivo \file{src/gtkwave.css} foi reescrito
para re-tintar as cores nomeadas do Adwaita para a paleta Surfer, definir fontes
(\enquote{Segoe UI} / \enquote{Cascadia Mono}) e espaçamentos. Ele é carregado em
prioridade de aplicação sobre o Adwaita-dark. Referência de paleta no cabeçalho do
CSS: \lstinline|foreground #d4d4d4|, \lstinline|primary_ui #171717|,
\lstinline|secondary_ui #0d1317|, \lstinline|canvas #0b151d|,
\lstinline|accent_info #007acc|, \lstinline|accent_error #d2302f|.

\paragraph{Ícones Phosphor (MIT).} 39 SVGs pré-recoloridos sob
\path{share/icons/phosphor}, organizados por código de cor (azul = estrutura, verde
= sinais, amarelo = controle, laranja/rosa = comportamental). Em \file{src/main.c} o
helper \lstinline|gw_phosphor_image(name)| carrega cada ícone diretamente do
\emph{gresource} pelo caminho
\path{/io/github/gtkwave/GTKWave/icons/phosphor/<name>.svg}.

\paragraph{Barra de título customizada (CSD).} No Windows com GTK3, o
\lstinline|gtk_header_bar_set_show_close_button()| frequentemente não renderiza os
controles, então o \emph{fork} substitui a barra de título nativa por um
\lstinline|GtkHeaderBar| com botões próprios de minimizar/maximizar/fechar. Em
\file{src/main.c} (verificado): cria-se o \lstinline|headerbar|, desliga-se o botão
de fechar nativo (\lstinline|set_show_close_button(..., FALSE)|), empacotam-se os
três botões (\lstinline|win-minimize|, \lstinline|win-maximize|,
\lstinline|win-close|) com \lstinline|gtk_header_bar_pack_end| e instala-se via
\lstinline|gtk_window_set_titlebar|. O \emph{callback} de maximizar alterna o ícone
entre \lstinline|win-maximize| e \lstinline|win-restore|.

\paragraph{Zoom animado.} Em \file{src/zoombuttons.c} o \emph{zoom in/out} passa a
suavizar com \emph{easing} (\emph{ease-out cubic}, 150\,ms) em vez de saltar
instantaneamente, dirigido por um \emph{tick callback} do \emph{frame clock} do
widget. Constantes verificadas: \lstinline|GW_ZOOM_ANIM_DURATION_US = 150000| (150\,ms)
e \lstinline|gw_zoom_animate_enabled = TRUE|. A interpolação aplica
\lstinline|gw_ease_out_cubic(t)| entre \lstinline|from| e \lstinline|to| a cada tick.

\subsection{Subsystem GUI e filhos sem console no Windows}
\label{sec:cap10-gtkwave-windows}

Dois ajustes específicos do Windows garantem uma experiência sem janelas de
console espúrias:

\begin{itemize}
  \item \textbf{Subsystem GUI.} Em \file{src/meson.build} o alvo \lstinline|gtkwave|
        é construído com \lstinline|win_subsystem: 'windows'|, de modo que nenhum
        console abre junto com a GUI. (O comentário do \emph{fork} registra que
        passar a flag via \lstinline|link_args| cru não funciona, pois o Meson
        anexa \lstinline|-Wl,--subsystem,console| \emph{depois} dos argumentos do
        usuário, e o \emph{ld} respeita o último \lstinline|--subsystem|.)
  \item \textbf{Filhos com \lstinline|CREATE_NO_WINDOW|.} Como o
        \path{gtkwave.exe} é GUI-subsystem, qualquer filho de subsistema console
        que ele lançasse abriria, por uma fração de segundo, uma janela tipo
        \lstinline|cmd|. As \emph{flags} de \lstinline|CreateProcess| em três
        pontos de chamada foram mudadas de \lstinline|0| para
        \lstinline|CREATE_NO_WINDOW|: o \emph{process filter}
        (\file{src/pipeio.c}, usado por filtros como \tool{comp2gtkw}), o
        \enquote{open a second viewer} (\file{src/menu.c}) e o \emph{reader} de
        \path{.stems}/rtlbrowse (\file{src/main.c}).
\end{itemize}

\subsection{Construção no Windows (Meson + MSYS2)}
\label{sec:cap10-gtkwave-build}

O \emph{fork} traz uma build de um clique para Windows nativo (sem WSL). O ponto de
entrada é \file{build-windows.bat} na raiz do repositório; ele localiza o MSYS2
(instalando-o via \lstinline|winget| se ausente), configura o ambiente
\lstinline|MINGW64| e delega para \file{tools/msys2-build.sh}.

\begin{lstlisting}[language=bash,caption={Uso do build-windows.bat (do README do fork).}]
build-windows.bat          # build + install (incremental)
build-windows.bat clean    # apaga build/ e refaz do zero
build-windows.bat run      # build, instala e roda gtkwave --dark
\end{lstlisting}

Detalhes verificados no \emph{script} de build:

\begin{itemize}
  \item \textbf{Prefixo de instalação.} O default é
        \path{C:/packs/gtkwave-bin} (variável \lstinline|GTKWAVE_PREFIX|); o
        MSYS2 pode ser sobrescrito por \lstinline|GTKWAVE_MSYS2|.
  \item \textbf{Configuração Meson.} \file{tools/msys2-build.sh} chama
        \lstinline|meson setup --prefix="$PREFIX" -Djudy=disabled -Dtests=false build|.
  \item \textbf{Launcher.} Gera um \path{gtkwave.cmd} no prefixo que ajusta o
        \lstinline|PATH| das DLLs do GTK antes de invocar o \path{gtkwave.exe}.
  \item \textbf{Bundle para a AURORA.} O passo final invoca
        \file{tools/deploy-aurora.sh}, que reempacota a build no \emph{package}
        portátil da AURORA (exe + DLLs do GTK via \lstinline|ldd| + \emph{pixbuf
        loaders} + \path{loaders.cache} portável).
  \item \textbf{Hooks só de Linux pulados.} Em \file{meson.build}, os
        \lstinline|gnome.post_install| de atualização de bancos
        \emph{desktop}/\emph{mime} são desativados no Windows
        (\lstinline|is_windows = host_machine.system() == 'windows'|).
\end{itemize}

\subsection{Empacotamento e \emph{bootstrap} pela AURORA}
\label{sec:cap10-gtkwave-bootstrap}

A AURORA não compila o \tool{GTKWave}: ela baixa um \emph{bundle} Windows portátil
publicado nas \emph{releases} de \repo{gtkwave-nipscern}. O \emph{script}
\file{components/Scripts/download-gtkwave-nipscern.js} roda no \emph{prestart},
depois do \texttt{download-toolchain} e antes do \texttt{copy-components}.

\begin{itemize}
  \item \textbf{Pinagem de versão.} Constantes no topo do \emph{script}:
        \lstinline|GTKWAVE_TAG = 'v0.1.2-nipscern'|, \emph{filename}
        \path{gtkwave-nipscern-v0.1.2-win-x64.zip}, \emph{owner}
        \lstinline|nipscernlab|, \emph{repo} \lstinline|gtkwave-nipscern|.
  \item \textbf{Destino.} Extraído em
        \path{components/Packages/gtkwave-nipscern/}; o arquivo-sentinela é
        \path{gtkwave.exe} nessa pasta (o ZIP não tem prefixo de pasta, então a
        extração vai direto para o diretório de instalação).
  \item \textbf{Bundle portátil.} Traz \path{gtkwave.exe} + \path{fst2vcd.exe} +
        todas as DLLs do \emph{runtime} GTK, sem depender de um MSYS2 instalado.
        Substitui o \tool{GTKWave} do \emph{toolchain} principal (o
        \tool{iverilog}/\tool{yosys} continuam vindo do outro bundle).
  \item \textbf{Integridade.} A verificação SHA-256 (\lstinline|EXPECTED_SHA256|)
        está como \lstinline|null| neste \emph{script} (computa e registra, sem
        impor). A falha de download é \emph{best-effort}: sai com código 0 para
        não bloquear o \lstinline|npm start| (a AURORA ainda compila/simula; só o
        botão Wave falha até o setup).
\end{itemize}

\subsection{Como a AURORA invoca o GTKWave}
\label{sec:cap10-gtkwave-invoca}

A linha de comando é montada em \file{js/compilation/builders/wave\_tools.ts}
(compilado para \path{wave\_tools.js}). O \emph{builder}
\lstinline|buildGtkwaveSpec| produz os argumentos \lstinline|--dark --zoom-fit
--left-justify <vcd>| e, quando há arquivo de \emph{save}, anexa \lstinline|-a
<gtkw>|. Como o SST já vem removido do \emph{fork}, não há necessidade de um
\lstinline|--rcvar hide_sst on|.

\begin{lstlisting}[language=Java,caption={buildGtkwaveSpec (wave\_tools.ts, verbatim).}]
export function buildGtkwaveSpec(ctx: GtkwaveBuilderCtx): CommandSpec {
  const args = ['--dark', '--zoom-fit', '--left-justify', ctx.vcdFile];
  if (ctx.gtkwSaveFile) {
    args.push('-a', ctx.gtkwSaveFile);
  }
  return {
    step: 'gtkwave',
    binary: ctx.gtkwaveBin,
    args,
    cwd: ctx.cwd,
    label: 'gtkwave --dark --zoom-fit --left-justify',
  };
}
\end{lstlisting}

O \emph{spawn} real ocorre no processo \emph{main} via IPC
\lstinline|launch-gtkwave-only| (\file{main/ipc/compile.js}). Pontos verificados:

\begin{itemize}
  \item O binário recebido do \emph{renderer} passa pelo \emph{allowlist} de
        \emph{toolchain} (\lstinline|isAllowed|): \path{gtkwave.exe} só é aceito a
        partir de \path{Packages/gtkwave-nipscern}.
  \item Lançado \emph{detached}, com \lstinline|stdio: 'ignore'| e
        \lstinline|shell: false|; é \lstinline|unref|'d para sobreviver a uma única
        execução, mas é rastreado (\lstinline|spawnTracked|) para morrer junto com
        a IDE.
  \item \lstinline|windowsHide: false| de propósito: um \emph{quirk} do
        \emph{fork} GUI-subsystem faz com que, sob \lstinline|CREATE_NO_WINDOW| e
        recebendo um VCD para carregar, o processo suba mas a janela \emph{top-level}
        não seja criada. Como o exe é GUI-subsystem, não há \emph{flash} de console
        de qualquer modo.
\end{itemize}

\subsection{O papel do \path{.gtkw} e do \tool{fst2vcd}}
\label{sec:cap10-gtkwave-gtkw-fst2vcd}

O \tool{GTKWave} recebe, além do dump, um arquivo de \emph{save} \path{.gtkw} que a
AURORA gera automaticamente (\lstinline|buildAuroraGtkw|, em
\file{js/wave/gtkw\_proc\_writer.ts}). Esse arquivo carrega a curadoria SAPHO:
seções por processador, cores, formatos e \emph{aliases}. Dois mecanismos do
próprio \path{.gtkw} são relevantes para o contraste com o Surfer (\autoref{sec:cap10-surfer}):

\begin{itemize}
  \item \textbf{\emph{File filter} (\path{trad\_*.txt}).} As trilhas de Assembly
        (\lstinline|valr2|) e de \cmm{} (\lstinline|linetabs|) usam tabelas
        \path{trad\_opcode.txt} / \path{trad\_cmm.txt} via a sintaxe
        \lstinline|^N <path>| do \path{.gtkw}.
  \item \textbf{\emph{Process filter} (\tool{comp2gtkw}).} Os números complexos
        (\lstinline|comp_me3_*| / \lstinline|comp_arr_me3_*|) são decodificados
        \emph{ao vivo} por um executável externo, via a sintaxe
        \lstinline|^>N <path>| (\emph{caret} + \emph{maior-que}). É justamente esse
        \emph{process filter} um dos filhos beneficiados pelo
        \lstinline|CREATE_NO_WINDOW| (\autoref{sec:cap10-gtkwave-windows}).
\end{itemize}

O \tool{fst2vcd} (que vem no \emph{bundle} do \emph{fork}, em
\path{components/Packages/gtkwave-nipscern/fst2vcd.exe}) é a única CLI do
\tool{GTKWave} que a AURORA usa diretamente. Ele converte o FST produzido pela
simulação em VCD-texto e — no caminho do Surfer — é \emph{streamado} para extrair só
o cabeçalho ou só os valores necessários (ver \autoref{sec:cap10-surfer-prepass}). A
resolução do caminho desses binários fica em
\file{js/compilation/wave\_toolchain.js} (\lstinline|resolveWaveToolchain|).

\section{O visualizador Surfer (opt-in)}
\label{sec:cap10-surfer}

O \tool{Surfer} (\url{https://surfer-project.org/}) é um visualizador de formas de
onda escrito em Rust/egui que lê o mesmo VCD/FST. A AURORA o trata como um
\path{surfer.exe} \emph{standalone} opcional sob \path{components/Packages/surfer/}.

\begin{nota}
\textbf{Estado real da integração (confirmado no código):} o Surfer é
\textbf{opt-in}. O default do visualizador é \lstinline|'gtkwave'|
(\file{viewer\_preference.js}), e o Surfer \textbf{não é empacotado por padrão} — o
comentário em \file{wave\_toolchain.js} é explícito: \enquote{NOT bundled by default
— \lstinline|_waveLaunchSurfer| degrades to GTKWave with a friendly message if it's
absent}. Quando o \path{surfer.exe} está ausente, o \emph{launch} reporta um
\enquote{not found} limpo e a AURORA cai de volta para o \tool{GTKWave}. A integração
está completa para o uso a que se propõe; o que não vem ligado/instalado é apenas o
\emph{default} e o binário.
\end{nota}

\subsection{Download e integração}
\label{sec:cap10-surfer-download}

O \emph{script} \file{components/Scripts/download-surfer.js} baixa o build Windows do
Surfer e extrai \path{surfer.exe} em \path{components/Packages/surfer/}. Roda no
\emph{bootstrap}, depois do \texttt{download-gtkwave-nipscern} e antes do
\texttt{copy-components}. Pontos verificados:

\begin{itemize}
  \item \textbf{Fonte e pinagem.} \lstinline|SURFER_TAG = 'v0.7.0'|, baixado do
        \emph{registro de pacotes genérico} do GitLab do projeto Surfer (id
        \lstinline|42073614|), arquivo \path{surfer\_win\_v0.7.0.zip}.
  \item \textbf{Integridade.} Diferente do \tool{GTKWave}, aqui o SHA-256 é
        \textbf{imposto}: \lstinline|EXPECTED_SHA256| traz o \emph{hash} pinado e
        uma divergência aborta a instalação (\lstinline|verifyChecksum| antes de
        extrair).
  \item \textbf{Sentinela e \emph{flatten}.} O arquivo-sentinela é
        \path{surfer.exe} no diretório de instalação; \lstinline|flattenIfNested|
        sobe os arquivos um nível caso o ZIP traga \path{surfer/surfer.exe}.
  \item \textbf{Best-effort.} Falha de download sai com código 0; o botão Wave
        (com Surfer) cai para o \tool{GTKWave} até o setup ser feito.
  \item \textbf{Licença.} O comentário registra EUPL-1.2 com atribuição no
        \path{LICENSE} da raiz; o \emph{spawn} \emph{arm's-length} (a AURORA só
        executa o \path{.exe}, não linka) não contamina a IDE.
\end{itemize}

\subsection{Lançamento do Surfer}
\label{sec:cap10-surfer-launch}

O método \lstinline|_waveLaunchSurfer| (\file{compilation\_module.js}) monta os
argumentos e delega ao IPC \lstinline|launch-surfer|. O VCD vai posicional; o
\emph{layout} ativo vem depois: um \emph{save state} \path{.surf.ron} via
\lstinline|-s|, ou um \emph{command file} \path{.sucl} via \lstinline|-c|. O VCD da
CLI tem precedência sobre qualquer caminho embutido no \emph{state}, de modo que um
\path{.surf.ron} registrado fica portável entre re-execuções (os itens re-vinculam por
nome).

\begin{lstlisting}[language=Java,caption={Montagem de argumentos do Surfer (\_waveLaunchSurfer).}]
const args = [vcdFile];
if (surferLayoutFile) {
    const flag = /\.sucl$/i.test(surferLayoutFile) ? '-c' : '-s';
    args.push(flag, surferLayoutFile);
}
const result = await electronAPI.launchSurfer({
    surferBin: tools.surferBin,
    args,
    workingDir: tools.tempBaseDir,
    multiWindow: getSurferMultiWindow(),
});
\end{lstlisting}

No \emph{main} (\file{main/ipc/compile.js}), o IPC \lstinline|launch-surfer|:

\begin{itemize}
  \item Pré-checa o binário com \lstinline|fs.existsSync| e o \emph{allowlist}
        (\path{surfer.exe} só de \path{Packages/surfer}), retornando um
        \enquote{not found} limpo quando ausente.
  \item \textbf{Uma janela por padrão.} Se \lstinline|multiWindow| é falso, mata o
        Surfer que a AURORA abriu antes (rastreado em \lstinline|lastSurferChild|,
        comparando o \emph{mesmo} objeto \lstinline|ChildProcess| para evitar reuso
        de PID), via \lstinline|killProcessSilently|. A preferência
        \lstinline|getSurferMultiWindow()| (chave \lstinline|aurora.surferMultiWindow|,
        default false) liga o modo \enquote{várias janelas} para comparar
        simulações lado a lado.
  \item \textbf{Janela centralizada.} Antes do \emph{spawn}, chama
        \lstinline|writeSurferCenteredWindowConfig|, que escreve geometria
        (\textasciitilde85\% da área de trabalho do display primário, centralizada)
        no \path{config.toml} \emph{global} do Surfer
        (\path{\%APPDATA\%/surfer-project/surfer/config/config.toml}). Um marcador
        \lstinline|# Managed by AURORA| protege um config escrito à mão (se o
        arquivo existe sem o marcador, a AURORA não o sobrescreve).
  \item \emph{Spawn} \emph{detached}, \lstinline|stdio: 'ignore'|,
        \lstinline|shell: false|, rastreado para morrer junto com a IDE.
\end{itemize}

\begin{nota}
O \emph{auto-reload} do Surfer (\lstinline|SurferConfig.autoreload_files|) não
dispara no Windows v0.7.0 — o \emph{watcher} de arquivo não detecta a reescrita do
FST. Por isso a AURORA fecha a janela anterior e reabre, em vez de depender do
reload automático. O campo \lstinline|autoreload_files| no \emph{state}
\path{.surf.ron} é inerte; quem controla o reload é o \path{config.toml}.
\end{nota}

\subsection{O layout declarativo \path{.surf.ron}}
\label{sec:cap10-surfer-ron}

A peça central da integração é o escritor de \emph{layout}
\file{js/wave/surfer\_layout\_writer.ts} (compilado para
\path{surfer\_layout\_writer.js}). Ele é a contraparte \emph{declarativa} do
\path{.gtkw}: enquanto o \path{.gtkw} é um \emph{savefile} do GTKWave, o
\path{.surf.ron} é um \emph{save state} do Surfer (formato RON), em que cada item
carrega sua própria cor/formato/análogo/nome.

O comentário de cabeçalho do módulo explica a escolha por \path{.surf.ron} em vez
de um \emph{command file} \path{.sucl}: o \path{.sucl} não monta grupos curados, não
seta \emph{display} analógico, e o \emph{targeting} por item (\lstinline|item_focus|
por um índice base-16 frágil) erra silenciosamente. O \path{.surf.ron} é
declarativo, e o Surfer \textbf{re-resolve} os itens por caminho hierárquico + nome
no \emph{load} (os IDs Wellen são apenas dica de performance, recomputados), de modo
que IDs \emph{placeholder} sequenciais ficam portáveis entre re-execuções.

O módulo separa duas camadas:

\begin{description}
  \item[Camada de formato — \lstinline|buildSurferState|] transforma uma lista
        \emph{ordenada} de \lstinline|SurferItem| em RON válido.
  \item[Camada de curadoria — \lstinline|buildSurferLayout|] decide \emph{quais}
        sinais e \emph{quais} cores/formatos, reusando a detecção de processador do
        \file{gtkw\_proc\_writer} (ver \autoref{sec:cap10-surfer-curadoria}).
\end{description}

\subsubsection{Tipos de item}

Os itens declarativos (\lstinline|SurferItem|) são quatro:

\begin{tabularx}{\linewidth}{l Y}
\toprule
\textbf{kind} & \textbf{Campos principais} \\
\midrule
\lstinline|variable| & \lstinline|scope[]|, \lstinline|name|, \lstinline|format|, \lstinline|color|, \lstinline|manualName| (alias), \lstinline|heightScale|, \lstinline|analog| \\
\lstinline|divider|  & \lstinline|name|, \lstinline|color| \\
\lstinline|timeline| & \lstinline|name| \\
\lstinline|group|    & \lstinline|name|, \lstinline|color|, \lstinline|isOpen|, \lstinline|children[]| \\
\bottomrule
\end{tabularx}

\subsubsection{Grupos colapsáveis e a árvore de itens}

Um \lstinline|group| é um \lstinline|DisplayedItem::Group| do Surfer. Detalhe
verificado contra o \emph{save} nativo do Surfer v0.7.0: os filhos \textbf{não} são
listados em \lstinline|content| (que serializa sempre como \lstinline|[]|); a
hierarquia é definida \emph{apenas} pelos níveis (\lstinline|level|) na
\lstinline|items_tree| — os filhos vêm logo após o nó do grupo, com
\lstinline|level + 1|. O campo \lstinline|is_open| persiste o estado dobrado
(\lstinline|true| = expandido).

A travessia em \lstinline|buildSurferState| é \emph{depth-first in-order}: os
\lstinline|ref| são sequenciais (1-based) na ordem de visita, o nó do grupo vem
\emph{antes} dos filhos, e cada item gera tanto uma linha na \lstinline|items_tree|
quanto uma entrada no \lstinline|displayed_items|.

\begin{lstlisting}[language=Java,caption={Travessia in-order que numera refs e níveis (buildSurferState).}]
const visit = (item, level) => {
  const ref = nextRef++;
  if (item.kind === 'group') {
    nodes.push({ ref, level, unfolded: item.isOpen ?? true });
    for (const child of item.children) visit(child, level + 1);
    entries.push(emitGroup(ref, item));
  } else if (item.kind === 'variable') {
    nodes.push({ ref, level, unfolded: true });
    /* scope id + var id placeholders, re-resolvidos por nome no load */
    entries.push(emitVariable(ref, item, sid, nextVarId++));
  } /* divider / timeline analogos */
  return ref;
};
\end{lstlisting}

\subsubsection{Esqueleto do estado}

O esqueleto externo do \path{.surf.ron} são os \emph{defaults} estáveis do
\lstinline|UserState| (todos os \emph{toggles} \lstinline|None|; \lstinline|frame_buffer|
e \lstinline|variable_filter| nos defaults). Só o bloco \lstinline|waves| é derivado:
a \lstinline|source: File(...)| (dica; o VCD posicional sobrescreve), o
\lstinline|format| (\lstinline|Vcd|/\lstinline|Fst|, magic-detectado de qualquer
modo), a \lstinline|items_tree|, o \lstinline|displayed_items|, os
\lstinline|viewports| e os \lstinline|markers|.

\subsubsection{Marcadores automáticos e janela de delta}

Quando há \lstinline|markerTimes| (tempos de evento), eles viram entradas
\lstinline|<id>: (1, [<tempo>])| no campo \lstinline|markers| (formato nativo do
Surfer v0.7.0), e \lstinline|show_cursor_window| abre a janela de cursor/delta para
medir o intervalo (latência). Os tempos são o \lstinline|#N| cru do dump (em
unidades do \emph{timescale}).

\subsection{Camada de curadoria: \lstinline|buildSurferLayout|}
\label{sec:cap10-surfer-curadoria}

A função \lstinline|buildSurferLayout| espelha o \lstinline|buildAuroraGtkw|:
mesma detecção de processador (\lstinline|detectProcessors|,
\lstinline|resolveScopeModules|, \lstinline|buildSignedSet|, importados \emph{verbatim}
de \file{gtkw\_proc\_writer}) e a mesma ordem de seções, mas emitindo
\lstinline|SurferItem| declarativos em vez de linhas \lstinline|@<hexflag>| do
GTKWave. A ordem por processador é: \textbf{Top-level} $\rightarrow$ por processador
(\emph{banner}/grupo $\rightarrow$ clk/rst/itr $\rightarrow$ I/O $\rightarrow$
Instructions $\rightarrow$ Variables $\rightarrow$ Flags).

\subsubsection{Grupos por processador e cores de seção}

Cada processador detectado vira um \textbf{grupo colapsável} (cabeçalho =
\lstinline|instanceName|), expandido por padrão (\lstinline|is_open: true|); em
\emph{designs} multi-processador o usuário pode dobrar os que não interessam. As
seções internas (\textbf{I/O}, \textbf{Instructions}, \textbf{Variables},
\textbf{Flags}) também são grupos. A cor de todos os \emph{headers} curados é
\lstinline|SECTION_COLOR = 'Red'| (\lstinline|mkGroup|), por contraste com as
trilhas (I/O em \lstinline|Yellow|, Variables em \lstinline|Orange|, Instructions em
\lstinline|Violet|). Blocos volumosos abrem fechados: \emph{arrays} e a seção
\textbf{Flags} têm \lstinline|isOpen = false|.

\subsubsection{Mapeamento de formato e cor (GTKWave $\rightarrow$ Surfer)}

A \autoref{tab:cap10-fmt} resume o mapeamento verificado no código.

\begin{table}[h]
\centering
\caption{Mapeamento de formato e cor do \path{.gtkw} para o \path{.surf.ron}.}
\label{tab:cap10-fmt}
\begin{tabularx}{\linewidth}{Y Y}
\toprule
\textbf{GTKWave} & \textbf{Surfer (item.format / color)} \\
\midrule
FMT\_BIN              & \lstinline|Binary| \\
FMT\_DEC             & \lstinline|Unsigned| \\
FMT\_SIGNED\_DEC      & \lstinline|Signed| \\
FMT\_REAL            & \lstinline|FP: 32-bit IEEE 754| \\
FMT\_ANALOG\_*        & formato base + \lstinline|analog| (Step) \\
Orange / Yellow      & \lstinline|Orange| / \lstinline|Yellow| \\
Indigo / Violet      & \lstinline|Violet| \\
NORMAL               & sem cor \\
\bottomrule
\end{tabularx}
\end{table}

\subsubsection{Seções de sinal}

\begin{description}
  \item[Top-level] sinais fora de qualquer processador; \lstinline|Binary| para
        escalares, \lstinline|Signed|/\lstinline|Unsigned| conforme o
        \lstinline|signedSet|.
  \item[clk/rst/itr] 1-bit, \lstinline|Binary|, com \lstinline|heightScale: 0.8|
        (levemente reduzidos para dar proeminência aos sinais de dado; o comentário
        registra que 0.5 era curto demais e o rótulo encavalava).
  \item[I/O] sinais \lstinline|req_in_sim|, \lstinline|in_sim|,
        \lstinline|out_en_sim|, \lstinline|out_sig| (regex), cor \lstinline|Yellow|,
        com \emph{aliases} legíveis (\enquote{req\_in 0}, \enquote{input 0},
        \enquote{out\_en 0}, \enquote{output 0}).
  \item[Instructions] as trilhas \lstinline|valr2| (Assembly) e \lstinline|linetabs|
        (\cmm{}), cor \lstinline|Violet|. São a \textbf{assinatura curada} do SAPHO e
        são \emph{sempre} emitidas quando existem — de propósito \textbf{não} passam
        pelo filtro do \emph{picker}. O \emph{alias} carrega o nome do processador
        (ex.: \enquote{Assembly (ProcDTW)}).
  \item[Variables] inteiros (\lstinline|me1_| $\rightarrow$ \lstinline|Signed|),
        floats (\lstinline|me2_| $\rightarrow$ \lstinline|FP: 32-bit IEEE 754|),
        complexos (\lstinline|comp_me3_|) e seus \emph{arrays}
        (\lstinline|arr_me1_| etc.), cor \lstinline|Orange|, com \emph{aliases} do
        tipo \enquote{int v in f()}. \emph{Arrays} viram subgrupos fechados.
  \item[Flags] grupos \textbf{Stack} (ponteiros/limites das pilhas de dados e de
        instrução, em \lstinline|sp|/\lstinline|isp|) e \textbf{ULA}
        (\lstinline|delta_int|/\lstinline|delta_float|, hexadecimal), com
        \emph{display} analógico \lstinline|ANALOG_STEP|
        (\lstinline|{ renderStyle: 'Step', yAxisScale: 'TypeLimits' }|).
\end{description}

\subsection{Tracks de Assembly/\cmm{}: \emph{mapping translators}}
\label{sec:cap10-surfer-mapping}

Onde o \path{.gtkw} usa \emph{file filters} (\path{trad\_opcode.txt} /
\path{trad\_cmm.txt}), o Surfer usa \emph{mapping translators}: arquivos de
mapeamento valor$\rightarrow$texto que o Surfer escaneia no \emph{startup} e
referencia pelo nome (campo \lstinline|format| do item). A AURORA converte os
\emph{trad files} do YANC nesses mappings.

\subsubsection{Conversão (\lstinline|convertTradToSurferMapping|)}

Um \emph{trad file} do YANC tem linhas \enquote{\lstinline|<chave-decimal> <texto>|}.
A conversão (regras validadas contra o binário v0.7.0):

\begin{itemize}
  \item emite um cabeçalho \lstinline|Name = <nome>| (e \lstinline|Bits = <n>|
        quando a largura é conhecida);
  \item o Surfer casa a chave pelo \emph{valor numérico} dos bits;
  \item chaves \textbf{negativas} (\emph{linetabs}: \lstinline|-1|/\lstinline|-2|/\lstinline|-3|)
        viram o padrão de bits \emph{unsigned} na largura do sinal (ex.: 20-bit
        \lstinline|-1| $\rightarrow$ \lstinline|0xFFFFF|), pois o Surfer não casa
        assinado;
  \item linhas \textbf{sem texto} são \textbf{puladas}: o Surfer rejeita
        (\enquote{Missing mapping}) e a falha derruba o arquivo inteiro, o que faria
        o Surfer entrar em \emph{panic} ao carregar um \path{.surf.ron} que
        referencia um \emph{translator} não carregado.
\end{itemize}

\subsubsection{Resolução e escrita}

\lstinline|resolveProcMappings| decide o \lstinline|format| das trilhas: com
\emph{trad file} presente, o \lstinline|format| é o nome do mapping (decode); sem
ele, o \emph{fallback} é decimal cru (\lstinline|Unsigned| para Assembly,
\lstinline|Signed| para \cmm{}). Os nomes de mapping são FS-safe e prefixados
(\lstinline|aurora_asm_<ns>_<procType>| etc.), com \emph{namespace} derivado do
caminho do projeto (FNV-1a de \lstinline|projectPath|), de modo que dois projetos
com o mesmo \emph{top} de testbench não se sobrescrevem.

A escrita dos mappings é feita pelo IPC \lstinline|write-surfer-mappings|
(\file{main/ipc/compile.js}), no diretório \textbf{global} do Surfer
(\path{\%APPDATA\%/surfer-project/surfer/config/mappings}) — o único lugar onde o
Surfer descobre os mappings de forma confiável no Windows (o \emph{walk-up} local
\path{.surfer/} está quebrado em v0.7.0). A escrita é \emph{atômica} (grava em
\path{.aurora.tmp} e renomeia) para o Surfer nunca ler um mapping pela metade; os
nomes são \lstinline|aurora_*|-prefixados para nunca tocar mappings do usuário, e
cada \emph{launch} sobrescreve o conjunto do projeto ativo (idempotente).

\subsection{Números complexos: pré-passo via \tool{comp2gtkw}}
\label{sec:cap10-surfer-prepass}

O Surfer não tem o \emph{process filter} externo do GTKWave, então os números
complexos (\lstinline|comp_me3_| / \lstinline|comp_arr_me3_|) são pré-decodificados
e \enquote{assados} em um mapping. O método
\lstinline|_buildSurferComplexMapping| (\file{compilation\_module.js}):

\begin{enumerate}[label=\arabic*.]
  \item faz \emph{stream} do \tool{fst2vcd} sobre o FST e coleta os valores
        \emph{distintos} dos sinais complexos (parando cedo ao atingir o \emph{cap});
  \item envia esses valores ao \tool{comp2gtkw} (o mesmo decoder do GTKWave) via o
        IPC \lstinline|decode-complex|, que os alimenta no \lstinline|stdin| (um
        token por linha) e devolve as strings \enquote{re imi} em ordem;
  \item monta um único mapping compartilhado por todos os sinais complexos (o
        decode depende só do \emph{bitpattern}).
\end{enumerate}

Tudo é \emph{gated}: projetos sem sinais complexos (\lstinline|hasComplexSignals|)
não pagam o \emph{stream}. Sem o decoder
(\path{components/bin/comp2gtkw.exe}) ou sem o \tool{fst2vcd}, o passo é pulado e os
complexos abrem em \lstinline|Binary| cru. O \lstinline|comp2gtkw.exe| só é aceito
pelo \emph{allowlist} a partir de \path{bin}.

\subsection{Marcadores de latência (eventos de I/O)}
\label{sec:cap10-surfer-markers}

\lstinline|_buildSurferEventMarkers| crava marcadores automáticos nos tempos do
primeiro \lstinline|req_in| (entrada) e do primeiro \lstinline|out_en| (saída) para
medir latência. Também é \emph{gated} (\lstinline|hasIoEventSignals|): só paga o
\emph{stream} do \tool{fst2vcd} quando há sinais de I/O, e para cedo ao achar os dois
eventos. Quando há marcadores, a janela de delta do Surfer abre
(\lstinline|showCursorWindow: markerTimes.length > 0|).

\subsection{Resolução do layout: usuário vs.\ auto-gerado}
\label{sec:cap10-surfer-resolve}

O \lstinline|_waveResolveSurferSaveFile| decide qual \emph{layout} carregar, em duas
fontes:

\begin{description}
  \item[Fonte 1 — curado pelo usuário] a entrada de \lstinline|surferFiles[]|
        marcada \lstinline|isActive| na \lstinline|WaveStore| (um \path{.surf.ron}
        ou \path{.sucl}). Se existir, é usada diretamente.
  \item[Fonte 2 — auto-gerado] quando não há entrada ativa, gera um
        \path{.surf.ron} curado via \lstinline|buildSurferLayout|, reusando a
        \emph{mesma} seleção do \emph{picker} e a mesma detecção de processador do
        \tool{GTKWave}. O arquivo vai para
        \path{<tempBaseDir>/<simTopModule>.surf.ron}.
\end{description}

Quando nenhuma fonte produz um \emph{layout}, retorna \lstinline|null| e o Surfer
abre o VCD cru. O \emph{parse} do cabeçalho prefere o irmão \path{.header.vcd}
(o FST binário não é \emph{parseável} como texto); se só houver FST e nenhum
\path{.header.vcd}, retorna \lstinline|null|.

\paragraph{Anti-\emph{staleness}.} Antes de assar os mappings, o método compara os
\emph{mtimes}: se os tradutores são \emph{mais novos} que o dump (recompilou o
\cmm{} sem re-simular), avisa no terminal que o decode pode estar desatualizado
(uma tabela nova casada com um dump velho geraria \enquote{lixo crível}), com margem
de 2\,s para a ordem normal compile$\rightarrow$simulate.

\subsection{Limites atuais da camada de curadoria}
\label{sec:cap10-surfer-limites}

O cabeçalho do próprio \file{surfer\_layout\_writer.ts} registra a paridade: cores,
formatos, \emph{analog}, \emph{aliases}, seções, \emph{arrays} e \emph{flags} têm
paridade total com o \path{.gtkw}. As trilhas de Assembly/\cmm{} e os complexos
dependem do mecanismo de \emph{mapping translator} / pré-passo descrito acima; quando
os \emph{trad files} ou o decoder não estão disponíveis, esses tracks abrem em
decimal/binário cru (nunca é fatal). A \autoref{tab:cap10-surfer-resumo} consolida
onde cada artefato é escrito.

\begin{table}[h]
\centering
\caption{Artefatos do caminho Surfer e onde são gravados.}
\label{tab:cap10-surfer-resumo}
\begin{tabularx}{\linewidth}{Y Y}
\toprule
\textbf{Artefato} & \textbf{Local} \\
\midrule
\path{.surf.ron} auto-gerado & \path{components/Temp/<top>.surf.ron} \\
\emph{Mapping translators}   & \path{\%APPDATA\%/surfer-project/surfer/config/mappings/} \\
Config de janela (centralizada) & \path{\%APPDATA\%/surfer-project/surfer/config/config.toml} \\
Binário do Surfer            & \path{components/Packages/surfer/surfer.exe} \\
\bottomrule
\end{tabularx}
\end{table}

\section{Resumo}
\label{sec:cap10-resumo}

A AURORA entrega dois visualizadores de onda sobre o mesmo \emph{pipeline} de
simulação. O \tool{GTKWave} (fork \repo{gtkwave-nipscern}) é o padrão empacotado,
com customizações verificadas no código: \lstinline|--zoom-fit| e
\lstinline|--left-justify| novas, painel/canvas escuros alinhados à paleta Surfer
(agora \emph{default}), SST removido, ícones Phosphor, barra de título CSD, \emph{zoom}
animado e build Windows nativo via Meson/MSYS2. O \tool{Surfer} é o viewer
\emph{opt-in}: baixado opcionalmente, lançado como janela externa, e alimentado por
um \emph{layout} declarativo \path{.surf.ron} gerado por \lstinline|buildSurferLayout|
— grupos colapsáveis por processador, cores, \emph{aliases}, \emph{display} analógico,
tracks de Assembly/\cmm{} via \emph{mapping translators} (conversão de
\path{trad\_*.txt}) e complexos via pré-passo \tool{comp2gtkw}. Quando o
\path{surfer.exe} está ausente, o botão Wave degrada limpa e automaticamente para o
\tool{GTKWave}.

% !TEX root = ../main.tex
\chapter{AURORA: o visualizador de RTL PRISM}
\label{cap:11-aurora-prism-rtl}

O \tool{PRISM} é o visualizador de RTL (\emph{register-transfer level}) embarcado
na \tool{AURORA}. A partir do Verilog sintetizável de um projeto SAPHO, ele
produz um \emph{esquemático} navegável do circuito: cada módulo da hierarquia
vira um diagrama vetorial (SVG) com células, portas e fios, no qual o usuário
pode entrar em submódulos com um clique, voltar ao código-fonte com duplo-clique
e --- desde a fase O9 --- alternar para uma \emph{simulação lógica interativa}
do mesmo \emph{netlist} via \tool{DigitalJS}.

Este capítulo documenta o subsistema completo: a janela e o canal de IPC
(\file{main/ipc/prism.js}), o \emph{preload} sandboxizado
(\file{js/app/preload\_prism.js}), o renderer da janela
(\file{html/prism/prism.js} e \file{html/prism/prism.html}), o \emph{pipeline}
Verilog~$\rightarrow$~Yosys~$\rightarrow$~netlistsvg-aurora~$\rightarrow$~SVG, o
fork \repo{netlistsvg-aurora} (o que ele muda em relação ao
\lstinline|@silimate/netlistsvg| de \emph{upstream}), o sistema de \emph{skins}
customizadas da família SAPHO, e o estado de integração de
\tool{DigitalJS}~+~\tool{yosys2digitaljs}. A síntese pela ferramenta \tool{Yosys}
e o \emph{toolchain} bundle são tratados em outros capítulos; aqui o foco é o que
o \tool{PRISM} faz com a saída do \tool{Yosys}.

\begin{nota}
Fonte da verdade: todas as afirmações deste capítulo foram verificadas lendo o
código real. Onde a documentação de prosa interna e o código divergem, vale o
código. Os arquivos efetivamente abertos estão listados ao final do estudo.
\end{nota}

\section{Visão geral e fluxo de acionamento}

O \tool{PRISM} é acionado pelo botão \enquote{PRISM} da \emph{toolbar} de
compilação da \tool{AURORA} (id \lstinline|prismcomp|). O \emph{handler} desse
botão vive em \file{js/compilation/compilation\_flow.js}, na função
\lstinline|handlePrismStep()|. Conforme a filosofia \enquote{cada botão é
self-contained} desse arquivo, o \tool{PRISM} é um \emph{superset} do botão
Verilog: ele primeiro roda \cmm{}~+~ASM para cada processador conhecido e a
verificação de sintaxe Verilog (top-level), e só então invoca o \tool{Yosys}
para a análise estrutural.

\begin{lstlisting}[language=bash,caption={Expansão dos botões (cabeçalho de compilation\_flow.js)}]
Verilog = cmm + asm + iverilog(-tnull, top)
PRISM   = Verilog                            + yosys
Wave    =     cmm + asm + iverilog(-o vvp, tb) + vvp + gtkwave
\end{lstlisting}

O botão \lstinline|prismcomp| só fica habilitado quando o \file{.spf} tem um
\emph{top-level} definido (\lstinline|setEnabled('prismcomp', hasTop)| em
\file{compilation\_flow.js}). O fluxo de alto nível é:

\begin{enumerate}[leftmargin=2.2em]
  \item \textbf{Renderer (Aurora principal):} \lstinline|handlePrismStep()| roda
        \cmm{}/ASM e \lstinline|verilogSyntaxCheck()|, monta o \emph{payload} de
        caminhos absolutos (\lstinline|buildPrismCompilationPaths|), consulta o
        \emph{store} de \emph{overrides} de IA para a etapa
        \lstinline|prism-yosys| e chama \lstinline|electronAPI.prismCompileWithPaths(paths)|.
  \item \textbf{Main:} o \emph{handler} IPC \lstinline|prism-compile-with-paths|
        executa o \emph{pipeline} de síntese e geração de SVG, e abre (ou foca) a
        janela do \tool{PRISM}.
  \item \textbf{Renderer (janela PRISM):} recebe \lstinline|compilation-complete|
        com o caminho do SVG do top-level, carrega o SVG, instrumenta cliques e
        navegação.
\end{enumerate}

O \emph{payload} de caminhos é montado em
\lstinline|buildPrismCompilationPaths(projectPath)| com todos os campos
absolutos e \emph{slashes} normalizadas para barra direta:

\begin{lstlisting}[caption={Campos do payload de caminhos (compilation\_flow.js)}]
{
  projectPath, componentsPath,
  hdlPath:      <components>/HDL,
  tempPath:     <components>/Temp/PRISM,
  yosysPath:    <components>/Packages/msys/mingw64/bin/yosys.exe,
  spfPath:      <currentSpfPath>,
  topLevelPath: <projectPath>/TopLevel,
  yosysOverride?: { appendArgs, prependArgs, removeArgs, envSet, envUnset }
}
\end{lstlisting}

O \emph{main} também expõe \lstinline|get-prism-compilation-paths|, que recompõe
esses mesmos campos a partir do \file{.spf} aberto
(\lstinline|state.currentOpenProjectPath|); a janela usa esse canal para
recompilar e para construir a simulação \tool{DigitalJS} por conta própria.

\section{A janela do PRISM}

A janela é criada por \lstinline|createPrismWindow(compilationData)| em
\file{main/ipc/prism.js}. É uma \lstinline|BrowserWindow| única
(\lstinline|state.prismWindow|): se já existe e não foi destruída, a chamada
apenas a foca e reenvia o payload de compilação, em vez de abrir uma segunda
janela.

\subsection{Configuração da BrowserWindow}

\begin{tabularx}{\linewidth}{Y Z}
\toprule
\textbf{Propriedade} & \textbf{Valor / efeito} \\
\midrule
\lstinline|width| / \lstinline|height| & 1400 $\times$ 900 (mín. 1000 $\times$ 700) \\
\lstinline|frame| & \lstinline|false| --- janela \emph{frameless} com \emph{titlebar} customizada desenhada pelo HTML do PRISM \\
\lstinline|thickFrame| & \lstinline|true| --- preserva \emph{snap} do Aero e redimensionamento por borda nativo no Windows \\
\lstinline|titleBarStyle| & \lstinline|'hidden'| \\
\lstinline|backgroundColor| & \lstinline|'#0A0D14'| (mesmo fundo escuro da \tool{AURORA}) \\
\lstinline|icon| & \path{assets/icons/sapho_aurora_icon.ico} \\
\lstinline|contextIsolation| & \lstinline|true| \\
\lstinline|sandbox| & \lstinline|true| \\
\lstinline|nodeIntegration| & \lstinline|false| \\
\lstinline|preload| & \path{js/app/preload_prism.js} (existência verificada antes de criar a janela) \\
\bottomrule
\end{tabularx}

A página em si é carregada por \lstinline|loadPage(prismWindow, 'html/prism/prism.html')|,
de \file{main/render\_loader.js}. Esse \emph{loader} é compartilhado por todas as
janelas de primeira parte e resolve a fonte do renderer em duas vias: servidor de
desenvolvimento Vite (quando \lstinline|AURORA_RENDERER_URL| está setado e o app
não está empacotado) ou o \emph{bundle} de produção em \path{dist/} via
\lstinline|file://|. \emph{Não há \emph{fallback}} para o fonte cru --- um
\emph{bundle} ausente é tratado como erro de \emph{build} (página de erro
explícita), porque pós-migração para \tool{Lit} o HTML cru carrega \emph{imports}
que só resolvem através do empacotador.

\subsection{Controles de janela e estado}

Como a janela é \emph{frameless}, os botões minimizar/maximizar/fechar são
desenhados em SVG no HTML e ligados via IPC genéricos do \emph{preload}
(\lstinline|window:minimize|, \lstinline|window:maximize-toggle|,
\lstinline|window:close|). O \emph{main} relê o estado de maximização/tela-cheia
para o renderer pelo canal \lstinline|prism:window-state| nos eventos
\lstinline|maximize|, \lstinline|unmaximize|, \lstinline|enter-full-screen| e
\lstinline|leave-full-screen|, para que o ícone \emph{maximize/restore} se
atualize. Ao carregar, a janela é maximizada e exibida; o \emph{main} também
notifica a janela principal com \lstinline|prism-status| (\lstinline|true| ao
abrir, \lstinline|false| ao fechar).

Erros de carregamento são tratados em três pontos: o \lstinline|try/catch| em
torno de \lstinline|loadPage| (mostra \lstinline|dialog.showMessageBox| e destrói
a janela), o \lstinline|did-fail-load| e o \lstinline|render-process-gone|
(ambos registram em \lstinline|electron-log| e exibem diálogo de erro).

\section{O preload sandboxizado}

\file{js/app/preload\_prism.js} expõe um \emph{único} objeto
\lstinline|window.electronAPI| via \lstinline|contextBridge|, com uma
\emph{allowlist enumerada} de canais. O comentário no próprio arquivo é
explícito: \emph{deliberadamente não há} um \lstinline|send|/\lstinline|removeAllListeners|
genérico, para que esse renderer nunca alcance um canal IPC arbitrário.

\begin{tabularx}{\linewidth}{Y Z}
\toprule
\textbf{Método exposto} & \textbf{Canal IPC / efeito} \\
\midrule
\lstinline|windowMinimize/MaximizeToggle/Close| & \lstinline|window:minimize| / \lstinline|maximize-toggle| / \lstinline|close| \\
\lstinline|onWindowState(cb)| & escuta \lstinline|prism:window-state| \\
\lstinline|readFile(path)| & \lstinline|read-file| (passa por \lstinline|safePath()| no main; evita \lstinline|fetch('file://')|) \\
\lstinline|joinPath(...segs)| & \lstinline|join-path| \\
\lstinline|getPrismCompilationPaths()| & \lstinline|get-prism-compilation-paths| \\
\lstinline|prismCompileWithPaths(paths)| & \lstinline|prism-compile-with-paths| \\
\lstinline|prismRecompile(paths)| & \lstinline|prism-recompile| \\
\lstinline|buildDigitalJS(paths, module)| & \lstinline|prism:build-digitaljs| (O9) \\
\lstinline|exportWave(payload)| & \lstinline|prism:export-wave| (grava o \lstinline|.vcd| do monitor no Temp do PRISM e manda a janela principal abri-lo, via \lstinline|aurora:open-wave|) \\
\lstinline|generateSVGFromModule(name, dir)| & \lstinline|generate-svg-from-module| \\
\lstinline|openSourceFile({filePath,line,column})| & \lstinline|prism:open-source-file| \\
\lstinline|onCompilationComplete(cb)| & escuta \lstinline|compilation-complete| \\
\lstinline|onTerminalLog| / \lstinline|onPrismStatus| & escuta \lstinline|terminal-log| / \lstinline|prism-status| \\
\bottomrule
\end{tabularx}

Além disso, o \emph{preload} reencaminha eventos \lstinline|terminal-log| do
\emph{main} para a página como \lstinline|window.postMessage|, de modo que a
janela do \tool{PRISM} possa exibir o progresso da síntese.

\section{O pipeline de compilação}

O coração do \tool{PRISM} é a função assíncrona
\lstinline|performPrismCompilationWithPaths(compilationPaths)|, que orquestra
três etapas: síntese com \tool{Yosys}, divisão do \emph{netlist} por módulo, e
geração do SVG do top-level. O top-level vem \emph{sempre} de
\lstinline|spf.structure.topLevelFile| (\lstinline|path.basename(...,'.v')|).

\begin{lstlisting}[caption={Diagrama de fluxo do pipeline (main/ipc/prism.js)}]
 +-------------------------+
 |  spf.structure.topLevel |  ->  topLevelModule (basename sem .v)
 +-----------+-------------+
             |
             v
 +-------------------------+    coleta unificada de .v + dedup por path
 |  runYosysCompilation    |    -> yosys_script.ys -> yosys.exe -> hierarchy.json
 +-----------+-------------+
             |
             v
 +-------------------------+    para cada modulo "clicavel":
 |  splitHierarchyJson     |    limpa nome, escreve <modulo>.json no tempDir
 +-----------+-------------+
             |
             v
 +-------------------------+    skin default + skins custom -> netlistsvg.render
 |  generateModuleSVG      |    -> injeta data-cell-type -> <modulo>.svg
 +-----------+-------------+
             |
             v
 +-------------------------+
 |  janela PRISM carrega   |  ->  compilation-complete { topLevelModule, svgPath, tempDir }
 +-------------------------+
\end{lstlisting}

\subsection{Etapa 1 --- coleta de arquivos e síntese com Yosys}

A função \lstinline|runYosysCompilationWithPaths(paths, topLevelModule, tempDir)|
coleta \emph{todos} os \file{.v} relevantes em um \lstinline|Set| (dedup por
caminho absoluto, evitando erros de módulo duplicado no \tool{Yosys}). As fontes,
nesta ordem, são:

\begin{enumerate}[leftmargin=2.2em]
  \item \path{components/HDL/*.v} --- a biblioteca SAPHO (processor, addr\_dec,
        core, ula, myFIFO, instr\_dec etc.), sempre incluída porque qualquer
        \emph{top} que seja um processador SAPHO depende dela. Arquivos com
        \lstinline|_tb| ou \lstinline|test| no nome são excluídos.
  \item \lstinline|spf.structure.synthesizableFiles[].path| --- a fonte canônica
        de todos os \file{.v} do projeto (auto-descobertos pela árvore de
        arquivos e/ou adicionados manualmente). Caminhos relativos são resolvidos
        contra \lstinline|spf.structure.basePath| (formato novo do \file{.spf};
        caminhos absolutos passam direto).
  \item \path{<projeto>/TopLevel/*.v}, se a pasta existir (\emph{wrapper}
        top-level, raro mas possível).
  \item \textbf{\emph{Fallback} defensivo:} para cada processador listado em
        \lstinline|spf.structure.processors|, varre
        \path{<projeto>/<proc>/Hardware/*.v}. No caminho feliz o \emph{dedup}
        torna isso \emph{no-op}, mas garante os módulos do processador em
        projetos antigos onde \lstinline|synthesizableFiles| esteja vazio.
\end{enumerate}

Se nenhum \file{.v} for encontrado, a etapa falha com \enquote{No Verilog files
found for compilation}.

O \emph{script} \tool{Yosys} é gerado em \path{<tempDir>/yosys_script.ys}. Cada
arquivo é lido com \lstinline|read_verilog -setattr src| --- a opção
\lstinline|-setattr src| faz o \tool{Yosys} gravar um atributo
\lstinline|src="arquivo.v:linha.col-linha.col"| em cada célula derivada do
\emph{source}, o que mais tarde habilita o duplo-clique \enquote{abrir no
editor}. O \emph{script} é exatamente:

\begin{lstlisting}[caption={Script Yosys do esquemático (gerado em main/ipc/prism.js)}]
read_verilog -setattr src "arquivo1.v"
read_verilog -setattr src "arquivo2.v"
...
hierarchy -top <topLevelModule>
proc
setundef -zero
opt_clean -purge
write_json "hierarchy.json"
\end{lstlisting}

\begin{tabularx}{\linewidth}{Y Z}
\toprule
\textbf{Comando} & \textbf{Função no PRISM} \\
\midrule
\lstinline|read_verilog -setattr src| & lê cada \file{.v} e anota a procedência (\lstinline|src|) em cada célula \\
\lstinline|hierarchy -top| & fixa o módulo top e resolve a árvore de instâncias/parametrizações \\
\lstinline|proc| & converte processos (\lstinline|always|) em células (mux, dff etc.) \\
\lstinline|setundef -zero| & troca \emph{don't-care} (\lstinline|x|) por constante 0, evitando \enquote{linhas fantasma} que o netlistsvg desenharia a partir de ramos \emph{full\_case} inalcançáveis \\
\lstinline|opt_clean -purge| & remove fios/células desconectados após a substituição \\
\lstinline|write_json| & emite o \emph{netlist} consolidado em \path{hierarchy.json} \\
\bottomrule
\end{tabularx}

Note que o \emph{script do esquemático não faz} \lstinline|techmap|/\lstinline|abc|:
o objetivo é um diagrama em nível de palavra (RTL), não \emph{gates}.

O \tool{Yosys} é executado com \lstinline|spawnTracked(yosysExe, finalArgs, ...)|.
O uso de \lstinline|spawnTracked| (de \file{main/process\_registry.js}) garante
que fechar a janela principal mata uma síntese do \tool{PRISM} em andamento ---
necessário porque o \tool{Yosys} roda do \path{mingw64/bin} bundle, fora de
\path{Temp/}. Os argumentos base são \lstinline|['-s', yosysScriptPath]|, mas
podem ser modificados por um \emph{override} de IA
(\lstinline|paths.yosysOverride|) com o mesmo contrato dos demais \emph{steps}:
\lstinline|removeArgs|, \lstinline|prependArgs|, \lstinline|appendArgs|,
\lstinline|envSet|, \lstinline|envUnset|. Sem \emph{override}, o comportamento é
idêntico ao padrão. O processo só resolve em sucesso se o código de saída for 0
\emph{e} \path{hierarchy.json} existir; caso contrário rejeita com o
\lstinline|stderr| capturado.

\subsection{Etapa 2 --- divisão do netlist por módulo}

\lstinline|splitHierarchyJson(hierarchyJsonPath, tempDir)| lê o
\path{hierarchy.json} consolidado e, para cada módulo \enquote{clicável}, escreve
um JSON \emph{por módulo} (\path{<sanitized>.json}) contendo apenas aquele
módulo. Isso é o que permite, em tempo de navegação, renderizar qualquer módulo
sob demanda sem reexecutar o \tool{Yosys}.

Dois \emph{helpers} regem o que entra:

\begin{description}[leftmargin=1.4em,style=nextline]
  \item[\lstinline|isClickableModule(name)|] descarta primitivas do \tool{Yosys}
        (qualquer \lstinline|$_|, \lstinline|$dff|, \lstinline|$mux|,
        \lstinline|$add|, \lstinline|$pmux|, \lstinline|$lut|, etc.) e aceita
        apenas \lstinline|$paramod...| ou identificadores Verilog legítimos
        (\lstinline|^[a-zA-Z_][a-zA-Z0-9_]*$|). É o mesmo conjunto de padrões
        replicado no renderer (\lstinline|_isClickableModuleType|), para que o
        que o \emph{main} grava como navegável coincida com o que o renderer
        torna clicável.
  \item[\lstinline|cleanModuleName(name)|] normaliza nomes gerados pelo
        \tool{Yosys}: desembrulha prefixos \lstinline|$paramod\...| de módulos
        parametrizados, remove \emph{hashes} (\lstinline|$<hex de 40+ dígitos>|),
        sufixos de parâmetro (\lstinline|\PARAM=...|), e prefixos de escopo de
        \emph{generate block} (\lstinline|genblk<N>.|), de modo que o usuário veja
        o nome real (ex.: \lstinline|my_f2i| em vez de
        \lstinline|genblk27.my_f2i|).
\end{description}

Para cada módulo aceito, o nome é limpo e \emph{sanitizado}
(\lstinline|sanitizeFileName|), as células internas também têm seus nomes e tipos
limpos, e o JSON resultante é escrito com a forma
\lstinline|{ creator, modules: { <cleanName>: <data> } }|.

\subsection{Etapa 3 --- renderização do SVG com netlistsvg-aurora}

\lstinline|generateModuleSVGWithPaths(moduleName, tempDir)| produz o SVG de um
módulo. Pontos-chave:

\begin{itemize}[leftmargin=1.4em]
  \item A biblioteca é carregada \emph{in-process} via
        \lstinline|require('@silimate/netlistsvg')| --- \emph{não há} \emph{spawn}
        de processo externo nem binário separado para o netlistsvg. (O comentário
        em \lstinline|get-prism-compilation-paths| confirma: \enquote{netlistsvg
        agora vem do node\_modules e roda in-process}.)
  \item A \emph{skin} é obtida por \lstinline|getDefaultSkinData()| (ver
        \autoref{sec:skins}), que combina a \path{lib/default.svg} do pacote com
        as \emph{skins} customizadas de \path{assets/prism-skins/}.
  \item Antes de renderizar, \lstinline|pruneNetlistToSkinPorts| poda do
        \emph{netlist} conexões a portas que a \emph{skin} customizada do tipo
        \emph{não} desenha (ver \autoref{sec:skin-ports}), evitando que o ELK
        aborte o \emph{layout}.
  \item A renderização usa \lstinline|netlistsvgLib.render(skinData, netlistJson, cb)|
        (estilo \emph{callback}), envolto em \lstinline|Promise|.
  \item Se a renderização não devolver uma \emph{string} não-vazia, a função
        falha com mensagem clara (em vez de estourar \lstinline|TypeError| na
        injeção seguinte).
  \item Por fim, \lstinline|injectCellTypesIntoSvg| insere
        \lstinline|data-cell-type=<tipo>| em cada \lstinline|<g id="cell_<inst>">|,
        e o SVG é gravado em \path{<sanitized>.svg}.
\end{itemize}

A injeção de \lstinline|data-cell-type| é necessária porque o netlistsvg rotula
cada célula com o \emph{nome da instância} (\lstinline|s:attribute="ref"|), não
com o \emph{tipo} de módulo. \lstinline|buildInstanceTypeMap| constrói o mapa
\lstinline|instância -> tipo| a partir do JSON do \tool{Yosys}, e
\lstinline|injectCellTypesIntoSvg| usa uma \emph{regex} (não um \emph{parser} XML)
para anotar cada \lstinline|<g>| de célula com o tipo --- assim o renderer sabe
para qual módulo navegar ao clicar. Valores de atributo são escapados por
\lstinline|xmlAttrEscape| (\lstinline|&|, \lstinline|"|, \lstinline|<|).

\subsection{Resultado e abertura da janela}

Em sucesso, \lstinline|performPrismCompilationWithPaths| retorna
\lstinline|{ success: true, topLevelModule, svgPath, tempDir }|, e o \emph{handler}
\lstinline|prism-compile-with-paths| chama \lstinline|createPrismWindow(result)|.
A janela recebe \lstinline|compilation-complete| com esse objeto, carrega o SVG
do top-level e inicializa a navegação. Durante todo o processo, mensagens de
progresso são enviadas à janela principal pelo canal \lstinline|terminal-log| no
terminal \lstinline|tveri| (\enquote{Top-level: \dots}, \enquote{Running Yosys
synthesis\dots}, \enquote{PRISM compilation completed successfully}, ou os erros
correspondentes).

\section{O fork netlistsvg-aurora}
\label{sec:fork}

O \file{package.json} da \tool{AURORA} aponta a dependência
\lstinline|@silimate/netlistsvg| para o fork do laboratório:

\begin{lstlisting}[caption={Dependência apontando para o fork (package.json da AURORA)}]
"@silimate/netlistsvg":
  "github:nipscernlab/netlistsvg-aurora#58a0703904a35841616db58716f40532b84fb018"
\end{lstlisting}

Ou seja, embora o nome do pacote permaneça \lstinline|@silimate/netlistsvg|
(versão \lstinline|1.1.4|, autor original Neil Turley), o código instalado vem de
\repo{netlistsvg-aurora}, fixado por \emph{commit hash}.

\subsection{O que o fork muda em relação ao upstream}

A diferença substantiva do fork é um \textbf{ajuste de \emph{tsconfig} para
destravar a compilação sob TypeScript~6}, seguida de uma reconstrução de
\path{built/}. O \file{tsconfig.json} do fork é:

\begin{lstlisting}[caption={tsconfig.json do fork netlistsvg-aurora}]
{
  "compilerOptions": {
    "outDir": "./built",
    "allowJs": true,
    "lib": [ "es6", "dom" ],
    "target": "es2015",
    "module": "commonjs",
    "moduleResolution": "node",
    "types": []
  },
  "include": [ "./lib/*", "./lib/@types" ]
}
\end{lstlisting}

A mensagem de \emph{commit} do fork descreve precisamente as mudanças:

\begin{itemize}[leftmargin=1.4em]
  \item \lstinline|target| passou de \lstinline|"es5"| (que o TS~6 marca como
        \emph{deprecated} e recusa compilar, erro \lstinline|TS5107|) para
        \lstinline|"es2015"|.
  \item Foram adicionados \lstinline|module: "commonjs"| e
        \lstinline|moduleResolution: "node"| explícitos, para que a sintaxe
        CommonJS existente (\lstinline|import x = require(...)|) continue
        funcionando sob o compilador moderno.
  \item \lstinline|types: []| foi fixado para que o \lstinline|tsc| não varra
        pacotes \lstinline|@types| transitivos e tropece em definições de
        \emph{babel}/\emph{jest}.
  \item O diretório \path{built/} foi reconstruído a partir dessas
        configurações; a mensagem afirma que a saída é funcionalmente idêntica
        (testada renderizando os exemplos digitais).
\end{itemize}

O \emph{layout} de pastas do fork (\path{lib/} com os \file{.ts} fonte ---
\file{index.ts}, \file{FlatModule.ts}, \file{Cell.ts}, \file{Port.ts},
\file{Skin.ts}, \file{drawModule.ts}, \file{elkGraph.ts}, \file{YosysModel.ts}
--- além de \file{lib/default.svg} e \file{lib/analog.svg}, e \path{built/} com os
\file{.js} compilados e \file{netlistsvg.bundle.js}) é o mesmo do \emph{upstream}.
A API usada pela \tool{AURORA} é a função
\lstinline|render(skinData, yosysNetlist, done)| de \file{lib/index.ts}, que
constrói o \lstinline|FlatModule|, monta o grafo ELK
(\lstinline|buildElkGraph|), roda o \emph{layout} via \lstinline|elkjs| e desenha
com \lstinline|drawModule| --- preservando o estilo de \emph{callback} legado que
o código do \tool{PRISM} consome.

\begin{nota}
As demais diferenças efetivas em relação ao netlistsvg de prateleira estão
\emph{fora} do fork da biblioteca: a \tool{AURORA} não edita
\path{lib/default.svg}; em vez disso \emph{mescla} suas próprias \emph{skins} em
tempo de execução (ver \autoref{sec:skins}). As \emph{skins} customizadas da
família SAPHO vivem no repositório \tool{AURORA}, em \path{assets/prism-skins/}.
\end{nota}

\section{Sistema de skins}
\label{sec:skins}

Uma \emph{skin} do netlistsvg é um SVG que declara, em blocos
\lstinline|<g s:type="...">|, como desenhar cada \emph{tipo} de célula (suas
portas, identificadas por \lstinline|s:pid|, e a geometria do símbolo). A
\path{lib/default.svg} do pacote traz os tipos genéricos: portas externas
(\lstinline|inputExt|, \lstinline|outputExt|), \emph{flip-flops} e \emph{latches}
(\lstinline|dff|, \lstinline|dffn|, \lstinline|dlatch| e variantes \lstinline|-bus|),
operadores aritméticos e lógicos (\lstinline|add|, \lstinline|sub|, \lstinline|mul|,
\lstinline|div|, \lstinline|mod|, \lstinline|and|, \lstinline|or|, \lstinline|not|,
\lstinline|mux|, \lstinline|eq|, \lstinline|lt|, \lstinline|shl|, \lstinline|shr|
etc.), \lstinline|split|/\lstinline|join|, \lstinline|constant| e o fallback
\lstinline|generic|.

\subsection{Mescla de skins customizadas em runtime}

\lstinline|getDefaultSkinData()| (em \file{main/ipc/prism.js}) não usa a
\path{default.svg} crua: ela a combina com as \emph{skins} de
\path{assets/prism-skins/}. O processo, executado a \emph{cada} compilação (sem
\emph{cache}, para que \enquote{Recompile} já reflita edições em \emph{skins}
custom):

\begin{enumerate}[leftmargin=2.2em]
  \item Lê a \path{lib/default.svg} resolvida via
        \lstinline|require.resolve('@silimate/netlistsvg')|.
  \item \lstinline|loadCustomSkinBlocks(customSkinDir)| varre
        \path{assets/prism-skins/*.svg} (ignorando arquivos começados por
        \lstinline|_|, como \file{\_template.svg}), e de cada um extrai os blocos
        \lstinline|<g s:type="...">| de topo via
        \lstinline|extractTopLevelGBlocks|. Esse extrator usa contagem de
        profundidade de \lstinline|<g>| (pilha) para suportar \lstinline|<g>|
        aninhados dentro de portas/labels --- uma \emph{regex} \emph{non-greedy}
        quebraria em blocos como o \lstinline|generic|.
  \item Para cada \lstinline|s:type| customizado, o bloco correspondente é
        \emph{removido} da \emph{skin} default e o customizado é injetado logo
        antes de \lstinline|</svg>|.
\end{enumerate}

O comentário no código explica a escolha de mesclar em tempo de execução em vez
de editar \path{node_modules/.../default.svg}: um \lstinline|npm install|
sobrescreveria a \path{default.svg}; mantendo as customizações em
\path{assets/prism-skins/} elas sobrevivem ao \emph{install} e ficam versionadas
no git.

\subsection{A família de símbolos SAPHO}

As \emph{skins} customizadas são geradas por
\file{scripts/prism-skin-standard.js}, que estabelece um \emph{padrão} de
símbolos: cada classe de módulo é desenhada como seu símbolo de livro-texto (mux
= pentágono apontando à direita; ALU = \emph{chevron} entalhado; decodificador =
trapézio \emph{fan-out}; registrador = cartão com \emph{clock} $\blacktriangleright$;
memória = grade; FIFO = \emph{chevrons} de fluxo; blocos hierárquicos/de controle
= cartão de \emph{chip} arredondado). Cada símbolo usa fundo escuro de \emph{chip},
fluxo da esquerda para a direita (dados a oeste, resultados a leste; sinais de
controle como \lstinline|clk|/\lstinline|rst|/\lstinline|op|/\lstinline|en| em
violeta), portas rotuladas no \emph{anchor} e cores temáticas via variáveis CSS
com \emph{fallbacks} em hexadecimal.

Há \emph{skins} para a biblioteca SAPHO inteira --- por exemplo
\file{processor.svg}, \file{core.svg}, \file{ula.svg}, \file{addr\_dec.svg},
\file{instr\_dec.svg}, \file{instr\_fetch.svg}, \file{mem\_ctrl.svg},
\file{mem\_data.svg}, \file{mem\_instr.svg}, \file{myFIFO.svg}, \file{pc.svg},
\file{stack.svg}, \file{prefetch.svg} --- além de uma família grande de variantes
de ALU (\file{ula\_add.svg}, \file{ula\_fadd.svg}, \file{ula\_fmlt.svg},
\file{ula\_i2f.svg}, \file{ula\_mux.svg} etc.). Cada arquivo declara seu
\lstinline|s:type| (o tipo de célula que casa com o nome do módulo Verilog) e um
\lstinline|<s:alias val="..."/>|; uma regra dura do padrão é que cada
\lstinline|s:pid| \emph{deve} ser igual ao nome da porta Verilog.

\begin{lstlisting}[language=Java,caption={Cabeçalho de uma skin SAPHO (assets/prism-skins/processor.svg)}]
<svg xmlns="..." xmlns:s="https://github.com/nturley/netlistsvg">
  <g s:type="processor" transform="translate(0, 0)"
     s:width="150" s:height="132">
    <s:alias val="processor"/>
    ...
    <g s:x="0"   s:y="66" s:pid="clk"/>
    <g s:x="0"   s:y="75" s:pid="rst"/>
    <g s:x="150" s:y="66" s:pid="io_out"/>
    ...
  </g>
</svg>
\end{lstlisting}

\subsection{Poda de portas ausentes na skin}
\label{sec:skin-ports}

Uma \emph{skin} customizada tem um conjunto \emph{fixo} de portas. Se o RTL tem
uma porta que a \emph{skin} não declara (ex.: o \lstinline|core| ganhou um sinal
novo mas \file{core.svg} ainda não tem o \emph{anchor}), o netlistsvg geraria uma
\emph{aresta} para uma \emph{shape} de porta que nunca desenhou, e o ELK abortaria
com \enquote{Referenced shape does not exist}, derrubando o \emph{render} inteiro
do módulo.

Para evitar isso, \lstinline|generateModuleSVGWithPaths| poda o \emph{netlist}
antes de renderizar:

\begin{itemize}[leftmargin=1.4em]
  \item \lstinline|buildSkinPortMap(skinData)| monta
        \lstinline|tipo -> Set(s:pid)| a partir dos blocos da \emph{skin}.
  \item \lstinline|pruneNetlistToSkinPorts(netlistJson, skinPortMap)| remove, de
        cada célula, as conexões cujo \emph{port} não está na \emph{skin} de seu
        tipo (e o respectivo \lstinline|port_directions|), registrando um aviso
        com \lstinline|log.warn| indicando que basta adicionar um
        \lstinline|<g s:pid="..."/>| à \emph{skin} para desenhá-la.
  \item Células de tipo \emph{genérico} (sem \emph{skin}) \emph{não} são tocadas:
        ali o netlistsvg cria as portas a partir das próprias conexões.
\end{itemize}

O resultado é que uma porta extra apenas não aparece (com aviso), em vez de
quebrar a tela.

\section{O renderer da janela}

\file{html/prism/prism.js} (carregado como \lstinline|<script type="module">|)
define a classe \lstinline|PRISMViewer|, instanciada em
\lstinline|DOMContentLoaded| e exposta como \lstinline|window.prismViewer|. Ela
cobre i18n, carregamento e interação do SVG, navegação, pan/zoom e a simulação
\tool{DigitalJS}.

\subsection{Internacionalização}

O renderer mantém um dicionário próprio \lstinline|PRISM_STRINGS| com locais
\lstinline|en| e \lstinline|pt| (não acessa a camada i18n da \tool{AURORA}
principal). \lstinline|detectLocale()| lê \lstinline|aurora-locale| (ou o legado
\lstinline|aurora-yanc-lang|) do \lstinline|localStorage|, com \emph{default}
\lstinline|pt|. As \emph{strings} resolvidas são aplicadas aos elementos da
\emph{toolbar} no carregamento.

\subsection{Carregamento e interação do SVG}

Em \lstinline|compilation-complete|, \lstinline|_onCompilationComplete(data)|
zera as pilhas de navegação, define o módulo corrente e chama
\lstinline|_loadSVG(svgPath, module)|. \lstinline|_loadSVG| lê o SVG via
\lstinline|electronAPI.readFile| (passando pelo \emph{main}, sem
\lstinline|fetch('file://')|), injeta-o no DOM com \emph{fade-in}, atualiza
\emph{breadcrumbs} e info do módulo, e instala as interações.

\lstinline|_setupSVGInteractions()| instala três comportamentos:

\begin{description}[leftmargin=1.4em,style=nextline]
  \item[Navegação por clique] Cada \lstinline|g[data-cell-type]| cujo tipo é
        clicável (mesma lista de \emph{skip} do \emph{main}, em
        \lstinline|_isClickableModuleType|) ganha cursor de ponteiro, realce ao
        passar o mouse e \emph{tooltip}. Um clique simples agenda
        \lstinline|navigateToModule(type)| após $\sim$250\,ms, dando janela para
        um eventual duplo-clique cancelar a navegação.
  \item[Abrir o código-fonte por duplo-clique] Cada
        \lstinline|g[onclick^="gotosrc"]| (atributo que o netlistsvg emite a
        partir do \lstinline|src| do \tool{Yosys}) tem seu \lstinline|onclick|
        removido e substituído por um \emph{listener} de \lstinline|dblclick| que
        extrai o localizador \lstinline|arquivo:linha.col| e chama
        \lstinline|electronAPI.openSourceFile(...)|. (O \emph{stub} global
        \lstinline|window.gotosrc = () => {}| neutraliza o \emph{handler} inline
        original.)
  \item[Realce de fios] Cliques em \lstinline|path|/\lstinline|line|/\lstinline|polyline|
        disparam \lstinline|_highlightWireConnection|, que faz \emph{flood-fill}
        geométrico de \emph{endpoints} (tolerância de 5 unidades) para realçar
        toda a conexão elétrica em \lstinline|var(--aurora-mint)|.
\end{description}

\subsection{Navegação hierárquica}

A entrada em submódulos é feita por \lstinline|navigateToModule(name)|, que chama
\lstinline|electronAPI.generateSVGFromModule(name, tempDir)| --- o \emph{main}
gera o SVG do \path{<modulo>.json} já existente no \emph{tempDir} (sem reexecutar
o \tool{Yosys}). O renderer mantém duas pilhas, \lstinline|navigationHistory| e
\lstinline|forwardHistory|, e oferece:

\begin{itemize}[leftmargin=1.4em]
  \item \lstinline|navigateBack()| / \lstinline|navigateForward()| (também ligados
        aos botões laterais do mouse X1/X2 e a \lstinline|Esc| para voltar);
  \item \emph{breadcrumbs} clicáveis que truncam o histórico até o nível
        escolhido;
  \item \lstinline|recompile()|, que reobtém os caminhos via
        \lstinline|getPrismCompilationPaths| e chama
        \lstinline|prismRecompile|.
\end{itemize}

\subsection{Pan, zoom e download}

O canvas é uma \lstinline|<div>| com \emph{transform} CSS
(\lstinline|translate(...) scale(...)|). O zoom é \enquote{suave}: a roda do mouse
ajusta um \lstinline|targetScale| (fator \lstinline|exp(-deltaY*0.0016)|, clamp
0{,}1--5) e \lstinline|currentScale| é interpolado por \emph{requestAnimationFrame},
ancorado sob o cursor. Há também \emph{pan} por arraste, gestos de toque
(\emph{pinch} e \emph{pan}), \lstinline|fitToScreen()| (ajusta a $\sim$90\% do
\emph{viewport}), \lstinline|resetView()| e atalhos de teclado
(\lstinline|Ctrl+=/-/0/F|, \lstinline|Ctrl+R| recompila, \lstinline|F11| tela
cheia). \lstinline|downloadSVG()| serializa o SVG atual (clonado, com
\emph{namespaces} XML adicionados) para um \lstinline|Blob| baixável, nomeado pelo
módulo corrente --- funciona em qualquer nível de navegação.

\section{Simulação interativa com DigitalJS (O9)}
\label{sec:digitaljs}

Desde a fase O9, o \tool{PRISM} oferece, além do esquemático estático, uma
\emph{simulação lógica interativa} do mesmo \emph{netlist}, via \tool{DigitalJS}.
O botão \enquote{Simular} (id \lstinline|simToggle|) alterna entre as duas
superfícies. \textbf{A integração está de fato funcional} --- não é apenas uma
dependência presente: há \emph{handler} IPC dedicado, conversão real do
\emph{netlist} e renderização da paper interativa, conforme o código abaixo.

As dependências relevantes no \file{package.json} da \tool{AURORA} são
\lstinline|digitaljs| \lstinline|^0.14.2| e \lstinline|yosys2digitaljs|
\lstinline|^0.10.3| (instalado com \lstinline|main: dist/index.js| e um
sub-\emph{entry} \lstinline|core|).

\subsection{Construção do circuito no main}

O \emph{handler} \lstinline|prism:build-digitaljs| (em \file{main/ipc/prism.js})
resolve o top-level do \file{.spf} e chama
\lstinline|buildDigitalJSCircuit(paths, topLevelModule, tempDir)|. Esse fluxo é
deliberadamente \emph{separado} do esquemático para nunca perturbá-lo:

\begin{enumerate}[leftmargin=2.2em]
  \item \lstinline|collectSynthFiles| coleta os \file{.v} com a mesma lógica de
        \lstinline|runYosysCompilationWithPaths| (HDL + \lstinline|synthesizableFiles|
        + \path{TopLevel/} + \path{<proc>/Hardware/}).
  \item Gera um \emph{script} \tool{Yosys} \enquote{amigável ao
        yosys2digitaljs}, em nível de palavra (sem \lstinline|abc|/\lstinline|techmap|,
        para manter \lstinline|$add|/\lstinline|$mux|/\lstinline|$dff| como
        células \tool{DigitalJS} em vez de \emph{gates}):
\end{enumerate}

\begin{lstlisting}[caption={Script Yosys da simulação (buildDigitalJSCircuit)}]
read_verilog "..."
...
hierarchy -top <topLevelModule>
proc
opt_clean
memory -nomap
wreduce -memx
opt_clean
write_json "digitaljs.json"
\end{lstlisting}

\begin{enumerate}[leftmargin=2.2em,start=3]
  \item Roda o \emph{próprio} \tool{Yosys} (allowlisted) via
        \lstinline|spawnTracked|, com \emph{timeout} de 45\,s (uma síntese
        descontrolada falha com mensagem clara em vez de \emph{spinner} infinito).
  \item Conta as células do \emph{netlist}; se passar de \lstinline|MAX_CELLS = 3000|,
        \emph{recusa} a simulação com orientação para usar o esquemático
        (\tool{DigitalJS} é um simulador didático para módulos pequenos; o
        conversor roda síncrono no processo principal e o \emph{layout} no
        navegador degrada além de alguns milhares de células).
  \item Carrega \emph{lazy} o conversor \lstinline|require('yosys2digitaljs/core')|
        e chama \lstinline|yosys2digitaljs(yosysJson, {})| --- uma função
        \emph{pura} que recebe o JSON do \tool{Yosys} já produzido (nenhum
        \emph{spawn} adicional de \tool{Yosys}) e devolve o \lstinline|TopModule|
        no formato \tool{DigitalJS}.
  \item \lstinline|stripYosysLabels(circuit)| zera \emph{labels} internos
        gerados pelo \tool{Yosys} (que começam com \lstinline|$|, como
        \lstinline|$xor$src:line$n|), preservando os úteis (\lstinline|a|,
        \lstinline|b|, \lstinline|sum| etc.).
\end{enumerate}

O \emph{handler} retorna \lstinline|{ ok: true, circuit, topLevelModule }|, com
\emph{logs} de progresso por fase (\enquote{synthesizing\dots}, \enquote{netlist
N cells --- converting\dots}, \enquote{converted to M devices\dots}).

\subsection{Renderização interativa no renderer}

No renderer, \lstinline|enterSimMode()| (em \file{html/prism/prism.js}):

\begin{itemize}[leftmargin=1.4em]
  \item Usa um \emph{guard} de reentrância (\lstinline|_simBusy|) para que um
        segundo clique não inicie um segundo \emph{build} concorrente.
  \item Obtém o circuito via \lstinline|electronAPI.buildDigitalJS(paths)|.
  \item \emph{Carrega \tool{DigitalJS} de forma \emph{lazy}}
        (\lstinline|_loadDigitalJS()|): expõe o \lstinline|jQuery| global,
        importa o \lstinline|jquery-ui| \emph{completo} (na ordem correta, para
        que \lstinline|$.widget|/\lstinline|$.fn.dialog| existam antes de o
        \emph{bundle} do \tool{DigitalJS} avaliar seu \lstinline|dialog.js|), e
        só então faz \lstinline|import('digitaljs')| para obter a classe
        \lstinline|Circuit|. Manter isso fora do \emph{import} de topo significa
        que o fluxo do esquemático nunca avalia o \tool{DigitalJS} (um erro de
        carga não derruba a janela inteira) e o \emph{chunk} de $\sim$2\,MB só é
        baixado sob demanda.
  \item Instancia \lstinline|new Circuit(res.circuit, { layoutEngine: 'dagre' })|
        e exibe com \lstinline|circuit.displayOn(paperHost)|, depois
        \lstinline|circuit.start()|. O motor de navegador \emph{síncrono} (padrão)
        com \emph{layout} \lstinline|dagre| evita qualquer \emph{Web Worker},
        permitindo rodar sob \lstinline|file://|.
\end{itemize}

A simulação tem seu próprio \emph{pan}/\emph{zoom} (transform CSS no
\lstinline|.djs-wrapper|, espelhando os valores do esquemático), \emph{pan} por
arraste de área vazia (\lstinline|blank:pointerdown|), \lstinline|_fitPaper()|
para centralizar/ajustar, e \lstinline|_buildValueOverlays()|, que sobrepõe um
dígito vivo \lstinline|0|/\lstinline|1|/\lstinline|x| em cada caixa de
entrada/saída de 1 bit (o \tool{DigitalJS} sozinho só pinta a caixa) e o atualiza
a cada mudança de sinal. \lstinline|exitSimMode()| e \lstinline|_destroyCircuit()|
param e descartam o circuito e a \emph{paper} JointJS, evitando vazamentos em
ciclos repetidos de entrar/sair. Uma recompilação do esquemático sai
automaticamente do modo de simulação (o circuito vivo fica obsoleto).

\section{Árvore de hierarquia pós-síntese}

A \enquote{árvore de hierarquia} que o \tool{PRISM} apresenta não é uma estrutura
em árvore separada, e sim a \emph{navegação} sobre os módulos individuais que
\lstinline|splitHierarchyJson| derivou do \path{hierarchy.json}. Como o
\emph{script} do esquemático roda \lstinline|hierarchy -top| seguido de
\lstinline|proc| (sem \emph{flatten}), o \path{hierarchy.json} preserva a
estrutura hierárquica: cada módulo guarda suas células, e células cujo
\lstinline|type| é outro módulo clicável tornam-se pontos de entrada.

Em tempo de execução, o usuário começa no top-level e \enquote{desce} a árvore
clicando em células de submódulo; os \emph{breadcrumbs} mostram o caminho atual e
permitem saltar para qualquer ancestral. Como cada nível corresponde a um
\path{<modulo>.json} já materializado, descer/subir é instantâneo --- nenhuma
reexecução de \tool{Yosys} ocorre durante a navegação (apenas uma nova chamada a
\lstinline|netlistsvg.render| por módulo visitado).

\section{Resumo de artefatos e responsabilidades}

\begin{tabularx}{\linewidth}{Y Z}
\toprule
\textbf{Arquivo} & \textbf{Responsabilidade} \\
\midrule
\file{main/ipc/prism.js} & janela, \emph{pipeline} Yosys, split por módulo, render SVG, skins, DigitalJS; \emph{handlers} IPC \\
\file{js/app/preload\_prism.js} & \emph{bridge} sandboxizada com \emph{allowlist} de canais \\
\file{html/prism/prism.html} & layout da janela (toolbar, breadcrumbs, canvas SVG, host DigitalJS) \\
\file{html/prism/prism.js} & classe \lstinline|PRISMViewer|: i18n, SVG, navegação, pan/zoom, sim \\
\file{js/compilation/compilation\_flow.js} & botão PRISM, pré-flight, payload de caminhos, \emph{overrides} \\
\file{main/render\_loader.js} & seleção dev/prod do renderer (\lstinline|loadPage|) \\
\repo{netlistsvg-aurora} & fork de \lstinline|@silimate/netlistsvg| (fix de tsconfig; skins default) \\
\file{assets/prism-skins/*.svg} & símbolos customizados da família SAPHO (mesclados em runtime) \\
\file{scripts/prism-skin-standard.js} & gerador do padrão de símbolos SAPHO \\
\bottomrule
\end{tabularx}

\begin{nota}
Em síntese: o \tool{PRISM} sintetiza o RTL do projeto com o \tool{Yosys} bundle
(\lstinline|read_verilog -setattr src| $\rightarrow$ \lstinline|hierarchy| $\rightarrow$
\lstinline|proc| $\rightarrow$ \lstinline|setundef -zero| $\rightarrow$
\lstinline|opt_clean| $\rightarrow$ \lstinline|write_json|), divide o
\emph{netlist} por módulo, e renderiza cada módulo \emph{in-process} com o fork
\repo{netlistsvg-aurora} usando \emph{skins} SAPHO mescladas em tempo de
execução, produzindo um esquemático SVG navegável. Sobre o mesmo desenho de
fluxo, a fase O9 acrescentou uma simulação lógica interativa via
\tool{yosys2digitaljs}~+~\tool{DigitalJS}, integrada e funcional, ativada pelo
botão \enquote{Simular}.
\end{nota}

% !TEX root = ../main.tex
\chapter{AURORA: o subsistema Aurora Intelligence (IA)}
\label{cap:12-aurora-intelligence-ia}

O \tool{Aurora Intelligence} e o assistente de IA integrado a IDE \tool{AURORA}.
Ele permite que o usuario converse com um modelo de linguagem sobre o projeto
aberto, sobre \cmm{}, Verilog e o ecossistema SAPHO, e que esse modelo
\emph{aja} sobre o ambiente de trabalho --- lendo arquivos, disparando
compilacoes, configurando formas de onda e mexendo no projeto --- por meio de um
conjunto curado de ferramentas (\emph{tools}). O subsistema esta
\textbf{completamente implementado}. O unico item ainda pendente e o modelo
proprio do laboratorio: o \emph{provider} e os modelos treinados pelo NIPSCERN
ainda \textbf{nao foram treinados}, sendo um item de roteiro. Enquanto isso, o
assistente opera contra provedores externos (descritos adiante), de modo que a
funcionalidade ja esta disponivel ao usuario final.

\begin{nota}
Este capitulo descreve o codigo real em \path{aurora/main/ai/*},
\path{aurora/main/ipc/ai.js} e \path{aurora/js/ai/*} mais
\file{js/ui/ai\_assistant\_manager.js}. Os assuntos vizinhos --- pipeline de
compilacao, terminais, formas de onda, PRISM --- sao tocados aqui apenas no
ponto em que as \emph{tools} os acionam; cada um e documentado em seu proprio
capitulo.
\end{nota}

\section{Arquitetura em duas camadas}

O Aurora Intelligence segue o modelo de processos do Electron: o
\textbf{processo principal} (\emph{main}) e dono de tudo o que e sensivel ou de
sistema (chaves de API, chamadas de rede aos provedores, subprocessos das CLIs,
servidor MCP local, persistencia em disco), e o \textbf{renderer} e dono da
interface (o painel lateral de chat) e da execucao concreta de cada ferramenta
contra o espaco de trabalho (\code{window.AuroraAPI}). As duas metades se
comunicam por IPC.

\begin{lstlisting}[caption={Fluxo geral de um turno de chat (ASCII).}]
  renderer (painel)                  main (processo principal)
  -----------------                  -------------------------
  envia mensagens  --- ai:chat-start --->  escolhe o transporte
                                           (SDK | claude-code | chatgpt)
                   <-- ai:chat-event ----  text-delta / tool-call /
                       (multicast)         tool-result / finish | aborted | error
   .                                       |  ao chamar uma tool:
   .            <--- ai:tool-exec --------'  tool_bridge.runTool
  tool_runner executa                       (renderer faz o trabalho real)
  AuroraAPI.<ns>.<fn>(...)
   --- ai:tool-result -------------------> resolveToolResult -> resultado volta
                                           ao modelo
\end{lstlisting}

Um ponto central de projeto e que \textbf{os dois transportes convergem nos
mesmos eventos de chat}. Independentemente de a resposta vir do Vercel AI SDK ou
de uma CLI por assinatura, o main emite sempre os mesmos pacotes
\code{ai:chat-event}, de modo que o laco de chat do renderer e
\emph{transporte-agnostico}.

\subsection{O ciclo de vida dos eventos de chat}

Todo turno termina com exatamente um evento de fechamento
(\code{finish}, \code{aborted} ou \code{error}). A sequencia tipica de eventos
emitidos por \file{main/ai/chat.js} (e replicada pelas pontes das CLIs) e:

\begin{tabularx}{\linewidth}{l Y}
\toprule
\textbf{Tipo} & \textbf{Carga e significado} \\
\midrule
\code{text-delta}   & \code{\{ delta \}} --- um fragmento de texto da resposta, transmitido token a token. \\
\code{tool-call}    & \code{\{ toolUseId, toolName, args \}} --- o modelo decidiu chamar uma ferramenta. \\
\code{tool-result}  & \code{\{ toolUseId, toolName, result \}} --- o resultado dessa chamada. \\
\code{finish}       & \code{\{ text, usage \}} --- fim normal do turno, com a contagem de tokens. \\
\code{aborted}      & \code{\{ text \}} --- o usuario interrompeu (ou um \emph{watchdog} agiu); devolve o texto parcial. \\
\code{error}        & \code{\{ message \}} --- falha; mensagem ja tratada para exibicao. \\
\bottomrule
\end{tabularx}

\section{Transporte 1: Vercel AI SDK (provedores com chave de API)}

A primeira via de transporte usa o \textbf{Vercel AI SDK} (pacote \code{ai},
funcoes \code{streamText} e \code{generateText}) para falar com provedores
comerciais por meio de uma chave de API fornecida pelo usuario (modelo
\emph{BYOK}, \emph{bring your own key}). A integracao vive em
\file{main/ai/provider.js} (montagem do provider e chamadas de um disparo) e em
\file{main/ai/chat.js} (laco de \emph{streaming} com ferramentas).

\subsection{Provedores suportados e suas dependencias}

\file{main/ai/provider.js} carrega cada SDK de provedor de forma defensiva (via
um \code{tryRequire}: se um pacote falhar ao carregar, apenas aquele provedor e
desabilitado em vez de derrubar o processo principal). Os provedores reconhecidos
correspondem as dependencias \code{@ai-sdk/*} declaradas em
\file{package.json}:

\begin{tabularx}{\linewidth}{l l Y}
\toprule
\textbf{Provedor} & \textbf{Pacote SDK} & \textbf{Modelo padrao (\code{DEFAULT\_MODELS})} \\
\midrule
openai    & \code{@ai-sdk/openai}    & \lstinline|gpt-4o-mini| \\
anthropic & \code{@ai-sdk/anthropic} & \lstinline|claude-haiku-4-5-20251001| \\
google    & \code{@ai-sdk/google}    & \lstinline|gemini-2.5-flash| \\
deepseek  & \code{@ai-sdk/deepseek}  & \lstinline|deepseek-chat| \\
groq      & \code{@ai-sdk/groq}      & \lstinline|llama-3.3-70b-versatile| \\
ollama    & (via \code{@ai-sdk/openai}) & \lstinline|llama3.1:8b| \\
\bottomrule
\end{tabularx}

\begin{nota}
O \tool{Ollama} e um caso especial: nao tem pacote dedicado. \code{getProvider}
o monta reaproveitando \code{createOpenAI} apontado para um \emph{base URL}
local (padrao \url{http://localhost:11434/v1}), no endpoint
\path{/v1/chat/completions}. Para o Ollama, a "chave" armazenada no cofre e, na
verdade, o \emph{base URL}; se nao houver nenhum, usa-se o padrao local, de modo
que ele funciona sem configuracao.
\end{nota}

\subsection{Governanca de modelos}

A escolha do modelo segue uma cadeia de resolucao deterministica em
\code{resolveModelId(provider, requested)}:

\begin{itemize}[leftmargin=1.4em]
  \item entrada vazia, \lstinline|'default'| ou \lstinline|'latest'| $\rightarrow$ o modelo padrao do provedor (\code{DEFAULT\_MODELS});
  \item um id sabidamente aposentado/renomeado $\rightarrow$ seu substituto, via o mapa \code{MODEL\_MIGRATIONS} (um mapa por provedor, semeado conforme ids sao aposentados);
  \item qualquer outro valor $\rightarrow$ inalterado.
\end{itemize}

A preferencia de modelo por provedor (a sobreposicao escolhida pelo usuario) e
persistida por \file{main/ai/prefs.js} em
\path{userData/aurora-ai-prefs.json} (dado nao sensivel). \code{getModelFor}
combina essa preferencia com a resolucao acima.

Quando um id de modelo esta indisponivel em tempo de execucao (aposentado, com
erro de digitacao ou nao habilitado para a chave), a heuristica
\code{isModelUnavailableError} reconhece o erro do SDK (por exemplo HTTP 404 ou
mensagens contendo \enquote{model \dots not found}) e o sistema recua uma vez
para o modelo padrao do provedor, em vez de encerrar o turno com uma mensagem
obscura. Esse recuo existe em \code{testConnection}, em \code{generateOneshot} e
no laco de chat.

\subsection{O laco de \emph{streaming} com ferramentas}

\file{main/ai/chat.js} mantem uma \code{Map} de sessoes ativas (uma por
\code{sessionId} gerado pelo renderer) e usa \code{streamText} consumindo
\code{result.fullStream} --- e nao apenas o fluxo de texto --- para que partes de
\code{tool-call} e \code{tool-result} aparecam no painel como
\emph{chips} de atividade. Caracteristicas relevantes:

\begin{description}[leftmargin=0pt, style=nextline]
  \item[Limite de passos] \code{stopWhen: stepCountIs(MAX\_STEPS)} com \code{MAX\_STEPS = 24} impede que o modelo entre em laco infinito de chamadas de ferramenta.
  \item[Anti-congelamento] um temporizador de inatividade (\code{STREAM\_INACTIVITY\_MS = 120000}) aborta o turno se o modelo ficar mudo por tempo demais \emph{sem} nenhuma ferramenta em voo. Enquanto ha ferramentas pendentes (controladas por um \code{Set} de ids), o temporizador fica pausado --- assim uma compilacao lenta nunca o dispara.
  \item[Conteudo multimodal] \code{toSdkMessage} expande anexos do compositor em partes nativas do SDK: o texto, cada imagem como parte \code{image} (\emph{data URL}) e cada arquivo como bloco de texto cercado.
  \item[\emph{Prompt caching} da Anthropic] quando o provedor e \code{anthropic} e o \emph{system prompt} e grande, ele e enviado com \code{cacheControl: \{ type: 'ephemeral' \}}, reduzindo o custo de tokens repetidos em turnos seguintes. Outros SDKs ignoram esse campo.
  \item[Compatibilidade] um caso conhecido de \code{item\_reference} (incompatibilidade do \emph{prompt caching} da Anthropic apos um resultado de ferramenta) e tratado como encerramento gracioso, com um eventual reenvio sem ferramentas para que o modelo resuma o que encontrou.
  \item[Ferramentas em texto] para modelos (tipicamente via Ollama) que emitem chamadas de ferramenta como JSON no proprio texto em vez de usar o campo de \emph{function calling}, ha um caminho de \emph{fallback} que extrai esse JSON, executa as ferramentas pela ponte e injeta os resultados numa chamada de continuacao. Artefatos textuais de chamada de ferramenta (tags \lstinline|<tool_call>|, JSON \enquote{name/arguments}) sao removidos do texto exibido.
\end{description}

\subsection{Geracao de um disparo}

Alem do chat com \emph{streaming}, \code{generateOneshot} oferece uma geracao
de texto unica (sem ferramentas, sem historico, sem \emph{streaming}) para
recursos que transformam uma entrada em uma saida num unico passo --- por
exemplo o gerador de \emph{harness} por IA. So provedores do SDK sao aceitos
nesse caminho (as pontes das CLIs sao apenas de chat, com excecao do caminho
proprio descrito na \autoref{sec:ia-claude-oneshot}). Para o Gemini, ele
desativa explicitamente os \enquote{tokens de raciocinio}
(\code{thinkingBudget: 0}) para que todo o orcamento de saida va para a resposta.

\section{Transporte 2: CLIs por assinatura (Claude Code e Codex)}

A segunda via permite usar a \emph{assinatura} do usuario (Claude Pro/MAX ou um
plano ChatGPT) em vez de uma chave de API metrada. Em vez de falar HTTP direto
com o provedor, a AURORA \textbf{invoca a CLI instalada localmente} em modo
nao interativo e traduz os eventos que ela emite para os mesmos pacotes
\code{ai:chat-event}. Sao duas pontes:

\begin{itemize}[leftmargin=1.4em]
  \item \file{main/ai/claude\_code.js} --- a CLI \tool{Claude Code} da Anthropic, provedor \enquote{claude-code} no painel;
  \item \file{main/ai/codex\_cli.js} --- a CLI \tool{Codex} da OpenAI, provedor \enquote{chatgpt} no painel.
\end{itemize}

Em ambas, a autenticacao e feita pelo \emph{login} OAuth da propria CLI
(\lstinline|claude login| / \lstinline|codex login|); a AURORA \textbf{nunca} ve
nem armazena um token. As pontes ainda \emph{removem} as variaveis de ambiente
de chave de API do processo filho (\code{ANTHROPIC\_API\_KEY},
\code{ANTHROPIC\_AUTH\_TOKEN} no Claude; \code{OPENAI\_API\_KEY} no Codex) para
que a CLI nao caia silenciosamente em cobranca metrada por API.

\subsection{Invocacao das CLIs}

\begin{description}[leftmargin=0pt, style=nextline]
  \item[Claude Code] e iniciada em modo \emph{print}: \lstinline|claude -p --output-format stream-json --verbose --include-partial-messages ...|. Os eventos \code{stream-json} (do tipo \code{system}, \code{stream\_event}, \code{assistant}, \code{user}, \code{result}, \code{rate\_limit\_event}) sao traduzidos para eventos de chat. O \emph{system prompt} grande viaja no corpo do \emph{prompt} (lido por \code{stdin}) quando ultrapassa um limiar conservador, contornando o limite de tamanho de linha de comando do Windows.
  \item[Codex] e iniciada como \lstinline|codex exec --json --skip-git-repo-check --dangerously-bypass-approvals-and-sandbox ...|. Seus eventos (modelo de \emph{thread}/\emph{turn}/\emph{item}: \code{thread.started}, \code{item.started}, \code{item.completed}, \code{turn.completed}, \code{turn.failed}) sao traduzidos da mesma forma.
\end{description}

Para nao perder contexto, cada ponte mantem um mapa de
\code{conversationId} para o id de sessao/\emph{thread} da CLI, de modo que
turnos seguintes usam \lstinline|--resume| (Claude) ou \lstinline|exec resume|
(Codex). Quando nao ha sessao a retomar mas o renderer entrega historico, os
turnos anteriores sao dobrados no \emph{prompt} dentro de um bloco
\lstinline|<conversation_context>|. Ambas as pontes tem \emph{watchdog} de
inatividade (mesma logica de \code{Set} de ferramentas pendentes do
transporte 1) e encerram subprocessos no fim da aplicacao
(\code{killAll}, via \lstinline|taskkill /T /F| no Windows).

\subsection{Deteccao, anexos e uso}

Cada ponte expoe \code{detect()} (instalado? autenticado? plano? versao?) e
\code{getUsage()} (uma contagem de tokens acumulada por execucao). O plano do
ChatGPT e extraido, somente leitura, da \emph{claim} \code{chatgpt\_plan\_type}
do \emph{id token} JWT das credenciais. Anexos sao tratados por
\file{main/ai/attachments.js}: arquivos viram blocos de texto cercados; imagens,
que as CLIs nao recebem inline, sao gravadas em arquivos temporarios e
referenciadas por caminho quando a CLI sabe le-las nativamente (Claude Code), ou
degradam para uma nota \enquote{este provedor nao ve imagens} (Codex). Os
temporarios sao limpos no inicio da aplicacao e podados por TTL a cada escrita.

\subsection{\emph{Download} das CLIs sob demanda}

O instalador da AURORA \textbf{nao empacota} os binarios nativos das CLIs
(eram cerca de 460 MB somados). Eles sao baixados na primeira vez que o usuario
de fato usa cada provedor:

\begin{itemize}[leftmargin=1.4em]
  \item \file{main/ai/cli\_manifest.js} fixa exatamente qual pacote de
        plataforma do npm buscar, sua versao e seu \emph{hash} de integridade
        (\emph{Subresource Integrity}, \lstinline|sha512-...|), por plataforma.
        As versoes (Claude \lstinline|2.1.144|, Codex \lstinline|0.131.0|) devem
        acompanhar as dependencias-base de \file{package.json}.
  \item \file{main/ai/cli\_downloader.js} baixa o \emph{tarball}, \textbf{verifica
        a integridade SHA-512 antes da extracao}, extrai com o \code{tar} do
        sistema para um cache por usuario em
        \path{userData/cli-cache/<pkg>@<versao>/}, e poda versoes antigas.
        Chamadas concorrentes para a mesma CLI compartilham um unico
        \emph{download}.
  \item \file{main/ai/cli\_locator.js} resolve o executavel em quatro \emph{layouts}
        de execucao, em ordem de prioridade: (0) o cache de \emph{download} sob
        demanda; (1) desenvolvimento (\path{node_modules/}); (2) build empacotado
        com a dependencia presente (reescrevendo \path{app.asar} para
        \path{app.asar.unpacked}, pois binarios nativos nao rodam de dentro do
        \emph{asar}); (3) uma instalacao global no \code{PATH}.
\end{itemize}

\section{O servidor MCP HTTP local}

As CLIs por assinatura conhecem apenas as \emph{suas} ferramentas nativas
(Read/Edit/Bash etc.). Para que elas dirijam o \emph{pipeline} da AURORA em vez
de invocar binarios da \emph{toolchain} por um \emph{shell}, a AURORA expoe sua
propria superficie de ferramentas a CLI por meio de um \textbf{servidor MCP
(\emph{Model Context Protocol}) HTTP local}, implementado em
\file{main/ai/aurora\_mcp\_server.js} sobre o
\code{@modelcontextprotocol/sdk}.

\subsection{Natureza e seguranca do servidor}

\begin{description}[leftmargin=0pt, style=nextline]
  \item[Transporte] \emph{Streamable HTTP} em modo \emph{stateless} (\code{sessionIdGenerator: undefined}): um \code{Server} e um \emph{transport} novos a cada requisicao. Como o ritmo de chamadas e humano (poucas por turno), a alocacao extra e desprezivel.
  \item[Endereco] vinculado a \textbf{127.0.0.1 numa porta efemera aleatoria} (\code{listen(0, '127.0.0.1')}). A URL final tem a forma \path{http://127.0.0.1:<porta>/mcp/<token>}.
  \item[Defesa contra \emph{DNS rebinding}] o cabecalho \code{Host} e reverificado a cada requisicao; so \code{127.0.0.1} ou \code{localhost} sao aceitos.
  \item[\emph{Token} de sessao] exige-se um \emph{token} de 256 bits (\code{crypto.randomBytes(32)}) em \emph{toda} requisicao, no caminho (\path{/mcp/<token>}) ou como \code{Authorization: Bearer}. Isso bloqueia uma pagina web maliciosa que alcance o \emph{loopback} mas nao consiga adivinhar a porta efemera e o \emph{token}.
  \item[Metodo] somente \code{POST}; corpo limitado a 4 MB.
\end{description}

O servidor registra cada entrada de \code{tools.TOOL\_MANIFEST} como uma
ferramenta MCP. As CLIs recebem um arquivo \lstinline|--mcp-config| apontando
para ele (Claude com \lstinline|--strict-mcp-config|; Codex via
\lstinline|-c mcp_servers.aurora.url=...|), de forma que as ferramentas da AURORA
aparecem a CLI como \lstinline|mcp__aurora__<nome>|. Cada chamada e reencaminhada
ao mesmo \code{tool\_bridge.runTool} usado pelo caminho do SDK --- ou seja, o
mesmo \emph{ask-before-write}, o mesmo \emph{log} de auditoria e o mesmo
\code{AuroraAPI} no renderer.

\subsection{Regras de uso das ferramentas}

Para garantir que o modelo \emph{prefira} as ferramentas \lstinline|mcp__aurora__*|
e nao recorra ao \emph{shell}, cada ponte injeta um texto de regras no primeiro
turno (\code{MCP\_TOOL\_RULES}) que instrui o modelo a nunca chamar binarios da
YANC (\tool{cmmcomp}, \tool{cppcomp}, \tool{asmcomp}, \dots), \tool{iverilog},
\tool{verilator} ou \tool{gtkwave} por um \emph{shell}, usando em vez disso as
ferramentas equivalentes da AURORA. Alem disso:

\begin{itemize}[leftmargin=1.4em]
  \item na CLI do Claude Code, ferramentas que escrevem ou usam \emph{shell}
        (\code{Bash}, \code{Edit}, \code{Write}, \code{MultiEdit},
        \code{NotebookEdit} \dots) sao \textbf{desabilitadas} via
        \lstinline|--disallowed-tools|, de modo que toda escrita passe pelas
        ferramentas da AURORA, que sao protegidas pelo cartao Permitir/Negar;
  \item na CLI do Codex, cuja ferramenta de \emph{shell} nao pode ser removida,
        a contencao e feita pelas regras do \emph{system prompt} (uma garantia
        mais suave, declaradamente uma limitacao aceita dessa integracao);
  \item o tempo limite de leitura de chamadas MCP nas CLIs e elevado (10 min) para
        ficar acima do teto da ponte de ferramentas, evitando que a CLI declare
        falha enquanto uma operacao lenta (como renomear projeto) ainda esta sendo
        concluida do lado do renderer.
\end{itemize}

\section{O conjunto de ferramentas (\emph{tools})}

O catalogo de ferramentas que o assistente pode chamar vive em
\file{main/ai/tools.js}, na constante \code{TOOL\_MANIFEST}. Trata-se de
\textbf{dados puros} (sem fechamentos), de modo que pode ser enviado por IPC
verbatim: o renderer puxa o mesmo manifesto e usa os campos \code{api},
\code{argStyle} e \code{argNames} para despachar a chamada contra
\code{window.AuroraAPI}, mantendo uma unica fonte de verdade.

\subsection{Contagem e estrutura}

O manifesto define \textbf{98 ferramentas}, das quais \textbf{41 sao de leitura}
(\code{access: 'read'}, executam sem confirmacao) e \textbf{57 sao de escrita}
(\code{access: 'write'}, exibem o cartao de confirmacao antes de agir). Cada
entrada declara:

\begin{tabularx}{\linewidth}{l Y}
\toprule
\textbf{Campo} & \textbf{Significado} \\
\midrule
\code{name}        & nome da ferramenta exposto ao modelo. \\
\code{description} & descricao em linguagem natural para o modelo. \\
\code{access}      & \lstinline|'read'| ou \lstinline|'write'|. \\
\code{api}         & par \code{[namespace, metodo]} em \code{window.AuroraAPI}. \\
\code{argStyle}    & \lstinline|'none'| / \lstinline|'positional'| / \lstinline|'object'|: como transformar os argumentos JSON na chamada. \\
\code{argNames}    & ordem dos argumentos posicionais, quando aplicavel. \\
\code{inputSchema} & esquema JSON dos argumentos. \\
\bottomrule
\end{tabularx}

\subsection{Categorias de ferramentas}

As ferramentas cobrem todo o fluxo de trabalho da IDE. A tabela a seguir agrupa
exemplos representativos por area funcional (a lista completa esta no codigo).

\begin{longtable}{l Y}
\caption{Categorias de ferramentas em \file{main/ai/tools.js}.}\\
\toprule
\textbf{Categoria} & \textbf{Exemplos de ferramentas} \\
\midrule
\endfirsthead
\toprule
\textbf{Categoria} & \textbf{Exemplos de ferramentas} \\
\midrule
\endhead
\bottomrule
\endfoot
Editor / leitura       & \code{get\_active\_file}, \code{get\_open\_files}, \code{read\_active\_file}, \code{get\_cursor} \\
Projeto / arquivos     & \code{get\_project\_tree}, \code{read\_file}, \code{create\_file}, \code{create\_folder}, \code{delete\_file}, \code{rename\_file}, \code{refresh\_file\_tree} \\
Edicao de texto        & \code{write\_active\_file}, \code{insert\_text}, \code{replace\_range}, \code{save\_file} \\
Compilacao / simulacao & \code{compile\_all}, \code{compile\_step}, \code{run\_in\_background}, \code{run\_verilator\_proc}, \code{run\_fast\_sim}, \code{cancel\_compilation} \\
Terminais              & \code{get\_terminal\_output}, \code{list\_terminals}, \code{read\_all\_terminals} \\
Processadores          & \code{list\_processors}, \code{create\_processor}, \code{delete\_processor}, \code{rename\_processor}, \code{get\_processor\_config}, \code{set\_processor\_config} \\
Topo / \emph{testbench}& \code{set\_top\_level}, \code{set\_testbench\_top} \\
Projetos (gestao)      & \code{create\_project}, \code{open\_project}, \code{rename\_project}, \code{backup\_project}, \code{list\_recent\_projects}, \code{get\_rename\_status} \\
Git                    & \code{git\_status}, \code{git\_log}, \code{git\_diff}, \code{git\_stage}, \code{git\_commit}, \code{git\_create\_branch}, \code{git\_pull}, \code{git\_push}, \code{git\_stash} \\
Formas de onda (GTKWave)& \code{list\_wave\_signals}, \code{select\_wave\_signals}, \code{open\_wave\_config}, \code{list\_gtkw\_files}, \code{add\_gtkw\_file}, \code{set\_active\_gtkw\_file} \\
Formas de onda (Surfer)& \code{list\_surfer\_files}, \code{add\_surfer\_file}, \code{set\_active\_surfer\_file} \\
Simulador / visualizador& \code{get\_simulator}, \code{set\_simulator}, \code{get\_waveform\_viewer}, \code{set\_waveform\_viewer} \\
Diretivas / opcodes    & \code{list\_directives}, \code{get\_directive}, \code{lookup\_compiler\_message}, \code{list\_opcodes}, \code{get\_opcode}, \code{analyze\_asm} \\
Configuracoes          & \code{get\_settings}, \code{set\_setting}, \code{get\_tree\_view}, \code{set\_tree\_view} \\
Sobreposicao de comando& \code{list\_compile\_steps}, \code{inspect\_compile\_command}, \code{preview\_compile\_command}, \code{list\_command\_overrides}, \code{set\_command\_override}, \code{clear\_command\_override} \\
Interacao com o usuario& \code{ask\_user\_question} \\
\end{longtable}

\subsection{Da definicao a chamada efetiva}

\code{buildTools(runToolFn)} converte o manifesto no objeto \code{tools} do
Vercel AI SDK, ligando o \code{execute} de cada ferramenta a uma funcao-ponte
fornecida pelo \file{chat.js}. \code{getManifest()} produz a versao serializavel
(sem fechamentos) que o renderer consome.

\begin{lstlisting}[caption={Construcao das ferramentas para o SDK (\file{main/ai/tools.js}).}]
function buildTools(runToolFn) {
  if (!tool || !jsonSchema) return {};
  const tools = {};
  for (const def of TOOL_MANIFEST) {
    tools[def.name] = tool({
      description: def.description,
      inputSchema: jsonSchema(def.inputSchema),
      execute: async (args) => runToolFn(def.name, args || {}),
    });
  }
  return tools;
}
\end{lstlisting}

\section{A ponte de ferramentas (\code{tool\_bridge})}

O SDK invoca o \code{execute()} de uma ferramenta dentro do \emph{processo
principal}, mas o espaco de trabalho a ser manipulado (Monaco, arvore de
arquivos, terminais) so existe no renderer. \file{main/ai/tool\_bridge.js}
preenche essa lacuna:

\begin{enumerate}[leftmargin=1.4em]
  \item \code{runTool(webContents, toolName, args)} gera um \code{requestId} e envia \code{ai:tool-exec} ao renderer;
  \item o \code{tool\_runner} do renderer executa o trabalho (incluindo a confirmacao \emph{ask-before-write}) e responde em \code{ai:tool-result};
  \item \code{resolveToolResult(requestId, result)} casa a resposta e resolve a \emph{promise} pendente exatamente uma vez.
\end{enumerate}

A conclusao e \textbf{dirigida por evento}, nao por relogio: a chamada se
resolve quando o renderer responde, ou imediatamente se o renderer morre no meio
(eventos \code{destroyed} / \code{render-process-gone}). O temporizador e apenas
um \emph{backstop} de ultimo recurso para o caso de um renderer vivo mas travado.
Os tempos limites sao escalonados conforme o tipo de ferramenta:

\begin{tabularx}{\linewidth}{l l Y}
\toprule
\textbf{Classe} & \textbf{Limite} & \textbf{Ferramentas} \\
\midrule
Padrao       & \code{TOOL\_TIMEOUT\_MS = 120000}   & a maioria das ferramentas. \\
Interativa   & \code{INTERACTIVE\_TIMEOUT\_MS = 600000} & \code{ask\_user\_question} (espera uma resposta humana deliberada). \\
Lenta        & \code{SLOW\_TIMEOUT\_MS = 300000}   & \code{rename\_project}, \code{rename\_processor} (mexem no disco). \\
\bottomrule
\end{tabularx}

A ponte \textbf{nunca rejeita}: falhas resolvem como
\code{\{ ok:false, error \}}, de modo que o SDK realimenta o erro ao modelo em
vez de abortar o fluxo.

\subsection{Execucao no renderer e o portao de permissao}

No renderer, \file{js/ai/tool\_runner.js} recebe cada \code{ai:tool-exec},
localiza a definicao no manifesto, aplica o portao de permissao
(\code{aiAssistantManager.confirmToolCall}), despacha a chamada para
\code{window.AuroraAPI[ns][fn](...)} e responde com \code{sendToolResult}. Uma
chamada negada resolve como um resultado normal
\code{\{ ok:false, error: 'The user denied this action.' \}}, de forma que o
modelo recebe \enquote{o usuario recusou} em vez de uma falha.

A decisao de permissao e pura, em \file{js/ai/tool\_permission.js}
(\code{decideToolPermission}), combinada com os modos definidos em
\file{js/ai/ai\_metadata.js}:

\begin{tabularx}{\linewidth}{l Y}
\toprule
\textbf{Modo} & \textbf{Comportamento} \\
\midrule
\lstinline|ask|    & toda chamada (leitura e escrita) pede um OK inline. \\
\lstinline|writes| & leituras correm livres; escritas pedem OK (\emph{padrao}). \\
\lstinline|allow|  & nada e perguntado; autonomia total. \\
\bottomrule
\end{tabularx}

Ha duas excecoes codificadas: \code{set\_command\_override} esta em
\code{ALWAYS\_CONFIRM} (sempre mostra o cartao, mesmo no modo \lstinline|allow|,
por reescrever uma linha de comando da \emph{toolchain}); e
\code{rename\_project}, \code{rename\_processor} e \code{get\_rename\_status}
estao em \code{PRE\_AUTHORIZED} (nunca passam pelo cartao bloqueante, pois rodam
desacoplados da IA). Alem dessas, o \code{tool\_runner} pre-autoriza o
\code{ask\_user\_question} pulando explicitamente o portao de permissao para ele
--- afinal ele \emph{e}, em si, um \emph{prompt} ao usuario.

\section{Cofre de chaves (\code{keystore})}

As chaves de API por provedor sao guardadas por \file{main/ai/keystore.js},
criptografadas com o \code{safeStorage} do Electron, que delega ao chaveiro do
sistema operacional: \textbf{DPAPI no Windows}, \emph{Keychain} no macOS e
\emph{libsecret} no Linux. O cofre em disco fica em
\path{userData/aurora-ai-keys.json} como um mapa de provedor para o
\emph{ciphertext} em base64.

\begin{nota}
\textbf{O texto puro nunca cruza a fronteira de IPC de volta ao renderer.} O
renderer so pode perguntar \emph{se} ha uma chave configurada
(\code{hasKey}, exposto por \code{ai:get-key-status}, que retorna apenas
booleanos) e \emph{testar} a chave; \textbf{ler os bytes e uma operacao exclusiva
do processo principal} --- nao existe canal \code{getKey} exposto.
\end{nota}

Detalhes de comportamento: \code{setKey} recusa-se a gravar se a criptografia do
SO nao estiver disponivel (lanca erro em vez de salvar texto puro silenciosamente);
\code{getKey} retorna \code{null} (e nao lanca) em falha de decifragem, para que o
chamador possa tratar um cofre corrompido sem quebrar; e os provedores aceitos
sao \code{openai}, \code{anthropic}, \code{google}, \code{deepseek}, \code{groq}
e \code{ollama} (\code{SUPPORTED\_PROVIDERS}).

\section{Persistencia: conversas, auditoria e preferencias}

\subsection{Conversas}

\file{main/ai/conversations.js} guarda \textbf{um arquivo JSON por conversa} em
\path{userData/aurora-intelligence-chats/<id>.json}. Esse desenho (um arquivo por
chat, em vez de um indice unico) faz com que adicionar uma conversa nunca
reescreva todo o historico e que um arquivo corrompido mate apenas aquela
conversa. \code{listAll()} devolve uma visao \emph{leve}
(id, titulo, provedor, modelo, datas, contagem de mensagens e tokens
acumulados), enquanto \code{read(id)} devolve o documento completo. Os canais IPC
correspondentes sao \code{ai:conv-list}, \code{ai:conv-read},
\code{ai:conv-save}, \code{ai:conv-delete}, \code{ai:conv-rename} e
\code{ai:conv-new-id}.

\subsection{Auditoria}

\file{main/ai/audit.js} mantem um \emph{log} \textbf{somente-acrescimo}, um objeto
JSON por linha, em \path{userData/aurora-intelligence-log.jsonl}. Cada chamada de
ferramenta e seu desfecho sao registrados (\code{kind: 'tool-call'} /
\code{'tool-result'}), dando ao usuario uma trilha revisavel do que o assistente
fez com seu espaco de trabalho. O registro e \emph{best-effort}: uma falha de
escrita e logada e engolida, pois a trilha de auditoria nunca pode quebrar um
turno de chat. O laco de chat acrescenta um par de entradas em torno de cada
ferramenta, e ha tambem \code{compile-override-applied} para registrar quando uma
sobreposicao de comando efetivamente disparou.

\subsection{Preferencias}

\file{main/ai/prefs.js} guarda preferencias \emph{nao} secretas (atualmente a
sobreposicao de modelo por provedor) em texto puro em
\path{userData/aurora-ai-prefs.json}.

\section{O painel no renderer}

O painel lateral de chat vive em \file{js/ui/ai\_assistant\_manager.js}, a classe
\code{AIAssistantManager}, apoiada por modulos puros extraidos em
\path{js/ai/*}:

\begin{tabularx}{\linewidth}{l Y}
\toprule
\textbf{Modulo} & \textbf{Papel} \\
\midrule
\file{js/ai/ai\_metadata.js}     & metadados de provedor/modelo, modos de permissao, formatadores. \\
\file{js/ai/provider\_view.js}   & marcacao (\emph{markup}) do seletor de provedor/modelo. \\
\file{js/ai/tool\_permission.js} & logica pura do portao de permissao. \\
\file{js/ai/tool\_runner.js}     & execucao de ferramentas contra \code{AuroraAPI}. \\
\file{js/ai/chat\_turn.js}       & montagem do \emph{payload} de cada turno. \\
\file{js/ai/chat\_render.js}     & renderizacao de Markdown, realce de codigo, \emph{links} de arquivo. \\
\file{js/ai/chat\_history.js}    & lista de conversas e serializacao. \\
\file{js/ai/system\_prompt.js}   & o \emph{system prompt} (constante imutavel). \\
\bottomrule
\end{tabularx}

\subsection{Comportamento da interface}

O painel mantem o historico em memoria, transmite cada turno via
\code{ai:chat-event} e renderiza o texto do assistente como Markdown ao vivo,
token a token. Pontos notaveis:

\begin{itemize}[leftmargin=1.4em]
  \item \textbf{Provedores sinteticos}: \enquote{Claude Code} e \enquote{ChatGPT} sao sempre listados (autenticacao por assinatura, sem chave); os provedores por chave so aparecem quando configurados. O modelo de cada um e persistido localmente.
  \item \textbf{Selecao a partir do editor}: \code{askAboutSelection(...)} abre o painel e semeia o compositor com um trecho destacado no Monaco, com intencoes pre-definidas (explicar, corrigir, melhorar, comentar, documentar).
  \item \textbf{Cartoes de atividade}: cada \code{tool-call} vira um \emph{chip} que se fecha quando chega o \code{tool-result} correspondente.
  \item \textbf{Rolagem inteligente}: a viewport so e empurrada quando o usuario ja esta no fim; rolar para cima libera a leitura, e uma pilula \enquote{Ir para o mais recente} reancorra.
  \item \textbf{\emph{Links} externos}: URLs vindas do modelo so abrem no navegador apos um aviso de redirecionamento (com opcao de confiar sempre); caminhos absolutos do sistema de arquivos abrem em Monaco, no Explorer ou no aplicativo padrao conforme o tipo.
  \item \textbf{\emph{Watchdog} de interface}: \code{STREAM\_STALL\_MS} e \code{STREAM\_STALL\_HARD\_MS} recuperam a UI para o estado ocioso se um turno travar.
\end{itemize}

\subsection{O \emph{system prompt} e o contexto de projeto}

\file{js/ai/system\_prompt.js} e uma string imutavel, montada a partir da
gramatica do compilador YANC, de exemplos reais de \cmm{}, do site do NIPSCERN e
do manifesto de ferramentas. Ela define a identidade do assistente
(\enquote{Aurora Intelligence}, sempre no feminino), o ecossistema SAPHO, a
linguagem \cmm{} e as regras de formatacao (sempre cercar codigo, citar
\lstinline|arquivo.ext:linha| para gerar \emph{links} clicaveis, responder no
idioma do usuario). A cada turno, \file{js/ai/chat\_turn.js}
(\code{buildProjectContext}) acrescenta o contexto do projeto ativo (caminho da
raiz e do arquivo \path{.spf}), poupando uma chamada de ferramenta e evitando que
o modelo alucine caminhos de projetos anteriores.

\section{Superficie IPC}

Todos os canais sao registrados em \file{main/ipc/ai.js}. Resumo:

\begin{longtable}{l Y}
\caption{Canais IPC do Aurora Intelligence.}\\
\toprule
\textbf{Canal} & \textbf{Funcao} \\
\midrule
\endfirsthead
\toprule
\textbf{Canal} & \textbf{Funcao} \\
\midrule
\endhead
\bottomrule
\endfoot
\code{ai:list-providers}    & lista provedores, modelo padrao e modelo efetivo. \\
\code{ai:get-key-status}    & quais provedores tem chave (apenas booleanos). \\
\code{ai:set-key} / \code{ai:clear-key} & grava / remove uma chave no cofre. \\
\code{ai:set-model}         & persiste a sobreposicao de modelo do provedor. \\
\code{ai:test-connection}   & testa a chave com uma chamada minima. \\
\code{ai:generate-oneshot}  & geracao de um disparo (SDK ou Claude Code). \\
\code{ai:claude-code-status} / \code{ai:claude-code-usage} & estado e uso da CLI Claude Code. \\
\code{ai:codex-status} / \code{ai:codex-usage} & estado e uso da CLI Codex. \\
\code{ai:chat-start}        & inicia um turno (roteia ao transporte correto). \\
\code{ai:chat-abort}        & aborta um turno em qualquer um dos transportes. \\
\code{ai:chat-event}        & (main $\rightarrow$ renderer) eventos de \emph{streaming}. \\
\code{ai:get-tool-manifest} & manifesto de ferramentas para o \code{tool\_runner}. \\
\code{ai:tool-exec}         & (main $\rightarrow$ renderer) pedido de execucao de ferramenta. \\
\code{ai:tool-result}       & (renderer $\rightarrow$ main) resposta da ferramenta. \\
\code{ai:audit-override-applied} & registro de sobreposicao de comando aplicada. \\
\code{ai:conv-*}            & lista, le, salva, apaga, renomeia conversas. \\
\bottomrule
\end{longtable}

O \code{ai:chat-start} ilustra o roteamento entre transportes: o provedor
\enquote{claude-code} vai para \file{claude\_code.js}, \enquote{chatgpt} para
\file{codex\_cli.js}, e todo o resto para o laco do Vercel AI SDK em
\file{chat.js}.

\section{Status do modelo proprio do laboratorio}

\label{sec:ia-claude-oneshot}
O Aurora Intelligence esta \textbf{totalmente implementado}: os dois transportes,
o servidor MCP, as 98 ferramentas, a ponte de ferramentas, o cofre criptografado,
a persistencia e o painel estao todos operacionais. A governanca de modelos
(\code{DEFAULT\_MODELS}, \code{MODEL\_MIGRATIONS}, resolucao e recuo) e factual e
ativa para os provedores externos hoje disponiveis.

\begin{nota}
\textbf{Item de roteiro --- modelo proprio do NIPSCERN.} O \emph{provider} e os
modelos \textbf{proprios} do laboratorio ainda \textbf{nao foram treinados}. Sao
um item planejado de roteiro. Ate la, o assistente funciona normalmente contra os
provedores externos descritos neste capitulo (provedores com chave de API e CLIs
por assinatura), de modo que toda a funcionalidade ja esta disponivel ao usuario.
\end{nota}

% !TEX root = ../main.tex
\chapter{AURORA: ferramentas de produtividade}
\label{cap:13-aurora-produtividade}

A \tool{AURORA} reúne, em torno do editor \tool{Monaco}, um conjunto de
recursos auxiliares que aproximam a IDE da experiência de um ambiente
profissional de desenvolvimento: um terminal de saída rico, controle de versão
embutido, busca global no projeto, servidores de linguagem (LSP) para
\textit{Verilog}/\textit{SystemVerilog}, realce sintático por
\textit{tree-sitter}, formatação de código, paleta de comandos,
internacionalização, painel de configurações, atualizador automático e um
\textit{overlay} de diagnóstico de desempenho para desenvolvimento.

Este capítulo descreve cada um desses recursos a partir do código-fonte: o que
faz, qual ferramenta o sustenta e como está ligado entre o processo principal
(\textit{main}) do \tool{Electron} e o processo de renderização (\textit{renderer}).
O fluxo de compilação propriamente dito é tratado em outros capítulos; aqui o
foco é na infraestrutura de produtividade que cerca a edição.

\begin{nota}
Um padrão recorrente atravessa todos os recursos: \textbf{best-effort}. Cada
módulo verifica se o binário ou recurso de que depende está presente e, quando
não está, vira um \textit{no-op} silencioso — o editor continua funcionando
exatamente como antes, sem erro visível ao usuário. Isso aparece em
\file{verible\_lsp.js}, \file{slang\_lsp.js}, \file{clang\_format.js},
\file{grammars.js} e no atualizador.
\end{nota}

\section{Arquitetura de ligação: \textit{main} $\leftrightarrow$ \textit{renderer}}
\label{sec:13-arquitetura-ligacao}

Quase todos os recursos deste capítulo seguem a mesma topologia em três
camadas:

\begin{enumerate}[leftmargin=*]
  \item Um módulo no processo \textit{main} (sob \path{aurora/main/}) que faz o
        trabalho privilegiado: gerar processos filhos, ler o sistema de
        arquivos, falar com a API do GitHub. Ele registra \textit{handlers} de
        IPC via \lstinline|ipcMain.handle(...)|.
  \item Uma ponte exposta no \file{preload.js} (em \path{js/app/preload.js})
        via \lstinline|contextBridge.exposeInMainWorld(...)|, que publica um
        objeto no \lstinline|window| do \textit{renderer} (por exemplo
        \lstinline|window.lspAPI|, \lstinline|window.gitAPI|,
        \lstinline|window.treeSitterAPI|, \lstinline|window.clangFormatAPI| e
        \lstinline|window.slangAPI|). Algumas operações não têm ponte própria e
        são agregadas em \lstinline|window.electronAPI| (por exemplo a busca, via
        \lstinline|electronAPI.searchInProject|).
  \item Um módulo no \textit{renderer} (sob \path{aurora/js/}) que consome essa
        ponte e a integra ao \tool{Monaco} ou à interface.
\end{enumerate}

A \autoref{tab:13-mapa-modulos} resume o mapeamento entre recurso, módulo
\textit{main}, ponte e módulo \textit{renderer}.

\begin{table}[h]
\centering
\caption{Mapa de módulos das ferramentas de produtividade.}
\label{tab:13-mapa-modulos}
\small
\begin{tabularx}{\linewidth}{Y Y Y}
\toprule
\textbf{Recurso} & \textbf{Módulo \textit{main}} & \textbf{Renderer / ponte} \\
\midrule
Terminal de saída    & ---                                & \file{js/terminal/terminal\_module.js} \\
Source control       & \file{main/ipc/git.js}             & \file{js/git/git\_panel.js} (\lstinline|gitAPI|) \\
Auth GitHub          & \file{main/ipc/github\_auth.js}    & \file{js/git/git\_panel.js} \\
Busca no projeto     & \file{main/ipc/search.js}          & \file{js/search/search\_panel.js} \\
LSP Verilog          & \file{main/lsp/verible\_lsp.js}    & \file{js/editor/lsp\_integration.js} (\lstinline|lspAPI|) \\
LSP semântico        & \file{main/lsp/slang\_lsp.js}      & \file{js/editor/slang\_integration.js} (\lstinline|slangAPI|) \\
Realce tree-sitter   & \file{main/treesitter/grammars.js} & \file{js/editor/treesitter\_highlight.js} (\lstinline|treeSitterAPI|) \\
Formatação C/C++     & \file{main/format/clang\_format.js}& \file{js/editor/clang\_format\_integration.js} \\
Paleta de comandos   & ---                                & \file{js/ui/command\_palette.js} \\
i18n                 & ---                                & \file{js/i18n/i18n.js} + \path{locales/} \\
Configurações        & ---                                & \file{js/processors/aurora\_settings.js} \\
Atualizador          & \file{main/updater.js}             & \path{html/update-notification.html} \\
Overlay de jank      & ---                                & \file{js/dev/jank\_overlay.js} \\
\bottomrule
\end{tabularx}
\end{table}

\section{Terminal de saída}
\label{sec:13-terminal}

O painel inferior da \tool{AURORA} é um console de \textit{scrollback} rico,
implementado pela classe \lstinline|TerminalManager| em
\file{js/terminal/terminal\_module.js}. Não é um terminal interativo (não há
\textit{shell}); é uma superfície de \textbf{saída} estruturada que recebe e
classifica as linhas emitidas pelas ferramentas do \textit{toolchain}.

\subsection{Abas e troca de terminal}

O painel é dividido em abas, cada uma com seu próprio corpo de rolagem,
mapeadas no construtor de \lstinline|TerminalManager|:

\begin{table}[h]
\centering
\caption{Abas do terminal de saída.}
\label{tab:13-terminal-abas}
\small
\begin{tabularx}{\linewidth}{Y Y}
\toprule
\textbf{Identificador} & \textbf{Conteúdo} \\
\midrule
\lstinline|tcmm|   & Compilação \cmm \\
\lstinline|tasm|   & Montagem (\textit{assembly}) \\
\lstinline|tveri|  & Síntese \textit{Verilog} / \tool{PRISM} \\
\lstinline|twave|  & Geração de \textit{waveform} \\
\lstinline|thtest| & \textit{Terminal Hardware Test} — etapas e barra de progresso \\
\lstinline|tcmd|   & Saída de comandos genéricos \\
\bottomrule
\end{tabularx}
\end{table}

A troca de aba é feita por \file{js/terminal/terminal.js}, que exporta a função
\lstinline|switchTerminal(targetId)|. Essa é a única implementação da troca
(consolidada de cópias que existiam em \file{compilation\_flow.js} e
\file{wave\_config\_manager.js}). Ela esconde todos os corpos
(\lstinline|.terminal-content|), revela o alvo e desliza um único indicador de
aba ativa (\lstinline|.terminal-tab-indicator|) via \lstinline|transform:
translateX(...)|, em vez de ligar/desligar uma barra por aba. Os botões de
compilação focam automaticamente a aba onde sua saída cai (por exemplo,
\lstinline|cmmcomp| foca \lstinline|terminal-tcmm|).

\subsection{Cartões por tipo de mensagem}

A saída não é texto cru: cada mensagem é classificada e agrupada em
\textbf{cartões} (\lstinline|.log-entry|). O \lstinline|TerminalManager| mantém
contadores por terminal nas categorias \lstinline|error|, \lstinline|warning|,
\lstinline|success| e \lstinline|tips|, exibidos como \textit{badges} sobre os
botões de filtro (\lstinline|createCounterBadges| /
\lstinline|updateCounterDisplay|).

Mensagens do mesmo tipo emitidas na mesma sessão são consolidadas num único
cartão agrupado (\lstinline|createGroupedCard| / \lstinline|addMessageToCard|):
o cartão recebe um \textit{timestamp} em formato \lstinline|pt-BR| e cada nova
linha vira um \lstinline|.grouped-message| dentro do contêiner. Botões de
filtro permitem mostrar/ocultar categorias (\lstinline|_cardMatchesFilter|,
\lstinline|applyFilter|); o filtro trata \lstinline|tips| e \lstinline|info|
como equivalentes. Há ainda um modo \textit{verbose} (persistido em
\lstinline|localStorage| sob \lstinline|terminal-verbose-mode|) que controla a
exibição das linhas \lstinline|plain| (não classificadas).

\begin{nota}
Para manter o consumo de memória e o desempenho de rolagem limitados num
\textit{build} que despeja milhares de linhas, há dois tetos: no máximo
\lstinline|MAX_TERMINAL_ENTRIES| (5000) nós \lstinline|.log-entry| por terminal
e no máximo \lstinline|MAX_GROUPED_MESSAGES| (5000) sub-mensagens por cartão
agrupado. Os mais antigos são descartados do topo, e os contadores são
recalculados a partir do DOM já podado.
\end{nota}

\subsection{Números de linha clicáveis}

O método \lstinline|makeLineNumbersClickable(text)| transforma referências de
diagnóstico em \textit{links} navegáveis. Há dois formatos reconhecidos:

\begin{description}[leftmargin=*]
  \item[Estilo \textit{toolchain} C] padrão \lstinline|<arquivo>:<linha>:|
        emitido por \tool{iverilog}, \tool{yosys}, \tool{gcc} e pelo
        \textit{appcomp}. A expressão regular captura o caminho (inclusive um
        prefixo de unidade do Windows como \lstinline|C:\foo\bar.v:15:|) e a
        linha; o caminho é guardado em \lstinline|data-file| para que o clique
        abra exatamente aquele arquivo.
  \item[Estilo \tool{AURORA}/\tool{YANC}] padrão \enquote{\lstinline|linha N|}
        ou \enquote{\lstinline|line N|}, sem prefixo de arquivo. Nesse caso o
        clique recai sobre o último \cmm{} compilado, resolvido pelo gerenciador
        de compilação.
\end{description}

Cada \textit{link} (\lstinline|.line-link|) tem seu clique conectado por
\lstinline|_attachLineLinkClicks|, e a navegação até a posição é feita por
\lstinline|goToLine(lineNumber)|, que comanda o \tool{Monaco} a posicionar e
revelar a linha-alvo (limitada ao número total de linhas do modelo). Há também
um \textit{link} de pasta (\lstinline|appendFolderLink|), construído com
\lstinline|textContent| (sem \lstinline|innerHTML|), que alterna a árvore de
arquivos para a visão de pastas e revela a pasta correspondente via
\lstinline|standardTreeRenderer.revealFolder|.

\subsection{Barra de progresso de teste de \textit{hardware}}

O fluxo de \textit{Terminal Hardware Test} (aba \lstinline|thtest|) usa
\lstinline|renderHardwareProgress(terminalId, p)| para exibir uma barra de
progresso real em DOM. Diferentemente dos cartões, esse é um \textbf{único}
elemento \lstinline|.hw-progress| que se atualiza no lugar: ele \textbf{não} é
um \lstinline|.log-entry|, portanto o filtro de \textit{verbose} e os
contadores o ignoram e ele permanece sempre visível. A cada atualização o nó é
reanexado ao fim do terminal, para acompanhar a última saída transmitida. O
preenchimento (\lstinline|.hw-progress-fill|) anima por
\lstinline|transform: scaleX(...)|, e a linha de metadados mostra
\lstinline|ciclo/total| mais o total de leituras de entrada (rotulado via
i18n por \lstinline|terminal.htest.reads|). O \textit{payload} \lstinline|p|
carrega \lstinline|pct|, \lstinline|cyc|, \lstinline|total|, \lstinline|reads|,
\lstinline|label| e \lstinline|done|.

\section{Controle de versão (Source Control) embutido}
\label{sec:13-git}

A \tool{AURORA} embute um painel de controle de versão completo. O lado
\textit{main} vive em \file{main/ipc/git.js} e é construído sobre a biblioteca
\textbf{simple-git}, um invólucro fino sobre o binário nativo \lstinline|git|.
Como ele dirige o \lstinline|git| real, \path{.gitignore}, \textit{diffs},
\textit{merges} e credenciais se comportam exatamente como na linha de comando.
A interface vive em \file{js/git/git\_panel.js}.

\subsection{Escopo das operações}

Todas as operações agem sobre o diretório do \textbf{projeto aberto}, derivado
de \lstinline|state.currentOpenProjectPath| (a fonte única de verdade): a função
\lstinline|projectDir()| devolve o diretório-pai do arquivo \path{.spf}. As
\textit{handlers} de leitura podem honrar um diretório explícito
(\lstinline|opts.dir|) para navegar um repositório recém-clonado que ainda não
tem \path{.spf} aberto; as mutações sempre agem sobre o projeto aberto. Toda
\textit{handler} é embrulhada por \lstinline|safe(fn)|, que converte exceções em
um objeto \lstinline|{ ok:false, error }| em vez de propagá-las pelo IPC.

\subsection{Inspeção e \textit{diffs}}

As \textit{handlers} de leitura cobrem o conjunto típico de inspeção:

\begin{longtable}{>{\raggedright\arraybackslash}p{0.34\linewidth} >{\raggedright\arraybackslash}p{0.6\linewidth}}
\caption{Handlers de inspeção do controle de versão.}\label{tab:13-git-read}\\
\toprule
\textbf{Canal IPC} & \textbf{Função} \\
\midrule
\endfirsthead
\toprule
\textbf{Canal IPC} & \textbf{Função} \\
\midrule
\endhead
\lstinline|git:is-repo|      & Verifica se o diretório é um repositório. \\
\lstinline|git:status|       & \textit{Status} com \textit{branch}, \textit{tracking}, \textit{ahead}/\textit{behind} e \texttt{flags} por arquivo; opcionalmente anexa contagens \texttt{+}/\texttt{-}. \\
\lstinline|git:diff|         & \textit{Diff} unificado de um arquivo ou de toda a árvore de trabalho. \\
\lstinline|git:commit-files| & Lista (apenas \textit{numstat}) os arquivos tocados por um \textit{commit}. \\
\lstinline|git:show|         & \textit{Diff} de um arquivo dentro de um \textit{commit} (carregado sob demanda). \\
\lstinline|git:log|          & Histórico de \textit{commits} (padrão: 50). \\
\lstinline|git:branches|     & \textit{Branches} locais e remotas (\lstinline|-a|), com remotas-só destacadas. \\
\lstinline|git:remotes|      & Lista de \textit{remotes}. \\
\lstinline|git:info|         & Nome derivado de \lstinline|origin| (\lstinline|owner/repo|) e URL. \\
\bottomrule
\end{longtable}

O texto dos \textit{diffs} é limitado por \lstinline|capDiff|: um \textit{diff}
de \textit{commit} pode ter megabytes (dados \path{.mif}/\path{.hex} gerados,
\textit{blobs} versionados), então a saída é cortada num limite de linha em
\lstinline|MAX_DIFF_BYTES| ($\approx$ 600 KB) e marcada como truncada. As
contagens por arquivo da árvore de trabalho (\lstinline|workNumstat|) combinam
dois \textit{diffs} \lstinline|--numstat| (\textit{staged} + \textit{unstaged}).
Na interface, os \textit{diffs} são renderizados por \textbf{diff2html}
(importado em \file{js/git/git\_panel.js} como \lstinline|html as renderDiff|),
com a CSS \path{diff2html/bundles/css/diff2html.min.css}, em modo
\lstinline|line-by-line| e \lstinline|colorScheme: 'dark'|. Há um teto adicional
de linhas no lado do \textit{renderer} antes de entregar o conteúdo ao
\lstinline|diff2html|.

\subsection{Mutações e operações remotas}

As mutações cobrem o ciclo completo de trabalho: \lstinline|git:init|,
\lstinline|git:stage| / \lstinline|git:stage-all| / \lstinline|git:unstage|,
\lstinline|git:discard|, \lstinline|git:commit| (com \lstinline|--amend|
opcional), \lstinline|git:undo-last-commit| (\textit{reset} \textit{soft}),
\lstinline|git:checkout| (incluindo criar \textit{branch} local e rastrear uma
remota com \lstinline|--track|), \lstinline|git:merge| (\lstinline|--no-edit|),
e o conjunto de \textit{stash} (\lstinline|git:stash|, \lstinline|git:stash-list|,
\lstinline|git:stash-pop|, \lstinline|git:stash-drop|), de modo que trocar de
\textit{branch} com a árvore suja simplesmente funciona.

As operações remotas — \lstinline|git:fetch|, \lstinline|git:pull|,
\lstinline|git:push| — usam o invólucro \lstinline|remoteGit()|. O
\lstinline|pull| roda com \lstinline|--no-edit --autostash| (nenhum
\lstinline|$EDITOR| abre e mudanças locais não impedem o \textit{merge}); o
\lstinline|push| define \textit{upstream} (\lstinline|-u origin <branch>|)
apenas quando ainda não há rastreamento. O \lstinline|git:clone| transmite
progresso ao vivo: o \textit{plugin} de progresso do \textbf{simple-git} analisa
o \texttt{stderr} do \lstinline|git --progress| e cada \textit{tick} é
encaminhado ao \textit{renderer} pelo canal \lstinline|git:clone-progress|, que
move a barra de progresso do painel. Após um \textit{clone}, \lstinline|git:scan-spf|
varre o diretório (profundidade $\le$ 5) em busca de arquivos \path{.spf} para
abrir o projeto recém-clonado.

\subsection{Autenticação GitHub}
\label{sec:13-github-auth}

O \enquote{conectar sua conta GitHub} vive em \file{main/ipc/github\_auth.js} e
oferece dois caminhos de autenticação:

\begin{enumerate}[leftmargin=*]
  \item \textbf{\textit{Personal Access Token} (PAT)} — o usuário cola um
        \textit{token} (clássico ou \textit{fine-grained} com escopo
        \lstinline|repo|); ele é validado contra \lstinline|/user| da API do
        GitHub e então armazenado.
  \item \textbf{OAuth \textit{Device Flow}} — \lstinline|deviceFlowLogin| pede
        um \textit{device}/\textit{user code}, abre a página de verificação no
        navegador (\lstinline|shell.openExternal|) e faz \textit{polling} até a
        autorização. O \textit{Client ID} é público (o \textit{device flow} não
        usa segredo) e os escopos são \lstinline|repo read:org|. Ao autorizar, a
        janela da \tool{AURORA} é trazida de volta ao primeiro plano.
\end{enumerate}

Em ambos os casos o \textit{token} é cifrado com
\lstinline|safeStorage.encryptString| (DPAPI no Windows) e gravado em
\path{userData/aurora-github.json}; o texto plano nunca toca o disco. O
\textit{token} é lido somente no lado \textit{main} — o \file{git.js} o injeta
como cabeçalho \lstinline|http.extraHeader| de uso único (\lstinline|-c|, nunca
escrito na configuração do repositório) para \textit{push}/\textit{pull}/
\textit{clone}; o \textit{renderer} só consegue perguntar \enquote{quem está
conectado?} (\lstinline|github:status|), nunca obter os \textit{bytes}. Outras
\textit{handlers}: \lstinline|github:create-repo| (cria repositório sob a conta),
\lstinline|github:list-repos| (lista repositórios próprios, de colaboração e de
organização, paginando) e \lstinline|github:disconnect|. O \textit{avatar} é
baixado e convertido em \textit{data URL} para passar pela CSP do
\textit{renderer} sem afrouxá-la.

\section{Busca no projeto (Find in Files)}
\label{sec:13-busca}

A busca global, equivalente ao painel \textit{Search} do VS Code, é
implementada em \file{main/ipc/search.js} e exposta pelo canal
\lstinline|search:in-project|. O \textit{renderer} é \file{js/search/search\_panel.js},
que agrupa os resultados por arquivo (cabeçalho recolhível + linhas de
ocorrência) e abre o arquivo na linha clicada via \lstinline|TabManager|.

\begin{nota}
O motor de busca \textbf{não é o ripgrep nem um binário externo}. É uma
varredura recursiva síncrona e \textbf{sem dependências}, escrita diretamente
sobre o módulo \lstinline|fs| do Node, enraizada no diretório do projeto aberto.
\end{nota}

\subsection{Consulta e modificadores}

O \textit{payload} aceita \lstinline|{ query, caseSensitive?, wholeWord?, regex? }|.
A função \lstinline|buildRegex| monta a expressão regular: em modo literal o
texto é escapado por \lstinline|escapeRegExp|; \lstinline|wholeWord| envolve o
corpo em \lstinline|\b...\b|; \lstinline|caseSensitive| controla a \textit{flag}
\lstinline|i|. Um padrão inválido em modo \textit{regex} é convertido em
\lstinline|{ ok:false, error }|. Os modificadores de caixa/palavra/\textit{regex}
são persistidos no \textit{renderer} sob \lstinline|aurora.search.toggles|.

\subsection{Limites da varredura}

Como a varredura é síncrona, há limites rígidos que impedem que uma árvore
grande congele o processo principal:

\begin{table}[h]
\centering
\caption{Limites da busca em \file{main/ipc/search.js}.}
\label{tab:13-busca-limites}
\small
\begin{tabularx}{\linewidth}{Y Y}
\toprule
\textbf{Constante} & \textbf{Valor / efeito} \\
\midrule
\lstinline|MAX_DEPTH|         & 12 níveis de profundidade. \\
\lstinline|MAX_FILE_BYTES|    & $\approx$ 1,5 MB — arquivos maiores são ignorados. \\
\lstinline|BINARY_SNIFF_BYTES|& 4 KB lidos para detectar byte \texttt{NUL} (binário). \\
\lstinline|PREVIEW_MAX|       & 240 caracteres por linha de \textit{preview}. \\
\lstinline|MAX_MATCHES|       & 2000 ocorrências no total. \\
\lstinline|MAX_FILES|         & 500 arquivos com ocorrência. \\
\bottomrule
\end{tabularx}
\end{table}

Diretórios de VCS, dependências e \textit{build} nunca são percorridos
(\lstinline|SKIP_DIRS|: \path{.git}, \path{node_modules}, \path{dist},
\path{build}, \path{components}, \path{.vite} entre outros). Arquivos binários
são detectados por uma heurística simples (\lstinline|looksBinary|: byte
\texttt{NUL} nos primeiros 4 KB). Ao atingir qualquer teto, a varredura para e o
resultado é marcado com \lstinline|truncated: true|. Cada ocorrência devolve
\lstinline|{ line, col, preview }|, e o caminho é relativo à raiz do projeto.

\section{Servidores de linguagem (LSP) para Verilog}
\label{sec:13-lsp}

A \tool{AURORA} integra dois servidores de linguagem distintos para
\textit{Verilog}/\textit{SystemVerilog}, com responsabilidades deliberadamente
divididas. Ambos falam JSON-RPC enquadrado por \lstinline|Content-Length| sobre
\textit{stdio}, por uma ponte de transporte escrita à mão — a \tool{AURORA} não
usa \textit{monaco-languageclient}; seu IPC já é próprio e um único \textit{backend}
de \textit{Verilog} precisa apenas de um transporte fino.

\subsection{Verible: o servidor sintático completo}

\file{main/lsp/verible\_lsp.js} gera um único processo de longa duração
\tool{verible-verilog-ls} (empacotado em
\path{components/Packages/verible/bin}), iniciado com
\lstinline|--lsp_enable_hover| e \lstinline|--rules_config_search| (para que o
\tool{Verible} procure, subindo a árvore, um \path{.rules.verible_lint} por
projeto). É o servidor \textbf{completo}: o \textit{renderer}
(\file{js/editor/lsp\_integration.js}) registra no \tool{Monaco} todo o conjunto
de provedores:

\begin{itemize}[leftmargin=*]
  \item \textbf{diagnósticos ao vivo} (\textit{lint} + sintaxe) como
        \textit{squiggles} e marcadores no painel \textit{Problems}, sob o dono
        \lstinline|verible|;
  \item \textbf{formatação} de documento (\textit{Format Document} /
        \lstinline|Shift+Alt+F|);
  \item \textbf{símbolos de \textit{outline}} (\textit{breadcrumbs} e visão
        \textit{Outline}), via \lstinline|provideDocumentSymbols|;
  \item \textbf{hover};
  \item \textbf{ir para definição} e \textbf{localizar referências}.
\end{itemize}

\begin{nota}
O \textbf{\textit{outline}} e o \textit{folding} simbólico vêm dos
\lstinline|documentSymbols| do \tool{Verible}, e \textbf{não} do
\textit{tree-sitter} (que é exclusivamente realce — ver
\autoref{sec:13-treesitter}). O \file{lsp\_integration.js} converte as
coordenadas LSP (base 0) para as do \tool{Monaco} (base 1) e mapeia
\lstinline|SymbolKind|/\lstinline|MarkerSeverity| entre os dois esquemas.
\end{nota}

O ciclo de vida é dirigido no nível do \textbf{modelo} do \tool{Monaco}: todo
modelo \path{.v}/\path{.sv} que aparece é \lstinline|didOpen|; suas edições são
debounced (350 ms) em \lstinline|didChange|; e o descarte dispara
\lstinline|didClose|. A sincronização é de documento inteiro (uma única mudança
sem \lstinline|range|). O \textit{bridge} é resiliente: se o servidor morre, ele
é respawnado na próxima requisição e os documentos abertos são re-enviados
(\lstinline|didOpen|) de forma transparente, de modo que os diagnósticos
continuam fluindo sem o \textit{renderer} perceber. Cada requisição
(\textit{format}/\textit{symbols}/\textit{hover}/...) tem \textit{timeout}
(\lstinline|REQUEST_TIMEOUT_MS| = 15 s) e, ao falhar, devolve um \textit{fallback}
vazio em vez de travar o editor.

\subsection{slang: o servidor semântico}

\file{main/lsp/slang\_lsp.js} gera o \tool{slang-server} (de
\textit{hudson-trading/slang-server}, em
\path{components/Packages/slang-server/bin}). Diferentemente do \tool{Verible}
(sintático e por arquivo), o \tool{slang} \textbf{elabora o projeto inteiro},
capturando erros semânticos que o \tool{Verible} não vê: identificadores não
declarados, incompatibilidades de tipo/porta, sinais não usados etc.

Pela divisão escolhida (\enquote{meio-termo}), o \textit{renderer}
(\file{js/editor/slang\_integration.js}) consome do \tool{slang} \textbf{apenas}:

\begin{itemize}[leftmargin=*]
  \item \textbf{diagnósticos semânticos} (sob o dono \lstinline|slang|,
        coexistindo com os marcadores \lstinline|verible|); e
  \item \textbf{completação} de símbolos (módulos, portas, membros de
        \textit{package}/\textit{struct}, macros), com caracteres de gatilho
        \lstinline| ` # . ( : [ |.
\end{itemize}

\textit{Hover}, definição, referências, \textit{outline} e formatação
permanecem com o \tool{Verible}. Como o \tool{slang} é acoplado ao
\textit{workspace} (indexa a árvore do projeto), o \textit{bridge} o inicia com
o diretório do projeto como \lstinline|rootUri| e o reinicia de forma
transparente quando o projeto muda (\lstinline|maybeRestartForProject|). Por
elaborar a cada mudança e poder ser ruidoso em projetos incompletos, o
\tool{slang} é \textbf{ligável/desligável}: \lstinline|window.AuroraSlang.toggle()|,
com estado persistido em \lstinline|localStorage| (\lstinline|slangEnabled|) e
sincronizado ao \textit{main}; há uma entrada na paleta de comandos para
alterná-lo. Ao desligar, os marcadores \lstinline|slang| são limpos e nenhuma
requisição é enviada.

\section{Realce sintático por tree-sitter}
\label{sec:13-treesitter}

O realce de sintaxe preciso é fornecido por \textit{tree-sitter} (via
\textbf{web-tree-sitter}, WASM), implementado em
\file{js/editor/treesitter\_highlight.js}. O \textit{parser} processa o
\textit{buffer}, a \textit{query} de \textit{highlights} da gramática produz
\textit{captures}, e o módulo os mapeia para \textbf{\textit{semantic tokens}}
do \tool{Monaco}, que sobrepõem o realce base (\textit{Monarch}) com cores
fiéis à gramática: nomes de módulo, instâncias, direção de portas, tipos,
macros etc.

\begin{nota}
O \textit{tree-sitter} na \tool{AURORA} faz \textbf{apenas realce} (semantic
tokens). Não provê \textit{folding} nem \textit{outline} — isso vem do LSP
(\autoref{sec:13-lsp}). As gramáticas cobertas são
\textit{Verilog}/\textit{SystemVerilog}, C e C++; \cmm{} e \textit{assembly}
não têm gramática e mantêm seus \textit{tokenizers} \textit{Monarch}.
\end{nota}

\subsection{Entrega das gramáticas WASM}

O \textit{renderer}, sob o \textit{sandbox} \lstinline|file://|, não consegue
buscar \path{.wasm} de forma confiável por URL. Por isso \file{main/treesitter/grammars.js}
serve os \textbf{bytes} crus de cada artefato pelo canal
\lstinline|treesitter:wasm|, e o \textit{renderer} os alimenta diretamente em
\lstinline|Parser.init({ wasmBinary })| e \lstinline|Language.load(bytes)|,
contornando todos os problemas de URL/\textit{fetch}/CSP. Apenas um conjunto fixo
de nomes lógicos é aceito (o \textit{renderer} não pode pedir caminhos
arbitrários):

\begin{table}[h]
\centering
\caption{Artefatos tree-sitter servidos por \file{grammars.js}.}
\label{tab:13-ts-artefatos}
\small
\begin{tabularx}{\linewidth}{Y Y}
\toprule
\textbf{Nome lógico} & \textbf{Arquivo} \\
\midrule
\lstinline|runtime|       & \path{web-tree-sitter.wasm} \\
\lstinline|systemverilog| & \path{tree-sitter-systemverilog.wasm} \\
\lstinline|c|             & \path{tree-sitter-c.wasm} \\
\lstinline|cpp|           & \path{tree-sitter-cpp.wasm} \\
\bottomrule
\end{tabularx}
\end{table}

Os artefatos ficam em \path{components/Packages/tree-sitter/} (baixados por
\file{download-tree-sitter-grammars.js} no \textit{bootstrap}). O canal
\lstinline|treesitter:status| informa se o \textit{runtime} e ao menos uma
gramática estão presentes.

\subsection{Detalhes de mapeamento}

As \textit{queries} (\path{js/editor/treesitter/queries/}) são incorporadas em
tempo de \textit{build} (\lstinline|?raw|); \textit{cpp} herda as \textit{queries}
de \textit{c}. Os nomes de \textit{capture} são mapeados para tipos de
\textit{token} semântico padrão do VS Code (\lstinline|CAPTURE_MAP|), de modo que
o tema ativo os colore sem configuração extra. Como as posições do
\textit{tree-sitter} são \textit{offsets} em \textbf{bytes} (UTF-8) e o
\tool{Monaco} usa colunas UTF-16, há conversão por linha (com caminho rápido
para ASCII), garantindo que comentários/\textit{strings} acentuados (por exemplo
em português) não desloquem o realce subsequente. Tudo é \textit{best-effort}:
se o \textit{runtime} ou os \textit{bytes} faltarem, o provedor devolve zero
\textit{tokens} e o realce \textit{Monarch} permanece, sem regressão. O
carregamento é \textbf{preguiçoso}: uma gramática só é buscada e compilada na
primeira vez que um \textit{buffer} daquela linguagem é realçado.

\section{Formatação com clang-format}
\label{sec:13-clang-format}

A formatação de C, C++ e \cmm{} é feita pelo \tool{clang-format} empacotado
(\path{components/Packages/clang-format/bin/clang-format.exe}), via
\file{main/format/clang\_format.js}. Diferentemente de um LSP, o
\tool{clang-format} é uma CLI de uma só passada: lê um \textit{buffer} no
\textit{stdin} e escreve o \textit{buffer} formatado no \textit{stdout}.

Por \lstinline|Shift+Alt+F|, o \textit{renderer}
(\file{js/editor/clang\_format\_integration.js}) registra um provedor de
formatação para \lstinline|c|, \lstinline|cpp| e \lstinline|cmm| e envia o
\textit{buffer} ao canal \lstinline|format:clang|. O \tool{Monaco} já despacha o
atalho para o provedor que casa com a linguagem do editor em foco — assim,
\textit{Verilog}/\textit{SystemVerilog} vão ao \tool{Verible} e
\lstinline|c|/\lstinline|cpp|/\lstinline|cmm| vão ao \tool{clang-format}. O
resultado é um \textbf{\textit{replace} de documento inteiro} (o
\tool{clang-format} devolve o \textit{buffer} completo, não um \textit{diff}).

O tratamento de linguagem usa \lstinline|-assume-filename|: o \cmm{} (um
subconjunto de C) é formatado com regras de C ao reivindicar um nome de arquivo
\path{.c}; C/C++ usam o caminho real, de modo que o \tool{clang-format} detecte o
dialeto \textbf{e} encontre, subindo a árvore, um \path{.clang-format} de
projeto (executado com \lstinline|-style=file -fallback-style=LLVM|). Há
\textit{timeout} de 15 s contra travamentos, verificação de \textit{allowlist}
de binário (\lstinline|isAllowed|, a mesma do executor do \textit{toolchain}) e
comportamento \textit{best-effort}: se o binário faltar ou houver erro,
\lstinline|format()| resolve \lstinline|null| e o \textit{buffer} fica
intocado.

\section{Paleta de comandos}
\label{sec:13-paleta}

A paleta de comandos (\file{js/ui/command\_palette.js}) é uma superfície única,
orientada a teclado, para as ações espalhadas pela barra de ferramentas e
menus. Abre com \lstinline|Ctrl/Cmd+Shift+K| (principal) ou
\lstinline|Ctrl/Cmd+Shift+P|; \lstinline|Esc| fecha, setas navegam,
\lstinline|Enter| executa. O \lstinline|Ctrl+K| puro é reservado para o painel
de IA e por isso não é vinculado. A captura é feita na fase de captura com
\lstinline|stopPropagation|, de forma que a paleta abra mesmo sobre um editor
\tool{Monaco} em foco (que liga \lstinline|Ctrl+Shift+K| a \enquote{apagar
linha}) sem disparar aquele comando.

A \textbf{visão} é o componente Lit \lstinline|<aurora-command-palette>| (Shadow
DOM + \textit{tokens} semânticos); o módulo é o dono do \textbf{registro}
(dados puros), da \textbf{pontuação} \textit{fuzzy} e da lógica de execução. Os
comandos preferem a API pública (\lstinline|window.AuroraAPI| / controlador da
árvore de arquivos) e, na ausência dela, clicam o botão de barra existente por
\lstinline|id| — assim um comando faz exatamente o que o botão faz (inclusive
ser \textit{no-op} quando o botão está desabilitado), sem lógica duplicada.

O registro \lstinline|COMMANDS| é agrupado em \lstinline|Compile|,
\lstinline|Project|, \lstinline|View|, \lstinline|Tools| e \lstinline|Dev| (a
ordem de exibição segue \lstinline|GROUP_ORDER|). Inclui, entre outros:
compilar \cmm, sintetizar \textit{Verilog}, gerar \textit{waveform}, execução
rápida (\tool{Verilator}), abrir o \tool{PRISM}, novo/abrir projeto,
configurações da \tool{AURORA}, alternar o \tool{slang} e alternar o
\textit{overlay} de \textit{jank} (este último por \textit{import} dinâmico, sem
custo de \textit{boot}). A função \lstinline|scoreCommand| pontua por
subsequência: cada termo da consulta precisa aparecer no \textit{haystack}
(título + \textit{keywords} + grupo); acertos no título (e prefixos) valem mais
que acertos em \textit{keywords}.

\section{Internacionalização (i18n)}
\label{sec:13-i18n}

A camada de internacionalização vive em \file{js/i18n/i18n.js}, com dois
catálogos em \path{locales/}: \path{en.json} e \path{pt.json}. São dois
\textit{locales} suportados — inglês e português do Brasil — e o padrão é
\lstinline|pt|. O DOM declara associações por dois atributos principais:

\begin{itemize}[leftmargin=*]
  \item \lstinline|data-i18n="caminho.da.chave"| $\rightarrow$
        \lstinline|element.textContent|;
  \item \lstinline|data-i18n-tooltip| $\rightarrow$ \lstinline|data-tooltip|;
        \lstinline|data-i18n-placeholder| $\rightarrow$ \lstinline|placeholder|;
        \lstinline|data-i18n-aria-label| $\rightarrow$ \lstinline|aria-label|;
        \lstinline|data-i18n-title| $\rightarrow$ \lstinline|title|.
\end{itemize}

A função \lstinline|applyDOM()| percorre o documento no \textit{boot} e a cada
troca de \textit{locale}, resolvendo cada associação por \lstinline|t(key)|. A
mesma \lstinline|t()| é exposta como \lstinline|window.t| para chamadas em JS
(mensagens do terminal, UI dinâmica). Há ainda um preenchimento automático de
\textit{tooltip} por \lstinline|id|: qualquer elemento com \lstinline|id="X"|
recebe \lstinline|data-tooltip| de \lstinline|tooltip.extended.X| quando a chave
existe.

\subsection{Resolução, \textit{fallback} e interpolação}

O \textit{locale} ativo é guardado em \lstinline|localStorage| sob
\lstinline|aurora-locale| (a chave legada \lstinline|aurora-yanc-lang| é migrada
silenciosamente na primeira leitura). A resolução tenta o \textit{locale} ativo;
se a chave faltar, recai em \lstinline|en|; se ainda faltar, devolve o próprio
caminho da chave — uma falha \enquote{barulhenta} que evidencia strings sem
tradução durante o desenvolvimento. A interpolação resolve marcadores
\lstinline|{{nome}}| no \textit{template}; parâmetros ausentes viram string
vazia (sem exceção).

Cada troca de \textit{locale} dispara o evento
\lstinline|aurora:locale-changed|, ao qual componentes que renderizam strings
imperativamente (terminal, notificações) reagem para re-renderizar. Por
compatibilidade com o compilador \tool{YANC}, \lstinline|window.getYancLang| é
um \textit{alias} de \lstinline|getLocale()|, de modo que a injeção do
\textit{flag} \lstinline|-pt|/\lstinline|-en| do compilador segue inalterada — a
preferência de idioma é, intencionalmente, unificada entre UI e compilador. O
controle visual é o \lstinline|<select>| de idioma em
\textbf{Configurações $\rightarrow$ Idioma}, ligado por
\file{js/ui/locale\_toggle.js}.

\section{Painel de configurações}
\label{sec:13-config}

O painel de configurações é montado por \file{js/processors/aurora\_settings.js},
acionado pelo botão \lstinline|aurora-settings| (e pelo comando homônimo na
paleta). Ele abre um modal (\lstinline|settings-modal|) com os controles de
preferências. O estado é persistido em \lstinline|localStorage|:

\begin{itemize}[leftmargin=*]
  \item \lstinline|aurora-settings| — as preferências gerais (por exemplo,
        \textit{tooltips} habilitados);
  \item \lstinline|aurora-shortcuts| — os atalhos de teclado personalizados;
  \item \lstinline|aurora-ai-trust-external-links| — chave compartilhada com a
        caixa de \enquote{confiar em links externos} do chat de IA, de modo que
        as duas fiquem ligadas.
\end{itemize}

Os atalhos padrão (\lstinline|defaultShortcuts|) cobrem ações como novo arquivo
(\lstinline|Ctrl+N|), compilar tudo (\lstinline|Ctrl+Shift+B|), fechar/reabrir
aba, salvar/salvar tudo e abrir configurações. Os rótulos dos atalhos são
resolvidos em tempo de renderização via chaves de i18n
(\lstinline|SHORTCUT_LABEL_KEYS|), não armazenados — assim uma troca de
\textit{locale} atualiza a UI sem tocar nos atalhos persistidos. O painel também
expõe o \textit{select} de idioma (\autoref{sec:13-i18n}) e o controle de
\textit{tooltips} (aplicado imediatamente por \lstinline|setTooltipsEnabled|).

\section{Atualizador automático (do ponto de vista do usuário)}
\label{sec:13-updater}

O \textit{auto-updater} vive em \file{main/updater.js} e é construído sobre
\textbf{electron-updater}, consultando os \textit{releases} do repositório
\repo{sapho} no GitHub. Os detalhes de \textit{release} e empacotamento são
tratados em \autoref{cap:16-release}; aqui descreve-se a experiência do usuário.

O fluxo (somente em produção — pulado em modo de desenvolvimento) é:

\begin{enumerate}[leftmargin=*]
  \item Cerca de 6 segundos após a janela principal aparecer, a \tool{AURORA}
        checa silenciosamente por atualizações
        (\lstinline|autoUpdater.autoDownload = false|).
  \item Em \lstinline|update-available|, abre-se uma janela de atualização
        personalizada (\path{html/update-notification.html}) mostrando o
        \textit{changelog} bilíngue e a opção de baixar — nunca um diálogo
        nativo. Se o \textit{electron-updater} não trouxer as notas, elas são
        buscadas diretamente na API do GitHub (\lstinline|fetchReleaseNotes|).
  \item Ao iniciar o \textit{download}, o evento \lstinline|download-progress|
        alimenta uma barra de progresso real (percentual, MB transferidos,
        velocidade, ETA) na mesma janela.
  \item Em \lstinline|update-downloaded|, a janela muda para o estado
        \enquote{Reiniciar e instalar}.
  \item \lstinline|quitAndInstall(false, true)| instala e relança diretamente
        na nova versão.
\end{enumerate}

Uma verificação manual (\enquote{checar atualizações}, via menu/configurações/
\textit{About}) usa \lstinline|checkForUpdates(true)| e, se nada for encontrado,
exibe um diálogo \enquote{você está atualizado}; a checagem silenciosa de
inicialização suprime esse aviso. Erros são traduzidos para mensagens curtas e
legíveis por \lstinline|friendlyError| (falha de rede, verificação de
assinatura, falta de espaço em disco, permissão negada). Há ainda a
\textbf{detecção pós-atualização}: a versão em execução é gravada em
\path{userData/aurora-version.json} a cada \textit{boot}; no \textit{boot}
seguinte, uma divergência entre a versão armazenada e a atual indica que o
usuário acabou de atualizar — o \textit{renderer} lê esse estado uma vez
(\lstinline|updates:post-update-status|) e mostra um \textit{toast} de
confirmação. Instalações novas não disparam falso \enquote{você foi
atualizado}, e o sinalizador é de uso único (uma segunda janela na mesma sessão
não repete o \textit{toast}).

\section{Overlay de desenvolvimento (Jank Overlay)}
\label{sec:13-jank}

Para diagnóstico de desempenho durante o desenvolvimento, há o \textit{Jank
Overlay} em \file{js/dev/jank\_overlay.js}, um HUD que mede e exibe em tempo
real:

\begin{itemize}[leftmargin=*]
  \item \textbf{FPS} instantâneo (último quadro);
  \item \textbf{p99} do tempo de quadro sobre uma janela deslizante de 300
        quadros;
  \item \textbf{taxa de \textit{jank}}: percentual de quadros cujo tempo excede
        o dobro do orçamento de 165 Hz --- isto é, $2 \times 6{,}06 \approx
        12$ ms (a constante \lstinline|BUDGET_MS| vale $1000/165 \approx 6{,}06$ ms);
  \item contagem de \textbf{\textit{longtasks}} (via \lstinline|PerformanceObserver|,
        tarefas $>$ 50 ms);
  \item aproximação de \textbf{TTI} (\textit{time to interactive}), a partir da
        marca \lstinline|aurora-interactive| ou do \lstinline|domInteractive| da
        entrada de navegação.
\end{itemize}

O \textit{overlay} tem \textbf{custo zero quando inativo}: o laço de
\lstinline|requestAnimationFrame| e o \lstinline|PerformanceObserver| só rodam
enquanto ele está visível. É alternado pela paleta de comandos
(\enquote{\textit{Toggle Jank Overlay}}, grupo \lstinline|Dev|), e o
\textit{import} dinâmico do módulo evita qualquer custo de inicialização da IDE.
Os valores são coloridos por limiares (verde/âmbar/vermelho) para leitura
rápida.

\section{Síntese}
\label{sec:13-sintese}

As ferramentas de produtividade da \tool{AURORA} compõem uma camada coesa em
torno do editor, conectando processos privilegiados do \textit{main} a
integrações no \tool{Monaco} e à interface. O terminal de saída estrutura e
torna navegável o que as ferramentas do \textit{toolchain} emitem; o controle de
versão embutido (sobre \textbf{simple-git} e \textbf{diff2html}) cobre o ciclo
completo de \lstinline|git| com autenticação GitHub cifrada; a busca global é um
varredor próprio sem dependências; o par \tool{Verible} (sintático completo) +
\tool{slang} (semântico, opcional) entrega diagnósticos, formatação,
\textit{outline}, \textit{hover}, navegação e completação; o
\textit{tree-sitter} acrescenta realce fiel à gramática; o \tool{clang-format}
formata C/C++/\cmm; e paleta de comandos, i18n, configurações, atualizador e
\textit{overlay} de \textit{jank} fecham a experiência. O fio condutor é o
\textit{best-effort}: cada recurso degrada graciosamente para um \textit{no-op}
quando seu binário ou dependência não está presente, sem nunca comprometer a
edição.

% !TEX root = ../main.tex
\chapter{O toolchain bundle (aurora-toolchain)}
\label{cap:14-toolchain-bundle}

A AURORA é uma IDE de desktop, mas o trabalho pesado de simulação, síntese e
compilação de hardware é feito por uma pilha de ferramentas FOSS de EDA
(\emph{Electronic Design Automation}) externas: \tool{iverilog}, \tool{verilator},
\tool{yosys}, além de um compilador C/C++ completo e um interpretador Python com
\tool{cocotb}. Em vez de exigir que o usuário instale e configure cada uma dessas
ferramentas no Windows, a AURORA baixa e desempacota um único artefato portátil:
o \emph{toolchain bundle}. Este capítulo documenta o repositório que produz esse
artefato --- \repo{aurora-toolchain} --- e a maneira como a AURORA o consome.

O repositório \repo{aurora-toolchain} é um \emph{projeto de build} em si: não
contém código que roda dentro da IDE, e sim os \emph{scripts} (\path{build/10..60 *.sh}),
o manifesto de versões (\path{manifest.txt}) e o \emph{workflow} de CI
(\path{.github/workflows/build.yml}) que montam, enxugam, testam e publicam o
arquivo \path{aurora-msys-vN.zip}. O artefato é publicado como um \emph{GitHub
Release} e a AURORA o busca em tempo de \emph{bootstrap}.

\begin{nota}
Toda afirmação deste capítulo foi verificada lendo os \emph{scripts} de build
(\path{build/10-install-packages.sh} a \path{build/60-package.sh}), o
\path{manifest.txt}, o \path{.github/workflows/build.yml} do repositório
\repo{aurora-toolchain} e o consumidor \file{components/Scripts/download-toolchain.js}
no repositório \repo{aurora}. Onde o \path{README.md} divergia do código, o
código prevaleceu.
\end{nota}

\section{Visão geral e papel na suíte SAPHO}

O bundle consolida, num único \emph{prefixo} \tool{MSYS2}/\tool{mingw64}, toda a
parte da pilha de ferramentas que pode ser empacotada como binários portáteis
para Windows. Nem tudo que a AURORA usa está no bundle: o que está fora vem de
outros repositórios, documentados em seus próprios capítulos.

\begin{tabularx}{\linewidth}{Y Y}
\toprule
\textbf{Componente} & \textbf{Onde vive} \\
\midrule
\tool{iverilog}, \tool{verilator}, \tool{yosys}, \tool{g++}, \tool{python}+\tool{cocotb} (ambos os VPIs) & \textbf{No bundle} (este capítulo) \\
\tool{gtkwave-nipscern} (exibição de ondas + \path{fst2vcd}) & \emph{Download} separado (ver \autoref{cap:10-aurora-visualizacao-gtkwave-surfer} e \autoref{cap:15-catalogo-terceiros}) \\
Compiladores YANC (\cmm{} $\rightarrow$ Verilog) & \emph{Download} separado (ver \autoref{cap:04-yanc-compiladores}) \\
\tool{netlistsvg} (renderização de RTL para o PRISM) & Pacote npm em \path{node_modules} (ver \autoref{cap:11-aurora-prism-rtl}) \\
\bottomrule
\end{tabularx}

\medskip

A divisão é deliberada: o bundle reúne apenas o que é grande, pinado por versão
e compartilhado por vários fluxos. O \path{README.md} resume o motivo de o
repositório existir separado da IDE: dar a esse trabalho de empacotamento um
histórico próprio, uma fonte única de verdade para as versões e um build
reproduzível, em vez de ele viver como \emph{scripts} espalhados dentro da AURORA.

\subsection{Layout do artefato}

O \path{aurora-msys-vN.zip} contém, a partir da sua raiz, um diretório \path{msys}
com duas subárvores. O \path{build/60-package.sh} exige explicitamente que o
diretório empacotado se chame \path{msys}, justamente para que a raiz interna do
zip seja \path{msys/...}; ao ser extraído pela AURORA, ele recai em
\file{components/Packages/msys/}.

\begin{tabularx}{\linewidth}{Y Z}
\toprule
\textbf{Caminho dentro do bundle} & \textbf{Conteúdo} \\
\midrule
\path{msys/mingw64/bin/} & Executáveis e DLLs: \path{iverilog.exe}, \path{vvp.exe}, \path{verilator}/\path{verilator_bin.exe}, \path{yosys.exe}, \path{g++}/\path{gcc}/\path{cc1plus}, \path{perl}, \path{make}, \path{ccache}, \path{python.exe} e as DLLs compartilhadas \\
\path{msys/mingw64/lib/ivl/} & \emph{Backends} do \tool{iverilog} \\
\path{msys/mingw64/lib/python3.12/} & Biblioteca padrão do Python + \path{site-packages} com \tool{cocotb} (e os dois VPIs) \\
\path{msys/mingw64/share/verilator/} & Cabeçalhos e \path{verilated.mk} do \tool{verilator} \\
\path{msys/mingw64/share/yosys/} & Dados de tempo de execução do \tool{yosys} \\
\path{msys/usr/bin/} & \path{bash} + \emph{coreutils} do \tool{MSYS} (necessários para o \path{verilated.mk} do \tool{verilator}) \\
\bottomrule
\end{tabularx}

\medskip

O \path{build/40-assemble-bundle.sh} copia o \path{/mingw64} inteiro e, do
\tool{MSYS}, apenas o \path{/usr/bin} (sem \path{usr/lib} nem \path{usr/share}),
porque o bundle enxuto comprovadamente funciona só com \path{usr/bin}, e arrastar
o resto apenas inflaria o artefato.

\section{O que o bundle contém, ferramenta a ferramenta}

\begin{tabularx}{\linewidth}{Y Z}
\toprule
\textbf{Ferramenta} & \textbf{Papel na AURORA} \\
\midrule
\path{iverilog} + \path{vvp} & Simulador Verilog padrão (botão \emph{Wave}) \\
\path{verilator} + \path{verilator_bin} & Simulador rápido opcional; compila um modelo em C++ \\
\path{yosys} & Síntese $\rightarrow$ hierarquia em JSON (consumida pelo PRISM) \\
\path{g++} / \path{gcc} / \path{cc1plus} & Compilador C++ que roda em \textbf{tempo de execução} para compilar o modelo gerado pelo \tool{verilator} \\
\path{perl}, \path{make}, \path{ccache} & \emph{Driver} de build do \tool{verilator} e cache de compilação \\
\path{python} 3.12 + \tool{cocotb} 2.0.1 & \emph{Testbenches} em \tool{cocotb}, com \textbf{os dois VPIs} \\
\bottomrule
\end{tabularx}

\medskip

Um ponto central, repetido em vários comentários dos \emph{scripts} e no
\path{manifest.txt}, é que o \tool{g++} não é só uma dependência de build do
repositório: ele \emph{roda em tempo de execução} na máquina do usuário. O fluxo
do \tool{verilator} emite C++ que precisa ser compilado a cada simulação; logo, o
bundle \emph{precisa} carregar um \tool{g++} funcional. Isso explica por que a
versão do GCC é tratada como um pino crítico (\autoref{sec:14-pins}) e por que o
\path{build/45-trim-bundle.sh} marca \path{g++}/\path{cc1plus} como
\enquote{intocáveis}.

\subsection{Os dois VPIs do cocotb}

O \tool{cocotb} comunica-se com o simulador através de uma \emph{VPI} (\emph{Verilog
Procedural Interface}) --- uma biblioteca de \emph{ponte} específica para cada
simulador. O bundle carrega \textbf{dois} VPIs, para que um \emph{único} Python
sirva aos dois \emph{backends}:

\begin{itemize}
  \item \path{libcocotbvpi_icarus.vpl} --- a VPI do \tool{iverilog}, que é
        \textbf{compilada automaticamente} pelo próprio \emph{build} do
        \tool{cocotb} durante a instalação a partir do \emph{sdist} (sua
        compilação \textbf{não} está sob o portão \code{posix}).
  \item \path{libcocotbvpi_verilator.a} --- a VPI do \tool{verilator}, cuja
        compilação o \emph{build} do \tool{cocotb} \textbf{pula} no Windows
        (está sob a condição \code{if os.name == "posix"}) e que por isso é
        compilada à mão por este repositório (\autoref{sec:14-recipe}).
\end{itemize}

\section{Os pinos de versão (manifest.txt)}
\label{sec:14-pins}

O \path{manifest.txt} é a fonte única de verdade para as versões do bundle. Como
o \tool{MSYS2} é de \emph{rolling release} (atualização contínua), sem esses pinos
um \emph{rebuild} pegaria o que estiver corrente no repositório e poderia quebrar.
Os \emph{scripts} de build leem o manifesto, e a CI instala exatamente essas
versões. Há dois grupos: \emph{pinados} (onde o mais novo é sabidamente quebrado)
e \emph{flutuantes} (acompanham o \tool{MSYS2} mais recente; registrados aqui
apenas para reprodutibilidade).

\subsection{Pinos críticos}

\begin{tabularx}{\linewidth}{Y Y Z}
\toprule
\textbf{Pacote} & \textbf{Versão pinada} & \textbf{Razão (verificada no manifesto e nos scripts)} \\
\midrule
\path{gcc} & \path{15.1.0-5} & A \path{16.1.0-5} traz uma libstdc++ quebrada (construtor de \emph{move} de \path{std::string} indefinido); a VPI do \tool{cocotb} não \emph{linka} contra ela. Como o \tool{g++} roda em tempo de execução, o bundle precisa de um \tool{g++} funcional. \\
\path{gcc-libs} & \path{15.1.0-5} & Bibliotecas de runtime do GCC; pinadas junto com o \path{gcc} pelo mesmo motivo. \\
\path{python} & \path{3.12.11-1} & O \tool{cocotb} \emph{linka} \code{-lpython3.X}; o \emph{minor} do Python do bundle precisa casar com aquele contra o qual o \tool{cocotb} foi construído. A 3.14 quebra o build do \tool{cocotb}. \\
\bottomrule
\end{tabularx}

\medskip

Esses três pacotes \textbf{já não estão} no repositório vivo do \tool{MSYS2} (que
serve gcc16/py3.14 agora). Por isso, em vez de instalá-los via \code{pacman -S},
o \path{build/10-install-packages.sh} os baixa de um \emph{Release} próprio do
repositório, com a \emph{tag} \path{pins-v1}, onde os arquivos
\path{.pkg.tar.zst} exatos ficam arquivados, e os instala com \code{pacman -U}.
Para mover um pino, sobe-se um novo \path{.pkg.tar.zst} num \emph{release}
\path{pins-vN} e aponta-se a variável \code{PINS_TAG} para ele.

\subsection{Pacotes flutuantes e dependências PyPI}

\begin{tabularx}{\linewidth}{Y Y Z}
\toprule
\textbf{Pacote} & \textbf{Versão registrada} & \textbf{Observação} \\
\midrule
\path{verilator} & \path{5.048-1} & Simulador rápido (modelo em C++). \\
\path{iverilog} & \path{1~13.0-2} & Era a 12 no bundle antigo; a 13 é mais estrita (exige HDL que declare antes de usar). \\
\path{yosys} & \path{0.56-4} & Instalado com \code{--assume-installed mingw-w64-x86_64-ghdl} para pular o \emph{frontend} VHDL (que arrastaria gcc16 via \path{gcc-ada}). \\
\path{ccache} & \path{4.13.6-2} & Cache de compilação. \\
\path{perl} & \path{5.38.5-1} & \emph{Wrapper} do \tool{verilator}. \\
\path{make} & \path{4.4.1-4} & \emph{Driver} de build do \tool{verilator}. \\
\path{cocotb} & \path{2.0.1} & Instalado a partir do \emph{sdist}; a VPI do \tool{verilator} é compilada à mão. \\
\path{find_libpython} & \path{0.5.1} & Dependência de runtime do \path{cocotb_tools.runner}. \\
\bottomrule
\end{tabularx}

\medskip

O artefato produzido também está registrado no manifesto: \code{bundle_tag = msys-v1}
(a \emph{tag} do \emph{GitHub Release}) e \code{bundle_filename = aurora-msys-v1.zip}
(o nome do asset, que extrai para \file{components/Packages/msys/}).

\section{A receita cocotb-on-Verilator (a parte difícil)}
\label{sec:14-recipe}

O ponto mais delicado --- e o que justifica a existência do repositório --- é
fazer o \tool{cocotb} rodar sobre o \tool{verilator} no \tool{mingw}. O
\emph{wheel} nativo do \tool{cocotb} para Windows \textbf{não} traz a VPI do
\tool{verilator}: no código-fonte do \tool{cocotb}, a compilação dessa VPI está
sob uma condição \code{if os.name == "posix"}, de modo que o Windows fica sem o
\path{libcocotbvpi_verilator}. A solução implementada em
\path{build/20-build-cocotb-vpi.sh} é construir essa VPI à mão.

\subsection{Por que a VPI é estática e por que \texttt{-DPLI\_DLLISPEC=}}

A biblioteca \path{libcocotbvpi_verilator.a} é construída como \textbf{estática}
(arquivo \path{.a}), e não como DLL, porque o \tool{verilator} \emph{linka} a VPI
diretamente no executável do modelo (não há \code{dlopen}). Ser estática deixa os
símbolos \path{vpi_*}, \path{gpi} e \path{gpilog} indefinidos no \path{.a}: eles
são resolvidos quando o executável \emph{verilado} faz o \emph{link} final --- os
\path{vpi_*} vêm do próprio \tool{verilator} (\code{--vpi}), e \path{gpi}/\path{gpilog}
vêm das DLLs \emph{core} do \tool{cocotb}.

A chave da compilação é a flag \code{-DPLI_DLLISPEC=} (definida como vazia),
combinada com o uso do \path{vpi_user.h} do \emph{próprio} \tool{verilator} (via
\code{-I} apontando para a raiz do \tool{verilator}). Isso faz com que os símbolos
\path{vpi_*} sejam \textbf{planos} --- sem o \code{__declspec(dllimport)}. Sem essa
flag, o \path{.a} referenciaria símbolos \path{__imp_vpi_*}, que o
\path{verilated_vpi.o} (que define os \path{vpi_*} planos) não satisfaria, e o
\emph{link} falharia.

O \path{build/20-build-cocotb-vpi.sh} compila cinco fontes do diretório
\path{src/cocotb/share/lib/vpi/} do \emph{sdist} do \tool{cocotb} --- \path{VpiImpl},
\path{VpiCbHdl}, \path{VpiObj}, \path{VpiIterator} e \path{VpiSignal} --- e as
arquiva com \path{ar} numa única \path{.a}:

\begin{lstlisting}[language=bash,caption={Compilação da VPI estática do Verilator (trecho de 20-build-cocotb-vpi.sh).}]
for f in VpiImpl VpiCbHdl VpiObj VpiIterator VpiSignal; do
  g++ -O2 -std=c++11 -fvisibility=hidden -fvisibility-inlines-hidden \
      -DCOCOTBVPI_EXPORTS= -DVERILATOR= -D__STDC_FORMAT_MACROS= -DWIN32= -DPLI_DLLISPEC= \
      -I"$VROOT/include" -I"$VROOT/include/vltstd" \
      -I"$INC" -I"$PKGDIR/src/cocotb" -I"$PYINC" \
      -c "$VPIDIR/$f.cpp" -o "$OBJ/$f.o"
done
ar rcs "$OBJ/libcocotbvpi_verilator.a" "$OBJ"/*.o
cp -f "$OBJ/libcocotbvpi_verilator.a" "$LIBS/"
\end{lstlisting}

As mesmas fontes e \emph{flags} do build do \tool{iverilog} são reusadas, trocando
\code{-DICARUS} por \code{-DVERILATOR}. Não há \code{-Werror}, para que um
\emph{warning} de versão nova do GCC não interrompa o build.

\subsection{O patch em cocotb\_tools/runner.py}

Compilar a VPI não basta: o \emph{runner} do \tool{cocotb} precisa saber linká-la
ao executável verilado e localizar o \tool{verilator}. O
\path{build/20-build-cocotb-vpi.sh} aplica dois \emph{patches} em
\path{src/cocotb_tools/runner.py} antes de instalar o \tool{cocotb}.

\subsubsection{Patch 1 --- LDFLAGS por token}

Como o \path{.a} é estático e não carrega suas dependências, o executável verilado
precisa também linkar as DLLs \emph{core} do \tool{cocotb}. Além disso, o
\emph{wrapper} \tool{perl} do \tool{verilator} no Windows \textbf{quebra} um
\code{-LDFLAGS} de várias palavras nos espaços (o \tool{verilator} interpreta
\code{-LC:/...} como opção solta e reclama de \enquote{Invalid option}). A
correção é emitir \emph{cada} flag como seu próprio par \code{-LDFLAGS <token>},
que o \tool{verilator} acumula:

\begin{lstlisting}[language=Java,caption={Substituição do par -LDFLAGS no runner (cada flag vira um token).}]
"-LDFLAGS", f"-L{cocotb_tools.config.libs_dir}",
"-LDFLAGS", "-lcocotbvpi_verilator",
"-LDFLAGS", "-lgpi",
"-LDFLAGS", "-lgpilog",
"-LDFLAGS", "-lcocotbutils",
"-LDFLAGS", "-static-libstdc++",
"-LDFLAGS", "-static-libgcc",
\end{lstlisting}

Detalhes do raciocínio registrados no \emph{script}: não se usa
\code{-Wl,-rpath,...} (inútil no Windows/PE, e o \tool{verilator} escaparia as
vírgulas, gerando \enquote{unrecognized option} no \tool{g++}); as DLLs \emph{core}
são achadas via \code{PATH} em tempo de execução. As flags
\code{-static-libstdc++} e \code{-static-libgcc} são necessárias porque o
\path{verilated.o} referencia símbolos de libstdc++ que não resolvem ao misturar
com as DLLs pré-construídas do \tool{cocotb}; linkar estático deixa tudo
consistente.

\subsubsection{Patch 2 --- resolução do verilator}

No \tool{mingw}, o \code{verilator} é um \emph{script} \tool{perl} \textbf{sem
extensão}, e o \code{shutil.which} do Python no Windows nunca encontra um nome
\enquote{puro} (sem \path{.exe}/\path{.bat}). O \emph{patch} acrescenta um
\emph{fallback} que varre o \code{PATH} procurando um arquivo chamado
\code{verilator}:

\begin{lstlisting}[language=Java,caption={Fallback de resolução do verilator (script perl sem extensão).}]
executable = shutil.which("verilator")
if executable is None:
    import os as _os
    for _d in _os.environ.get("PATH", "").split(_os.pathsep):
        _p = _os.path.join(_d, "verilator")
        if _os.path.isfile(_p):
            executable = _p
            break
if executable is None:
    raise SystemExit("ERROR: verilator executable not found!")
\end{lstlisting}

\subsection{Construção a partir do sdist e versão do Python}

O \tool{cocotb} é instalado a partir do \emph{sdist} (não do \emph{wheel}), porque
são necessárias as \emph{fontes} da VPI (\path{src/cocotb/share/lib/vpi/*.cpp})
para compilar o \path{libcocotbvpi_verilator.a} à mão. O \emph{script} cria um
\emph{venv} (o \tool{MSYS2} marca o Python do sistema como \emph{externally
managed} pela PEP 668, o que bloqueia \code{pip install}), baixa o \emph{sdist}
com \code{pip download --no-binary :all:}, aplica os \emph{patches} no
\path{runner.py} e instala. A VPI do \tool{iverilog} (\path{.vpl}) é compilada
automaticamente pelo \emph{build} do \tool{cocotb} durante essa instalação;
apenas a do \tool{verilator} é gerada à mão num passo posterior do \emph{script}.

O \path{build/20-build-cocotb-vpi.sh} também valida a versão do Python: o
\tool{cocotb} 2.0.x só suporta até a 3.13. Como o \tool{pacman} tende a instalar
a mais nova, há uma checagem que aborta o build em Python $>$ 3.13 (com uma saída
de escape opcional \code{ALLOW_PY314=1} que exporta
\code{COCOTB_IGNORE_PYTHON_REQUIRES}, sem garantia). No fluxo oficial, isso não é
exercitado: o \path{build/10-install-packages.sh} já pina o Python em 3.12.11-1.

\section{O pipeline de build, em ordem}
\label{sec:14-pipeline}

Os \emph{scripts} numerados em \path{build/} formam o \emph{pipeline}, executado
numa \emph{shell} \tool{MSYS2 MINGW64}, em ordem. O \path{<out>} típico é
\path{dist/msys}.

\begin{tabularx}{\linewidth}{Y Z}
\toprule
\textbf{Script} & \textbf{O que faz} \\
\midrule
\path{10-install-packages.sh} & Instala os pacotes pinados (gcc15/python3.12) via \code{pacman -U} a partir do \emph{release} \path{pins-v1}, fixa-os em \code{IgnorePkg} no \path{/etc/pacman.conf} e puxa as ferramentas flutuantes (\tool{verilator}/\tool{iverilog}/\tool{ccache}/\tool{perl}/\tool{make}) do repositório vivo. Instala o \tool{yosys} sem \tool{ghdl}. \\
\path{20-build-cocotb-vpi.sh} & Cria o \emph{venv}, baixa o \emph{sdist} do \tool{cocotb}, aplica os \emph{patches} no \path{runner.py}, instala o \tool{cocotb} e compila o \path{libcocotbvpi_verilator.a} (a receita da \autoref{sec:14-recipe}). \\
\path{40-assemble-bundle.sh} & Copia o \path{/mingw64} inteiro e o \path{/usr/bin} do \tool{MSYS} para \path{<out>}, e \enquote{assa} o \tool{cocotb} (com os dois VPIs) no \path{site-packages} do Python do bundle. \\
\path{45-trim-bundle.sh} & Enxuga o bundle em cerca de 40\% (de $\approx$1077\,MB para $\approx$660\,MB descompactado), removendo o que nenhum fluxo da AURORA toca. \\
\path{50-smoke.sh} & Executa os 4 fluxos da AURORA contra o bundle. É o \emph{gate}: precisa dar 4/4 PASS. \\
\path{60-package.sh} & Compacta o bundle (preferindo \tool{7zip}/\tool{zip} a \emph{Compress-Archive}) em \path{aurora-msys-vN.zip}, com a raiz interna \path{msys/...}. \\
\bottomrule
\end{tabularx}

\medskip

Os \emph{scripts} \path{30-package-cocotb.sh} e \path{21-rebuild-vpi.sh} são
auxiliares originais que operam sobre um bundle \emph{já existente}
(\path{21} recompila apenas a \path{.a} estática quando só as \emph{flags} da VPI
mudaram; \path{30} monta o \tool{cocotb} dentro de um bundle pronto e o testa de
forma \emph{standalone}). O \path{40-assemble-bundle.sh} formaliza o caminho
\enquote{do zero}.

\subsection{O \emph{trim} (45-trim-bundle.sh)}

O \path{build/45-trim-bundle.sh} remove, em blocos, partes que nenhum fluxo da
AURORA carrega. Cada bloco foi validado empiricamente: após cada remoção, o
\emph{smoke} de 4 fluxos precisa continuar passando.

\begin{tabularx}{\linewidth}{Y Z}
\toprule
\textbf{Bloco} & \textbf{O que remove} \\
\midrule
1 & Suíte de testes do Python ($\approx$132\,MB), \path{idlelib}, \path{tkinter}, \path{lib2to3}, \path{ensurepip} e dados de GUI em \path{share} (\path{gtk-3.0}, \path{icons}, \path{glib-2.0} etc.). \\
2 & Bibliotecas estáticas \path{.a} de terceiros que a AURORA nunca \emph{linka} estaticamente (gnutls, glib, archive, isl/gmp/mpfr/mpc, bfd/opcodes etc.). \\
3 & Cabeçalhos de bibliotecas (\path{include/}) de GTK/GUI/fontes que o \tool{g++} sobre a saída do \tool{verilator} não precisa. \\
4 & DLLs da pilha GTK/GUI/fontes (nenhuma ferramenta de EDA \emph{headless} as carrega). \\
5 & \tool{perl} duplicado em \path{usr/} (o \tool{verilator} usa o \tool{perl} do \path{mingw64}). \\
\bottomrule
\end{tabularx}

\medskip

O \emph{script} documenta o \enquote{piso intocável}: \path{g++}/\path{cc1plus}
(que roda em tempo de execução para compilar o modelo do \tool{verilator}), o
núcleo da \emph{stdlib} do Python + \tool{cocotb}, os três simuladores/sintetizador,
os \emph{backends} do \tool{iverilog} (\path{lib/ivl}), \path{share/verilator} +
\path{share/yosys}, o \path{bash}+\emph{coreutils} para o \path{verilated.mk}, e as
DLLs que esses executáveis carregam.

\subsection{O smoke obrigatório (50-smoke.sh)}

O \path{build/50-smoke.sh} é o \emph{gate} de cada bloco de \emph{trim} e de cada
\emph{release}. Ele monta arquivos de teste num diretório temporário e executa
quatro fluxos, todos apontando o \code{PATH} para o \emph{próprio} bundle (provando
que ele não depende do \tool{MSYS2} do sistema):

\begin{enumerate}
  \item \textbf{Simulação \tool{iverilog}}: compila \path{dff.v}+\path{dff_tb.v}
        com \path{iverilog.exe}, roda com \path{vvp.exe} e confere a saída
        \code{IVOK}.
  \item \textbf{Hierarquia \tool{yosys} $\rightarrow$ JSON}: lê dois módulos,
        roda \code{hierarchy -top top}, \code{proc} e \code{write_json}, e
        confere que o JSON foi escrito.
  \item \textbf{cocotb-icarus}: \code{get_runner("icarus")} com a VPI \path{.vpl}
        gerada pelo \emph{build} do \tool{cocotb}; confere \code{COCOTB_icarus_OK}.
  \item \textbf{cocotb-verilator}: \code{get_runner("verilator")} com a VPI
        \path{.a} construída à mão; confere \code{COCOTB_verilator_OK}.
\end{enumerate}

O \emph{script} sai com código não-zero se qualquer fluxo falhar; ele imprime
\code{--- PASS=$pass FAIL=$fail ---} e só termina com sucesso se
\code{FAIL=0}. Os dois fluxos do \tool{cocotb} usam \emph{o mesmo} Python do
bundle, com \code{PYTHONHOME} apontando para \path{mingw64} e o \code{PATH}
restrito a \path{mingw64/bin} + \path{usr/bin}.

\section{Build reproduzível via CI}

O \path{.github/workflows/build.yml} é a versão \enquote{de botão} do
\emph{pipeline}: monta o \path{aurora-msys-<tag>.zip} num \emph{runner}
\tool{windows-latest} limpo e o publica como \emph{GitHub Release}, removendo a
dependência da máquina de qualquer desenvolvedor. É acionado manualmente
(\emph{workflow\_dispatch}), recebendo a \emph{tag} e se é \emph{pre-release}.

As etapas, na ordem do \emph{workflow}: configurar o \tool{MSYS2} (MINGW64),
instalar a base mínima (\path{tar}, \path{unzip}, \path{zip}, \tool{7zip} ---
\textbf{sem} \path{git}, para não arrastar \tool{openssh}/módulos CPAN que vazariam
para o bundle), rodar \path{10-install-packages.sh}, \textbf{verificar} que as
versões instaladas batem com o \path{manifest.txt} (um \emph{guardrail} que falha
se um pacote pinado tiver \emph{driftado}), rodar
\path{20}/\path{40}/\path{45}/\path{50}/\path{60}, arquivar os
\path{.pkg.tar.zst} exatos usados (reprodutibilidade) e publicar o \emph{Release}.

\begin{lstlisting}[language=bash,caption={Guardrail de versões pinadas no CI (trecho de build.yml).}]
for kv in "gcc:gcc" "gcc-libs:gcc-libs" "python:python"; do
  key="${kv%%:*}"; pkg="${kv##*:}"
  want="$(sed -nE "s/^${key}[[:space:]]*=[[:space:]]*([^[:space:]#]+).*/\1/p" manifest.txt | head -1)"
  got="$(pacman -Q mingw-w64-x86_64-$pkg | awk '{print $2}')"
  [ "$want" = "$got" ] || { echo "::error::$pkg drifted ($got != $want)"; exit 1; }
done
\end{lstlisting}

\subsection{Como evoluir um release}

O \path{README.md} e os \emph{scripts} descrevem um único caminho de evolução:
editar o \path{manifest.txt} (subir uma versão e a \code{bundle_tag}), rodar o
\emph{pipeline} (ou deixar a CI fazer), confirmar o \emph{smoke} 4/4 e publicar com
\code{gh release create}. Do lado da AURORA, ajusta-se \code{MSYS_TAG} /
\code{MSYS_FILENAME} em \file{components/Scripts/download-toolchain.js} para a nova
\emph{tag}. Em resumo: \emph{um bump de manifesto $\rightarrow$ rebuild $\rightarrow$
smoke $\rightarrow$ publish}.

\section{Como a AURORA consome o bundle}
\label{sec:14-consumo}

O consumidor é \file{components/Scripts/download-toolchain.js} no repositório
\repo{aurora}. Ele baixa o \path{aurora-msys-vN.zip} do \emph{GitHub Release}
pinado e o extrai em \file{components/Packages/}. As constantes que fixam o
artefato:

\begin{tabularx}{\linewidth}{Y Z}
\toprule
\textbf{Constante} & \textbf{Valor} \\
\midrule
\code{MSYS_TAG} & \path{msys-v1} \\
\code{MSYS_FILENAME} & \path{aurora-msys-v1.zip} \\
\code{GITHUB_OWNER} & \path{nipscernlab} \\
\code{GITHUB_REPO} & \path{aurora-toolchain} \\
\bottomrule
\end{tabularx}

\medskip

A URL final é montada como
\path{https://github.com/<owner>/<repo>/releases/download/<tag>/<filename>}. O
download segue até 5 redirecionamentos, e a escrita do zip só é resolvida após o
arquivo estar totalmente \emph{flushed} e fechado --- caso contrário, o
\emph{Expand-Archive} do PowerShell correria contra o \emph{writer} e atingiria o
zip ainda travado (corrupção silenciosa). A extração usa
\code{Expand-Archive} com \code{$ErrorActionPreference='Stop'}, para que uma falha
do \emph{cmdlet} se propague como código de saída não-zero.

\subsection{As sentinelas verificadas}

O \emph{script} usa duas \emph{sentinelas} --- arquivos que comprovam a presença e
a completude da extração:

\begin{itemize}
  \item \textbf{Bundle presente}: \file{components/Packages/msys/mingw64/bin/verilator_bin.exe}
        (constante \code{MSYS_SENTINEL}), um binário fundo no bundle que prova uma
        extração completa.
  \item \textbf{cocotb presente}: a VPI estática do \tool{verilator}. A função
        \code{cocotbInstalled()} varre \path{msys/mingw64/lib/} por qualquer
        diretório \path{python3.*} e checa
        \path{site-packages/cocotb/libs/libcocotbvpi_verilator.a}. Varrer por
        \path{python3.*} (em vez de fixar \path{python3.12}) faz a checagem
        sobreviver a um futuro \emph{bump} de \emph{minor} do Python.
\end{itemize}

A lógica de \code{main()} combina as duas: se o bundle e o \tool{cocotb} estão
presentes (e não há \code{--force}), pula o download; se o bundle está presente
mas falta o suporte a \tool{cocotb} (por exemplo, um \emph{snapshot} mais antigo),
faz o download para \emph{atualizar no lugar}. Após extrair, reconfere a sentinela
do bundle (e sai com erro se faltar) e avisa, sem ser fatal, se o suporte a
\tool{cocotb} não veio.

\begin{lstlisting}[language=Java,caption={Detecção do cocotb via varredura de python3.* (download-toolchain.js).}]
function cocotbInstalled() {
    const libBase = path.join(PACKAGES_DIR, 'msys', 'mingw64', 'lib');
    let entries;
    try { entries = fs.readdirSync(libBase); }
    catch (_e) { return false; }
    return entries.some((name) =>
        /^python3\./i.test(name) &&
        fs.existsSync(path.join(
            libBase, name, 'site-packages', 'cocotb', 'libs',
            'libcocotbvpi_verilator.a')));
}
\end{lstlisting}

\begin{nota}
Uma falha de download não bloqueia o \code{npm start}: o \emph{script} sai com
código 0 após imprimir instruções de download manual, e a própria IDE mostra um
erro claro (\enquote{run npm run bootstrap}) quando uma ferramenta for de fato
invocada sem o bundle presente. Há ainda um gancho de verificação de
\emph{checksum} (\code{EXPECTED_SHA256}), hoje em \code{null} (calcula e registra,
sem impor), para fixar o SHA-256 do artefato no futuro.
\end{nota}

\section{Resumo}

O \repo{aurora-toolchain} encapsula toda a pilha FOSS de EDA num único
\path{aurora-msys-vN.zip}: um prefixo \tool{mingw64} pinado com \tool{iverilog},
\tool{verilator}, \tool{yosys}, \tool{g++} (que roda em tempo de execução),
\tool{perl}/\tool{make}/\tool{ccache} e um Python 3.12 com \tool{cocotb} 2.0.1
carregando os dois VPIs. A peça inédita é a VPI estática
\path{libcocotbvpi_verilator.a}, compilada à mão com \code{-DPLI_DLLISPEC=} e
acompanhada de \emph{patches} cirúrgicos no \path{cocotb_tools/runner.py}, que
permitem rodar \tool{cocotb} sobre \tool{verilator} no Windows. O \emph{pipeline}
numerado (instalar pinos $\rightarrow$ build da VPI $\rightarrow$ \emph{assemble}
$\rightarrow$ \emph{trim} $\rightarrow$ \emph{smoke} 4/4 $\rightarrow$ \emph{package})
é reproduzível na CI e protegido por um \emph{guardrail} de versões; do lado da
AURORA, o \file{download-toolchain.js} busca o artefato pela \emph{tag} pinada e
confirma a instalação por duas sentinelas.

% !TEX root = ../main.tex
\chapter{Catálogo de ferramentas de terceiros}
\label{cap:15-catalogo-terceiros}

Este capítulo é o catálogo de referência das ferramentas de terceiros que o
SAPHO usa, orquestra ou empacota. Cada componente aparece com seu nome, a versão
\emph{exata} verificada na fonte primária, a licença confirmada e o papel que
desempenha na plataforma. As fontes primárias deste levantamento são o
\file{package.json} da AURORA (campos \code{dependencies} e \code{devDependencies}),
o \file{manifest.txt} do repositório \repo{aurora-toolchain}, os scripts
\file{components/Scripts/download-*.js} (que baixam e fixam cada binário externo),
o \file{main/ai/cli\_manifest.js} (manifesto das CLIs de IA baixadas sob demanda)
e os arquivos \file{LICENSE} e \file{THIRD\_PARTY\_NOTICES.md} da raiz.

As versões aqui registradas foram confirmadas lendo diretamente esses arquivos e,
quando instaladas, os \file{package.json} de cada pacote em \path{node_modules/}.
Os caminhos de instalação dos binários nativos (toolchain, viewers, LSPs) e o
fluxo de \emph{bootstrap} que os baixa são detalhados em
\autoref{cap:10-toolchain} e \autoref{cap:14-build-empacotamento}; aqui o foco é
o \emph{inventário} de cada peça.

\begin{nota}
A AURORA é distribuída sob licença MIT (\copyright{} 2024--presente NIPS-CERN).
Cada componente de terceiros permanece sob sua própria licença, que se aplica aos
respectivos arquivos em adição à licença da AURORA. Os componentes nativos sob
GPL (Icarus Verilog, GTKWave) e o Surfer sob EUPL são invocados como processos
externos, separados e em \emph{arm's-length} --- a AURORA os executa, mas não
faz \emph{link} contra eles. Os componentes LGPL são usados sem modificação.
\end{nota}

\section{Como ler este catálogo}

\subsection{Modos de empacotamento}

As ferramentas chegam ao usuário por quatro mecanismos distintos, e a coluna
\enquote{origem} de cada tabela indica qual se aplica.

\begin{description}[leftmargin=2.2em]
  \item[npm (\code{dependencies}).] Bibliotecas JavaScript/TypeScript instaladas
  em \path{node_modules/} e empacotadas dentro do instalador pelo
  \tool{electron-builder}. São as bibliotecas de runtime da aplicação.
  \item[npm (\code{devDependencies}).] Ferramentas usadas apenas no
  desenvolvimento e na CI (build, testes, lint). Não vão para o instalador final.
  \item[Bootstrap (download fixado).] Binários nativos (\code{.exe}, \code{.wasm},
  fontes, \emph{bundle} MSYS2) baixados de releases do GitHub/GitLab por um script
  \file{download-*.js} durante o \code{npm run bootstrap}. Cada um fixa uma
  \emph{tag} de versão e, na maioria dos casos, um \code{SHA-256} verificado antes
  da extração. São empacotados no instalador via \code{extraResources}.
  \item[On-demand (runtime).] As CLIs de IA por assinatura (Claude Code, Codex)
  \emph{não} são empacotadas: são baixadas do registro npm na primeira utilização,
  a partir de um manifesto fixado em \file{main/ai/cli\_manifest.js}.
\end{description}

\subsection{Fixação de versões e integridade}

O fluxo de bootstrap é descrito em
\autoref{cap:14-build-empacotamento}; aqui basta registrar os mecanismos que
garantem que as versões deste catálogo são as que efetivamente chegam ao usuário.

\begin{itemize}[leftmargin=1.4em]
  \item \textbf{Pacotes npm fixados.} O script \file{scripts/check-pinned-versions.js}
  trata qualquer dependência declarada com \emph{semver} puro (sem \lstinline|^|,
  \lstinline|~| ou outro operador de intervalo) como \enquote{fixada} e
  falha o build se a versão instalada divergir. Hoje os pacotes com versão exata
  são \lstinline|monaco-editor| (\code{0.52.2}), \lstinline|web-tree-sitter|
  (\code{0.26.9}), \lstinline|@anthropic-ai/claude-code| (\code{2.1.144}),
  \lstinline|@openai/codex| (\code{0.131.0}) e \lstinline|knip| (\code{6.17.1}).
  \item \textbf{Binários de bootstrap.} Cada \file{download-*.js} carrega uma
  constante de \emph{tag} e, na maioria, um \code{EXPECTED_SHA256} verificado
  via \file{components/Scripts/lib/checksum.js} antes da extração. Uma divergência
  de hash aborta a instalação.
  \item \textbf{CLIs de IA on-demand.} O \file{cli\_manifest.js} fixa versão,
  URL do \emph{tarball} e \emph{Subresource-Integrity} (\code{sha512-...}) por
  plataforma; o \file{check-pinned-versions.js} cruza essas versões com o
  \file{package.json} e os hashes com o \file{package-lock.json}.
\end{itemize}

\section{Runtime da aplicação}

A base de execução da AURORA é o Electron, que combina o Chromium (processo
\emph{renderer}, onde vive a UI) e o Node.js (processo \emph{main}, onde vivem
o sistema de arquivos, os spawns de toolchain e a IPC). A versão de Node mínima
declarada em \code{engines} é \code{>=18.0.0}.

\begin{longtable}{>{\raggedright\arraybackslash}p{0.20\linewidth} >{\raggedright\arraybackslash}p{0.12\linewidth} >{\raggedright\arraybackslash}p{0.13\linewidth} >{\raggedright\arraybackslash}p{0.43\linewidth}}
\caption{Runtime da aplicação}\label{tab:cat-runtime}\\
\toprule
\textbf{Componente} & \textbf{Versão} & \textbf{Licença} & \textbf{Papel no SAPHO} \\
\midrule
\endfirsthead
\toprule
\textbf{Componente} & \textbf{Versão} & \textbf{Licença} & \textbf{Papel no SAPHO} \\
\midrule
\endhead
\bottomrule
\endfoot
\bottomrule
\endlastfoot
\tool{electron} & \code{39.8.10} & MIT & Runtime desktop (Chromium + Node). Hospeda toda a aplicação; é \code{devDependency} porque o \tool{electron-builder} a embute no instalador. \\
Node.js & \code{>=18} & MIT & Ambiente do processo \emph{main}; declarado em \code{engines}. \\
\end{longtable}

\section{Editor e interface}

A camada de edição e UI combina o Monaco (o editor de código do VS Code),
componentes web baseados em Lit, ícones Phosphor, renderização de fórmulas com
KaTeX, jQuery/jQuery-UI para alguns painéis legados, \emph{diffs} HTML com
diff2html e as fontes tipográficas. O destaque semântico de sintaxe usa o
runtime \lstinline|web-tree-sitter| com gramáticas \code{.wasm} (catalogadas em
\autoref{tab:cat-lsp}).

\begin{longtable}{>{\raggedright\arraybackslash}p{0.21\linewidth} >{\raggedright\arraybackslash}p{0.10\linewidth} >{\raggedright\arraybackslash}p{0.13\linewidth} >{\raggedright\arraybackslash}p{0.44\linewidth}}
\caption{Editor e interface (npm)}\label{tab:cat-ui}\\
\toprule
\textbf{Componente} & \textbf{Versão} & \textbf{Licença} & \textbf{Papel no SAPHO} \\
\midrule
\endfirsthead
\toprule
\textbf{Componente} & \textbf{Versão} & \textbf{Licença} & \textbf{Papel no SAPHO} \\
\midrule
\endhead
\bottomrule
\endfoot
\bottomrule
\endlastfoot
\lstinline|monaco-editor| & \code{0.52.2} & MIT & Editor de código (núcleo do VS Code). Fixado em versão exata: a \code{0.53.x} regride na inicialização. Hospeda \cmm{}, C/C++, ASM, Verilog/SystemVerilog. \\
\lstinline|lit| & \code{3.3.3} & BSD-3-Clause & Biblioteca de Web Components reativos usada nos elementos de UI próprios (\emph{shell}, abas, painéis). \\
\lstinline|@phosphor-icons/web| & \code{2.1.2} & MIT & Conjunto de ícones da interface. \\
\lstinline|katex| & \code{0.17.0} & MIT & Renderização de fórmulas matemáticas (LaTeX) em superfícies de UI/documentação. \\
\lstinline|jquery| & \code{3.7.1} & MIT & Manipulação de DOM em componentes de UI legados. \\
\lstinline|jquery-ui| & \code{1.14.2} & MIT & Widgets de interface (sobre jQuery). \\
\lstinline|diff2html| & \code{3.4.56} & MIT & Renderização HTML de \emph{diffs} unificados (visualização de controle de versão). \\
\end{longtable}

\subsection{Fontes tipográficas}

As fontes ficam em \path{assets/fonts/} e todas são versionadas no repositório
como \code{.woff2}: Inter e JetBrains Mono na interface, Metamorphous e Noto
Sans Runic no letreiro \enquote{Dagr} do controle de versão. As quatro estão sob
SIL OFL 1.1 e são buscadas pelo \file{scripts/fetch-fonts.js}, que não faz parte
do build --- roda à mão apenas para atualizá-las.

Até 08/08/2026 o letreiro usava a fonte Norse, de Joël Carrouché, que não era
versionada porque sua licença proíbe redistribuição. O arranjo parecia correto,
mas não era: o Vite copiava a fonte para \path{dist/assets/} e o
\file{electron-builder} não exclui \path{assets/}, então o instalador publicado
levava o arquivo dentro. Manter uma fonte fora do repositório não é o mesmo que
mantê-la fora da distribuição. Sob OFL o problema deixa de existir, e com ele
sumiram o \file{download-norse-font.js} e o passo correspondente no CI.

\begin{longtable}{>{\raggedright\arraybackslash}p{0.22\linewidth} >{\raggedright\arraybackslash}p{0.18\linewidth} >{\raggedright\arraybackslash}p{0.50\linewidth}}
\caption{Fontes}\label{tab:cat-fonts}\\
\toprule
\textbf{Fonte} & \textbf{Licença} & \textbf{Papel no SAPHO} \\
\midrule
\endfirsthead
\toprule
\textbf{Fonte} & \textbf{Licença} & \textbf{Papel no SAPHO} \\
\midrule
\endhead
\bottomrule
\endfoot
\bottomrule
\endlastfoot
Inter & SIL OFL 1.1 & Fonte de interface (texto da UI). Versionada (\file{inter-latin.woff2}, \file{inter-latin-ext.woff2}). \\
JetBrains Mono & SIL OFL 1.1 & Fonte monoespaçada (editor/terminal). Versionada (\file{jetbrains-mono-latin.woff2}, \file{jetbrains-mono-latin-ext.woff2}). \\
Metamorphous (James Grieshaber) & SIL OFL 1.1 & Letreiro \enquote{Dagr} do painel de controle de versão. Versionada como \code{.woff2}. \\
Noto Sans Runic (Google) & SIL OFL 1.1 & A runa Dagaz (\code{U+16DE}) da mesma marca; é a única das quatro que cobre o bloco rúnico. Versionada como \code{.woff2}. \\
\end{longtable}

\section{Toolchain EDA}

O coração do fluxo de hardware é o \emph{bundle} unificado MSYS2/MinGW64 do
repositório \repo{aurora-toolchain}, baixado pelo \file{download-toolchain.js}
como \file{aurora-msys-v1.zip} (\emph{tag} \code{msys-v1}) e extraído em
\path{components/Packages/msys/}. Esse \emph{bundle} carrega, em um só lugar,
Icarus Verilog, Verilator, Yosys, GCC/G++, Python e cocotb com suas DLLs. As
versões são fixadas no \file{manifest.txt} do \repo{aurora-toolchain}, que é a
fonte única de verdade para a montagem do \emph{bundle} (MSYS2 é
\emph{rolling-release}; sem os pinos, uma reconstrução puxa versões correntes que
podem quebrar o fluxo). Cada linha do manifesto é uma versão exata de pacote
MSYS2.

O GTKWave não vem do \emph{bundle} principal: é um \emph{fork} próprio
(\repo{gtkwave-nipscern}) baixado à parte. O netlistsvg deixou de ser um binário
e roda \emph{in-process} a partir de \path{node_modules/} (ver
\autoref{sec:cat-prism}). DigitalJS e yosys2digitaljs são bibliotecas npm
empregadas pelo visualizador PRISM (\autoref{cap:12-prism}).

\begin{nota}
O \code{7-Zip} não é uma dependência ativa. O empacotamento de \emph{backups}
de projeto usa o \code{Compress-Archive} nativo do PowerShell
(\file{main/ipc/files.js}); a referência a \code{7z} no código é um comentário
documentando que essa dependência foi removida.
\end{nota}

\subsection{Bundle MSYS2/MinGW e ferramentas nativas}

\begin{longtable}{>{\raggedright\arraybackslash}p{0.20\linewidth} >{\raggedright\arraybackslash}p{0.14\linewidth} >{\raggedright\arraybackslash}p{0.18\linewidth} >{\raggedright\arraybackslash}p{0.36\linewidth}}
\caption{Toolchain EDA (bundle e binários de bootstrap)}\label{tab:cat-eda}\\
\toprule
\textbf{Componente} & \textbf{Versão} & \textbf{Licença} & \textbf{Papel no SAPHO} \\
\midrule
\endfirsthead
\toprule
\textbf{Componente} & \textbf{Versão} & \textbf{Licença} & \textbf{Papel no SAPHO} \\
\midrule
\endhead
\bottomrule
\endfoot
\bottomrule
\endlastfoot
Icarus Verilog (\code{iverilog}/\code{vvp}) & \code{13.0} (pkg \lstinline|1~13.0-2|) & GNU GPL v2 & Compilação e simulação Verilog. A v13 é mais estrita (exige declaração antes do uso). No \emph{bundle}. \\
Verilator & \code{5.048} (pkg \code{5.048-1}) & LGPL-3.0 ou Artistic-2.0 & Simulação Verilog de alto desempenho (compila o modelo via G++ em runtime). No \emph{bundle}. \\
Yosys (+ ABC) & \code{0.56} (pkg \code{0.56-4}) & ISC (Yosys); MIT-style (ABC, UC Berkeley) & Síntese RTL para o PRISM. Instalado sem o \emph{frontend} VHDL (GHDL). No \emph{bundle}. \\
GTKWave (\emph{fork} nipscern) & \code{v0.1.2-nipscern} & GNU GPL v2 & Visualizador de formas de onda + \code{fst2vcd}/\code{vcd2fst}. \emph{Fork} portátil (DLLs GTK incluídas) com \code{--zoom-fit}, painel escuro, \code{--left-justify}. Baixado por \file{download-gtkwave-nipscern.js}. \\
cocotb & \code{2.0.1} & BSD-3-Clause & Framework de testbench em Python (VPIs Icarus + Verilator compiladas no \emph{bundle}). \\
\code{find_libpython} & \code{0.5.1} & MIT & Dependência de runtime de \code{cocotb_tools.runner} (no \emph{bundle}). \\
Python (MinGW) & \code{3.12.11} (pkg \code{3.12.11-1}) & PSF & Interpretador que hospeda o cocotb. Fixado: a 3.14 quebra o build do cocotb. No \emph{bundle}. \\
GCC / G++ & \code{15.1.0} (pkg \code{15.1.0-5}) & GNU GPL v3 (runtime: GCC Runtime Library Exception) & Compilador que o Verilator usa em runtime para compilar o modelo. Fixado: a 16.x quebra o link do VPI do cocotb. No \emph{bundle}. \\
\code{ccache} & \code{4.13.6} (pkg \code{4.13.6-2}) & GNU GPL v3 & Cache de compilação (acelera o Verilator). No \emph{bundle}. \\
Perl (MinGW) & \code{5.38.5} (pkg \code{5.38.5-1}) & Artistic / GPL & \emph{Wrapper} do Verilator. No \emph{bundle}. \\
\code{make} (MinGW) & \code{4.4.1} (pkg \code{4.4.1-4}) & GNU GPL v3 & Orquestra a compilação do modelo Verilator (\file{verilated.mk}). No \emph{bundle}. \\
MSYS2 / MinGW runtime & --- & Diversas (GPL/LGPL/MIT por pacote) & Hospeda os binários GTK/Yosys no Windows (bash + coreutils em \path{usr/bin}). No \emph{bundle}. \\
\end{longtable}

\subsection{PRISM: visualização de RTL via npm}
\label{sec:cat-prism}

O PRISM (\autoref{cap:12-prism}) sintetiza o RTL com o Yosys do \emph{bundle} e
renderiza o esquemático. O netlistsvg deixou de ser um \code{.exe} empacotado e
hoje é carregado \emph{in-process} a partir de \path{node_modules/} em
\file{main/ipc/prism.js} (\lstinline|require('@silimate/netlistsvg')|), o que
elimina a necessidade de resolver caminho ou de entrada na \emph{allowlist} de
binários. O pacote é um \emph{fork} do laboratório, fixado a um commit Git
específico no \file{package.json}.

\begin{longtable}{>{\raggedright\arraybackslash}p{0.21\linewidth} >{\raggedright\arraybackslash}p{0.12\linewidth} >{\raggedright\arraybackslash}p{0.12\linewidth} >{\raggedright\arraybackslash}p{0.43\linewidth}}
\caption{PRISM e visualização de circuitos (npm)}\label{tab:cat-prism}\\
\toprule
\textbf{Componente} & \textbf{Versão} & \textbf{Licença} & \textbf{Papel no SAPHO} \\
\midrule
\endfirsthead
\toprule
\textbf{Componente} & \textbf{Versão} & \textbf{Licença} & \textbf{Papel no SAPHO} \\
\midrule
\endhead
\bottomrule
\endfoot
\bottomrule
\endlastfoot
\lstinline|@silimate/netlistsvg| (\emph{fork} nipscern) & \code{1.1.4} (commit \code{58a0703}) & MIT & Renderiza o JSON do Yosys em esquemático SVG, em processo. Fixado por commit em \repo{netlistsvg-aurora}. \\
\lstinline|digitaljs| & \code{0.14.2} & BSD-2-Clause & Simulador/visualizador de circuitos digitais no navegador (PRISM interativo). \\
\lstinline|yosys2digitaljs| & \code{0.10.3} & BSD-2-Clause & Converte a saída do Yosys para o formato DigitalJS. \\
\end{longtable}

\section{Linguagem, LSP e formatação}

A inteligência de edição combina dois \emph{language servers} de Verilog/SystemVerilog
(Verible para lint/sintaxe; slang-server para análise semântica completa),
gramáticas \code{tree-sitter} para destaque semântico de Verilog/C/C++, e o
clang-format para formatação de C/C++/\cmm{}. Esses motores são detalhados em
\autoref{cap:09-editor-lsp}; abaixo está apenas o inventário.

Os três binários de LSP/formatação são baixados no bootstrap, cada um com tag e
\code{SHA-256} fixados, e instalados em \path{components/Packages/}. As gramáticas
\code{.wasm} são baixadas pelo \file{download-tree-sitter-grammars.js} (cada uma
com hash fixado), e o runtime \lstinline|web-tree-sitter.wasm| é copiado de
\path{node_modules/} para casar exatamente com a versão do JavaScript empacotado.

\begin{longtable}{>{\raggedright\arraybackslash}p{0.24\linewidth} >{\raggedright\arraybackslash}p{0.18\linewidth} >{\raggedright\arraybackslash}p{0.14\linewidth} >{\raggedright\arraybackslash}p{0.30\linewidth}}
\caption{Linguagem, LSP e formatação}\label{tab:cat-lsp}\\
\toprule
\textbf{Componente} & \textbf{Versão} & \textbf{Licença} & \textbf{Papel no SAPHO} \\
\midrule
\endfirsthead
\toprule
\textbf{Componente} & \textbf{Versão} & \textbf{Licença} & \textbf{Papel no SAPHO} \\
\midrule
\endhead
\bottomrule
\endfoot
\bottomrule
\endlastfoot
Verible (\code{verible-verilog-ls.exe}) & \code{v0.0-4080-ga0a8d8eb} & Apache-2.0 & LSP de Verilog: diagnóstico (lint+sintaxe), formatação, \emph{outline}, hover, ir-para-definição. Só o \emph{language server} é extraído. Baixado por \file{download-verible.js}. \\
slang-server (\code{slang-server.exe}) & \code{v0.2.7} & MIT & LSP semântico de SystemVerilog (elaboração completa): erros que o lint sintático não vê + autocompletar. Hudson River Trading. Baixado por \file{download-slang-server.js}. \\
clang-format (\code{clang-format.exe}) & LLVM \code{20} (tag \code{master-796e77c}) & Apache-2.0 com exceção LLVM & Formatação de C/C++/\cmm{} no editor (\cmm{} usa regras de C). \emph{Repack} estático de \code{muttleyxd/clang-tools-static-binaries}. Baixado por \file{download-clang-format.js}. \\
\lstinline|web-tree-sitter| (runtime) & \code{0.26.9} & MIT & Runtime WASM do tree-sitter; gera \emph{semantic tokens} no Monaco. Fixado em versão exata; copiado para \path{components/Packages/tree-sitter/}. \\
\code{tree-sitter-systemverilog.wasm} & \code{v0.3.1} (gmlarumbe) & MIT & Gramática para destaque semântico de \code{.v}/\code{.sv}. \\
\code{tree-sitter-c.wasm} & \code{v0.24.2} & MIT & Gramática para destaque semântico de C. \\
\code{tree-sitter-cpp.wasm} & \code{v0.23.4} & MIT & Gramática para destaque semântico de C++. \\
\end{longtable}

\section{Visualização de formas de onda (Surfer)}

O Surfer é o visualizador de formas de onda \emph{opt-in} (alternância
GTKWave\,$\leftrightarrow$\,Surfer); na sua ausência, o botão de onda recai sobre
o GTKWave. É um binário Rust baixado do registro de pacotes genérico do GitLab do
projeto Surfer, com \code{SHA-256} fixado, e instalado em
\path{components/Packages/surfer/}.

\begin{longtable}{>{\raggedright\arraybackslash}p{0.20\linewidth} >{\raggedright\arraybackslash}p{0.12\linewidth} >{\raggedright\arraybackslash}p{0.13\linewidth} >{\raggedright\arraybackslash}p{0.43\linewidth}}
\caption{Visualização de formas de onda}\label{tab:cat-surfer}\\
\toprule
\textbf{Componente} & \textbf{Versão} & \textbf{Licença} & \textbf{Papel no SAPHO} \\
\midrule
\endfirsthead
\toprule
\textbf{Componente} & \textbf{Versão} & \textbf{Licença} & \textbf{Papel no SAPHO} \\
\midrule
\endhead
\bottomrule
\endfoot
\bottomrule
\endlastfoot
Surfer (\code{surfer.exe}) & \code{v0.7.0} & EUPL-1.2 & Visualizador de formas de onda em Rust (alternativa ao GTKWave). Spawn em \emph{arm's-length}. Baixado por \file{download-surfer.js}. \\
\end{longtable}

\section{Inteligência artificial}

O subsistema Aurora Intelligence (\autoref{cap:13-aurora-intelligence}) é
construído sobre o AI SDK (Vercel) com adaptadores para vários provedores, o SDK
do Model Context Protocol e o validador de esquemas Zod. As CLIs por assinatura
da Anthropic (Claude Code) e da OpenAI (Codex) são integradas, mas não
empacotadas no instalador: são baixadas sob demanda na primeira utilização, a
partir do manifesto fixado em \file{main/ai/cli\_manifest.js} (versão, URL do
\emph{tarball} e \emph{Subresource-Integrity} por plataforma). O
\file{package.json} declara essas CLIs como dependências para fixar as versões
base, mas as exclui explicitamente do \emph{bundle} (\code{build.files}).

\begin{nota}
O subsistema Aurora Intelligence está completamente implementado. O \emph{provider}
e os modelos \emph{próprios} do laboratório ainda não foram treinados (item de
\emph{roadmap}); até lá, a inteligência opera por meio dos provedores externos
listados abaixo.
\end{nota}

\begin{longtable}{>{\raggedright\arraybackslash}p{0.27\linewidth} >{\raggedright\arraybackslash}p{0.10\linewidth} >{\raggedright\arraybackslash}p{0.16\linewidth} >{\raggedright\arraybackslash}p{0.35\linewidth}}
\caption{Inteligência artificial}\label{tab:cat-ai}\\
\toprule
\textbf{Componente} & \textbf{Versão} & \textbf{Licença} & \textbf{Papel no SAPHO} \\
\midrule
\endfirsthead
\toprule
\textbf{Componente} & \textbf{Versão} & \textbf{Licença} & \textbf{Papel no SAPHO} \\
\midrule
\endhead
\bottomrule
\endfoot
\bottomrule
\endlastfoot
\lstinline|ai| & \code{6.0.184} & Apache-2.0 & Núcleo do AI SDK (chamadas de modelo, streaming, ferramentas). \\
\lstinline|@ai-sdk/anthropic| & \code{3.0.85} & Apache-2.0 & Adaptador do provedor Anthropic. \\
\lstinline|@ai-sdk/openai| & \code{3.0.64} & Apache-2.0 & Adaptador do provedor OpenAI. \\
\lstinline|@ai-sdk/google| & \code{3.0.83} & Apache-2.0 & Adaptador do provedor Google. \\
\lstinline|@ai-sdk/groq| & \code{3.0.39} & Apache-2.0 & Adaptador do provedor Groq. \\
\lstinline|@ai-sdk/deepseek| & \code{2.0.39} & Apache-2.0 & Adaptador do provedor DeepSeek. \\
\lstinline|@modelcontextprotocol/sdk| & \code{1.29.0} & MIT & SDK do Model Context Protocol (servidores/clientes de ferramentas). \\
\lstinline|zod| & \code{3.25.76} & MIT & Validação de esquemas (formato de saídas e ferramentas da IA). \\
\lstinline|@anthropic-ai/claude-code| & \code{2.1.144} & Termos comerciais da Anthropic & CLI Claude Code. On-demand: baixada na 1ª utilização; não empacotada. \\
\lstinline|@openai/codex| & \code{0.131.0} & Apache-2.0 & CLI Codex. On-demand: baixada na 1ª utilização; não empacotada. \\
\end{longtable}

\section{Build e infraestrutura}

O empacotamento usa Vite (build do \emph{renderer}) e electron-builder (instalador
NSIS), com atualização automática via electron-updater (publicação em releases do
GitHub no repositório \repo{sapho}). O log, a observação de arquivos, as operações
de arquivo e o Git são cobertos por bibliotecas dedicadas. O processo de build
completo está em \autoref{cap:14-build-empacotamento}.

\begin{longtable}{>{\raggedright\arraybackslash}p{0.23\linewidth} >{\raggedright\arraybackslash}p{0.11\linewidth} >{\raggedright\arraybackslash}p{0.11\linewidth} >{\raggedright\arraybackslash}p{0.43\linewidth}}
\caption{Build e infraestrutura}\label{tab:cat-build}\\
\toprule
\textbf{Componente} & \textbf{Versão} & \textbf{Licença} & \textbf{Papel no SAPHO} \\
\midrule
\endfirsthead
\toprule
\textbf{Componente} & \textbf{Versão} & \textbf{Licença} & \textbf{Papel no SAPHO} \\
\midrule
\endhead
\bottomrule
\endfoot
\bottomrule
\endlastfoot
\lstinline|vite| & \code{8.0.16} & MIT & Bundler do \emph{renderer} (\code{build:renderer}). \emph{devDependency}. \\
\lstinline|vite-plugin-static-copy| & \code{4.1.1} & MIT & Cópia de ativos estáticos no build do Vite. \emph{devDependency}. \\
\lstinline|electron-builder| & \code{26.15.3} & MIT & Geração do instalador NSIS e empacotamento. \emph{devDependency}. \\
\lstinline|electron-updater| & \code{6.8.9} & MIT & Atualização automática (releases do GitHub). \\
\lstinline|electron-log| & \code{5.4.3} & MIT & Sistema de log de \emph{main} e \emph{renderer}. \\
\lstinline|chokidar| & \code{4.0.3} & MIT & Observação de mudanças no sistema de arquivos (árvore do projeto). \\
\lstinline|fs-extra| & \code{11.3.2} & MIT & Operações de arquivo estendidas (cópia, remoção recursiva). \\
\lstinline|simple-git| & \code{3.36.0} & MIT & Operações Git (painel de controle de versão \enquote{Dagr}). \\
\end{longtable}

\section{Qualidade e ferramentas de desenvolvimento}

As ferramentas de qualidade e desenvolvimento (\code{devDependencies}, não
empacotadas) cobrem tipagem TypeScript, lint com ESLint, testes unitários com
Vitest, testes \emph{end-to-end} com Playwright (sobre Electron, com DOM
simulado via happy-dom), \emph{hooks} de Git com Husky + lint-staged, validação
de mensagens de commit com commitlint e detecção de código morto com Knip. O
versionamento automatizado de releases é feito pelo \tool{release-please}, que é
uma \emph{GitHub Action} (workflow \file{.github/workflows/release-please.yml}),
não um pacote npm.

\begin{longtable}{>{\raggedright\arraybackslash}p{0.27\linewidth} >{\raggedright\arraybackslash}p{0.11\linewidth} >{\raggedright\arraybackslash}p{0.10\linewidth} >{\raggedright\arraybackslash}p{0.40\linewidth}}
\caption{Qualidade e desenvolvimento (devDependencies)}\label{tab:cat-dev}\\
\toprule
\textbf{Componente} & \textbf{Versão} & \textbf{Licença} & \textbf{Papel no SAPHO} \\
\midrule
\endfirsthead
\toprule
\textbf{Componente} & \textbf{Versão} & \textbf{Licença} & \textbf{Papel no SAPHO} \\
\midrule
\endhead
\bottomrule
\endfoot
\bottomrule
\endlastfoot
\lstinline|typescript| & \code{6.0.3} & Apache-2.0 & Compilador TypeScript (\code{build:ts}). \\
\lstinline|@types/node| & \code{25.9.3} & MIT & Tipos do Node.js para o processo \emph{main}. \\
\lstinline|eslint| & \code{9.39.1} & MIT & Linter de JavaScript. \\
\lstinline|@eslint/js| & \code{9.39.1} & MIT & Configuração base recomendada do ESLint. \\
\lstinline|globals| & \code{16.5.0} & MIT & Definições de variáveis globais para o ESLint. \\
\lstinline|vitest| & \code{4.1.9} & MIT & Framework de testes unitários (\code{test}). \\
\lstinline|@vitest/coverage-v8| & \code{4.1.9} & MIT & Cobertura de testes (V8). \\
\lstinline|happy-dom| & \code{20.10.6} & MIT & DOM simulado para os testes. \\
\lstinline|playwright| & \code{1.61.0} & Apache-2.0 & Testes \emph{end-to-end} da aplicação Electron. \\
\lstinline|husky| & \code{9.1.7} & MIT & Gestão de \emph{hooks} de Git (\code{prepare}). \\
\lstinline|lint-staged| & \code{16.4.0} & MIT & Roda o linter apenas nos arquivos \emph{staged}. \\
\lstinline|@commitlint/cli| & \code{19.8.1} & MIT & Validação de mensagens de commit. \\
\lstinline|@commitlint/config-conventional| & \code{19.8.1} & MIT & Regras de \emph{Conventional Commits}. \\
\lstinline|knip| & \code{6.17.1} & ISC & Detecção de código/dependências mortos (\code{deadcode}). Fixado em versão exata. \\
\end{longtable}

\section{Componentes próprios do laboratório (YANC)}

Embora os compiladores YANC (\autoref{cap:05-yanc-pipeline}) sejam desenvolvidos
internamente pelo NIPS-CERN e não sejam \enquote{terceiros}, eles chegam à AURORA
pelo mesmo mecanismo de bootstrap dos binários externos --- o
\file{download-yanc.js} baixa \file{yanc-bin-v5.2.zip} (\emph{tag} \code{v5.2}) e o
extrai em \path{components/}. Registra-se aqui apenas por completude do inventário
de artefatos baixados; o conteúdo (6 executáveis em \path{bin/} --- os
compiladores/pré-processadores \code{cmmcomp}, \code{cppcomp}, \code{asmcomp},
\code{cpppp} e \code{appcomp}, mais o utilitário \code{comp2gtkw} de decodificação
de números complexos para formas de onda --- além da biblioteca Verilog HDL,
\emph{headers} C++ e macros de ponto flutuante) é descrito no capítulo do
\emph{pipeline} YANC. Tudo é licenciado segundo o projeto YANC.

\section{Síntese de licenças}

A tabela a seguir consolida as famílias de licença presentes na distribuição. A
AURORA em si é MIT; os componentes mais permissivos (MIT, BSD, ISC, Apache-2.0)
são empacotados sem restrição prática; os componentes \emph{copyleft} fortes (GPL)
e o EUPL são invocados como processos externos separados (sem \emph{link});
o Verilator (LGPL) é usado sem modificação.

\begin{longtable}{>{\raggedright\arraybackslash}p{0.26\linewidth} >{\raggedright\arraybackslash}p{0.30\linewidth} >{\raggedright\arraybackslash}p{0.36\linewidth}}
\caption{Famílias de licença na distribuição}\label{tab:cat-licencas}\\
\toprule
\textbf{Licença} & \textbf{Componentes (exemplos)} & \textbf{Forma de uso} \\
\midrule
\endfirsthead
\toprule
\textbf{Licença} & \textbf{Componentes (exemplos)} & \textbf{Forma de uso} \\
\midrule
\endhead
\bottomrule
\endfoot
\bottomrule
\endlastfoot
MIT & AURORA; Monaco; jQuery/jQuery-UI; KaTeX; Phosphor; netlistsvg; MCP SDK; Zod; web-tree-sitter; slang-server; gramáticas tree-sitter; Electron e a maioria das libs de build & Empacotamento direto. \\
BSD-2/3-Clause & Lit (3); DigitalJS, yosys2digitaljs (2); cocotb (3) & Empacotamento direto / processo externo (cocotb). \\
Apache-2.0 & AI SDK (\lstinline|ai|, \code{@ai-sdk/*}); Codex; TypeScript; Playwright; Verible; clang-format (com exceção LLVM) & Empacotamento direto / spawn (Verible, clang-format, Codex). \\
ISC & Yosys; Knip & Processo externo (Yosys) / dev. \\
LGPL-3.0 ou Artistic-2.0 & Verilator & Processo externo, sem modificação. \\
GNU GPL v2 & Icarus Verilog; GTKWave (\emph{fork}) & Processo externo separado, sem \emph{link}. \\
GNU GPL v3 & GCC/G++ (runtime com exceção), \code{ccache}, \code{make} & Toolchain no \emph{bundle}; runtime de compilação. \\
EUPL-1.2 & Surfer & Processo externo separado, sem \emph{link}. \\
PSF & Python & Interpretador no \emph{bundle}. \\
SIL OFL 1.1 & Inter; JetBrains Mono; Metamorphous; Noto Sans Runic & Embutimento de fonte. \\
Termos comerciais & Claude Code (Anthropic) & CLI on-demand; não empacotada. \\
\end{longtable}

% !TEX root = ../main.tex
\chapter{Build, empacotamento e distribuição}
\label{cap:16-build-distribuicao}

Este capítulo descreve, de ponta a ponta, como a \tool{AURORA} é construída a
partir do código-fonte e como é empacotada e distribuída ao usuário final. A
descrição segue o que está efetivamente codificado em \file{package.json}
(scripts e bloco \texttt{build}), \file{vite.config.mjs}, nos scripts de
\path{scripts/} e \path{components/Scripts/}, nos workflows de
\path{.github/workflows/}, no instalador \file{build/installer.nsh} e no
auto-updater em \path{main/updater.js}. Onde a documentação histórica do
repositório divergia do código, a descrição aqui reflete o código.

A cadeia completa tem cinco estágios encadeados: (1) \emph{bootstrap}, que baixa
o \emph{toolchain} e os artefatos não versionados; (2) compilação \texttt{tsc}
do processo principal e dos \emph{preloads}; (3) \emph{build} do \emph{renderer}
com Vite; (4) empacotamento com \tool{electron-builder} (alvo NSIS x64); e (5)
publicação no canal de \emph{release}, consumido pelo auto-updater. A
\autoref{sec:16-fluxo} traz o diagrama geral.

\section{Identidade do pacote e visão geral dos scripts}
\label{sec:16-identidade}

O \file{package.json} define a identidade do projeto e o conjunto de
\emph{scripts} npm que orquestram todo o ciclo. Os campos principais:

\begin{tabularx}{\linewidth}{l Y}
\toprule
\textbf{Campo} & \textbf{Valor} \\
\midrule
\texttt{name}        & \texttt{sapho} \\
\texttt{productName} & \texttt{SAPHO} \\
\texttt{version}     & \texttt{6.3.2} (no momento da leitura) \\
\texttt{main}        & \texttt{main.js} (ponto de entrada do processo principal) \\
\texttt{license}     & \texttt{MIT} \\
\texttt{repository}  & \lstinline|github.com/nipscernlab/aurora| \\
\texttt{engines.node}& \lstinline|>=18.0.0| \\
\bottomrule
\end{tabularx}

\noindent
Há uma distinção importante entre o \texttt{name}/\texttt{productName} do pacote
npm (\enquote{sapho}/\enquote{SAPHO}) e o \texttt{productName} do bloco
\texttt{build} do \tool{electron-builder} (também \enquote{SAPHO} desde
06/08/2026, antes \enquote{Aurora-IDE}): o primeiro identifica o projeto npm; o
segundo nomeia o aplicativo instalado.

Cada um desses campos aterrissa num lugar diferente do disco, e trocá-los move
estado que cópias já instaladas usam. O \texttt{name} é o mais delicado: dele
sai o diretório de cache do atualizador
(\path{%LOCALAPPDATA%/sapho-updater}), que guarda o \path{installer.exe} usado
como base da atualização incremental. Renomeá-lo faz cada máquina procurar essa
base num diretório inexistente, sem erro e sem aviso, e cair para download
completo. O \texttt{build.productName} move o diretório de instalação, então
mudá-lo depois de uma implantação instala a versão nova \emph{ao lado} da
antiga. E o \texttt{appId} precisa coincidir com o \texttt{setAppUserModelId}
do \path{main.js}, ou o Windows trata a janela em execução e o atalho fixado
como aplicativos distintos.

\subsection{Os scripts npm}
\label{sec:16-scripts}

A tabela a seguir lista os \emph{scripts} declarados e o que cada um executa.
Os \emph{hooks} \texttt{pre*} do npm disparam automaticamente antes do
\emph{script} de mesmo nome.

\begin{longtable}{l Y}
\toprule
\textbf{Script} & \textbf{Comando / efeito} \\
\midrule
\endhead
\texttt{bootstrap}      & encadeia o \emph{check} de versões pinadas, os 9 \emph{downloaders} e o \texttt{copy-components} (\autoref{sec:16-bootstrap}). \\
\texttt{build:ts}        & \lstinline|tsc| — compila TypeScript \emph{in-place} (processo principal e \emph{preloads}). \\
\texttt{build:renderer}  & \lstinline|vite build| — empacota o \emph{renderer} em \path{dist/} (\autoref{sec:16-vite}). \\
\texttt{predev}          & \lstinline|npm run build:ts| (roda antes de \texttt{dev}). \\
\texttt{dev}             & \lstinline|node scripts/dev.js| — Vite \emph{dev server} + Electron com HMR (\autoref{sec:16-dev}). \\
\texttt{prestart}        & \lstinline|bootstrap && build:ts && build:renderer|. \\
\texttt{start}           & \lstinline|node scripts/launch-electron.js| — inicia o Electron sobre o \path{dist/} já construído. \\
\texttt{test}            & \lstinline|vitest run| (testes unitários). \\
\texttt{pretest:e2e}     & \lstinline|npm run build:renderer|. \\
\texttt{test:e2e}        & \lstinline|vitest run --config vitest.config.e2e.js| (\emph{smoke} Electron). \\
\texttt{rules:sync}      & \lstinline|node scripts/sync-sapho-rules.js| (\autoref{sec:16-rules-sync}). \\
\texttt{prebuild}        & \lstinline|bootstrap && build:ts && build:renderer| (idêntico a \texttt{prestart}). \\
\texttt{build}           & \lstinline|electron-builder| — empacota o instalador (\autoref{sec:16-electron-builder}). \\
\texttt{prepare}         & \lstinline|husky| (instala os \emph{git hooks}). \\
\bottomrule
\end{longtable}

\begin{nota}
O \emph{hook} \texttt{prebuild} garante que \lstinline|npm run build| sempre
parte de um \emph{toolchain} baixado, de um \texttt{tsc} fresco e de um
\path{dist/} recém-gerado. Importante: os \emph{workflows} de CI e de
\emph{release} \textbf{não} invocam \lstinline|npm run build| --- eles chamam
\lstinline|electron-builder| diretamente, de modo que \texttt{prebuild} não
dispara nesses contextos; por isso os \emph{workflows} repetem os passos de
\emph{bootstrap}, \texttt{build:ts} e \texttt{build:renderer} explicitamente
(\autoref{sec:16-ci}).
\end{nota}

\section{O bootstrap}
\label{sec:16-bootstrap}

O \emph{toolchain} e a maior parte dos binários de terceiros \textbf{não são
versionados} no repositório \repo{aurora}; eles são baixados em tempo de
\emph{setup} pelo \texttt{npm run bootstrap}. O \emph{script} executa, em ordem:

\begin{enumerate}[nosep]
  \item \file{scripts/check-pinned-versions.js} --- verifica versões pinadas;
  \item nove \emph{downloaders} em \path{components/Scripts/};
  \item \file{components/Scripts/copy-components.js} --- liga a pasta
        \path{components/} ao Electron de desenvolvimento.
\end{enumerate}

\subsection{Verificação de versões pinadas}
\label{sec:16-pinned}

O \file{check-pinned-versions.js} é uma guarda contra a deriva de versões.
A convenção: uma dependência declarada em \file{package.json} com
\emph{semver} exato (sem \lstinline|^|, \lstinline|~| ou outro operador de
faixa) entra na checagem estrita; qualquer faixa fica fora. O regex de
reconhecimento é \lstinline|/^\d+\.\d+\.\d+(-[0-9A-Za-z.-]+)?$/|.

Para cada dependência pinada, o \emph{script} lê a versão realmente instalada em
\path{node_modules/<pkg>/package.json} e compara com o valor declarado; uma
divergência (ou pacote ausente) faz o \emph{script} sair com código 1. No estado
lido, o único pacote pinado é o \texttt{monaco-editor} (\texttt{0.52.2}).

O \emph{script} faz ainda duas verificações cruzadas relativas às CLIs de IA
(que \textbf{não} são empacotadas no instalador, mas baixadas sob demanda em
tempo de execução --- ver \autoref{cap:12-aurora-intelligence-ia}):

\begin{itemize}[nosep]
  \item compara as versões em \file{main/ai/cli\_manifest.js}
        (\texttt{CLAUDE\_VERSION}, \texttt{CODEX\_VERSION}) com as versões-base
        declaradas em \file{package.json} para \lstinline|@anthropic-ai/claude-code|
        e \lstinline|@openai/codex|;
  \item confere \texttt{integrity} e \texttt{tarball} do manifesto contra o
        \file{package-lock.json}, evitando que um \emph{bump} de versão fique com
        \emph{hash} de integridade defasado.
\end{itemize}

Este \emph{script} roda no \texttt{prestart}/\texttt{prebuild} (via
\texttt{bootstrap}) e também explicitamente no \emph{workflow} de CI logo após o
\lstinline|npm ci| (\autoref{sec:16-ci}).

\subsection{Os nove downloaders}
\label{sec:16-downloaders}

Cada \emph{downloader} é um \emph{script} Node autossuficiente em
\path{components/Scripts/}. Todos compartilham o mesmo padrão:

\begin{itemize}[nosep]
  \item \textbf{HTTPS com \emph{redirect}}: usam o módulo \lstinline|https|,
        seguindo redirecionamentos 301/302/307/308 até cinco saltos
        (os \emph{assets} de \emph{release} do GitHub redirecionam para um CDN);
  \item \textbf{\emph{sentinel}}: um arquivo-marcador prova a presença do pacote;
        se já existe (e sem \lstinline|--force|), o \emph{download} é pulado;
  \item \textbf{verificação de \emph{checksum}}: via
        \file{components/Scripts/lib/checksum.js} antes de extrair
        (\autoref{sec:16-checksum});
  \item \textbf{extração}: arquivos \texttt{.zip} são expandidos com
        \lstinline|powershell -NoProfile -Command "$ErrorActionPreference='Stop'; Expand-Archive ..."|
        (disponível em todo Windows 10+);
  \item \textbf{\emph{best-effort}}: em caso de falha, a maioria sai com código 0
        (com instruções de \emph{download} manual), para não bloquear o
        \lstinline|npm start| de um desenvolvedor \emph{offline} --- a
        funcionalidade correspondente fica indisponível até o \emph{setup}.
\end{itemize}

A tabela resume cada \emph{downloader}, sua origem e seu destino.

\begin{longtable}{l Y l}
\toprule
\textbf{Downloader / artefato} & \textbf{Origem (pinada)} & \textbf{Destino} \\
\midrule
\endhead
\texttt{download-toolchain.js} \newline (\file{aurora-msys-v1.zip}, \emph{tag} \texttt{msys-v1})
  & \emph{release} de \repo{aurora-toolchain}
  & \path{components/Packages/msys/} \\
\addlinespace
\texttt{download-yanc.js} \newline (\file{yanc-bin-v5.2.zip})
  & \emph{release} de \repo{yanc}
  & \path{components/} (gera \path{bin/}, \path{HDL/}, \path{Header/}, \path{Macros/}) \\
\addlinespace
\texttt{download-gtkwave-nipscern.js} \newline (\emph{tag} \texttt{v0.1.2-nipscern})
  & \emph{release} de \repo{gtkwave-nipscern}
  & \path{components/Packages/gtkwave-nipscern/} \\
\addlinespace
  & dafont (\lstinline|dl.dafont.com|)
  & \path{assets/fonts/} \\
\addlinespace
\texttt{download-surfer.js} \newline (\file{surfer\_win\_v0.7.0.zip})
  & registro genérico do GitLab do Surfer
  & \path{components/Packages/surfer/} \\
\addlinespace
\texttt{download-verible.js} \newline (\emph{LS} \texttt{v0.0-4080-...})
  & \emph{release} de \lstinline|chipsalliance/verible|
  & \path{components/Packages/verible/bin/} \\
\addlinespace
\texttt{download-clang-format.js} \newline (\file{clang-format-20...exe})
  & \lstinline|muttleyxd/clang-tools-static-binaries|
  & \path{components/Packages/clang-format/bin/} \\
\addlinespace
\texttt{download-slang-server.js} \newline (\emph{LS} \texttt{v0.2.7})
  & \lstinline|hudson-trading/slang-server|
  & \path{components/Packages/slang-server/bin/} \\
\addlinespace
\texttt{download-tree-sitter-grammars.js} \newline (3 \texttt{.wasm} + \emph{runtime})
  & \emph{releases} GitHub das gramáticas
  & \path{components/Packages/tree-sitter/} \\
\bottomrule
\end{longtable}

\subsubsection{Toolchain MSYS unificado}

O \file{download-toolchain.js} baixa o \emph{bundle} portátil
\file{aurora-msys-v1.zip} e o extrai em \path{components/Packages/msys/}. Esse
\emph{snapshot} MSYS2/mingw64 reúne, em um só lugar, toda a cadeia FOSS de
síntese e simulação:

\begin{itemize}[nosep]
  \item \path{mingw64/bin/} --- \tool{iverilog}, \tool{vvp}, \tool{yosys},
        \tool{verilator}, \tool{perl}, \texttt{g++}, \texttt{make},
        \texttt{ccache}, \texttt{python} e DLLs compartilhadas;
  \item \path{mingw64/lib/} --- \emph{backends} do iverilog (\path{ivl/}),
        \path{python3.12/} com \tool{cocotb} (VPI do Icarus
        \file{libcocotbvpi\_icarus.vpl} e o estático do Verilator
        \file{libcocotbvpi\_verilator.a}), \texttt{gcc/};
  \item \path{mingw64/share/} --- \tool{verilator}, \tool{yosys};
  \item \path{usr/bin/} --- \texttt{bash} + \emph{coreutils}.
\end{itemize}

\noindent
O \emph{sentinel} é \file{components/Packages/msys/mingw64/bin/verilator\_bin.exe}.
Há um segundo marcador, o \tool{cocotb} estático
(\file{...libs/libcocotbvpi\_verilator.a}, varrido sob qualquer
\path{python3.*}): um \emph{bundle} antigo que não o contenha dispara um
\emph{re-download} para atualizar no lugar. O \emph{bundle} é construído e
publicado pelo \emph{pipeline} dedicado do \repo{aurora-toolchain}; o
\tool{netlistsvg} \textbf{não} vem nele --- roda \emph{in-process} a partir de
\path{node_modules} (\lstinline|@silimate/netlistsvg|).

\subsubsection{Compiladores YANC}

O \file{download-yanc.js} baixa \file{yanc-bin-v5.2.zip} do \repo{yanc} e o
extrai em \path{components/} (de modo que \path{bin/cmmcomp.exe} etc.\ caiam em
\path{components/bin/}). O \texttt{.zip} contém os seis binários (desde a v4):
\texttt{cmmcomp.exe}, \texttt{cppcomp.exe}, \texttt{asmcomp.exe},
\texttt{cpppp.exe}, \texttt{appcomp.exe} e \texttt{comp2gtkw.exe}, além das
pastas \path{HDL/} (bibliotecas Verilog do SAPHO), \path{Macros/} (\emph{helpers}
de ponto flutuante \file{float\_*.asm}) e \path{Header/} (\emph{shims} de C++).

Este \emph{downloader} tem o controle de versão mais elaborado:

\begin{itemize}[nosep]
  \item \emph{sentinels} de presença: \file{components/bin/cppcomp.exe},
        \file{components/HDL/core.v}, \file{components/Header/cmath} e
        \file{components/Macros/float\_sin.asm} (um por pasta, pois HDL/Header/
        Macros não são mais versionados);
  \item \emph{sentinel} de versão: \file{components/bin/.yanc-version}, escrito
        após uma extração bem-sucedida e comparado contra \texttt{YANC\_TAG}.
        Só pula o \emph{download} quando \emph{ambos} valem (presença \emph{e}
        \emph{tag} correta) --- assim um \emph{bump} de \texttt{YANC\_TAG}
        alcança quem já tinha a versão anterior.
\end{itemize}

\subsubsection{Visualizadores de forma de onda e ferramentas de editor}

Os demais \emph{downloaders} buscam ferramentas que a \tool{AURORA} executa de
forma \emph{arm's-length} (invoca o \texttt{.exe}, não \emph{linka} a
biblioteca), o que evita contaminação de licença:

\begin{description}[nosep]
  \item[\tool{gtkwave-nipscern}] \emph{fork} portátil do GTKWave (com
        \texttt{gtkwave.exe} + \texttt{fst2vcd.exe} + DLLs do GTK), visualizador
        de formas de onda padrão.
  \item[\tool{Surfer}] visualizador \emph{opt-in} (alternável com o GTKWave),
        baixado do registro genérico do GitLab; licença EUPL-1.2. Tem
        \texttt{EXPECTED\_SHA256} pinado e \emph{flatten} de subpasta aninhada.
  \item[\tool{Verible}] apenas o \emph{language server}
        \file{verible-verilog-ls.exe} é mantido (os demais \textasciitilde{}10 executáveis do
        \texttt{.zip} são descartados); motor de diagnóstico/formatação/outline
        de Verilog (O2, LSP). Licença Apache-2.0.
  \item[\tool{clang-format}] \emph{build} estático do LLVM (um \texttt{.exe}
        avulso); formatação de C/C++/\cmm{} no editor. Licença Apache-2.0 com
        exceção LLVM.
  \item[\tool{slang-server}] \emph{LS} de SystemVerilog (análise semântica,
        O11); apenas o \texttt{.exe} é mantido. Licença MIT.
  \item[gramáticas \tool{tree-sitter}] três \texttt{.wasm}
        (SystemVerilog, C, C++) baixados com \texttt{SHA-256} por arquivo, mais
        o \emph{runtime} \file{web-tree-sitter.wasm} copiado de
        \path{node_modules} para casar exatamente com a versão do JS empacotado.
\end{description}

As fontes não são caso à parte: as quatro estão sob SIL OFL 1.1, que permite
redistribuir e embutir, e são versionadas como \code{.woff2}. Foram caso à parte
até 08/08/2026, quando o letreiro \enquote{Dagr} usava a fonte Norse, cuja
licença proíbe redistribuição; ver \autoref{cap:catalogo-terceiros}.

\subsection{Verificação de integridade (checksum)}
\label{sec:16-checksum}

O \file{components/Scripts/lib/checksum.js} expõe \lstinline|verifyChecksum()|.
Após o \texttt{.zip} chegar e \emph{antes} de extrair:

\begin{itemize}[nosep]
  \item se houver \texttt{EXPECTED\_SHA256} pinado, o \texttt{SHA-256} é
        recomputado e comparado; uma divergência \emph{lança erro} e o arquivo
        potencialmente adulterado nunca é extraído;
  \item se o \emph{hash} for \lstinline|null|, o \texttt{SHA-256} é apenas
        computado e registrado no \emph{log} (para um mantenedor pinar depois).
\end{itemize}

\noindent
No estado lido, têm \emph{hash} pinado e enforçado: Surfer, Verible,
clang-format, slang-server e as três gramáticas tree-sitter. O \emph{toolchain}
MSYS, o YANC e o gtkwave-nipscern usam \lstinline|EXPECTED_SHA256 = null|
(computa e registra, sem enforçar).

\subsection{Ligação dos componentes (copy-components)}
\label{sec:16-copy-components}

O último passo do \emph{bootstrap}, \file{components/Scripts/copy-components.js},
torna a pasta \path{components/} alcançável a partir do Electron de
desenvolvimento, no caminho \path{node_modules/electron/dist/components}.

Como o \emph{toolchain} pesa cerca de 1\,GB, o \emph{script} \textbf{prefere uma
junção de diretório} (\emph{junction}, via \lstinline|fse.ensureSymlink(..., 'junction')|)
em vez de copiar: não duplica, não fica desatualizado. Se já existir uma junção
apontando para a origem correta, nada é feito; se a junção não puder ser criada
(volume/permissões), há \emph{fallback} para cópia física com
\lstinline|fse.copy|.

\section{Build do renderer com Vite}
\label{sec:16-vite}

O \file{vite.config.mjs} é a configuração \emph{só do renderer}. O processo
principal e os \emph{preloads} \textbf{não} são empacotados pelo Vite --- ficam
em CommonJS cru, carregados diretamente pelo Electron; o \texttt{tsc} apenas os
compila \emph{in-place} (\autoref{sec:16-tsc}). O Vite é dono de
\file{index.html} e do seu grafo de módulos.

\subsection{base relativa para file://}

A diretiva \lstinline|base: './'| é obrigatória: o aplicativo empacotado carrega
\file{dist/index.html} por \texttt{file://}. A \texttt{base} padrão
(\lstinline|'/'|) emitiria URLs absolutas \lstinline|/assets/...| que, sob
\texttt{file://}, resolveriam à raiz do sistema de arquivos e dariam 404. Com
\lstinline|'./'|, toda URL de \emph{asset} é relativa ao HTML.

\subsection{Saída multi-página}

O \emph{build} é multi-página: a janela principal mais janelas secundárias.
O \lstinline|rollupOptions.input| declara cinco entradas:

\begin{tabularx}{\linewidth}{l Y}
\toprule
\textbf{Entrada} & \textbf{Arquivo-fonte} \\
\midrule
\texttt{index}       & \file{index.html} (janela principal) \\
\texttt{splash}      & \file{html/splash.html} \\
\texttt{update}      & \file{html/update-notification.html} \\
\texttt{prism}       & \file{html/prism/prism.html} (visualizador PRISM) \\
\texttt{design-lab}  & \file{html/design-lab.html} (galeria interna de componentes) \\
\bottomrule
\end{tabularx}

\noindent
Cada página é emitida no caminho relativo à fonte (\path{dist/index.html},
\path{dist/html/splash.html}, \dots) e tem suas referências de \emph{asset}
reescritas conforme a profundidade (graças a \lstinline|base: './'|). Outros
parâmetros de \emph{build}: \lstinline|outDir: 'dist'|,
\lstinline|emptyOutDir: true|, \lstinline|target: 'chrome130'| e
\lstinline|cssMinify: false| (o JS continua minificado; o CSS é mantido sem
minificar para preservar cada \texttt{@font-face} exatamente como na página
crua, evitando que a deduplicação/minificação altere a resolução de pesos das
fontes, já que o CSS vem de disco local/asar).

\subsection{Vendorização de Monaco, KaTeX e Phosphor}

Três dependências de grande porte são \enquote{vendorizadas} verbatim para
\path{dist/vendor/} pelo \emph{plugin} \lstinline|vite-plugin-static-copy|, em
vez de passar pelo grafo de \emph{import} do Vite (são módulos AMD/UMD que o Vite
não consegue analisar):

\begin{tabularx}{\linewidth}{l l}
\toprule
\textbf{Origem (\path{node_modules})} & \textbf{Destino} \\
\midrule
\path{monaco-editor/min/vs/}        & \path{dist/vendor/vs/} \\
\path{katex/dist/}                  & \path{dist/vendor/katex/dist/} \\
\path{@phosphor-icons/web/src/}     & \path{dist/vendor/phosphor/src/} \\
\bottomrule
\end{tabularx}

\noindent
A cópia usa \lstinline|rename.stripBase| para descartar os segmentos iniciais do
caminho e aterrissar as árvores exatamente nos caminhos que \file{index.html}
referencia. O mesmo \emph{plugin} ainda espelha para \path{dist/} três conjuntos
de recursos que o \emph{renderer} busca por caminho relativo ao documento em
tempo de execução (fora do grafo de \emph{import}): \path{locales/} (i18n),
\path{resources/} (incluindo \file{sapho\_rules.json}) e \path{assets/icons/}.

Um \emph{plugin} adicional, \lstinline|rewriteVendorPaths()|
(\lstinline|order: 'pre'|), reescreve no HTML os caminhos
\path{node_modules/...} para \path{vendor/...} antes da análise do Vite. O
\emph{logger} também é customizado para suprimir o aviso cosmético
\enquote{without type=module} dos \emph{loaders} AMD/UMD.

\subsection{Seleção da fonte do renderer em runtime}
\label{sec:16-render-loader}

O \file{main/render\_loader.js} centraliza, via \lstinline|loadPage(win, relPath)|,
a escolha de onde cada janela de primeira parte carrega seu HTML:

\begin{itemize}[nosep]
  \item \textbf{Dev (Vite)}: quando \texttt{AURORA\_RENDERER\_URL} está definido
        \emph{e} o app não está empacotado, carrega a página do \emph{dev
        server} (com HMR); o caminho é anexado à origem do servidor;
  \item \textbf{Produção/empacotado}: carrega a página de \path{dist/} por
        \texttt{file://} (\lstinline|win.loadFile(...)|).
\end{itemize}

\begin{nota}
\textbf{Não existe \emph{fallback} para a fonte crua.} Após a migração para Lit,
o \file{index.html} cru carrega \emph{imports} \lstinline|lit| que só resolvem
pelo \emph{bundler}; um \path{dist/} ausente é, portanto, \emph{erro de build} ---
o \emph{loader} mostra uma página de erro explícita em vez de carregar uma UI
quebrada. Todo fluxo real constrói \path{dist/} antes: \texttt{npm run dev}
(Vite), \texttt{npm start} (via \texttt{prestart}), o arnês e2e
(\texttt{pretest:e2e}) e o empacotamento.
\end{nota}

\subsection{Compilação TypeScript in-place}
\label{sec:16-tsc}

O \texttt{build:ts} executa \lstinline|tsc| (configurado por \file{tsconfig.json})
e compila o TypeScript do processo principal e dos \emph{preloads}
\emph{no lugar} (gerando \texttt{.js} ao lado dos \texttt{.ts}). Esses
arquivos não passam pelo Vite; o Electron os carrega como CommonJS. Em CI/release
há ainda a guarda \file{scripts/check-no-generated-js.js}, que falha o
\emph{build} se algum \texttt{.js} gerado for commitado por engano.

\subsection{Servidor de desenvolvimento (npm run dev)}
\label{sec:16-dev}

O \file{scripts/dev.js} sobe o \emph{dev server} do Vite na porta fixa
\texttt{5273} (com \lstinline|--strictPort|), espera-o responder por
\emph{polling} HTTP (até \textasciitilde{}30\,s) e então lança o Electron apontado para ele via
\texttt{AURORA\_RENDERER\_URL}. O Vite é \emph{spawnado} como filho Node direto
(via seu \emph{bin} \file{node\_modules/vite/bin/vite.js}), não por \texttt{npx}
ou \emph{shell}, o que dá encerramento limpo. Se o Vite morre, o Electron é
derrubado junto. Tanto \file{dev.js} quanto \file{scripts/launch-electron.js}
\textbf{removem} \texttt{ELECTRON\_RUN\_AS\_NODE} do ambiente antes de lançar o
Electron --- alguns \emph{shells} (terminal do VS Code, Claude Code) exportam
essa variável, o que faria o Electron iniciar como Node puro e falhar em
\lstinline|app.getAppPath()|.

\section{Empacotamento com electron-builder}
\label{sec:16-electron-builder}

O bloco \texttt{build} do \file{package.json} configura o \tool{electron-builder}
(invocado por \lstinline|npm run build|).

\subsection{Identidade do aplicativo e alvo}

\begin{tabularx}{\linewidth}{l Y}
\toprule
\textbf{Chave} & \textbf{Valor} \\
\midrule
\texttt{appId}        & \lstinline|com.nipscern.sapho| \\
\texttt{productName}  & \texttt{SAPHO} \\
\texttt{win.target}   & \texttt{nsis}, \texttt{arch: ["x64"]} \\
\texttt{win.icon}     & \path{./assets/icons/sapho_aurora_icon.ico} \\
\texttt{win.artifactName} & \path{sapho-aurora-Setup-v${version}.${ext}} \\
\texttt{directories.buildResources} & \path{build} \\
\bottomrule
\end{tabularx}

\noindent
O nome do artefato resultante é, portanto,
\file{sapho-aurora-Setup-vX.Y.Z.exe}.

\subsection{Associação de arquivo e protocolo}

\begin{description}[nosep]
  \item[fileAssociations] registra a extensão \texttt{.spf}
        (\enquote{SAPHO Project File}), papel \texttt{Editor}, com o ícone
        \file{sapho\_aurora\_icon.ico} --- dar duplo-clique em um \texttt{.spf}
        abre a \tool{AURORA};
  \item[protocols] registra o esquema de URL \lstinline|sapho://| (\enquote{SAPHO
        Project Protocol}), permitindo \emph{deep links} \lstinline|sapho://...|.
\end{description}

\subsection{Listas files e extraResources}
\label{sec:16-files}

A chave \texttt{files} define o que entra no \emph{asar} do aplicativo. Inclui
\lstinline|**/*| e \path{dist/**}, e \textbf{exclui} (prefixo \lstinline|!|) o
que não pertence ao pacote: \path{components/**} (vai como \emph{extraResource},
ver adiante), \path{tests/**}, \path{docs/**}, \path{scripts/**}, todos os
\texttt{*.md}, arquivos de configuração de \emph{release} e variantes pesadas do
Monaco (\path{node_modules/monaco-editor/dev|esm|min-maps/**}). Também exclui
explicitamente as CLIs de IA empacotáveis
(\lstinline|node_modules/@anthropic-ai/claude-code*| e
\lstinline|node_modules/@openai/codex*|), coerente com o modelo de
\emph{download} sob demanda.

A chave \texttt{extraResources} copia a pasta \path{components/} (o
\emph{toolchain} de \textasciitilde{}1\,GB) para \path{resources/components\_tmp} no aplicativo
instalado, em vez de embuti-la no \emph{asar}, aplicando um \emph{filter} que
remove projetos do usuário, temporários e \emph{backups}
(\path{Projeto/**}, \path{Temp/**}, \path{Backup/**}, \texttt{*.bak},
\texttt{*.aurora\_backup}, \texttt{*.spf.bak}, \texttt{*.log}).

\begin{nota}
O destino é \path{components\_tmp} (não \path{components}) de propósito: o
instalador NSIS o renomeia para \path{components} no passo de pós-instalação
(\autoref{sec:16-nsis}). Em produção, o processo principal procura os componentes
ao lado do executável: \file{main/paths.js} computa \texttt{componentsPath} como
\path{<dir do exe>/components} quando não está em \emph{dev}.
\end{nota}

\subsection{asarUnpack}

O empacotamento padrão coloca o código em um \emph{asar}; binários e recursos que
precisam existir como arquivos reais em disco (em vez de dentro do \emph{archive})
são desempacotados via \texttt{asarUnpack}. No estado lido, o \emph{toolchain}
viaja como \texttt{extraResource} (fora do \emph{asar} por construção), e o bloco
\texttt{build} não declara uma lista \texttt{asarUnpack} explícita adicional ---
os executáveis ficam em \path{resources/components/} graças ao
\texttt{extraResources}/renomeação NSIS.

\section{O instalador NSIS}
\label{sec:16-nsis}

A configuração \texttt{nsis} do \file{package.json}:

\begin{tabularx}{\linewidth}{l Y}
\toprule
\textbf{Opção} & \textbf{Valor} \\
\midrule
\texttt{oneClick}                       & \lstinline|true| (instalação de um clique) \\
\texttt{allowToChangeInstallationDirectory} & \lstinline|false| \\
\texttt{createDesktopShortcut}          & \lstinline|true| \\
\texttt{createStartMenuShortcut}        & \lstinline|true| \\
\texttt{include}                        & \path{build/installer.nsh} \\
\bottomrule
\end{tabularx}

\noindent
O \emph{script} customizado \file{build/installer.nsh} contém apenas uma macro de
pós-instalação que renomeia a pasta de componentes para o nome final:

\begin{lstlisting}[language=bash]
!macro customInstall
  Rename "$INSTDIR\resources\components_tmp" "$INSTDIR\components"
!macroend
\end{lstlisting}

\noindent
É esse passo que materializa o casamento entre o \texttt{extraResources}
(\path{components\_tmp}) e o caminho que \file{main/paths.js} espera em produção
(\path{<INSTDIR>/components}).

\section{Auto-updater}
\label{sec:16-updater}

A atualização automática usa \lstinline|electron-updater| e está implementada em
\file{main/updater.js}. O fluxo (apenas em produção --- é pulado em \emph{dev}):

\begin{enumerate}[nosep]
  \item cerca de 6\,s após a janela principal aparecer, faz uma verificação
        silenciosa nos \emph{releases};
  \item em \texttt{update-available}, abre a janela customizada
        \file{html/update-notification.html} (nunca um diálogo nativo) com o
        \emph{changelog} bilíngue e a opção de baixar;
  \item o \emph{download} é iniciado pelo usuário; \texttt{download-progress}
        alimenta a barra de progresso;
  \item em \texttt{update-downloaded}, a janela passa ao estado
        \enquote{Reiniciar e instalar};
  \item \lstinline|quitAndInstall(false, true)| instala e relança direto na nova
        versão.
\end{enumerate}

\subsection{Feed apontando para o canal de release}

O \emph{feed} é configurado explicitamente em
\lstinline|initializeUpdateSystem()|:

\begin{lstlisting}
autoUpdater.setFeedURL({
  provider: 'github',
  owner: 'nipscernlab',
  repo: 'sapho',
  releaseType: 'release',
});
\end{lstlisting}

\noindent
Note o destino: o repositório \repo{sapho} (não \repo{aurora}). É o canal de
distribuição estável (\autoref{sec:16-split}). As constantes
\texttt{REPO\_OWNER}/\texttt{REPO\_NAME} também são usadas no \emph{fallback} que
busca as notas de \emph{release} pela API do GitHub quando o
\lstinline|electron-updater| não as expõe. As opções
\lstinline|autoDownload = false| e \lstinline|autoInstallOnAppQuit = false|
garantem que nada é baixado ou instalado sem ação do usuário.

\subsection{latest.yml e blockmap}

O \tool{electron-builder} publica, junto do \texttt{.exe}, um \file{latest.yml}
(o manifesto que o \lstinline|electron-updater| lê: versão, caminho do arquivo,
\texttt{sha512}, tamanho) e um \texttt{.blockmap} (mapa de blocos para
atualizações diferenciais/\emph{delta}). O \lstinline|electron-updater| valida o
\texttt{sha512} do \file{latest.yml} contra os bytes baixados antes de instalar.

Após a primeira inicialização, o \file{main/updater.js} também persiste a versão
corrente em \file{aurora-version.json} (no \texttt{userData}); na inicialização
seguinte, uma divergência entre a versão armazenada e a atual indica que o
usuário acabou de atualizar, e o \emph{renderer} surge um \emph{toast} de
confirmação (uma única vez por atualização).

\section{Workflows de CI e de release}
\label{sec:16-ci}

\subsection{CI (ci.yml)}

O \emph{workflow} \file{.github/workflows/ci.yml} roda em \texttt{push} e
\emph{pull request} para \texttt{main}, no \texttt{windows-latest}. Após
\lstinline|npm ci|, executa, em sequência:

\begin{enumerate}[nosep]
  \item \file{check-pinned-versions.js} (versões pinadas);
  \item ESLint (\lstinline|--max-warnings=0|);
  \item \emph{typecheck} (\lstinline|tsc --noEmit|);
  \item guarda de \texttt{.js} gerado (\file{check-no-generated-js.js});
  \item consistência de i18n (\file{check-i18n.js});
  \item \emph{dead code} (\lstinline|npm run deadcode|);
  \item \lstinline|node --check| em todo \texttt{.js} de \path{js/} e \path{main/};
  \item testes unitários com cobertura + \emph{upload} ao Codecov
        (informativo, nunca falha o CI);
  \item testes e2e (\emph{smoke} Electron);
  \item \texttt{build:renderer} (Vite);
  \item \emph{build} de \emph{smoke} com
        \lstinline|electron-builder --win --x64 --publish never|.
\end{enumerate}

\noindent
O CI invoca \lstinline|electron-builder| diretamente (não por \texttt{npm run
build}), por isso constrói o \emph{renderer} explicitamente antes --- caso
contrário o \emph{smoke build} validaria um \path{dist/} obsoleto/ausente.

\subsection{Release (release.yml)}

O \file{.github/workflows/release.yml} dispara quando um \emph{release} do GitHub
é \emph{publicado} (e tem \texttt{workflow\_dispatch} como gatilho manual). No
\texttt{windows-latest}, com permissão \lstinline|contents: write|, executa:

\begin{enumerate}[nosep]
  \item \lstinline|npm ci|;
  \item \lstinline|npm run bootstrap| (baixa todo o \emph{toolchain});
  \item \textbf{verificação de \emph{sentinels}}: ao contrário dos
        \emph{downloaders} (que saem 0 em falha para não bloquear o
        desenvolvedor), o \emph{release} \textbf{falha alto} se faltar qualquer
        \emph{sentinel}
        (\file{.../verilator\_bin.exe}, \file{components/bin/cppcomp.exe},
        \file{.../gtkwave.exe}, \file{.../surfer.exe}) --- um \emph{release}
        nunca pode sair sem o \emph{toolchain};
  \item \texttt{build:ts} e \texttt{build:renderer};
  \item \lstinline|electron-builder --win --x64 --publish always|, usando
        \texttt{GH\_TOKEN} para criar o \emph{release} e subir os \emph{assets}.
\end{enumerate}

\subsection{release-please}

O \file{.github/workflows/release-please.yml} roda em \texttt{push} para
\texttt{main} e usa a \lstinline|googleapis/release-please-action@v4|. Ele mantém
um \emph{pull request} permanente de \emph{release} que agrega todos os
\emph{Conventional Commits} desde o último \emph{release}; ao ser mesclado, gera
o \emph{commit} de \emph{version bump}, cria a \emph{tag}, atualiza o
\file{CHANGELOG.md} e publica o \emph{release} do GitHub --- o que, por sua vez,
dispara o \file{release.yml}.

A configuração vive em \file{release-please-config.json}
(\lstinline|release-type: node|, \lstinline|include-v-in-tag: true|,
\lstinline|bump-minor-pre-major: true| e o mapeamento de seções do
\emph{changelog} por tipo de \emph{commit}) e em
\file{.release-please-manifest.json}, que registra a versão corrente
(\lstinline|{".": "6.3.2"}|).

\section{Assinatura de código}
\label{sec:16-signing}

Segundo o \file{docs/CODE\_SIGNING.md} e as configurações lidas, o instalador
NSIS e o auto-updater são \textbf{atualmente não assinados}: o SmartScreen do
Windows pode alertar (\enquote{Windows protegeu seu PC}) na primeira execução, e
o \emph{updater} verifica apenas o \texttt{sha512} do \file{latest.yml}. O passo
\enquote{Build \& publish} do \file{release.yml} traz um comentário registrando
esse estado e apontando para o \emph{runbook}.

O \emph{runbook} documenta o caminho recomendado: como o Aurora-IDE é MIT, ele se
qualifica ao programa gratuito da \textbf{SignPath Foundation} (OSS). O fluxo de
CI proposto (ainda não aplicado ao \file{release.yml}) é \emph{gated} na
existência do \emph{secret} \texttt{SIGNPATH\_API\_TOKEN}: enquanto ele não
existir, o passo de \emph{build} não assinado roda sem alteração; quando
presente, o \emph{build} gera o instalador sem publicar, o envia para assinatura
e então republica.

Como a assinatura externa (SignPath) devolve \emph{novos bytes}, o
\file{latest.yml} precisa ser re-\emph{hasheado}, ou todo auto-update falharia.
Esse re-\emph{hash} está pronto em \file{scripts/patch-latest-yml.js}: ele
recomputa \texttt{sha512} (base64) e tamanho do \texttt{.exe} assinado, reescreve
as entradas correspondentes em \file{latest.yml} e remove o \texttt{.blockmap}
defasado (forçando o \emph{updater} a um \emph{download} completo, sempre válido,
em vez de um \emph{delta} quebrado). Há ainda a alternativa de assinar
\emph{durante} o \emph{build} (\emph{hook} \lstinline|win.sign| do
\tool{electron-builder}, ou Azure Trusted Signing via
\texttt{CSC\_LINK}/\texttt{CSC\_KEY\_PASSWORD}), que dispensaria o re-\emph{hash}.

\section{Split intencional de repositórios}
\label{sec:16-split}

A separação entre \repo{aurora} e \repo{sapho} é \textbf{intencional}, não um
defeito:

\begin{description}
  \item[\repo{aurora}] repositório de \textbf{desenvolvimento} da IDE. É onde o
        código vive, onde a CI roda e onde o \tool{electron-builder} executa.
  \item[\repo{sapho}] \textbf{canal de distribuição estável}. O instalador
        \file{sapho-aurora-Setup-vX.Y.Z.exe} e o \file{latest.yml} são
        publicados como \emph{assets} de \emph{release} \textbf{do repositório
        sapho}.
\end{description}

\noindent
Isso é coerente em todo o código: o bloco \texttt{publish} do \file{package.json}
aponta \lstinline|owner: nipscernlab|, \lstinline|repo: sapho|; o auto-updater em
\file{main/updater.js} configura o \emph{feed} para
\lstinline|owner: 'nipscernlab', repo: 'sapho'|; e o \emph{runbook} de assinatura
nota que publicar de \texttt{aurora} para \texttt{sapho} é uma operação
\emph{cross-repo} (exige um \emph{token} com \lstinline|contents:write| em
\texttt{sapho}). Em resumo: o usuário final baixa de \repo{sapho} e recebe
atualizações de \repo{sapho}; toda a engenharia acontece em \repo{aurora}.

\section{Fluxo de build (visão geral)}
\label{sec:16-fluxo}

O diagrama a seguir resume a cadeia, do código-fonte ao usuário final.

\begin{lstlisting}
                     repo: nipscernlab/aurora  (DESENVOLVIMENTO)
  +---------------------------------------------------------------+
  |  npm ci                                                       |
  |    |                                                          |
  |    v                                                          |
  |  npm run bootstrap                                            |
  |    |-- check-pinned-versions.js  (monaco 0.52.2; manifest IA) |
  |    |-- download-toolchain (aurora-msys-v1.zip) -> Packages/msys|
  |    |-- download-yanc (yanc-bin-v5.2.zip)       -> components/  |
  |    |-- download-gtkwave-nipscern               -> Packages/   |
  |    |-- download-norse-font (dafont)            -> assets/fonts|
  |    |-- download-surfer / verible / clang-format / slang-server|
  |    |-- download-tree-sitter-grammars (3x .wasm + runtime)     |
  |    +-- copy-components (junction -> electron/dist/components)  |
  |    |                                                          |
  |    v                                                          |
  |  build:ts  (tsc)         ----> main/preload .js in-place      |
  |  build:renderer (vite)   ----> dist/ (index + 4 paginas;      |
  |                                 vendor/ Monaco/KaTeX/Phosphor) |
  |    |                                                          |
  |    v                                                          |
  |  electron-builder --win --x64                                 |
  |    |-- asar (files: !components !tests !docs ...)             |
  |    |-- extraResources: components -> resources/components_tmp |
  |    |-- NSIS (oneClick) + installer.nsh                        |
  |    |     customInstall: rename components_tmp -> components   |
  |    +-- artefatos: sapho-aurora-Setup-vX.Y.Z.exe               |
  |                   latest.yml + .blockmap                      |
  +-------------------------------|-------------------------------+
            release-please PR     |  --publish always (cross-repo)
            (Conventional Commits)|
                                  v
                     repo: nipscernlab/sapho  (RELEASE ESTAVEL)
  +---------------------------------------------------------------+
  |  Release assets: Setup .exe + latest.yml + .blockmap          |
  +-------------------------------|-------------------------------+
                                  |  electron-updater (feed: sapho)
                                  v
                          USUARIO FINAL  (instala / auto-update)
\end{lstlisting}

\noindent
Para o detalhamento do \emph{toolchain} de simulação/síntese baixado no
\emph{bootstrap}, ver \autoref{cap:14-toolchain-bundle}; para o subsistema de IA
cujas CLIs são baixadas sob demanda (e não empacotadas), ver
\autoref{cap:12-aurora-intelligence-ia}.

% !TEX root = ../main.tex
\chapter{O pipeline de criação ponta a ponta}
\label{cap:17-pipeline-ponta-a-ponta}

Os capítulos anteriores descreveram cada subsistema da plataforma SAPHO de forma
isolada: a arquitetura do soft-processor (\autoref{cap:02-arquitetura-soft-processor-sapho}),
a linguagem \cmm{} (\autoref{cap:03-linguagem-cmm}), os compiladores \tool{YANC}
(\autoref{cap:04-yanc-compiladores}), o motor de compilação da AURORA
(\autoref{cap:08-aurora-compilacao}), a simulação e instrumentação de ondas
(\autoref{cap:09-aurora-simulacao-ondas}), os visualizadores
(\autoref{cap:10-aurora-visualizacao-gtkwave-surfer}) e o \tool{PRISM}
(\autoref{cap:11-aurora-prism-rtl}). Este capítulo é de \emph{síntese}: ele costura
todos esses subsistemas em um único fio condutor, do projeto vazio ao resultado
visualizado, mostrando \emph{quais artefatos} cada etapa produz e \emph{quem invoca
o quê}.

A descrição é construída sobre o que o código realmente faz. Do lado da AURORA, o
orquestrador é o renderer (\file{js/compilation/compilation\_flow.js} e
\file{js/compilation/compilation\_module.js}, apoiados em
\file{js/compilation/processor\_compiler.js}); a execução dos binários acontece no
processo \emph{main}, atrás da \emph{allowlist} de
\file{main/compile/binary\_allowlist.js}. Os mesmos passos, em forma de script de
referência, existem no repositório \tool{YANC} em
\file{Scripts/single\_proc.sh} (um processador) e \file{Scripts/multi\_proc.sh}
(vários processadores sob um top-level). Onde a documentação anterior e o código
divergem, vale o código.

\section{Visão geral: o fio condutor}

O caminho de criação tem sete estágios. Cada um consome os artefatos do anterior e
produz artefatos novos; nenhum estágio precisa adivinhar nada que o anterior não
tenha gravado em disco ou no \file{.spf}.

\begin{enumerate}[leftmargin=2.2em]
  \item \textbf{Projeto e processadores} — cria-se o arquivo de projeto \file{.spf}
        e a estrutura de pastas dos processadores na AURORA.
  \item \textbf{Código-fonte} — escreve-se o programa em \cmm{} (\file{.cmm}) ou
        em C++ (\file{.cpp}), além do Verilog de top-level e do testbench quando
        houver.
  \item \textbf{Compilação} — \tool{cmmcomp} gera o \file{.asm}; \tool{appcomp}
        pré-processa; \tool{asmcomp} gera o \file{.v} sintetizável, as imagens de
        memória e o testbench padrão \file{<proc>\_tb.v}.
  \item \textbf{Simulação} — Icarus (\tool{iverilog}\,+\,\tool{vvp}), \tool{Verilator}
        ou \tool{cocotb} (Python) elabora o design e produz o \emph{dump} de ondas.
  \item \textbf{Instrumentação do dump} — o testbench recebe \lstinline|$dumpfile|/%
        \lstinline|$dumpvars| de acordo com a seleção de sinais escolhida.
  \item \textbf{Visualização de ondas} — \tool{GTKWave} (fork) ou \tool{Surfer}
        abre o waveform com as faixas de \cmm{}, assembly e variáveis em
        \emph{lockstep} com o clock.
  \item \textbf{Inspeção do RTL} — o \tool{PRISM} converte o design via \tool{Yosys}
        e \tool{netlistsvg} em um diagrama navegável.
\end{enumerate}

\begin{nota}
Os estágios 3 a 6 não têm desvio estrutural ``tem processador?''. Quando o
\file{.spf} não lista nenhum processador (projeto Verilog puro), os laços sobre o
array de processadores são no-op e o pipeline segue idêntico — o cabeçalho de
\file{compilation\_module.js} documenta essa decisão de \emph{pipeline único}. Por
isso este capítulo descreve o caminho com processador, que é o caso mais completo;
o caso puro é o mesmo fluxo com os laços vazios.
\end{nota}

\section{Quem invoca o quê: as duas fronteiras}
\label{sec:17-fronteiras}

Antes do passo a passo, é preciso fixar a topologia de invocação, porque ela se
repete em todos os estágios que tocam binários.

\subsection{Fronteira 1: orquestração no renderer}

Os botões da UI (Verilog, Wave, ASM, PRISM) entram por
\file{compilation\_flow.js}. Esse módulo prepara a lista de processadores a
compilar e delega para um \code{CompilationModule} (instanciado por compilação).
Os métodos públicos do \code{CompilationModule} são o ``backend'' de cada etapa:

\begin{itemize}[leftmargin=1.4em]
  \item \code{loadConfig()} lê o \file{.spf} para \code{this.projectConfig} via
        \code{SpfStore.read}.
  \item \code{cmmCompilation(proc)} e \code{asmCompilation(proc)} delegam para
        \file{processor\_compiler.js} (passos por processador).
  \item \code{verilogSyntaxCheck()} roda o \emph{check} de sintaxe
        (\lstinline|iverilog -tnull|) e regenera a hierarquia via \tool{Yosys}.
  \item \code{waveBuildVvp()} compila o \file{.vvp} com o testbench instrumentado.
  \item \code{runGtkWave()} orquestra todo o fluxo de ondas em fases
        privadas \lstinline|_wave*|.
\end{itemize}

A função \code{precompileAllProcessors} (em \file{compilation\_flow.js}) é o ponto
comum: antes de qualquer simulação ou check, ela percorre os processadores do
\file{.spf} e chama \code{cmmCompilation} + \code{asmCompilation} em cada um. Os
defaults numéricos de cada processador (\lstinline|clk = 100| e
\lstinline|numClocks = 2000|) ficam embutidos em \code{PROC\_DEFAULTS} (objeto
congelado em \file{compilation\_flow.js}) e são aplicados quando o \file{.spf} não
os define; o nome do \file{.cmm}, por sua vez, deriva do nome do processador
(\lstinline|`${proc.name}.cmm`|).

\subsection{Fronteira 2: execução no main atrás da allowlist}

Nenhum binário é executado diretamente pelo renderer. Cada etapa monta uma
\emph{CommandSpec} (um \code{builder} em \file{js/compilation/builders/}), e a
\code{runSpec} (\file{js/compilation/spec\_runner.js}) é o \emph{único} ponto de
passagem: ela aplica os \emph{overrides} persistidos/efêmeros (o canal pelo qual a
Aurora Intelligence ajusta linhas de comando, \autoref{cap:12-aurora-intelligence-ia}),
loga a linha resultante e dispara o IPC \lstinline|exec-spec| / \lstinline|exec-spec-streamed|.

Do lado do main, a spec é \emph{re-validada}: o campo \code{binary} precisa bater
com o \emph{basename} e com um dos diretórios legais listados em
\file{main/compile/binary\_allowlist.js}. Essa tabela é a fronteira de confiança —
um renderer comprometido não consegue spawnar um executável fora dela. A tabela
inclui \file{cmmcomp.exe}, \file{appcomp.exe}, \file{asmcomp.exe} e
\file{comp2gtkw.exe} (em \path{components/bin}); \file{iverilog.exe},
\file{vvp.exe}, \file{verilator.exe}, \file{yosys.exe}, \file{g++.exe},
\file{make.exe} e \file{perl.exe} (em \path{Packages/msys/mingw64/bin});
\file{gtkwave.exe} e \file{fst2vcd.exe} (do fork \tool{gtkwave-nipscern}); e
\file{surfer.exe}. O \tool{netlistsvg} do \tool{PRISM} não está na tabela porque
roda \emph{in-process} (não é um \file{.exe} bundled).

\begin{nota}
Resumo da regra: \textbf{o renderer orquestra; o main executa.} Toda etapa que
chama um compilador, simulador ou visualizador passa por
\code{runSpec} $\rightarrow$ \lstinline|exec-spec| $\rightarrow$ allowlist. Detalhes
de cada \code{builder} e do executor estão em \autoref{cap:08-aurora-compilacao}.
\end{nota}

\section{Estágio 1 — Criar o projeto e os processadores}
\label{sec:17-projeto}

O artefato raiz é o arquivo de projeto \file{.spf} (\emph{SAPHO Project File}). Ele
é a fonte canônica única de configuração: as listas \code{synthesizableFiles} e
\code{testbenchFiles}, o \code{topLevelFile}/\code{testbenchFile}, e o array
\code{processors} (cada um com \code{name}, \code{cmmFile}, \code{clk},
\code{numClocks}, \code{showArrays}). Formatos legados (\path{projectOriented.json},
\path{processorConfig.json}) foram consolidados no \file{.spf} — ver o cabeçalho de
\file{compilation\_module.js} e \autoref{cap:07-aurora-projeto-arvore}.

Para cada processador, a estrutura de pastas é \path{<proj>/<proc>/Software/}
(onde mora o \file{.cmm}) e \path{<proj>/<proc>/Hardware/} (onde o \file{.v}
gerado será gravado). A AURORA mantém ainda a área de \emph{scratch}
\path{components/Temp/<proc>/}, criada por \code{ensureDirectories} /
\code{createDirectory}, que serve de diretório temporário (\lstinline|-t|) dos
compiladores e de CWD da simulação. A árvore de arquivos e a leitura/escrita do
\file{.spf} são detalhadas em \autoref{cap:07-aurora-projeto-arvore}.

\section{Estágio 2 — Escrever o código-fonte}
\label{sec:17-codigo}

O insumo primário é o programa em \cmm{} (\file{<proc>.cmm}), descrito em
\autoref{cap:03-linguagem-cmm}. Alternativamente pode-se escrever C++
(\file{.cpp}), que o \tool{YANC} também aceita como entrada
(\autoref{cap:04-yanc-compiladores}). O editor Monaco (\autoref{cap:06-aurora-editor})
fornece destaque, formatação e os \emph{language servers}.

Além do \file{.cmm}, um projeto típico tem: o Verilog de top-level marcado como
\code{isTopLevel} em \code{synthesizableFiles}, e um testbench. O testbench pode
ser o padrão auto-gerado pelo \tool{asmcomp} (\file{<proc>\_tb.v}, que exercita um
único processador isolado) ou um testbench de usuário (por exemplo
\file{top\_level\_tb.v}, que instancia o top-level com vários processadores).

\begin{nota}
\textbf{Instrumentação \lstinline|#TOAQUI| (botão Verilator).} Quando a simulação
escolhida é Verilator, \code{ensureChegueiToaqui} (em
\file{processor\_compiler.js}) garante que o \file{.cmm} tenha a diretiva
\lstinline|#TOAQUI| antes do \lstinline|}| de \code{main()} — sem ela, o pino
\code{cheguei} não vira porta do \file{<proc>.v} e o harness não detecta o fim do
programa. A operação é idempotente e ocorre \emph{depois} de salvar os buffers,
sincronizando o editor para que um Ctrl+S posterior não derrube a instrumentação.
\end{nota}

\section{Estágio 3 — Compilar: \texorpdfstring{\cmm{}}{C+-} até Verilog}
\label{sec:17-compilar}

Este é o coração do pipeline. São três binários do \tool{YANC} em sequência,
um por processador, com a AURORA salvando os buffers do editor antes de cada
passo (\code{TabManager.saveAllFiles()}).

\subsection{3a — \texttt{cmmcomp}: \texorpdfstring{\cmm{}}{C+-} → assembly}

\code{cmmCompilation} (\file{processor\_compiler.js}) monta a spec via
\code{buildCmmSpec} e roda \file{components/bin/cmmcomp.exe}. As opções nomeadas
do \tool{YANC} v4 são \lstinline|-i input -n name -p proc-dir -m macros-dir -t temp-dir|,
com \lstinline|-A|/\lstinline|--array| habilitando o dump de arrays
(\code{showArrays} do \file{.spf}) e a flag de locale (\lstinline|-pt|/\lstinline|-en|)
indo \emph{primeiro}. O script de referência mostra a mesma chamada:

\begin{lstlisting}[language=bash]
"$BIN_DIR/cmmcomp" -i "$FNAM" -n "$PROC" -p "$PROC_DIR" -m "$MAC_DIR" -t "$TMP_PRO"
\end{lstlisting}

\textbf{Artefatos produzidos:} \file{<proj>/<proc>/Software/<base>.asm} (o assembly
SAPHO), as imagens de memória do processador (\file{<proc>\_data.mif} /
\file{<proc>\_inst.mif} no fluxo single, e \file{pc\_<proc>\_mem.txt} no
\path{Temp/<proc>/}) e o \file{cmm\_log.txt}. A AURORA cacheia o caminho do
\file{.cmm} corrente (\code{lastCompiledCmmPath}) \emph{antes} de rodar, para que
o clique em ``linha N'' no terminal resolva o arquivo mesmo após falha.

\subsection{3b — \texttt{appcomp}: pré-montagem}

\code{asmCompilation} chama primeiro \file{components/bin/appcomp.exe} (spec por
\code{buildAsmPreSpec}), que pré-processa o assembly resolvendo parâmetros e
endereços. As opções são \lstinline|-i input -t temp-dir|:

\begin{lstlisting}[language=bash]
"$BIN_DIR/appcomp" -i "$ASM_FILE" -t "$TMP_PRO"
\end{lstlisting}

\subsection{3c — \texttt{asmcomp}: assembly → Verilog sintetizável}

Em seguida \code{asmCompilation} roda \file{components/bin/asmcomp.exe} (spec por
\code{buildAsmSpec}), com as opções nomeadas
\lstinline|-i -p -d -m -t -f -c|, onde \lstinline|-f| é a frequência de clock
(\code{clk}) e \lstinline|-c| o número de clocks (\code{numClocks}) — ambos
obrigatoriamente inteiros:

\begin{lstlisting}[language=bash]
"$BIN_DIR/asmcomp" -i "$ASM_FILE" -p "$PROC_DIR" -d "$HDL_DIR" \
    -m "$MAC_DIR" -t "$TMP_PRO" -f "$FRE_CLK" -c "$NUM_CLK"
\end{lstlisting}

\textbf{Artefatos produzidos:} \file{<proj>/<proc>/Hardware/<proc>.v} (o RTL
sintetizável do processador), as imagens de memória \file{.mif} e o testbench
padrão \file{<proc>\_tb.v}. Quando o processador usa o testbench ``standard''
(sem testbench customizado no \file{.spf}), a AURORA copia o \file{<base>\_tb.v}
gerado de \path{Temp/<proc>/} para \path{<proj>/<proc>/Simulation/}.

\begin{nota}
\textbf{Biblioteca SAPHO via \lstinline|-y|.} O \file{<proc>.v} gerado instancia
módulos da biblioteca SAPHO (\file{processor.v}, \file{core.v}, \file{ula.v},
\file{addr\_dec.v}, \file{instr\_dec.v}, \file{myFIFO.v}) que vivem em
\path{components/HDL}. Em todos os passos seguintes que chamam \tool{iverilog} ou
\tool{Yosys}, a AURORA acrescenta \path{components/HDL} como diretório de busca
(\lstinline|-y| para o iverilog, \code{read\_verilog} para o Yosys) para resolver
esses módulos sem o usuário precisar listá-los. Ver
\autoref{cap:02-arquitetura-soft-processor-sapho}.
\end{nota}

\section{Estágio 4 — Simular}
\label{sec:17-simular}

Com o \file{<proc>.v} e o testbench prontos, a simulação elabora o design e produz
o dump de ondas. A escolha do simulador vem de \code{getSimulator()}
(preferência persistida; \tool{iverilog} é o default, \tool{Verilator} é opt-in) e
o tipo de testbench decide se o caminho é HDL ou \tool{cocotb} (Python). O
orquestrador é \code{runGtkWave()} em \file{compilation\_module.js}; a etapa de
pré-compilação dos processadores (\code{precompileAllProcessors}) já rodou antes.

\subsection{4a — Icarus (iverilog + vvp)}

\code{waveBuildVvp} monta a spec via \code{buildIverilogBuildSpec} e compila o
\file{.vvp} a partir dos \code{synthesizableFiles} + testbench instrumentado, com
\lstinline|-s <tbModule>| (top de simulação) e \lstinline|-y components/HDL|. A
fase \code{\_waveRunVvpSimulation} então roda \tool{vvp} com CWD em
\path{components/Temp/}, depois de \code{stageProcessorMemoryFiles} ter copiado os
\file{pc\_<proc>\_mem.txt} para a raiz desse diretório (é onde o
\lstinline|$readmemb| do \file{<proc>.v} procura). A linha de referência:

\begin{lstlisting}[language=bash]
"$IVERILOG" -s "$TB_MOD" -o "$TMP_PRO/$PROC.vvp" \
    "$TMP_PRO/${PROC}_tb.v" "$UPROC.v" addr_dec.v instr_dec.v processor.v core.v ula.v
"$VVP" "$PROC.vvp" -fst
\end{lstlisting}

\textbf{Artefato produzido:} o dump de ondas (FST no fluxo da AURORA).

\subsection{4b — Verilator}

Quando \code{getSimulator()} retorna \code{'verilator'}, \code{\_waveBuildVerilator}
(specs por \code{buildVerilatorBuildSpec} / \code{buildVerilatorTbBuildSpec}) faz a
verilation + \tool{g++} num executável nativo sob \path{components/Temp/}, e
\code{\_waveRunVerilatorSimulation} o executa. O script de referência usa
\lstinline|--binary --timing --trace +define+YANC_TRACE| reaproveitando o
\file{<proc>\_tb.v} como top.

\subsection{4c — cocotb (Python)}

Se o testbench é um \file{.py} (\code{isPythonFile}), \code{runGtkWave} entra no
ramo cocotb: \code{\_waveValidateCocotbConfig} resolve o módulo HDL de topo (via
diretiva \lstinline|TOPLEVEL| no \file{.py} ou fallback para o top-level do
\file{.spf}, com aviso), e \code{\_waveRunCocotbSimulation} (spec por
\code{buildCocotbRunSpec}) roda o teste usando, por baixo, \tool{iverilog} ou
\tool{Verilator}. Detalhes em \autoref{cap:09-aurora-simulacao-ondas}.

\section{Estágio 5 — Instrumentar o dump}
\label{sec:17-instrumentar}

Antes de spawnar o simulador, a AURORA decide \emph{quais sinais} entram no dump e
\emph{como}. \code{\_resolveWaveSelection} (em \file{wave\_signal\_validator.js})
aplica a precedência: \textbf{.gtkw ativo} > \textbf{Wave Config (seleção do
usuário)} > \textbf{testbench com \lstinline|$dumpvars| escrito à mão} >
\textbf{default} (\lstinline|$dumpvars(1, <tb>)|). O resultado é
\code{\{ signalsToDump, overrideUserDumpvars, source \}}.

\code{instrumentTestbench} então lê o testbench original e, via
\code{instrumentTestbenchSource} (\file{testbench\_instrumenter.js}), injeta
\lstinline|$dumpfile|/\lstinline|$dumpvars| conforme a seleção, gravando uma
\emph{cópia} em \path{components/Temp/instr_<tb>.v}. \textbf{O testbench original
nunca é tocado.} A escrita é idempotente por conteúdo (só regrava se mudou), o que
evita recompilações desnecessárias no caminho do Verilator. O conjunto final de
fontes passado ao simulador é \code{synthesizableFiles} + esse testbench
instrumentado.

\section{Estágio 6 — Visualizar as ondas}
\label{sec:17-visualizar}

Concluída a simulação, \code{runGtkWave} faz uma captura única de cabeçalho:
\code{\_extractFstHeaderVcd} usa \tool{fst2vcd} para extrair, do FST, o VCD de
cabeçalho (escopos e sinais) que alimenta o seletor de sinais e o gerador
automático de layout. Em seguida, o ramo depende de \code{getViewer()}:

\begin{itemize}[leftmargin=1.4em]
  \item \textbf{GTKWave} (default): \code{\_waveResolveGtkwSaveFile} resolve ou
        gera o arquivo de layout \file{.gtkw} (via \code{buildAuroraGtkw} em
        \file{gtkw\_proc\_writer.js}), e \code{\_waveLaunchGtkwave} (spec por
        \code{buildGtkwaveSpec}) abre o \tool{gtkwave.exe} do fork. O processo é
        monitorado por \code{monitorGtkwaveProcess}.
  \item \textbf{Surfer} (opt-in): \code{\_waveResolveSurferSaveFile} gera o layout
        (\code{buildSurferLayout} em \file{surfer\_layout\_writer.js}) e
        \code{\_waveLaunchSurfer} abre o \tool{surfer.exe}.
\end{itemize}

O diferencial do layout SAPHO é mostrar as faixas de \cmm{}, assembly e variáveis
em \emph{lockstep} com o clock do processador — o usuário vê o programa executando
linha a linha junto às ondas. O script de referência delega a construção do layout
ao \tool{gen\_gtkw} do \tool{YANC}; na AURORA, os writers
\file{gtkw\_proc\_writer.js}/\file{surfer\_layout\_writer.js} desempenham esse
papel. Toda a mecânica de layout, decodificação de complexos
(\file{complex\_decode.js}) e marcadores de evento (\file{event\_markers.js}) está
em \autoref{cap:10-aurora-visualizacao-gtkwave-surfer}.

\section{Estágio 7 — Inspecionar o RTL no PRISM}
\label{sec:17-prism}

Paralelamente à simulação, o usuário pode inspecionar a estrutura do RTL. O botão
PRISM dispara, primeiro, o \code{verilogSyntaxCheck}, que além do check de sintaxe
(\lstinline|iverilog -tnull|) regenera a hierarquia: \code{generateProjectHierarchy}
monta um script \tool{Yosys} (\code{read\_verilog} de \path{components/HDL/*} + dos
\code{synthesizableFiles}, \code{hierarchy -top}, \code{proc}, \code{write\_json})
e o roda via \code{buildYosysHierarchySpec}, produzindo
\path{Temp/project_hierarchy.json}, que \code{parseYosysHierarchy} converte na
árvore de módulos da file-tree.

O \tool{PRISM} propriamente (\file{main/ipc/prism.js}) usa o mesmo \tool{Yosys}
allowlisted para gerar o JSON da netlist e o entrega ao \tool{netlistsvg}
(rodando in-process) para produzir o SVG navegável do RTL. Os detalhes do
conversor e da renderização estão em \autoref{cap:11-aurora-prism-rtl}.

\section{Diagrama ASCII do pipeline completo}
\label{sec:17-diagrama}

O diagrama a seguir reúne os sete estágios, os binários envolvidos e os artefatos
de cada transição. A coluna da esquerda é o que o renderer orquestra; os binários
(em colchetes) são executados pelo main atrás da allowlist (\autoref{sec:17-fronteiras}).

\begin{lstlisting}
==============================================================================
   PIPELINE SAPHO PONTA A PONTA  (renderer orquestra | [bin] = main + allowlist)
==============================================================================

 (1) PROJETO
     AURORA: novo projeto + processadores
        |  grava
        v
     <proj>.spf  ........... fonte canonica (synthesizableFiles, testbenchFiles,
        |                     processors[]: name/cmmFile/clk/numClocks/showArrays)
        |  cria pastas:  <proj>/<proc>/{Software,Hardware,Simulation}
        |                components/Temp/<proc>/   (scratch -t / CWD da sim)
        v
 (2) CODIGO-FONTE  (editor Monaco)
     <proc>.cmm  (ou .cpp)   + top-level.v  + testbench (<proc>_tb.v auto ou
        |                                     top_level_tb.v / *.py do usuario)
        |   [Verilator] ensureChegueiToaqui -> injeta #TOAQUI no .cmm
        v
 (3) COMPILACAO   (precompileAllProcessors -> por processador)
     +--------------------------------------------------------------+
     | [cmmcomp.exe]  -i .cmm -n -p -m -t  [--array] [-pt|-en]      |
     |     -> <proc>.asm  +  cmm_log.txt  +  pc_<proc>_mem.txt      |
     |                       +  <proc>_data.mif / <proc>_inst.mif   |
     | [appcomp.exe]  -i .asm -t                                    |
     |     -> .asm pre-processado (params + enderecos)              |
     | [asmcomp.exe]  -i .asm -p -d HDL -m -t -f<clk> -c<numClk>    |
     |     -> Hardware/<proc>.v  (RTL sintetizavel)                 |
     |     -> <proc>_tb.v  (testbench padrao) -> Simulation/        |
     +--------------------------------------------------------------+
        |   (-y components/HDL resolve processor/core/ula/addr_dec/instr_dec)
        |
        +-------------------------------+----------------------------------+
        |                               |                                  |
        v (4a Icarus)                   v (4b Verilator)        v (4c cocotb/Python)
 (5) INSTRUMENTAR DUMP  (comum aos tres ramos, ANTES de simular)
     _resolveWaveSelection: .gtkw ativo > Wave Config > $dumpvars do tb > default
     instrumentTestbench -> Temp/instr_<tb>.v  (copia; original intacto)
        |
        v
     [iverilog -s tb -o .vvp     [verilator --binary       [cocotb: iverilog
       -y HDL] ; [vvp -fst]       --timing --trace]          ou verilator por baixo]
        |   stageProcessorMemoryFiles: pc_*_mem.txt -> Temp/ (readmemb)
        v
     DUMP de ondas (FST)
        |
        |  _extractFstHeaderVcd  -> [fst2vcd] -> <tb>.header.vcd (escopos/sinais)
        v
 (6) VISUALIZAR   (getViewer)
        +--- gtkwave -> buildAuroraGtkw -> <tb>.gtkw -> [gtkwave.exe --dark ...]
        +--- surfer  -> buildSurferLayout -> .surf -> [surfer.exe]
              faixas C+- / assembly / variaveis em LOCKSTEP com o clock

 (7) INSPECIONAR RTL  (botao PRISM, em paralelo)
     verilogSyntaxCheck -> [iverilog -tnull] (ok) ->
     generateProjectHierarchy -> [yosys] read_verilog + hierarchy -top + proc +
        write_json -> Temp/project_hierarchy.json -> arvore de modulos
     PRISM -> [yosys] netlist JSON -> netlistsvg (in-process) -> SVG navegavel
==============================================================================
\end{lstlisting}

\section{Os dois scripts de referência do YANC}
\label{sec:17-scripts}

O repositório \tool{YANC} encapsula o pipeline em dois scripts que servem de
especificação executável e batem, passo a passo, com o que a AURORA faz no main.

\begin{description}[leftmargin=1.2em, style=nextline]
  \item[\file{Scripts/single\_proc.sh}] Um único processador. Exemplo embutido:
        \code{proc\_fft} (FFT de ordem 8 com endereçamento bit-reverso do \cmm{}).
        Usa o testbench auto-gerado \file{<proc>\_tb.v} como top de simulação,
        roda Icarus (default) ou Verilator via \lstinline|--sim|, e abre o
        \tool{GTKWave} com layout do \tool{gen\_gtkw}.
  \item[\file{Scripts/multi\_proc.sh}] Vários processadores sob um top-level
        escrito pelo usuário. Exemplo embutido: projeto \code{DTW} com
        \code{ProcDTW} + \code{ZeroCross} (já testado em FPGA). Como os
        processadores rodam juntos, o testbench é o \file{top\_level\_tb.v} do
        usuário (o \file{<proc>\_tb.v} auto-gerado só dirige um processador
        isolado). A view mostra os sinais do top-level lado a lado com a execução
        de software dos dois processadores em \emph{lockstep}.
\end{description}

A diferença essencial entre os dois é qual testbench dirige a simulação — o
auto-gerado por processador (single) ou o top-level de usuário (multi) — e o
\emph{staging} extra de \file{pc\_<proc>\_mem.txt} e \file{.mif} para o CWD da
simulação no caso multi. A AURORA reproduz exatamente essa lógica: o staging vive
em \code{stageProcessorMemoryFiles} e a escolha do top de simulação em
\code{\_waveDeriveSimTopModule} / \code{validateForWave}.

\section{Síntese final}

O pipeline SAPHO é uma cadeia determinística de transformações sobre artefatos em
disco, ancorada em um único arquivo de configuração (\file{.spf}). Cada estágio
tem um responsável claro no renderer (um método de \code{CompilationModule} ou um
helper de \file{processor\_compiler.js}) e, quando precisa de um binário, atravessa
a fronteira \code{runSpec} $\rightarrow$ \lstinline|exec-spec| $\rightarrow$
allowlist do main. Do \cmm{} ao SVG do RTL, passando pelo \file{.v} sintetizável,
pelo dump de ondas e pela visualização em lockstep, o caminho é único e sem
ramificações estruturais — o que muda entre os cenários (com/sem processador,
single/multi, Icarus/Verilator/cocotb, GTKWave/Surfer) são escolhas de dados e de
preferência, não desvios de fluxo. Os capítulos referenciados ao longo deste texto
detalham cada subsistema; aqui o objetivo foi mostrar como todos se encaixam.

% !TEX root = ../main.tex
\chapter{Engenharia, qualidade e evolução do projeto}
\label{cap:18-engenharia-evolucao}

Este capítulo descreve dois aspectos complementares da plataforma \tool{SAPHO}.
A \autoref{sec:18-engenharia} (Parte~A) documenta a \emph{infraestrutura de
engenharia e qualidade} do repositório \repo{aurora}: como o código é testado,
verificado, formatado, versionado e integrado de forma contínua. A
\autoref{sec:18-evolucao} (Parte~B) reconstrói a \emph{linha do tempo de
evolução} dos repositórios \repo{aurora} e \repo{yanc} a partir do histórico
real de \emph{commits} do Git --- marcos técnicos, progressão de versões e
contribuidores.

Todas as afirmações desta seção foram verificadas diretamente nos arquivos de
configuração e no histórico de \emph{commits}; onde a documentação textual
divergir do que está descrito aqui, prevalece o que foi lido do repositório.

\section{Parte A --- Engenharia e qualidade}
\label{sec:18-engenharia}

A IDE \tool{AURORA} é um aplicativo Electron escrito majoritariamente em
JavaScript, com uma fatia crescente de TypeScript compilado \emph{in place}. A
disciplina de qualidade do repositório é construída sobre uma cadeia de
ferramentas que atua em três momentos: no \emph{editor} do desenvolvedor
(\file{.editorconfig}, \file{tsconfig.json}), no momento do \emph{commit}
(\tool{husky} + \tool{lint-staged} + \tool{commitlint}) e na
\emph{integração contínua} (GitHub Actions). Cada um desses estágios é
descrito a seguir.

\subsection{Visão geral dos scripts npm}
\label{sub:18-scripts-npm}

O ponto de entrada de toda a automação é a seção \code{scripts} de
\file{package.json}. A \autoref{tab:18-npm-scripts} resume os comandos
relevantes para teste e qualidade.

\begin{longtable}{Y Z}
  \caption{Scripts npm de qualidade e ciclo de vida (\file{package.json}).}
  \label{tab:18-npm-scripts}\\
  \toprule
  \textbf{Script} & \textbf{Ação} \\
  \midrule
  \endfirsthead
  \toprule
  \textbf{Script} & \textbf{Ação} \\
  \midrule
  \endhead
  \bottomrule
  \endfoot
  \code{build:ts}      & \lstinline|tsc| --- compila os módulos TypeScript declarados em \file{tsconfig.json} para \code{.js} ao lado do fonte. \\
  \code{build:renderer}& \lstinline|vite build| --- empacota o \emph{renderer} em \path{dist/}. \\
  \code{predev} / \code{dev} & \code{predev} roda \code{build:ts}; \code{dev} executa \file{scripts/dev.js}. \\
  \code{prestart} / \code{start} & \code{prestart} encadeia \code{bootstrap}, \code{build:ts} e \code{build:renderer}; \code{start} chama \file{scripts/launch-electron.js}. \\
  \code{test}          & \lstinline|vitest run| --- suíte unitária (sem Electron). \\
  \code{test:watch}    & \lstinline|vitest| em modo observação. \\
  \code{test:coverage} & \lstinline|vitest run --coverage| --- cobertura via provedor v8. \\
  \code{test:e2e} (com \code{pretest:e2e}) & compila o \emph{renderer} e roda os testes E2E sob \file{vitest.config.e2e.js}. \\
  \code{deadcode} / \code{deadcode:all} & \tool{knip} para detecção de código/dependências mortas. \\
  \code{prebuild} / \code{build} & \code{prebuild} encadeia \code{bootstrap}, \code{build:ts}, \code{build:renderer}; \code{build} invoca \lstinline|electron-builder|. \\
  \code{prepare}       & \lstinline|husky| --- instala os \emph{git hooks} ao rodar \lstinline|npm install|. \\
\end{longtable}

O campo \code{engines} exige \lstinline|node >=18.0.0|. As versões pinadas de
ferramentas (em \code{devDependencies}) incluem \tool{vitest}, \tool{playwright},
\tool{eslint}, \tool{knip} (fixado em \code{6.17.1}), \tool{husky},
\tool{lint-staged}, \tool{commitlint}, \tool{typescript} e \tool{electron-builder}.

\subsection{Estratégia de testes}
\label{sub:18-testes}

A suíte é dividida em duas camadas com configurações independentes, refletindo
custos de execução muito diferentes.

\subsubsection{Testes unitários (Vitest)}

A configuração padrão fica em \file{vitest.config.js}. Ela inclui apenas
\path{tests/unit/**/*.test.js} e exclui \path{node_modules}. São testes
\enquote{puros}: exercitam módulos isolados de lógica (parsing, normalização,
validação, escrita de formatos) sem subir o Electron, o que os torna rápidos.

No diretório \path{tests/unit/} existem dezenas de arquivos de teste. A
\autoref{tab:18-unit-tests} agrupa amostras por domínio para ilustrar a
cobertura.

\begin{longtable}{Y Z}
  \caption{Amostra de testes unitários por domínio (\path{tests/unit/}).}
  \label{tab:18-unit-tests}\\
  \toprule
  \textbf{Domínio} & \textbf{Arquivos de teste (exemplos)} \\
  \midrule
  \endfirsthead
  \toprule
  \textbf{Domínio} & \textbf{Arquivos de teste (exemplos)} \\
  \midrule
  \endhead
  \bottomrule
  \endfoot
  Segurança / caminhos & \file{safePath.test.js}, \file{sanitizeFileName.test.js} \\
  Forma de ondas (\emph{wave}) & \file{vcdParser.test.js}, \file{signalParser.test.js}, \file{gtkwWriter.test.js}, \file{gtkwProcWriter.test.js}, \file{surferLayoutWriter.test.js}, \file{waveStateStore.test.js}, \file{complexDecode.test.js}, \file{eventMarkers.test.js}, \file{selectionValidator.test.js} \\
  Compilação & \file{compilationHelpers.test.js}, \file{processor\_compiler.test.js}, \file{wave\_toolchain.test.js} \\
  Projeto / árvore & \file{spfStore.test.js}, \file{processorList.test.js}, \file{verilogClassifier.test.js}, \file{projectTreeRender.test.js}, \file{fileTreeViewController.test.js} \\
  Aurora Intelligence & \file{ai\_system\_prompt.test.js}, \file{ai\_chat\_render.test.js}, \file{ai\_metadata.test.js}, \file{chat\_history.test.js}, \file{tool\_permission.test.js}, \file{modelGovernance.test.js}, \file{provider\_view.test.js} \\
  Editor / componentes & \file{documentTypeDetector.test.js}, \file{monacoPin.test.js}, \file{aurora\_tabs.test.js}, \file{aurora\_panel.test.js}, \file{tab\_utils.test.js} \\
  Git / decorações & \file{gitDecorations.test.js}, \file{gitNs.test.js} \\
  Diversos & \file{i18n.test.js}, \file{utils.test.js}, \file{formatTimestamp.test.js}, \file{cliDownloader.test.js}, \file{download-toolchain.test.js} \\
\end{longtable}

Como exemplo concreto, \file{safePath.test.js} valida o auxiliar defensivo
\lstinline|safePath| de \file{main/utils.js} usado por todos os \emph{handlers}
de IPC de sistema de arquivos: rejeita não-strings (\code{TypeError}), strings
vazias e injeção de \emph{null byte}, e normaliza separadores. O próprio teste
documenta que \lstinline|safePath| é \emph{normalização}, não
\emph{confinamento} de diretório.

\paragraph{Cobertura.} A seção \code{coverage} de \file{vitest.config.js}
configura o provedor \code{v8}, com relatórios \code{text-summary} e \code{lcov}
gravados em \path{./coverage}. A opção \lstinline|all: false| restringe o
relatório aos módulos efetivamente importados pelos testes (em vez de diluir o
número com toda a árvore do \emph{renderer}/\emph{main} a 0\%). Listas de
exclusão removem arquivos de teste, \path{node_modules}, \path{dist/},
\path{components/} e arquivos de configuração.

\subsubsection{Testes ponta a ponta (E2E) com Playwright + Electron}

Os testes E2E ficam em \path{tests/e2e/} e usam a API \lstinline|_electron| do
Playwright para lançar o aplicativo Electron completo (processos \emph{main},
\emph{renderer} e GPU). A configuração \file{vitest.config.e2e.js} inclui apenas
\path{tests/e2e/**/*.test.js} e ajusta o ambiente para o custo de inicialização
do Electron:

\begin{itemize}[nosep]
  \item \lstinline|fileParallelism: false| --- os arquivos rodam \emph{um por
        vez}. Cada arquivo faz \emph{cold start} de um Electron completo; rodar
        vários em paralelo estourava o orçamento de tempo em \emph{runners} de
        2 núcleos do CI. Serializar mantém apenas um Electron vivo de cada vez.
  \item \lstinline|testTimeout: 60_000| e \lstinline|hookTimeout: 120_000| ---
        margens generosas, dimensionadas para o pior caso observado no
        \emph{runner} \code{windows-latest}, onde o \emph{cold start} mais a
        cadeia de IPC de carga de \code{.spf} chega a 10--15~s.
\end{itemize}

Os três arquivos E2E são \file{smoke.test.js}, \file{edit-flow.test.js} e
\file{split-pane.test.js}. O teste de fumaça (\file{smoke.test.js}) lança o
aplicativo com um diretório de dados de usuário isolado
(\lstinline|--user-data-dir|), define a variável
\lstinline|SAPHO_SKIP_SINGLE_INSTANCE| para contornar o \emph{lock} de instância
única e remove \lstinline|ELECTRON_RUN_AS_NODE| do ambiente. Em seguida espera a
janela principal (identificada por \path{index.html}) e verifica:

\begin{itemize}[nosep]
  \item que o \emph{container} \lstinline|#monaco-editor| está montado;
  \item que os módulos AMD do Monaco carregaram (\lstinline|window.monaco| existe);
  \item que nenhum marcador conhecido de falha de inicialização do Monaco
        apareceu no console ou em \emph{page errors} (a lista
        \code{MONACO_FAILURE_MARKERS} inclui
        \enquote{EditorManager has not been initialized},
        \enquote{Property description must be an object} e
        \lstinline|monaco.contribution.js|);
  \item que o \emph{renderer} alcança o estado interativo dentro de um orçamento
        de 30~s, medido pela marca de desempenho
        \lstinline|performance.mark('aurora-interactive')| emitida ao fim da
        inicialização. Trata-se de uma \emph{guarda de regressão}, não de um
        \emph{benchmark}.
\end{itemize}

\subsection{Lint com ESLint (flat config)}
\label{sub:18-eslint}

O \emph{linting} é configurado em \file{eslint.config.mjs} no formato
\emph{flat config} do ESLint~9, exportando um array de blocos via
\lstinline|defineConfig|. O desenho reconhece que o \emph{renderer} da AURORA
\textbf{não} é um grafo de \emph{imports} empacotado: \file{index.html} carrega
\emph{scripts} clássicos via tags \lstinline|<script>| que compartilham estado
por globais de \lstinline|window|. Por isso a configuração declara um conjunto de
globais compartilhados (\code{rendererSharedGlobals}: \lstinline|editor|,
\lstinline|monaco|, \lstinline|TabManager|, \lstinline|CodeFormatter| etc.) como
\code{readonly}.

Os blocos escopam o ambiente correto para cada tipo de código:

\begin{itemize}[nosep]
  \item \textbf{Ignorados globais}: \path{components/**} (\emph{toolchain}
        baixado, artefatos de terceiros), exceto \path{components/Scripts/}
        (código próprio); \path{dist/} e \path{release/} (saídas de \emph{build}).
  \item \textbf{Regra base} \code{js/recommended} com \code{no-unused-vars}
        ajustada: \code{args: "none"}, \code{caughtErrors: "none"} e o padrão de
        escape \code{varsIgnorePattern: "^_"} (prefixo
        sublinhado marca variáveis intencionalmente não usadas, como
        \lstinline|_event|, \lstinline|_info|).
  \item \textbf{Processo \emph{main} e scripts Node} (\file{main.js},
        \path{main/**}, \path{components/Scripts/**}, \path{scripts/**}):
        \code{sourceType: "commonjs"} com globais do Node.
  \item \textbf{Testes E2E e unitários}: \code{sourceType: "module"}, globais do
        Node mais \lstinline|window|/\lstinline|document| em \code{readonly}.
  \item \textbf{Preloads} (\file{js/app/preload*.js}): CommonJS com globais de
        Node \emph{e} \emph{browser} (a ponte \lstinline|contextBridge| vive nos
        dois mundos).
  \item \textbf{Renderer} (\path{js/**}, \path{html/**}, exceto os preloads):
        globais de \emph{browser} mais os globais compartilhados.
\end{itemize}

Na CI o \emph{lint} roda com tolerância zero a avisos
(\lstinline|npx eslint . --max-warnings=0|).

\subsection{Detecção de código morto (knip)}
\label{sub:18-knip}

\tool{knip} é configurado em \file{knip.config.js} (CommonJS) como portão de
código morto. Como o \emph{renderer} não é um grafo de \emph{imports}
rastreável, a configuração informa explicitamente todos os pontos de entrada
reais:

\begin{itemize}[nosep]
  \item \file{main.js} (processo \emph{main}); \file{js/app/preload*.js} (4
        \emph{preloads} carregados por caminho); \file{html/prism/prism.js};
        \file{js/components/design-lab.js};
        \path{scripts/*.js}; \path{components/Scripts/*.js};
  \item e --- de forma engenhosa --- todas as entradas
        \lstinline|<script src="/foo.js">| \textbf{lidas de
        \file{index.html} em tempo de execução}, pela função
        \lstinline|rendererEntriesFromIndexHtml()|. Assim, adicionar ou remover
        uma tag de \emph{script} mantém a configuração correta sem edição.
\end{itemize}

O array \code{project} restringe a análise ao código próprio
(\file{main.js}, \path{main/**}, \path{js/**}, \path{scripts/*.js},
\path{components/Scripts/*.js}, \file{html/prism/prism.js}). A lista
\code{ignoreDependencies} cobre pacotes alcançados de formas que a análise
estática não enxerga: \tool{monaco-editor} (carregado pelo próprio \emph{loader}
AMD), os provedores \lstinline|@ai-sdk/*| (resolvidos por string em
\file{main/ai/provider.js}), \tool{@openai/codex} (CLI \emph{spawned} como
subprocesso), \tool{@phosphor-icons/web} e \tool{katex} (referenciados via
\lstinline|<link>|/\lstinline|<script>| diretos a \path{node_modules}).

A CI roda apenas \lstinline|npm run deadcode|, que limita o portão a
\code{files} (módulos não usados) e \code{dependencies} (pacotes npm não usados)
--- as classes que o autor considera confiáveis para falhar a \emph{build}.

\subsection{Hooks de Git (husky + lint-staged)}
\label{sub:18-husky}

Os \emph{git hooks} são gerenciados por \tool{husky} (instalado pelo script
\code{prepare}). Há dois \emph{hooks} em \path{.husky/}:

\begin{itemize}[nosep]
  \item \file{pre-commit}: executa \lstinline|npx lint-staged|. A configuração
        \code{lint-staged} em \file{package.json} roda
        \lstinline|eslint --max-warnings=0 --no-warn-ignored| sobre os arquivos
        \lstinline|*.{js,mjs,cjs}| preparados (\emph{staged}).
  \item \file{commit-msg}: executa
        \lstinline|npx --no-install commitlint --edit "$1"|, validando a
        mensagem de \emph{commit}.
\end{itemize}

\subsection{Conventional Commits (commitlint)}
\label{sub:18-commitlint}

A convenção de mensagens é definida em \file{commitlint.config.mjs}, que estende
\code{@commitlint/config-conventional} com ajustes para o estilo bilíngue e
escopado da AURORA:

\begin{itemize}[nosep]
  \item \code{type-enum} acrescenta tipos próprios do projeto aos padrões:
        além de \code{feat}, \code{fix}, \code{perf}, \code{refactor},
        \code{chore}, \code{docs}, \code{test}, \code{build}, \code{ci},
        \code{style}, \code{revert}, inclui \code{hardening}, \code{ux},
        \code{ui} e \code{diag}.
  \item \code{header-max-length} elevado para 120 (assuntos costumam estar em
        português e referenciar seções com o sinal de parágrafo).
  \item \code{subject-case} e \code{body-max-line-length} desativados
        (valor 0), para não conflitar com o registro bilíngue.
\end{itemize}

Essa convenção pareia com o \tool{release-please} (ver
\autoref{sub:18-ci}), que agrega \emph{Conventional Commits} para gerar
notas de versão e \emph{bumps} semânticos.

\subsection{TypeScript (compilação \emph{in place})}
\label{sub:18-typescript}

A adoção de TypeScript é \emph{incremental e in place}: \file{tsconfig.json}
define \lstinline|rootDir: "."| e \lstinline|outDir: "."|, de modo que
\lstinline|tsc| emite cada \code{.js} ao lado do seu \code{.ts}. As opções
relevantes são \code{target: "ES2022"}, \code{module: "ESNext"},
\code{moduleResolution: "Bundler"}, \lstinline|strict: true|,
\code{isolatedModules: true} e \code{esModuleInterop: true}.

A lista \code{include} é \emph{explícita} --- apenas os módulos já migrados são
compilados, concentrados em três áreas: utilidades de caminho e tipos globais
(\file{js/utils/path\_utils.ts}, \file{js/types/aurora-globals.d.ts}), o
subsistema de forma de ondas (\path{js/wave/*.ts}: \emph{parsers} de sinais e
VCD, escritores \code{.gtkw}, instrumentação de \emph{testbench}, decodificação
de complexos, marcadores de evento) e o subsistema de compilação
(\path{js/compilation/*.ts} e \path{js/compilation/builders/*.ts}: \emph{specs}
de comando, \emph{builders} para \cmm, \emph{asm}, \tool{cocotb},
\tool{iverilog}, \tool{vvp}, \tool{yosys}, \tool{verilator} e ferramentas de
\emph{wave}).

\begin{nota}
Os \code{.js} gerados por \lstinline|tsc| são \textbf{artefatos de build} e estão
no \file{.gitignore}, de modo que a árvore versionada nunca diverge do fonte.
A CI mantém duas guardas complementares: \lstinline|tsc --noEmit| prova que o
\code{.ts} tipa corretamente, e o script \file{scripts/check-no-generated-js.js}
falha a \emph{build} se algum \code{.js} gerado (um \code{.js} versionado com um
\code{.ts} irmão também versionado) reaparecer no Git.
\end{nota}

\subsection{Convenções de formatação (.editorconfig)}
\label{sub:18-editorconfig}

O arquivo \file{.editorconfig} fixa as convenções básicas para qualquer editor
compatível: \code{charset = utf-8}, \code{end_of_line = lf},
\code{insert_final_newline = true}, \code{trim_trailing_whitespace = true},
indentação por espaços com largura 4 por padrão. Sobrescrições por tipo: largura
2 para \code{json}/\code{yml}/\code{yaml} e para \code{html}/\code{css}/\code{scss};
preservação de espaços finais em \code{md}/\code{markdown}; e indentação por
\emph{tab} no \code{Makefile}.

\subsection{Integração contínua (GitHub Actions)}
\label{sub:18-ci}

O diretório \path{.github/workflows/} contém cinco \emph{workflows}, descritos na
\autoref{tab:18-workflows}.

\begin{longtable}{Y Z}
  \caption{Workflows de GitHub Actions (\path{.github/workflows/}).}
  \label{tab:18-workflows}\\
  \toprule
  \textbf{Arquivo} & \textbf{Papel} \\
  \midrule
  \endfirsthead
  \toprule
  \textbf{Arquivo} & \textbf{Papel} \\
  \midrule
  \endhead
  \bottomrule
  \endfoot
  \file{ci.yml} & Lint + \emph{typecheck} + testes + \emph{smoke build} em \code{windows-latest}, em \emph{push} e \emph{pull request} para \code{main}. \\
  \file{codeql.yml} & Análise CodeQL (\code{javascript-typescript}) em \emph{push}, PR e agendamento semanal (segunda 06:00 UTC), ignorando \path{dist}, \path{components}, \path{node_modules} e \path{assets/prism-skins}. \\
  \file{release-please.yml} & Mantém uma PR de \enquote{release} em \code{main} que agrega os \emph{Conventional Commits}; ao mesclar, cria o \emph{commit} de \emph{bump}, a \emph{tag} e a GitHub Release. \\
  \file{release.yml} & Constrói o instalador Windows e publica nos \emph{assets} da Release (consumidos pelo \tool{electron-updater}); dispara quando uma Release é \emph{published} (mais \code{workflow_dispatch} manual). \\
  \file{dependabot-automerge.yml} & Faz \emph{auto-merge} (squash) de PRs do Dependabot do tipo \emph{patch}/\emph{minor} após a CI passar. \\
\end{longtable}

\subsubsection{Pipeline de CI (\file{ci.yml})}

O \emph{job} \code{lint-and-build} roda em \code{windows-latest} com Node~20 e
\lstinline|npm ci|. A sequência de passos é uma cadeia de portões:

\begin{enumerate}[nosep]
  \item \textbf{Verificar versões pinadas}: \file{scripts/check-pinned-versions.js}
        garante que pacotes com semver exata (sem \code{^}/\code{~})
        não derivem da versão declarada --- mecanismo criado por causa de uma
        regressão do \tool{monaco-editor} 0.53.0 que quebrava a inicialização do
        editor; \tool{monaco-editor} está fixado em \code{0.52.2}.
  \item \textbf{ESLint} com \lstinline|--max-warnings=0|.
  \item \textbf{Typecheck} via \lstinline|npx tsc --noEmit|.
  \item \textbf{Sem \code{.js} gerado versionado}: \file{scripts/check-no-generated-js.js}.
  \item \textbf{Consistência de i18n}: \file{scripts/check-i18n.js} (sincronia
        en/pt e ausência de chaves indefinidas).
  \item \textbf{Código morto}: \lstinline|npm run deadcode| (\tool{knip}).
  \item \textbf{Sanidade dos módulos JS}: \lstinline|node --check| sobre cada
        \code{.js} em \path{js} e \path{main}.
  \item \textbf{Testes unitários com cobertura}: \lstinline|npm run test:coverage|.
  \item \textbf{Upload de cobertura ao Codecov} (informativo;
        \lstinline|fail_ci_if_error: false|).
  \item \textbf{Testes E2E} (\emph{smoke} do Electron).
  \item \textbf{Build do renderer} (Vite) e \textbf{build sem publicação}
        (\lstinline|electron-builder --win --x64 --publish never|).
\end{enumerate}

\subsubsection{Fluxo de release}

O versionamento é automatizado por \tool{release-please}: cada \emph{push} em
\code{main} atualiza uma PR de release acumulada. Mesclá-la publica a GitHub
Release, o que dispara \file{release.yml}. Esse \emph{workflow} roda
\lstinline|npm run bootstrap| (baixando o \emph{bundle} mingw unificado ---
\tool{iverilog}/\tool{yosys}/\tool{verilator}/Python+\tool{cocotb} ---, o fork
\repo{gtkwave-nipscern}, os compiladores \tool{YANC} e o \tool{Surfer}),
verifica \emph{sentinelas} obrigatórias do \emph{toolchain} (falhando se alguma
faltar, pois uma Release não pode sair sem o \emph{toolchain}), compila
\code{.ts} e o \emph{renderer}, e finalmente publica o instalador com
\lstinline|electron-builder --win --x64 --publish always|. O \tool{electron-updater}
lê \file{latest.yml} e o \code{.exe} diretamente desses \emph{assets} (ver
\autoref{cap:11-build-empacotamento} para detalhes de empacotamento e
auto-atualização). A configuração \code{build} de \file{package.json} publica no
repositório \repo{sapho}, refletindo o \emph{split} intencional entre o repo de
desenvolvimento (\repo{aurora}) e o canal de release estável (\repo{sapho}).

\section{Parte B --- Evolução via histórico de Git}
\label{sec:18-evolucao}

Esta seção reconstrói a história dos repositórios a partir do log de
\emph{commits} real (e não do \file{CHANGELOG.md}). Os números e datas vêm de
\lstinline|git log|, \lstinline|git shortlog -sn| e \lstinline|git tag|
executados em ambos os repositórios.

\subsection{Panorama quantitativo}
\label{sub:18-panorama}

\begin{tabularx}{\linewidth}{Y Y Y}
  \toprule
  \textbf{Métrica} & \repo{aurora} & \repo{yanc} \\
  \midrule
  \emph{Commits} totais & 1269 & 665 \\
  Primeiro \emph{commit} & 04/02/2025 & 09/10/2024 \\
  \emph{Commit} mais recente (deste estudo) & 19/06/2026 & 14/06/2026 \\
  Tags de versão & até \code{v6.3.2} & até \code{v5.2} \\
  \bottomrule
\end{tabularx}

\medskip

A atividade concentra-se fortemente em 2026, com picos em maio e junho de 2026.
No \repo{aurora}, os meses de maior volume foram 2026-06 (498 \emph{commits}) e
2026-05 (447); no \repo{yanc}, 2026-05 (280) e 2026-06 (125). Esse adensamento
coincide com a fase de profissionalização da plataforma (introdução da cadeia de
qualidade, do \emph{bundle} de \emph{toolchain} e da Aurora Intelligence).

\subsection{Contribuidores}
\label{sub:18-contribuidores}

A \autoref{tab:18-autores-aurora} mostra a contagem de \emph{commits} por autor
em cada repositório (de \lstinline|git shortlog -sn --all|). Vários nomes
correspondem à mesma pessoa em identidades Git distintas: \enquote{Chrysthofer},
\enquote{chrysthofer}, \enquote{Kristoffer} e a conta \enquote{nipscern}
(\code{nipscernlab}) são todos Chrysthofer Afonso; \enquote{Luciano} é o
responsável pelo projeto (Luciano Andrade).

\begin{tabularx}{\linewidth}{Y Y}
  \toprule
  \multicolumn{2}{c}{\textbf{\repo{aurora} --- commits por autor}} \\
  \midrule
  \textbf{Autor} & \textbf{Commits} \\
  \midrule
  Chrysthofer & 611 \\
  Luciano & 473 \\
  nipscern & 228 \\
  dependabot[bot] & 21 \\
  Kristoffer & 10 \\
  ART\_AM (Arthur) & 8 \\
  chrysthofer & 7 \\
  LucasAredesCruz & 2 \\
  github-actions[bot] & 2 \\
  Claude & 1 \\
  robertaluchini & 1 \\
  \bottomrule
\end{tabularx}

\begin{tabularx}{\linewidth}{Y Y}
  \toprule
  \multicolumn{2}{c}{\textbf{\repo{yanc} --- commits por autor}} \\
  \midrule
  \textbf{Autor} & \textbf{Commits} \\
  \midrule
  Luciano & 399 \\
  nipscern & 243 \\
  Claude & 21 \\
  Chrysthofer & 3 \\
  Kristoffer & 2 \\
  \bottomrule
\end{tabularx}

\medskip

\paragraph{Colaboração via \emph{branches}.} O \emph{branch} \code{main} é o de
integração. O histórico de \emph{branches} (remotos) evidencia colaboração e
automação por ramos: \code{feature/aurora-revamp} e \code{ide-ux-fixes}
(trabalho de funcionalidade), \code{claude/wave-button-pipeline-V4zTD} (sessão
assistida por IA), além de \emph{branches} automáticos do Dependabot
(\code{dependabot/npm_and_yarn/*}, \code{dependabot/github_actions/*}) e do
\tool{release-please} (\code{release-please--branches--main--*}). Arthur
(\code{ART_AM}) contribui pontualmente via \emph{branches} de correção. As
entradas \code{dependabot[bot]}, \code{github-actions[bot]} e \code{Claude}
refletem, respectivamente, atualizações automáticas de dependências, \emph{commits}
de automação e sessões de assistência por IA.

\subsection{Progressão de versões (tags)}
\label{sub:18-versoes}

A \autoref{tab:18-tags-aurora} lista as \emph{tags} do \repo{aurora} com a data
e o assunto do \emph{commit} apontado.

\begin{longtable}{Y Y Z}
  \caption{Tags do repositório \repo{aurora}.}
  \label{tab:18-tags-aurora}\\
  \toprule
  \textbf{Tag} & \textbf{Data} & \textbf{Commit apontado} \\
  \midrule
  \endfirsthead
  \toprule
  \textbf{Tag} & \textbf{Data} & \textbf{Commit apontado} \\
  \midrule
  \endhead
  \bottomrule
  \endfoot
  \code{v4.1.13} & 02/05/2026 & \emph{build}: corrige caminho de \code{yosys.exe} em \code{extraResources} \\
  \code{v6.0.0}  & 16/05/2026 & \emph{feat}(intelligence,toolbar): Aurora Intelligence ponta a ponta + marca \\
  \code{v6.2.0}  & 17/05/2026 & \emph{chore}(release): v6.2.0 \\
  \code{v6.3.1}  & 19/05/2026 & \emph{fix}(wave): PATH para a passada 2 do Verilator \\
  \code{v6.3.2}  & 19/05/2026 & \emph{chore}(release): 6.3.2 --- sincroniza regras do YANC \\
  \code{toolchain-v1} & 06/05/2026 & \emph{feat}(toolchain): \emph{bootstrap} de auto-download + remoção de \code{Packages} do Git \\
  \code{toolchain-v2} & 07/05/2026 & \emph{chore}(deps): Electron 38 \textrightarrow{} 39.8.10 \\
  \code{msys-v1} & 03/06/2026 & \emph{refactor}(toolchain): consolida o \emph{toolchain} FOSS em um \emph{bundle} mingw \\
  \code{verilator-v1} & 19/05/2026 & \emph{fix}(ai): desbloqueia Claude Code no Windows + persiste tool calls \\
  \code{verilator-v2} & 02/06/2026 & marco do fluxo Verilator \\
  \code{main-pre-revamp-20260617} & 15/06/2026 & \emph{snapshot} pré-\emph{revamp} \\
  \code{shelf-a2-godfiles-2026-06-18} & 15/06/2026 & \emph{snapshot} de prateleira (refatoração A2) \\
  \bottomrule
\end{longtable}

A versão declarada em \file{package.json} no momento deste estudo é
\code{6.3.2}. Observa-se que parte das \emph{tags} é semântica de versão
(\code{v4.1.13}, \code{v6.x}) e parte é de \emph{snapshot}/componente
(\code{toolchain-v*}, \code{msys-v1}, \code{verilator-v*}, e tags datadas de
\enquote{prateleira}). As versões anteriores a \code{v4.1.13} aparecem no
\emph{assunto} de muitos \emph{commits} (\code{v2.x} em 2025, \code{v3.0.0} e
\code{v3.4.0} em maio de 2025) ainda na fase em que o \emph{bump} de versão era
manual, antes da adoção do \tool{release-please}.

Para o \repo{yanc}, as \emph{tags} são \code{v1}, \code{v2}, \code{v3},
\code{v4}, \code{v4.1}, \code{v4.2}, \code{v4.3}, \code{v4.4}, \code{v4.4.1},
\code{v5.0}, \code{v5.1} e \code{v5.2}, todas concentradas entre 10/05/2026
(\code{v1}, \enquote{\emph{build} e release dos binários do YANC em \emph{push}
de tag}) e 08/06/2026 (\code{v5.2}, \enquote{paridade de GTKWave no pipeline
C++ + otimização de \emph{asm} no \code{cppcomp}}).

\subsection{Linha do tempo do \repo{aurora}}
\label{sub:18-timeline-aurora}

A \autoref{tab:18-timeline} consolida os marcos verificados no histórico. As
datas são dos \emph{commits} reais.

\begin{longtable}{Y Z}
  \caption{Linha do tempo de marcos do \repo{aurora}.}
  \label{tab:18-timeline}\\
  \toprule
  \textbf{Período} & \textbf{Marco} \\
  \midrule
  \endfirsthead
  \toprule
  \textbf{Período} & \textbf{Marco} \\
  \midrule
  \endhead
  \bottomrule
  \endfoot
  Fev/2025 & Início do histórico público (\code{v2.0.0}, 04/02). Cadência rápida de \emph{tags} \code{v2.x} no \emph{assunto} dos \emph{commits}; versionamento manual. \\
  Mar/2025 & Introdução do PRISM (visualizador de RTL) na faixa \code{v2.21.0} (23/03). \\
  Abr/2025 & Editor Monaco (10/04); \enquote{Consolidation of CSS} (consolidação de CSS); \code{v2.38.0} \enquote{Huge UI Changes!} (26/04). \\
  Mai/2025 & \code{v3.0.0} (20/05) e \code{v3.4.0} (23/05); primeiros trabalhos de \emph{updater}. \\
  Jul--Nov/2025 & Refatorações do fluxo de compilação e do terminal; gerência de modos (\enquote{Project Mode}, \enquote{Verilog Mode}) com \emph{file manager}, \emph{drag \& drop} e hierarquia (nov/2025). \\
  Mar/2026 & \enquote{Split editor} (30/03) --- base de painéis divididos; primeiros \emph{commits} no formato \emph{Conventional} (\enquote{Fix(ui)} em 03/04). \\
  26/04/2026 & \enquote{SAPHO IDE: full UI overhaul + frameless window} --- \emph{revamp} visual maior, janela sem moldura. \\
  01--03/05/2026 & Reorganização do código em pastas de domínio (CSS e JS); \emph{split} de \file{main.js} em submódulos \path{main/}; LICENSE, README, CONTRIBUTING, SECURITY, templates; documentação de release + auto-\emph{updater}. \\
  06--07/05/2026 & \tool{husky} + \tool{lint-staged} (07/05); \emph{auto-download} de \emph{toolchain} e remoção de \code{Packages} do Git (\code{toolchain-v1}, 06/05); endurecimento de segurança (desabilita \code{nodeIntegration}, IPC de FS/exec). \\
  07/05/2026 & \tool{Vitest} adotado (\emph{bootstrap} com \file{safePath} + cobertura de download); \emph{opt-in} de TypeScript via \lstinline|// @ts-check|; \emph{release-please} automatiza releases. \\
  07--09/05/2026 & Primeiro teste E2E do Electron (\emph{smoke}: abre + Monaco inicializado, 07/05). \\
  10--11/05/2026 & \textbf{Consolidação para modo único}: fase 1 esconde o \emph{toggle} de modo e trava em \enquote{project} (11/05); fase 2.3 deleta os métodos órfãos de \enquote{Processor/Verilog Mode}; reestruturação \enquote{filosofia modo único}. \\
  13--14/05/2026 & \emph{Bump} para \code{5.0.0} (renomeia instalador, poda componentes empacotados). \\
  16--18/05/2026 & \textbf{Aurora Intelligence} ponta a ponta + marca (\code{v6.0.0}, 16/05); \emph{overhaul} de provedor Ollama + \emph{multi-step tool calling}; redesenho do PRISM (\emph{titlebar} própria, tema Aurora). \\
  19/05/2026 & Desbloqueio do Claude Code no Windows + persistência de \emph{tool calls} + novas \emph{tools} de IDE (\code{verilator-v1}); \code{v6.3.x}. \\
  02--03/06/2026 & Consolidação do \emph{toolchain} FOSS em um único \emph{bundle} mingw (\code{msys-v1}, 03/06). \\
  06--07/06/2026 & Conversão de \emph{builders} de compilação para TypeScript (\code{spec_factory}, \code{spec_runner}, \code{command_overrides}). \\
  13/06/2026 & \textbf{Vite adotado como \emph{bundler} do \emph{renderer}} (A1, primeiro \emph{flag-gated}, 21:39; default \emph{flipado} para \emph{renderer} empacotado às 22:36); \textbf{fundação Lit} + \enquote{Design Lab} + primeiro componente \lstinline|<aurora-statusbar>|. \\
  14/06/2026 & \textbf{Surfer} como \emph{viewer} \emph{opt-in} ao lado do GTKWave (\emph{toggle} + API da IA); migração de telas para componentes Lit (\emph{welcome}, \emph{command palette}); CSP + \emph{sandbox} em todas as janelas. \\
  15--16/06/2026 & \emph{Snapshots} \code{main-pre-revamp} e \code{shelf-a2-godfiles}; profissionalização do repo (CODEOWNERS, \tool{commitlint}, CodeQL, CITATION, README \textrightarrow{} \repo{sapho}); integridade de download e fluxo de release (\enquote{Wave B}). \\
  16--17/06/2026 & \textbf{Source Control embutido} (\enquote{Dagr}, runa Dagaz, G2); APIs de Git completas para a IA (G1, 14 \emph{tools}); decorações de status Git na árvore de arquivos (estilo VS Code); cobertura de testes \tool{vitest} v8 + badge Codecov. \\
  18--19/06/2026 & Simulação DigitalJS interativa no PRISM (O9, modo \enquote{Simular}); \emph{language servers}: \textbf{Verible} (O2: diagnósticos, \emph{format}, \emph{outline}, \emph{hover}, def/refs) e \textbf{slang-server} (O11); \textbf{tree-sitter} para \emph{highlight} preciso (O7); \emph{format} C/C++/\cmm{} via \tool{clang-format}; refatorações A2/A3 (extração de domínios, migração de globais \lstinline|electronAPI| para \emph{imports}); \emph{shells} semânticos \lstinline|<aurora-editor>|, \lstinline|<aurora-panel>|, \lstinline|<aurora-tabs>|; \emph{runbook} de assinatura de código. \\
  \bottomrule
\end{longtable}

\subsubsection{Marcos de arquitetura em destaque}

\begin{description}[style=nextline,leftmargin=0pt]
  \item[Adoção do Vite] O \emph{commit} \enquote{build(vite): adopt Vite as the
    renderer bundler (A1, renderer-only, flag-gated)} (13/06/2026) introduziu o
    Vite empacotando apenas o \emph{renderer}, atrás de uma \emph{flag}; horas
    depois o default foi invertido para o \emph{renderer} empacotado, com todas
    as janelas passando pelo Vite. É o ponto em que \code{build:renderer}
    (\lstinline|vite build| \textrightarrow{} \path{dist/}) entrou no ciclo de
    \emph{build}.
  \item[Revamp visual] O \emph{commit} \enquote{SAPHO IDE: full UI overhaul +
    frameless window} (26/04/2026) marca o \emph{revamp} maior; o trabalho
    continuou no \emph{branch} \code{feature/aurora-revamp}. A consolidação
    visual prossegue nas fases Lit (\autoref{cap:05-lit-componentes}, se
    aplicável).
  \item[Consolidação para modo único] Em 10--11/05/2026, uma série de
    \emph{commits} \emph{refactor(mode)} removeu os modos \enquote{Processor
    Mode} e \enquote{Verilog Mode} (que tinham \emph{file managers} próprios em
    nov/2025): a fase 1 escondeu o \emph{toggle} e travou em \enquote{project},
    e a fase 2.3 apagou os métodos órfãos. O resultado é a \enquote{filosofia
    modo único} da IDE atual.
  \item[Bundle de toolchain e split de repos] A remoção do diretório
    \code{Packages} do Git e o \emph{auto-download} via \code{bootstrap}
    (\code{toolchain-v1}, 06/05/2026), seguidos da consolidação em um único
    \emph{bundle} mingw (\code{msys-v1}, 03/06/2026), separaram o código-fonte
    do \emph{toolchain} pesado. As releases passaram a publicar em \repo{sapho}.
  \item[Source Control, Surfer, Verible, tree-sitter] A faixa final (14--19/06/2026)
    embutiu o controle de versão (\enquote{Dagr}), adicionou o \tool{Surfer}
    como \emph{viewer} de ondas \emph{opt-in}, e trouxe \emph{language servers}
    (Verible, slang) e \emph{highlight} via \tool{tree-sitter} ---
    aproximando a IDE de uma experiência de \emph{language tooling} completa.
\end{description}

\subsection{Linha do tempo do \repo{yanc}}
\label{sub:18-timeline-yanc}

O repositório dos compiladores (\enquote{Yet Another Compiler}) é anterior ao
\repo{aurora} (início em out/2024). A \autoref{tab:18-timeline-yanc} resume seus
marcos.

\begin{longtable}{Y Z}
  \caption{Linha do tempo de marcos do \repo{yanc}.}
  \label{tab:18-timeline-yanc}\\
  \toprule
  \textbf{Período} & \textbf{Marco} \\
  \midrule
  \endfirsthead
  \toprule
  \textbf{Período} & \textbf{Marco} \\
  \midrule
  \endhead
  \bottomrule
  \endfoot
  Out--Dez/2024 & \emph{Initial commit} (09/10/2024); trabalho em FFT, números complexos, geração de imagens de memória (\code{.mif}); \emph{Makefiles} para \code{CMMComp} e \code{Assembler}. \\
  Jan/2025 & Renomeação \code{Assembler} \textrightarrow{} \code{ASMComp}; variáveis \emph{float} no GTKWave. \\
  Fev--Mar/2025 & Limpeza de código; remoção da criação de Verilog para simulação (30/03). \\
  Mai/2026 & Profissionalização intensa: CI que constrói e publica binários do YANC em \emph{push} de tag (\code{v1}, 10/05); empacotamento de \code{HDL/} e \code{Macros/} no zip de release (\code{v2}); \emph{smoke} cobrindo o pipeline \code{cmmcomp} \textrightarrow{} \code{appcomp} \textrightarrow{} \code{asmcomp} (13/05); tradução dos comentários PT \textrightarrow{} EN; compiladores bilíngues (PT/EN). \\
  28--31/05/2026 & Série de releases \code{v3} (args nomeados de CLI, tabelas que crescem sob demanda, \code{--help}/\code{--version}), \code{v4}, \code{v4.1}, \code{v4.2}, \code{v4.3}, \code{v4.4}; \emph{harness} de regressão (\code{regress.sh}) que roda toolchain + simulação e compara arquivos de saída. \\
  01--03/06/2026 & \code{v4.4.1}; \code{v5.0} (correção em \code{gen_gtkw} sob \code{-Werror=format-truncation}). \\
  07--08/06/2026 & \code{v5.1} (biblioteca matemática + números complexos); \code{v5.2} (paridade de GTKWave no pipeline C++ + otimização de \emph{asm} no \code{cppcomp}). \\
  14/06/2026 & Estudo de viabilidade e projeto-piloto de prova formal em Lean~4 (\emph{feasibility study} + corpus de análise \code{f2mf}); correção de \emph{overflow} de mantissa em \code{f2mf}. \\
  \bottomrule
\end{longtable}

\paragraph{Automação e assistência por IA no YANC.} Assim como no \repo{aurora},
o \repo{yanc} recebeu uma cadeia de CI (\emph{workflows} de \emph{build} por
\emph{push} e release por \emph{tag}) e \emph{commits} de sessão de IA
(autor \enquote{Claude}, 21 \emph{commits}), incluindo um documento de
\emph{handoff} de sessão com sumário, \emph{bugs} encontrados e tarefas Windows.
O detalhamento dos compiladores em si está em
\autoref{cap:02-yanc-compiladores}.

\begin{nota}
Os dois repositórios refletem o \emph{split} intencional da plataforma: o
\repo{yanc} entrega binários de compilador versionados (\code{v1}--\code{v5.2})
que o \repo{aurora} baixa via \code{bootstrap}, enquanto o \repo{aurora}
desenvolve a IDE e publica releases estáveis no canal \repo{sapho}. A explosão
de atividade em maio--junho de 2026 nos dois repositórios corresponde à fase em
que a plataforma adquiriu sua infraestrutura de engenharia (testes, CI, releases
automatizadas) e seus subsistemas mais recentes (Aurora Intelligence, Source
Control embutido, \emph{language tooling}).
\end{nota}


% --------------------------- APÊNDICES ---------------------------
\appendix
% !TEX root = ../main.tex
\chapter{Glossário}
\label{ap:glossario}

Este apêndice reúne os termos, siglas e nomes de artefato recorrentes ao longo do documento.

\begin{description}[style=nextline,leftmargin=0pt,font=\bfseries\ttfamily]

\item[SAPHO] \textit{Scalable-Architecture Processor for Hardware Optimization}. A plataforma/suíte
desenvolvida pelo NIPSCERN que reúne a IDE \tool{AURORA}, os compiladores \tool{YANC} e o
\textit{toolchain bundle}. Também nomeia o canal de distribuição (repositório \repo{sapho}) e o
arquivo de projeto (\file{.spf}).

\item[AURORA] A IDE desktop (Electron + Monaco) da plataforma SAPHO, onde o usuário projeta,
compila, simula e visualiza \textit{soft-processors}. Repositório de desenvolvimento: \repo{aurora}.

\item[YANC] \textit{Yet Another Compiler}. O conjunto de compiladores e utilitários que traduzem
\cmm{} (ou C++) até Verilog sintetizável, imagens de memória e \textit{testbench}. Repositório:
\repo{yanc}.

\item[\cmm{} (CMM)] A linguagem de alto nível, semelhante a C, com suporte de primeira classe a
ponto-fixo, ponto-flutuante e números complexos, usada para descrever o programa de um processador
SAPHO. Arquivos com extensão \file{.cmm}.

\item[PRISM] O visualizador de RTL embutido na AURORA: converte Verilog em um diagrama de netlist
(via \tool{Yosys} e \tool{netlistsvg}).

\item[Aurora Intelligence] O subsistema de inteligência artificial da AURORA (assistente de chat e
ferramentas), com transportes via SDK (provedores com chave) e via CLIs por assinatura.

\item[POLARIS] Tentativa de sucessor multiplataforma da AURORA, \textbf{descontinuada}. Citada
apenas por contexto histórico; não faz parte do escopo deste documento.

\item[\file{.spf}] \textit{SAPHO Project File}. O arquivo canônico de configuração de um projeto
(processadores, arquivos sintetizáveis, \textit{testbenches}, \textit{top-level}, metadados).

\item[\file{.asm}] O assembly intermediário emitido pelos compiladores de \textit{front-end}
(\tool{cmmcomp}/\tool{cppcomp}) e consumido pelo \textit{back-end} (\tool{asmcomp}).

\item[\file{.mif}] \textit{Memory Initialization File}. Imagem de memória (programa e/ou dados)
gerada pelo \tool{asmcomp} para inicializar as memórias do core.

\item[testbench (\file{\_tb.v})] Banco de testes em Verilog gerado pelo \tool{asmcomp} para exercitar
o processador na simulação. Também pode ser um \textit{testbench} Python via \tool{cocotb}.

\item[VCD / FST] Formatos de \textit{dump} de simulação (\textit{Value Change Dump} e o formato
binário compacto do \tool{GTKWave}) que registram a evolução dos sinais ao longo do tempo.

\item[\file{.gtkw}] Arquivo de \textit{save} do \tool{GTKWave} que descreve a visualização
(sinais, cores, agrupamentos) a ser aplicada sobre um \textit{dump}.

\item[cmmcomp / cppcomp / asmcomp] Os três compiladores do YANC: \textit{front-ends} de \cmm{} e
C++ para assembly, e o \textit{back-end} de assembly para Verilog + memórias + \textit{testbench}.

\item[appcomp / cpppp] Pré-processadores do YANC: \tool{appcomp} faz a primeira passada do assembly
(parâmetros do processador e endereços de variáveis/labels); \tool{cpppp} é o pré-processador de C++.

\item[comp2gtkw / gen\_gtkw] Utilitários do YANC: tradução do padrão de bits de números complexos
para o \tool{GTKWave}, e geração de uma \textit{view} \file{.gtkw} a partir do cabeçalho de um VCD.

\item[Icarus Verilog (iverilog / vvp)] Simulador Verilog padrão da AURORA; \tool{iverilog} compila e
\tool{vvp} executa a simulação.

\item[Verilator] Simulador Verilog de alto desempenho (opt-in) que compila o modelo para C++.

\item[cocotb] \textit{Framework} Python para \textit{testbenches} de hardware, executável sobre
\tool{Icarus} ou \tool{Verilator}.

\item[Yosys] Ferramenta de síntese de RTL usada pelo PRISM para gerar a netlist em JSON.

\item[GTKWave] Visualizador de formas de onda; a AURORA empacota um \textit{fork} customizado
(\repo{gtkwave-nipscern}).

\item[Surfer] Visualizador de formas de onda alternativo, integrado de forma opcional à AURORA.

\item[netlistsvg] Conversor de netlist JSON em diagrama SVG; a AURORA usa um \textit{fork}
(\repo{netlistsvg-aurora}).

\item[Verible / slang-server] Servidores de linguagem (LSP) para Verilog/SystemVerilog que fornecem
diagnósticos e análise ao editor.

\item[tree-sitter] Biblioteca de \textit{parsing} incremental (via \tool{web-tree-sitter}) usada
para recursos de estrutura de código no editor.

\item[clang-format] Formatador de código usado pela IDE.

\item[Monaco Editor] O componente de editor de texto (o mesmo do VS Code) usado pela AURORA, fixado
na versão 0.52.2.

\item[Electron] O \textit{framework} que combina Chromium e Node.js para construir a aplicação
\textit{desktop}.

\item[Vite] O empacotador (\textit{bundler}) usado para construir o \textit{renderer} da AURORA.

\item[electron-builder / NSIS] Ferramenta de empacotamento e o sistema de instalador (Windows)
usados para produzir o executável distribuível.

\item[electron-updater] Biblioteca responsável pelas atualizações automáticas via GitHub Releases.

\item[MCP] \textit{Model Context Protocol}. Protocolo pelo qual as CLIs de IA acessam as ferramentas
da AURORA por meio de um servidor HTTP local.

\item[DPAPI / safeStorage] Mecanismo de criptografia do sistema operacional (via \tool{safeStorage}
do Electron) usado para proteger as chaves de API armazenadas.

\item[release-please] Automação que mantém o versionamento e o \textit{changelog} a partir de
\textit{Conventional Commits}.

\item[toolchain bundle] O pacote portátil \file{aurora-msys-vN.zip} (prefixo \texttt{msys/mingw64})
com simuladores, síntese e \tool{cocotb}, produzido pelo repositório \repo{aurora-toolchain}.

\item[numBits / MABits / MIBits] Diretivas de hardware do \cmm{} que parametrizam as larguras de
palavra e de endereçamento de memória do core gerado.

\end{description}

% !TEX root = ../main.tex
\chapter{Referências e fontes}
\label{ap:referencias}

Este documento foi construído a partir do estudo direto do \textbf{código-fonte} dos repositórios
da organização \texttt{nipscernlab} no GitHub. As referências abaixo apontam para esses
repositórios, para as ferramentas de terceiros orquestradas pela plataforma e para os padrões
relevantes.

\section{Repositórios da plataforma SAPHO}

\begin{description}[style=nextline,leftmargin=0pt]
\item[\repo{aurora}] A IDE AURORA (código estudado).
\url{https://github.com/nipscernlab/aurora}
\item[\repo{yanc}] Os compiladores YANC.
\url{https://github.com/nipscernlab/yanc}
\item[\repo{sapho}] O canal de distribuição (instalador da suíte).
\url{https://github.com/nipscernlab/sapho}
\item[\repo{aurora-toolchain}] O \textit{toolchain bundle} reprodutível (msys/mingw64).
\url{https://github.com/nipscernlab/aurora-toolchain}
\item[\repo{gtkwave-nipscern}] O \textit{fork} customizado do GTKWave.
\url{https://github.com/nipscernlab/gtkwave-nipscern}
\item[\repo{netlistsvg-aurora}] O \textit{fork} do netlistsvg usado pelo PRISM.
\url{https://github.com/nipscernlab/netlistsvg-aurora}
\item[\repo{nipscernweb}] O sítio do laboratório.
\url{https://github.com/nipscernlab/nipscernweb}
\end{description}

\section{Ferramentas de terceiros (projetos de origem)}

\begin{description}[style=nextline,leftmargin=0pt]
\item[Electron] \url{https://www.electronjs.org}
\item[Monaco Editor] \url{https://microsoft.github.io/monaco-editor/}
\item[Vite] \url{https://vitejs.dev}
\item[Icarus Verilog] \url{https://steveicarus.github.io/iverilog/}
\item[Verilator] \url{https://www.veripool.org/verilator/}
\item[Yosys] \url{https://yosyshq.net/yosys/}
\item[GTKWave] \url{https://gtkwave.sourceforge.net}
\item[cocotb] \url{https://www.cocotb.org}
\item[netlistsvg] \url{https://github.com/nturley/netlistsvg}
\item[DigitalJS] \url{https://github.com/tilk/digitaljs}
\item[Surfer] \url{https://surfer-project.org}
\item[Verible] \url{https://github.com/chipsalliance/verible}
\item[slang] \url{https://github.com/MikePopoloski/slang}
\item[tree-sitter] \url{https://tree-sitter.github.io/tree-sitter/}
\item[clang-format] \url{https://clang.llvm.org/docs/ClangFormat.html}
\item[MSYS2] \url{https://www.msys2.org}
\item[KaTeX] \url{https://katex.org}
\item[Lit] \url{https://lit.dev}
\item[Model Context Protocol] \url{https://modelcontextprotocol.io}
\item[Vercel AI SDK] \url{https://sdk.vercel.ai}
\item[electron-builder] \url{https://www.electron.build}
\item[electron-updater] \url{https://www.electron.build/auto-update}
\item[release-please] \url{https://github.com/googleapis/release-please}
\end{description}

\section{Padrões, formatos e convenções}

\begin{description}[style=nextline,leftmargin=0pt]
\item[Verilog HDL] IEEE Std 1364 --- linguagem de descrição de hardware alvo da síntese SAPHO.
\item[VCD] \textit{Value Change Dump} --- formato textual de \textit{dump} de simulação (parte do
padrão Verilog).
\item[FST] \textit{Fast Signal Trace} --- formato binário de \textit{dump} do GTKWave.
\item[Conventional Commits] Convenção de mensagens de \textit{commit} adotada no repositório.
\url{https://www.conventionalcommits.org}
\item[NSIS] \textit{Nullsoft Scriptable Install System} --- sistema de instalador para Windows.
\end{description}

\section{Documentos internos do repositório}

Os documentos abaixo existem nos repositórios e foram consultados apenas como pistas para localizar
o código; todas as afirmações deste relatório foram verificadas contra a implementação:
\file{README.md}, \file{ARCHITECTURE.md}, \file{CHANGELOG.md}, \file{ROADMAP.md}, \file{RELEASE.md},
\file{CONTRIBUTING.md}, \file{THIRD_PARTY_NOTICES.md}, \file{CITATION.cff} e os documentos em
\file{docs/} (entre eles \file{docs/CODE_SIGNING.md}, \file{docs/DESIGN.md} e os estudos internos).


\end{document}
